% GENERATED FILE -- DO NOT EDIT.
%
% Produced by harness/assemble_contract.py from module.yaml and claims/*.tex.
% Edit the claim files and regenerate; `--check` fails if this file is stale.
\documentclass[11pt,a4paper]{article}
\pdftrailerid{60be5848636939ca2ca280f7f0febda7}
%% --- 1. Mathematics (order matters: physics after braket wins \ket, \tr) ---
\usepackage{amsmath}
\usepackage{amssymb}
\usepackage{amsfonts}
\usepackage{amsthm}
\usepackage{mathtools}
\usepackage{dsfont}
\usepackage{braket}
\usepackage{physics}
%% --- 2. Page furniture ------------------------------------------------------
\usepackage{xcolor}
\usepackage{array}    % >{...} column prefixes in the wide tables of some claims
\usepackage{enumitem}
\usepackage{graphicx}
\usepackage{booktabs}
\usepackage{longtable}
\usepackage{xltabular}              % front- and back-matter tables across pages
\usepackage{needspace}
\usepackage{xurl}
\usepackage[noautomatic]{imakeidx}  % defined-terms index; latexmk runs makeindex
\makeindex[title=Index of defined terms,intoc]
\usepackage{tcolorbox}
\tcbuselibrary{skins,breakable}
%% --- 3. Links and navigation (hyperref late; bookmark, refcount after it) ---
\usepackage{hyperref}
\usepackage{bookmark}
\usepackage{refcount}
%% --- 4. Colours: every distinction also has a shape (rule, caps, size, weight)
\definecolor{nnGrey}{gray}{0.40}      % non-normative text
\definecolor{metaGrey}{gray}{0.30}    % the generated meta line
\definecolor{techGrey}{gray}{0.45}    % heading of a technical claim
\definecolor{mainRule}{RGB}{0,110,140}
\definecolor{reqAccent}{RGB}{165,70,0}
\definecolor{linkBlue}{RGB}{0,62,150}
\definecolor{extLink}{RGB}{120,40,140}  % links into another PDF
\hypersetup{colorlinks=true, linkcolor=linkBlue, citecolor=linkBlue,
  urlcolor=linkBlue, filecolor=extLink, bookmarksnumbered=true,
  bookmarksopen=true, bookmarksopenlevel=1, pdfdisplaydoctitle=true}
% A link into another PDF (\href to a .pdf, \hyperref{file}..., xr-hyper) is a
% relative URI, <file>#nameddest=<destination>, in the file-link colour:
% browser viewers (Edge, Chrome) ignore hyperref's default GoToR action.
\makeatletter
\renewcommand*{\hyper@linkfile}[3]{% anchor text, file name, destination
  \begingroup
    \ifx\\#3\\\def\CQ@target{#2}\else\def\CQ@target{#2\hyper@hash nameddest=#3}\fi
    \let\@urlcolor\@filecolor
    \hyper@linkurl{#1}{\CQ@target}%
  \endgroup}
\makeatother
%% --- 5. The mapped Unicode set (harness/contract_navigation.py, UNICODE_TEX)
\DeclareUnicodeCharacter{03B1}{\ensuremath{\alpha}}
\DeclareUnicodeCharacter{03B2}{\ensuremath{\beta}}
\DeclareUnicodeCharacter{03B3}{\ensuremath{\gamma}}
\DeclareUnicodeCharacter{03B4}{\ensuremath{\delta}}
\DeclareUnicodeCharacter{03B5}{\ensuremath{\varepsilon}}
\DeclareUnicodeCharacter{03F5}{\ensuremath{\epsilon}}
\DeclareUnicodeCharacter{03B6}{\ensuremath{\zeta}}
\DeclareUnicodeCharacter{03B7}{\ensuremath{\eta}}
\DeclareUnicodeCharacter{03B8}{\ensuremath{\theta}}
\DeclareUnicodeCharacter{03D1}{\ensuremath{\vartheta}}
\DeclareUnicodeCharacter{03B9}{\ensuremath{\iota}}
\DeclareUnicodeCharacter{03BA}{\ensuremath{\kappa}}
\DeclareUnicodeCharacter{03BB}{\ensuremath{\lambda}}
\DeclareUnicodeCharacter{03BC}{\ensuremath{\mu}}
\DeclareUnicodeCharacter{03BD}{\ensuremath{\nu}}
\DeclareUnicodeCharacter{03BE}{\ensuremath{\xi}}
\DeclareUnicodeCharacter{03C0}{\ensuremath{\pi}}
\DeclareUnicodeCharacter{03C1}{\ensuremath{\rho}}
\DeclareUnicodeCharacter{03F1}{\ensuremath{\varrho}}
\DeclareUnicodeCharacter{03C3}{\ensuremath{\sigma}}
\DeclareUnicodeCharacter{03C2}{\ensuremath{\varsigma}}
\DeclareUnicodeCharacter{03C4}{\ensuremath{\tau}}
\DeclareUnicodeCharacter{03C5}{\ensuremath{\upsilon}}
\DeclareUnicodeCharacter{03C6}{\ensuremath{\varphi}}
\DeclareUnicodeCharacter{03D5}{\ensuremath{\phi}}
\DeclareUnicodeCharacter{03C7}{\ensuremath{\chi}}
\DeclareUnicodeCharacter{03C8}{\ensuremath{\psi}}
\DeclareUnicodeCharacter{03C9}{\ensuremath{\omega}}
\DeclareUnicodeCharacter{0393}{\ensuremath{\Gamma}}
\DeclareUnicodeCharacter{0394}{\ensuremath{\Delta}}
\DeclareUnicodeCharacter{0398}{\ensuremath{\Theta}}
\DeclareUnicodeCharacter{039B}{\ensuremath{\Lambda}}
\DeclareUnicodeCharacter{039E}{\ensuremath{\Xi}}
\DeclareUnicodeCharacter{03A0}{\ensuremath{\Pi}}
\DeclareUnicodeCharacter{03A3}{\ensuremath{\Sigma}}
\DeclareUnicodeCharacter{03A5}{\ensuremath{\Upsilon}}
\DeclareUnicodeCharacter{03A6}{\ensuremath{\Phi}}
\DeclareUnicodeCharacter{03A8}{\ensuremath{\Psi}}
\DeclareUnicodeCharacter{03A9}{\ensuremath{\Omega}}
\DeclareUnicodeCharacter{00E1}{\'a}
\DeclareUnicodeCharacter{00E0}{\`a}
\DeclareUnicodeCharacter{00E2}{\^a}
\DeclareUnicodeCharacter{00E4}{\"a}
\DeclareUnicodeCharacter{00E5}{\aa{}}
\DeclareUnicodeCharacter{00E9}{\'e}
\DeclareUnicodeCharacter{00E8}{\`e}
\DeclareUnicodeCharacter{00EA}{\^e}
\DeclareUnicodeCharacter{00EB}{\"e}
\DeclareUnicodeCharacter{00ED}{\'{\i}}
\DeclareUnicodeCharacter{00EF}{\"{\i}}
\DeclareUnicodeCharacter{00F3}{\'o}
\DeclareUnicodeCharacter{00F2}{\`o}
\DeclareUnicodeCharacter{00F4}{\^o}
\DeclareUnicodeCharacter{00F6}{\"o}
\DeclareUnicodeCharacter{00F8}{\o{}}
\DeclareUnicodeCharacter{00FA}{\'u}
\DeclareUnicodeCharacter{00F9}{\`u}
\DeclareUnicodeCharacter{00FB}{\^u}
\DeclareUnicodeCharacter{00FC}{\"u}
\DeclareUnicodeCharacter{00F1}{\~n}
\DeclareUnicodeCharacter{00E7}{\c{c}}
\DeclareUnicodeCharacter{00DF}{\ss{}}
\DeclareUnicodeCharacter{00C4}{\"A}
\DeclareUnicodeCharacter{00C5}{\AA{}}
\DeclareUnicodeCharacter{00C9}{\'E}
\DeclareUnicodeCharacter{00D6}{\"O}
\DeclareUnicodeCharacter{00DC}{\"U}
\DeclareUnicodeCharacter{2013}{\textendash{}}
\DeclareUnicodeCharacter{2014}{\textemdash{}}
\DeclareUnicodeCharacter{2018}{`}
\DeclareUnicodeCharacter{2019}{'}
\DeclareUnicodeCharacter{201C}{``}
\DeclareUnicodeCharacter{201D}{''}
\DeclareUnicodeCharacter{2026}{\ldots{}}
\DeclareUnicodeCharacter{00A0}{~}
\DeclareUnicodeCharacter{00A7}{\S{}}
\DeclareUnicodeCharacter{2264}{\ensuremath{\le}}
\DeclareUnicodeCharacter{2265}{\ensuremath{\ge}}
\DeclareUnicodeCharacter{2260}{\ensuremath{\neq}}
\DeclareUnicodeCharacter{2248}{\ensuremath{\approx}}
\DeclareUnicodeCharacter{2261}{\ensuremath{\equiv}}
\DeclareUnicodeCharacter{223C}{\ensuremath{\sim}}
\DeclareUnicodeCharacter{2243}{\ensuremath{\simeq}}
\DeclareUnicodeCharacter{2245}{\ensuremath{\cong}}
\DeclareUnicodeCharacter{221D}{\ensuremath{\propto}}
\DeclareUnicodeCharacter{226A}{\ensuremath{\ll}}
\DeclareUnicodeCharacter{226B}{\ensuremath{\gg}}
\DeclareUnicodeCharacter{2192}{\ensuremath{\to}}
\DeclareUnicodeCharacter{2190}{\ensuremath{\leftarrow}}
\DeclareUnicodeCharacter{2194}{\ensuremath{\leftrightarrow}}
\DeclareUnicodeCharacter{21D2}{\ensuremath{\Rightarrow}}
\DeclareUnicodeCharacter{21D0}{\ensuremath{\Leftarrow}}
\DeclareUnicodeCharacter{21D4}{\ensuremath{\Leftrightarrow}}
\DeclareUnicodeCharacter{21A6}{\ensuremath{\mapsto}}
\DeclareUnicodeCharacter{2208}{\ensuremath{\in}}
\DeclareUnicodeCharacter{2209}{\ensuremath{\notin}}
\DeclareUnicodeCharacter{2286}{\ensuremath{\subseteq}}
\DeclareUnicodeCharacter{2282}{\ensuremath{\subset}}
\DeclareUnicodeCharacter{2287}{\ensuremath{\supseteq}}
\DeclareUnicodeCharacter{2283}{\ensuremath{\supset}}
\DeclareUnicodeCharacter{2229}{\ensuremath{\cap}}
\DeclareUnicodeCharacter{222A}{\ensuremath{\cup}}
\DeclareUnicodeCharacter{2216}{\ensuremath{\setminus}}
\DeclareUnicodeCharacter{2205}{\ensuremath{\emptyset}}
\DeclareUnicodeCharacter{2200}{\ensuremath{\forall}}
\DeclareUnicodeCharacter{2203}{\ensuremath{\exists}}
\DeclareUnicodeCharacter{2115}{\ensuremath{\mathbb{N}}}
\DeclareUnicodeCharacter{2124}{\ensuremath{\mathbb{Z}}}
\DeclareUnicodeCharacter{211A}{\ensuremath{\mathbb{Q}}}
\DeclareUnicodeCharacter{211D}{\ensuremath{\mathbb{R}}}
\DeclareUnicodeCharacter{2102}{\ensuremath{\mathbb{C}}}
\DeclareUnicodeCharacter{00B7}{\ensuremath{\cdot}}
\DeclareUnicodeCharacter{00D7}{\ensuremath{\times}}
\DeclareUnicodeCharacter{2212}{\ensuremath{-}}
\DeclareUnicodeCharacter{00B1}{\ensuremath{\pm}}
\DeclareUnicodeCharacter{2213}{\ensuremath{\mp}}
\DeclareUnicodeCharacter{221A}{\ensuremath{\surd}}
\DeclareUnicodeCharacter{2211}{\ensuremath{\sum}}
\DeclareUnicodeCharacter{220F}{\ensuremath{\prod}}
\DeclareUnicodeCharacter{222B}{\ensuremath{\int}}
\DeclareUnicodeCharacter{221E}{\ensuremath{\infty}}
\DeclareUnicodeCharacter{2202}{\ensuremath{\partial}}
\DeclareUnicodeCharacter{2207}{\ensuremath{\nabla}}
\DeclareUnicodeCharacter{2297}{\ensuremath{\otimes}}
\DeclareUnicodeCharacter{2295}{\ensuremath{\oplus}}
\DeclareUnicodeCharacter{2218}{\ensuremath{\circ}}
\DeclareUnicodeCharacter{2308}{\ensuremath{\lceil}}
\DeclareUnicodeCharacter{2309}{\ensuremath{\rceil}}
\DeclareUnicodeCharacter{230A}{\ensuremath{\lfloor}}
\DeclareUnicodeCharacter{230B}{\ensuremath{\rfloor}}
\DeclareUnicodeCharacter{27E8}{\ensuremath{\langle}}
\DeclareUnicodeCharacter{27E9}{\ensuremath{\rangle}}
\DeclareUnicodeCharacter{2016}{\ensuremath{\|}}
\DeclareUnicodeCharacter{2113}{\ensuremath{\ell}}
\DeclareUnicodeCharacter{00B0}{\ensuremath{^{\circ}}}
\DeclareUnicodeCharacter{00B2}{\ensuremath{^{2}}}
\DeclareUnicodeCharacter{00B3}{\ensuremath{^{3}}}
\DeclareUnicodeCharacter{00B9}{\ensuremath{^{1}}}
\DeclareUnicodeCharacter{2080}{\ensuremath{_{0}}}
\DeclareUnicodeCharacter{2081}{\ensuremath{_{1}}}
\DeclareUnicodeCharacter{2082}{\ensuremath{_{2}}}
\DeclareUnicodeCharacter{207F}{\ensuremath{^{n}}}
\DeclareUnicodeCharacter{2096}{\ensuremath{_{k}}}
%% --- 6. Numbering: one counter per section -----------------------
\numberwithin{equation}{section}
\setcounter{tocdepth}{1}
%% --- 7. The project macro set, then the claim environments -----------------
% spec/macros.tex, read relative to this module's directory (latexmk -cd).
% =============================================================================
%  spec/macros.tex -- the one macro set every generated contract inputs.
%
%  Two parts: part A, the macros spec/notation.md catalogues, and part B,
%  the macros local to one module, grouped by module.
%  harness/assemble_contract.py emits % =============================================================================
%  spec/macros.tex -- the one macro set every generated contract inputs.
%
%  Two parts: part A, the macros spec/notation.md catalogues, and part B,
%  the macros local to one module, grouped by module.
%  harness/assemble_contract.py emits % =============================================================================
%  spec/macros.tex -- the one macro set every generated contract inputs.
%
%  Two parts: part A, the macros spec/notation.md catalogues, and part B,
%  the macros local to one module, grouped by module.
%  harness/assemble_contract.py emits \input{../../macros} for it, so the
%  file is read from spec/modules/<id>/.
%
%  RULES
%    * A macro that spec/notation.md catalogues is a verbatim copy of it
%      (part A). Changing one is a notation change: bump spec/notation.md.
%    * Part B holds the module-local macros. A macro defined by two modules
%      with the same expansion appears once, under the first module (in the
%      order below) that defines it. Two macros with the same name and
%      DIFFERENT expansions are both kept under a module prefix, and the
%      claims write the prefixed names:
%        \Chain  -> \redChain (reductions, \mathfrak{C}: the reduction chain)
%                   \pcChain  (protocol-compiler, \mathfrak{R}: the two-segment
%                             chain; \Chsrc and \Chrec are built on it)
%        \Nmsm   -> \redNmsm  (reductions, \mathcal{N}: a POVM measurement
%                             channel)
%                   \pcNmsm   (\mathcal{N}_{\mathsf A}: Alice's measurement
%                             channel, written by
%                             transmitter/thm:certificate-sufficiency,
%                             transmitter/rem:xreg and
%                             reductions/lem:source-transfer)
%    * Deliberately omitted: a \provenance{tag}{details} macro. Provenance
%      is a header field of every claim file (spec/claim-format.md section
%      3); a claim body that calls such a macro fails to compile, which is
%      the intended signal.
%
%  REQUIRED PACKAGES, loaded by the assembler's preamble before this file:
%    amsmath, amssymb, amsfonts, mathtools, dsfont, braket, physics, xcolor
%    (braket before physics, so that physics wins for \ket, \bra, \tr, ...).
%
%  This file contains only definitions: no \usepackage, no document body.
% =============================================================================


% =============================================================================
%  PART A -- notation macros, the ones spec/notation.md catalogues, each a
%  verbatim copy of its catalogue row; grouped under numbered headings.
% =============================================================================

% -----------------------------------------------------------------------------
% 1. State spaces and operator sets (a "system label" -> spaces of operators)
% -----------------------------------------------------------------------------
\newcommand{\Hs}[1]{\mathcal{H}(#1)}        % Hilbert space of a system label
\newcommand{\Lin}[1]{\mathcal{L}(#1)}       % bounded linear operators on H(#1)
\newcommand{\Linn}[2]{\mathcal{L}(#1,#2)}   % linear maps H(#1) -> H(#2)
\newcommand{\Pos}[1]{\mathrm{Pos}(#1)}      % positive semidefinite operators on H(#1)
\newcommand{\Density}{\mathrm{D}}           % density operators: \Density(A) = D(A)
\newcommand{\Seq}[1]{S_{=}(#1)}             % normalized states (trace = 1)
\newcommand{\Sleq}[1]{S_{\leq}(#1)}         % subnormalized states (0 <= trace <= 1)
\newcommand{\id}{\mathds{1}}                % identity operator
\newcommand{\idon}[1]{\mathds{1}_{#1}}      % identity on a named system, e.g. \idon{C}

% DEPRECATED: applies \mathcal to both arguments itself, so it takes bare
% letters, not system labels. Prefer \Linn; do not use in new contracts.
\providecommand{\Lop}[2]{\mathcal{L}(\mathcal{#1}, \mathcal{#2})}


% -----------------------------------------------------------------------------
% 2. Composite systems
% -----------------------------------------------------------------------------
\newcommand{\dsum}{\oplus}                  % direct sum of system labels
% Tensor product: use \otimes directly.
% n-fold direct sum of an operator space: \Lin{A}^{\oplus n}


% -----------------------------------------------------------------------------
% 3. Events, conditioning and classical registers (for STATES)
%    Convention:  \wedge  -> subnormalized "and-event" component, e.g. \rho_{A\wedge x}
%                 |        -> normalized conditional state,        e.g. \rho_{A|x}
% (Write these directly in subscripts; no extra macros needed.)
% -----------------------------------------------------------------------------


% -----------------------------------------------------------------------------
% 4. Photon-number subspaces / direct-sum blocks
%    Convention:  operators carry the restriction "on top", in brackets;
%                 states carry it as a \wedge subscript (see Sec. 3).
% -----------------------------------------------------------------------------
\newcommand{\Pnum}[1]{P_{#1}}               % photon-number projector (system from context)
\newcommand{\Pnumon}[2]{P_{#1}^{#2}}        % with explicit system superscript, e.g. P_n^C
\newcommand{\blk}[2]{#1^{(#2)}}             % operator block, e.g. \blk{M}{\leqN} = M^{(<=N)}
\newcommand{\leqN}{\leq N}                  % use as \blk{M}{\leqN}
\newcommand{\gtN}{> N}                      % use as \blk{M}{\gtN}
\newcommand{\eqn}{= n}                      % use as \blk{M}{\eqn}


% -----------------------------------------------------------------------------
% 5. System labels used in the protocol
% -----------------------------------------------------------------------------
\newcommand{\Chan}{C}                       % quantum-channel (Eve) input to Bob
\newcommand{\Prb}{P}                        % probe input
\newcommand{\Flag}{{\mathrm{Flg}}}          % flag system (EXISTS ONLY AFTER squashing)
\newcommand{\Loss}{{\mathrm{Lss}}}          % loss system
\newcommand{\Vc}{{\mathrm{Vcm}}}            % vacuum system
\newcommand{\Nout}{N_{\mathrm{out}}}
\newcommand{\Nin}{N_{\mathrm{in}}}


% -----------------------------------------------------------------------------
% 6. Alphabets
% -----------------------------------------------------------------------------
\newcommand{\Yset}{\mathcal{Y}}             % Bob's (coarse-grained) detection alphabet
\newcommand{\Yraw}{\mathcal{Y}_{\mathrm{raw}}}% raw click patterns {0,1}^4
\newcommand{\Yscset}{\mathcal{Y}_{\singlec}}% single-click outcomes
\newcommand{\Ymcset}{\mathcal{Y}_{\mc}}     % multi-click outcomes
\newcommand{\Tset}{\mathcal{T}_{\mathrm{set}}}% private setting set (the probe settings in the optical instance)
\newcommand{\Xset}{\mathcal{X}}             % Alice's alphabet
\newcommand{\Dset}{\{H,V,+,-\}}             % polarization / detector labels
\newcommand{\pass}{\top}                    % pass sifting


% -----------------------------------------------------------------------------
% 7. Mode / outcome abbreviations
% -----------------------------------------------------------------------------
\newcommand{\on}{{\mathrm{on}}}
\newcommand{\off}{{\mathrm{off}}}
\newcommand{\dc}{{\mathrm{dc}}}             % dark count
\newcommand{\sq}{\mathrm{sq}}               % squashed
\newcommand{\mc}{\mathrm{mc}}               % multi-click
\newcommand{\nc}{\mathrm{nc}}               % no-click
\newcommand{\singlec}{\mathrm{sc}}          % single-click
\newcommand{\vac}{\mathrm{vac}}             % vacuum
\newcommand{\kout}{{k_\mathrm{out}}}
\newcommand{\kin}{{k_\mathrm{in}}}


% -----------------------------------------------------------------------------
% 8. Maps, channels and miscellaneous operators
% -----------------------------------------------------------------------------
\newcommand{\Prob}{\mathrm{Prob}}           % probability (use \Prob[\cdot] or \Prob(\cdot))
\newcommand{\lmin}{\lambda_{\min}}          % minimum eigenvalue
\newcommand{\Wflag}{W_{\Flag}}              % flag-space weight
\newcommand{\transpose}{\mathsf{T}}
\newcommand{\innerp}[2]{\left\langle #1 \vert #2 \right\rangle}
\renewcommand{\vec}[1]{\boldsymbol{#1}}     % NB: overrides the LaTeX arrow accent


% -----------------------------------------------------------------------------
% 9. Key-rate numerics
% -----------------------------------------------------------------------------
\newcommand{\cG}{\mathcal{G}}               % protocol map
\newcommand{\cZ}{\mathcal{Z}}               % pinching map
\newcommand{\Amap}{\mathcal{A}_{\mathrm{ann}}} % announcement + measurement map
\newcommand{\Amapnb}{\mathcal{A}_{\mathrm{ann}|\mathrm{nb}}}
\newcommand{\simpl}{}                       % empty suffix for simplified block maps
\newcommand{\acSet}{\mathcal{S}_{\mathrm{acc}}} % acceptance set
\newcommand{\testE}{\mathcal{E}_{\mathrm{test}}} % test-statistics map
\newcommand{\ecLeak}{\mathrm{leak}^{\mathrm{EC}}_{\mathrm{obs}}}
\newcommand{\pgen}{p_{\mathrm{gen}}}
\newcommand{\askeyr}{R_{\infty}}             % asymptotic key rate
\newcommand{\ppass}{p_{\mathrm{pass}}}       % sifting pass probability


% -----------------------------------------------------------------------------
% 10. Finite-size analysis / Renyi entropies
% -----------------------------------------------------------------------------

% --- Round types: test vs. key generation (register T, alphabet \Fset) --------
\newcommand{\tst}{\mathrm{test}}            % test round
\newcommand{\gen}{\mathrm{gen}}             % key-generation round
\newcommand{\Fset}{\mathcal{T}}             % round-type alphabet {test, gen} (not \Tset)
\newcommand{\Chat}{\widehat{C}}             % acceptance-test register
\newcommand{\Chatset}{\widehat{\mathcal{C}}}% its alphabet

% --- Basis / key extraction functions ----------------------------------------
\newcommand{\bA}{b_{A}}                     % Alice's basis-extraction function
\newcommand{\bB}{b_{B}}                     % Bob's basis-extraction function
\newcommand{\kB}{k_{B}}                     % Bob's key-extraction function

% --- Maps --------------------------------------------------------------------
\newcommand{\Mprot}{\mathcal{M}}            % single-round protocol map
\newcommand{\Phimap}{\boldsymbol{\Phi}}     % test-statistics map
\newcommand{\Pdeph}{\mathcal{P}}            % dephasing (pinching) map on \Blreg

% --- Isometries --------------------------------------------------------------
\newcommand{\Vmsm}{V_{\mathrm{msm}}}        % measurement isometry
\newcommand{\Vann}{V_{\mathrm{ann}}}        % announcement isometry
\newcommand{\VI}{V_{I}}                     % Eve's copy of the announcements
\newcommand{\VChat}{V_{\Chat}}              % acceptance-data isometry
\newcommand{\VR}{V_{R}}                     % sifting + key map -> raw key R
\newcommand{\VZ}{V_{Z}}                     % raw key -> key register Z
\newcommand{\Wprot}{W}                      % full protocol isometry

% --- Renyi entropies ---------------------------------------------------------
% SECURITY-CRITICAL: see spec/notation.md §4. Never change an arrow to
% make an inequality work out; never convert silently between the sandwiched
% (tilde) and Petz (overline) families.
\newcommand{\Hsandup}{\widetilde{H}^{\uparrow}_{\alpha}}   % sandwiched, up
\newcommand{\Hsanddown}{\widetilde{H}^{\downarrow}_{\alpha}}% sandwiched, down
\newcommand{\Hpetzup}{\overline{H}^{\uparrow}_{\alpha}}    % Petz, up
\newcommand{\Hpetzdown}{\overline{H}^{\downarrow}_{\alpha}}% Petz, down
\newcommand{\hup}{h^{\uparrow}_{\alpha}}                   % single-round rate
\newcommand{\hqkd}{h^{\uparrow,\mathrm{QKD}}_{\alpha}}     % single-round QKD rate
\newcommand{\hlow}{\underline{h}_{\alpha}}                 % lower bound on \hqkd
% --- spec/notation.md 3.5 and 4.5 -------------------------------------------
% \hlb  is the first-order linearization lower bound on \hlow: the headline
%       key rate (the `first-order` tier of foundations/def:tiers). NOT \hlow
%       itself.
% \tnum is the instance's one declared numerical tolerance. Its
%       near-collision with the declared marginal \tau_A is disclosed in
%       spec/notation.md 3.5; bare \tau is reserved and denotes neither.
\newcommand{\hlb}{\underline{h}^{\mathrm{FO}}_{\alpha}}    % first-order lower bound on \hlow
\newcommand{\tnum}{\tau_{\mathrm{num}}}                    % declared numerical tolerance


% =============================================================================
%  PART B -- module-local macros, grouped by module. Deprecated aliases,
%  defined so that a text written with them compiles, used in no new contract: \Lop,
%  \testE, \Pdeph, \pass, \ecLeak (spec/notation.md section 3.9).
% =============================================================================

% --- meat-core -----------------------------------------------------------------
\newcommand{\SecSet}{\mathbb{S}}            % sector index set of Q
\newcommand{\Qsys}{Q}                       % declared sector-carrying measured system
\newcommand{\Vset}{\mathcal{V}}             % declared joint outcome alphabet of the Q-POVM
\newcommand{\CAset}{\mathcal{C}_{A}}        % Alice announcement alphabet
\newcommand{\CBset}{\mathcal{C}_{B}}        % Bob announcement alphabet
\newcommand{\Party}{\mathfrak{P}}           % declared party partition
\newcommand{\Oset}[1]{\mathcal{O}_{#1}}     % readout alphabet of party #1
\newcommand{\rd}[1]{\pi_{#1}}               % readout map of party #1
\newcommand{\Sadm}{\mathcal{S}_{A}}         % admissible Alice subspace
\newcommand{\nomarg}{\mathsf{none}}         % "no trusted marginal declared"
\newcommand{\Zset}{\mathcal{Z}_{\perp}}     % raw-key alphabet incl. \perp
\newcommand{\Zval}{\mathcal{Z}}             % key alphabet without \perp
\newcommand{\Etil}{\widetilde{E}}           % single-round purifying register
\newcommand{\Ehat}{\widehat{E}}             % Eve's n-round register
\newcommand{\OmAT}{\Omega_{\mathrm{AT}}}    % acceptance event
\newcommand{\epa}{\varepsilon_{\mathrm{PA}}}
\newcommand{\eev}{\varepsilon_{\mathrm{EV}}}
\newcommand{\lEC}{\lambda_{\mathrm{EC}}}
\newcommand{\Decl}{\mathfrak{D}}            % protocol declaration
\newcommand{\CList}{\mathfrak{L}}           % declared constraint list
\newcommand{\PinS}{\mathcal{P}_{\SecSet}}   % sector pinching on Q
\newcommand{\Feas}{\mathcal{W}}             % feasible set of block states
\newcommand{\Chkey}{\Chatset_{\mathrm{key}}}% key-live acceptance values
\newcommand{\Kann}{\mathcal{K}}             % announcement-class alphabet
\newcommand{\hold}{\mathsf{h}}              % declared key holder (A or B)
\newcommand{\sfD}{D}                        % classical KL divergence (base 2)
\newcommand{\Dtwo}{D}                       % quantum relative entropy (base 2)
\newcommand{\Dsand}{\widetilde{D}_{\alpha}} % sandwiched Renyi divergence
\newcommand{\idmap}{\mathrm{id}}            % identity channel
\newcommand{\Aev}{\mathcal{A}}              % Eve's per-round attack channel
\newcommand{\AttSet}{\mathfrak{A}}          % licensed attack-channel set
\newcommand{\Chacc}{\Chatset_{\mathrm{acc}}}% acceptance support of \acSet
\newcommand{\epsstar}{\varepsilon^{\star}}      % distance to an exact POVM
\newcommand{\dfH}{\Delta_{\mathrm{herm}}}       % Hermiticity defect
\newcommand{\dfP}{\Delta_{\mathrm{pos}}}        % negative part
\newcommand{\dfC}{\Delta_{\mathrm{cpl}}}        % completeness defect
\newcommand{\lmoq}{g}                       % acceptance-side gradient block
\newcommand{\lmoJ}{G}                       % feasible-side gradient block
\newcommand{\bacc}{\beta_{\mathrm{acc}}}    % certified acceptance-side minimum
\newcommand{\Itset}{\mathfrak{I}}           % finite index set of iterates (C22)

% --- reductions ----------------------------------------------------------------
\newcommand{\Bsys}{B}                        % physical (pre-squashing) system
\newcommand{\Tcl}[1]{\mathcal{T}(#1)}        % trace-class operators on H(#1)
\newcommand{\Ins}{\mathcal{I}}               % generalized-instrument branch map
\newcommand{\redNmsm}{\mathcal{N}}              % POVM measurement channel
\newcommand{\Eprot}{\mathcal{E}}             % n-round protocol channel
\newcommand{\Grnd}{\mathsf{G}}               % one round map of a protocol
\newcommand{\Kfin}{\mathsf{K}}               % end-of-protocol processing
\newcommand{\Crit}{\mathsf{sec}}             % value of a security criterion
\newcommand{\Pflag}{\Pi_{\Flag}}             % flag-sector projector on Q
\newcommand{\Pleq}{\Pi_{\leqN}}              % preserved-block projector
\newcommand{\Pgt}{\Pi_{\gtN}}                % above-cutoff projector
\newcommand{\Sph}{S^{\mathrm{ph}}}           % family of physical inputs
\newcommand{\Srch}{S^{\mathrm{rch}}}         % reachable squashed states
\newcommand{\Ssq}{S^{\sq}}                   % squashed admissible set
\newcommand{\Ffss}{F^{\mathrm{FSS}}}         % flag-state-squasher target
\newcommand{\Fwp}{F^{\mathrm{WP}}}           % multi-bound WPFSS target
\newcommand{\Mwp}{M^{\mathrm{WP}}}           % renormalized above-cutoff POVM
\newcommand{\Mdec}{{M'}^{\mathrm{ph}}}       % declared physical family on B
\newcommand{\Ppvm}{P^{\Flag}}                % flag PVM
\newcommand{\Lmp}{{\Lambda_{\mathrm{mp}}}}   % measure-and-prepare channel
\newcommand{\Lwp}{{\Lambda^{\mathrm{WP}}}}   % multi-bound WPFSS channel
\newcommand{\Lfss}{{\Lambda^{\mathrm{FSS}}}} % flag-state-squasher channel
\newcommand{\lmax}{\lambda_{\max}}           % maximum eigenvalue
\newcommand{\Lnoise}{\widetilde{\Lambda}}    % flag-noise channel
\newcommand{\cg}{\varpi}                     % declared coarse-graining map
\newcommand{\MCK}{\mathsf{MC}}               % multiclick raw-pattern set
\newcommand{\redChain}{\mathfrak{C}}            % reduction chain
\newcommand{\Lnk}{\ell}                      % a single typed link
\newcommand{\dsq}{\preceq}                      % the squashing relation
\newcommand{\opn}[1]{\left\lVert #1\right\rVert}           % operator (spectral) norm
\newcommand{\mxn}[1]{\left\lVert #1\right\rVert_{\max}}    % entrywise max modulus
\newcommand{\Fld}{\mathbb{F}}                   % exact ordered substrate field
\newcommand{\Fldi}{\mathbb{F}_{i}}              % its Gaussian extension F + iF
\newcommand{\lat}[1]{\big\lfloor #1\big\rfloor_{\kappa}} % round DOWN onto the lattice

% --- protocol-compiler ---------------------------------------------------------
\newcommand{\Spec}{\mathfrak{S}}             % protocol specification (compiler input)
\newcommand{\pcChain}{\mathfrak{R}}            % reduction chain (two typed segments)
\newcommand{\Chsrc}{\pcChain_{\mathrm{src}}}   % its Alice-side source-map segment
\newcommand{\Chrec}{\pcChain_{\mathrm{rec}}}   % its receiver-side squashing segment
\newcommand{\rvirt}{\rho^{\mathrm{v}}}       % virtual (PRE-source-map) signal state
\newcommand{\Fph}{F^{\mathrm{ph}}}           % physical (pre-reduction) POVM element
\newcommand{\Vph}{\mathcal{V}^{\mathrm{ph}}} % physical outcome alphabet
\newcommand{\Tpub}{\mathcal{T}_{\mathrm{pub}}}% public-choice alphabet
\newcommand{\Uset}{\mathcal{U}}              % residual outcome alphabet
\newcommand{\idxm}{\mathsf{idx}}             % label-to-index map
\newcommand{\Asig}{A'}                       % transmitted (signal) system -- the REAL one
\newcommand{\Avirt}{A''}                     % virtual transmitted system (source descriptor of Xi_1)
\newcommand{\Abar}{\bar{A}}                  % Alice's measured register
\newcommand{\Ahat}{\hat{A}}                  % Alice's shield register
\newcommand{\pcNmsm}{\mathcal{N}_{\mathsf A}}  % Alice's measurement channel (see the RULES block)
\newcommand{\Phip}{\Phi^{+}}                 % maximally entangled vector label
\newcommand{\Smap}{\mathfrak{s}}             % the emitted sector map (fraktur: a compile-time object)

% --- receiver ------------------------------------------------------------------
\newcommand{\Modes}{\mathbb{M}}              % declared ordered mode set
\newcommand{\Inp}{\mathbb{I}}                % declared input modes
\newcommand{\InpC}{\mathbb{I}_{\Chan}}       % channel input modes
\newcommand{\InpS}{\mathbb{I}_{\Prb}}        % probe (source) input modes
\newcommand{\Env}{\mathbb{E}}                % dilation (environment) modes
\newcommand{\Port}[1]{\mathbb{P}_{#1}}       % mode set of detector port #1
\newcommand{\Emd}[1]{E_{#1}}                 % coordinate projector on modes #1
\newcommand{\Gmat}{G}                        % instrument (transfer) matrix
\newcommand{\Fock}{\mathfrak{F}}             % Fock / second-quantization functor
\newcommand{\Symn}[1]{\mathrm{Sym}^{#1}}     % symmetric tensor power
\newcommand{\normord}[1]{{:}\mspace{2mu}#1\mspace{2mu}{:}}   % normal ordering
\newcommand{\rk}{\operatorname{rank}}        % rank (physics.sty owns \rank)
\newcommand{\SCK}{\mathsf{SC}}               % single-click pattern set
\newcommand{\Recv}{\mathfrak{R}}             % receiver instrument tuple
\newcommand{\Cat}{\mathfrak{C}}              % component catalogue
\newcommand{\Grph}{\mathfrak{E}}             % declared element graph
\newcommand{\defic}{\mathbb{L}}              % loss deficiency operator
\newcommand{\Pvac}[1]{\Pi_{\vac}(#1)}        % joint no-click element on ports #1
\newcommand{\pmc}{p_{\mc}}                   % uniform high-photon mc weight
\newcommand{\Bnd}{\beta}                     % the scalar bound function beta(n)
\newcommand{\Rot}{R}                         % planar rotation matrix
\newcommand{\Wplate}{W}                      % retarder (wave-plate) matrix
\newcommand{\Bsp}{B}                         % beamsplitter core matrix

% --- statistics ----------------------------------------------------------------
\newcommand{\honest}{p^{\mathrm{hon}}}          % honest per-round distribution
\newcommand{\ecomp}{\varepsilon_{\mathrm{comp}}}% completeness budget
\newcommand{\simplex}[1]{\Delta(#1)}            % probability simplex
\newcommand{\Chatcon}{\Chatset_{\mathrm{con}}}  % constrained outcomes
\newcommand{\Ireg}[3]{\mathcal{I}_{#1}(#2,#3)}  % regularized incomplete beta
\newcommand{\Ncount}[1]{\mathsf{N}_{#1}}        % per-outcome count
\newcommand{\qlo}{\mathbf{q}^{-}}               % box lower bounds
\newcommand{\qhi}{\mathbf{q}^{+}}               % box upper bounds
 for it, so the
%  file is read from spec/modules/<id>/.
%
%  RULES
%    * A macro that spec/notation.md catalogues is a verbatim copy of it
%      (part A). Changing one is a notation change: bump spec/notation.md.
%    * Part B holds the module-local macros. A macro defined by two modules
%      with the same expansion appears once, under the first module (in the
%      order below) that defines it. Two macros with the same name and
%      DIFFERENT expansions are both kept under a module prefix, and the
%      claims write the prefixed names:
%        \Chain  -> \redChain (reductions, \mathfrak{C}: the reduction chain)
%                   \pcChain  (protocol-compiler, \mathfrak{R}: the two-segment
%                             chain; \Chsrc and \Chrec are built on it)
%        \Nmsm   -> \redNmsm  (reductions, \mathcal{N}: a POVM measurement
%                             channel)
%                   \pcNmsm   (\mathcal{N}_{\mathsf A}: Alice's measurement
%                             channel, written by
%                             transmitter/thm:certificate-sufficiency,
%                             transmitter/rem:xreg and
%                             reductions/lem:source-transfer)
%    * Deliberately omitted: a \provenance{tag}{details} macro. Provenance
%      is a header field of every claim file (spec/claim-format.md section
%      3); a claim body that calls such a macro fails to compile, which is
%      the intended signal.
%
%  REQUIRED PACKAGES, loaded by the assembler's preamble before this file:
%    amsmath, amssymb, amsfonts, mathtools, dsfont, braket, physics, xcolor
%    (braket before physics, so that physics wins for \ket, \bra, \tr, ...).
%
%  This file contains only definitions: no \usepackage, no document body.
% =============================================================================


% =============================================================================
%  PART A -- notation macros, the ones spec/notation.md catalogues, each a
%  verbatim copy of its catalogue row; grouped under numbered headings.
% =============================================================================

% -----------------------------------------------------------------------------
% 1. State spaces and operator sets (a "system label" -> spaces of operators)
% -----------------------------------------------------------------------------
\newcommand{\Hs}[1]{\mathcal{H}(#1)}        % Hilbert space of a system label
\newcommand{\Lin}[1]{\mathcal{L}(#1)}       % bounded linear operators on H(#1)
\newcommand{\Linn}[2]{\mathcal{L}(#1,#2)}   % linear maps H(#1) -> H(#2)
\newcommand{\Pos}[1]{\mathrm{Pos}(#1)}      % positive semidefinite operators on H(#1)
\newcommand{\Density}{\mathrm{D}}           % density operators: \Density(A) = D(A)
\newcommand{\Seq}[1]{S_{=}(#1)}             % normalized states (trace = 1)
\newcommand{\Sleq}[1]{S_{\leq}(#1)}         % subnormalized states (0 <= trace <= 1)
\newcommand{\id}{\mathds{1}}                % identity operator
\newcommand{\idon}[1]{\mathds{1}_{#1}}      % identity on a named system, e.g. \idon{C}

% DEPRECATED: applies \mathcal to both arguments itself, so it takes bare
% letters, not system labels. Prefer \Linn; do not use in new contracts.
\providecommand{\Lop}[2]{\mathcal{L}(\mathcal{#1}, \mathcal{#2})}


% -----------------------------------------------------------------------------
% 2. Composite systems
% -----------------------------------------------------------------------------
\newcommand{\dsum}{\oplus}                  % direct sum of system labels
% Tensor product: use \otimes directly.
% n-fold direct sum of an operator space: \Lin{A}^{\oplus n}


% -----------------------------------------------------------------------------
% 3. Events, conditioning and classical registers (for STATES)
%    Convention:  \wedge  -> subnormalized "and-event" component, e.g. \rho_{A\wedge x}
%                 |        -> normalized conditional state,        e.g. \rho_{A|x}
% (Write these directly in subscripts; no extra macros needed.)
% -----------------------------------------------------------------------------


% -----------------------------------------------------------------------------
% 4. Photon-number subspaces / direct-sum blocks
%    Convention:  operators carry the restriction "on top", in brackets;
%                 states carry it as a \wedge subscript (see Sec. 3).
% -----------------------------------------------------------------------------
\newcommand{\Pnum}[1]{P_{#1}}               % photon-number projector (system from context)
\newcommand{\Pnumon}[2]{P_{#1}^{#2}}        % with explicit system superscript, e.g. P_n^C
\newcommand{\blk}[2]{#1^{(#2)}}             % operator block, e.g. \blk{M}{\leqN} = M^{(<=N)}
\newcommand{\leqN}{\leq N}                  % use as \blk{M}{\leqN}
\newcommand{\gtN}{> N}                      % use as \blk{M}{\gtN}
\newcommand{\eqn}{= n}                      % use as \blk{M}{\eqn}


% -----------------------------------------------------------------------------
% 5. System labels used in the protocol
% -----------------------------------------------------------------------------
\newcommand{\Chan}{C}                       % quantum-channel (Eve) input to Bob
\newcommand{\Prb}{P}                        % probe input
\newcommand{\Flag}{{\mathrm{Flg}}}          % flag system (EXISTS ONLY AFTER squashing)
\newcommand{\Loss}{{\mathrm{Lss}}}          % loss system
\newcommand{\Vc}{{\mathrm{Vcm}}}            % vacuum system
\newcommand{\Nout}{N_{\mathrm{out}}}
\newcommand{\Nin}{N_{\mathrm{in}}}


% -----------------------------------------------------------------------------
% 6. Alphabets
% -----------------------------------------------------------------------------
\newcommand{\Yset}{\mathcal{Y}}             % Bob's (coarse-grained) detection alphabet
\newcommand{\Yraw}{\mathcal{Y}_{\mathrm{raw}}}% raw click patterns {0,1}^4
\newcommand{\Yscset}{\mathcal{Y}_{\singlec}}% single-click outcomes
\newcommand{\Ymcset}{\mathcal{Y}_{\mc}}     % multi-click outcomes
\newcommand{\Tset}{\mathcal{T}_{\mathrm{set}}}% private setting set (the probe settings in the optical instance)
\newcommand{\Xset}{\mathcal{X}}             % Alice's alphabet
\newcommand{\Dset}{\{H,V,+,-\}}             % polarization / detector labels
\newcommand{\pass}{\top}                    % pass sifting


% -----------------------------------------------------------------------------
% 7. Mode / outcome abbreviations
% -----------------------------------------------------------------------------
\newcommand{\on}{{\mathrm{on}}}
\newcommand{\off}{{\mathrm{off}}}
\newcommand{\dc}{{\mathrm{dc}}}             % dark count
\newcommand{\sq}{\mathrm{sq}}               % squashed
\newcommand{\mc}{\mathrm{mc}}               % multi-click
\newcommand{\nc}{\mathrm{nc}}               % no-click
\newcommand{\singlec}{\mathrm{sc}}          % single-click
\newcommand{\vac}{\mathrm{vac}}             % vacuum
\newcommand{\kout}{{k_\mathrm{out}}}
\newcommand{\kin}{{k_\mathrm{in}}}


% -----------------------------------------------------------------------------
% 8. Maps, channels and miscellaneous operators
% -----------------------------------------------------------------------------
\newcommand{\Prob}{\mathrm{Prob}}           % probability (use \Prob[\cdot] or \Prob(\cdot))
\newcommand{\lmin}{\lambda_{\min}}          % minimum eigenvalue
\newcommand{\Wflag}{W_{\Flag}}              % flag-space weight
\newcommand{\transpose}{\mathsf{T}}
\newcommand{\innerp}[2]{\left\langle #1 \vert #2 \right\rangle}
\renewcommand{\vec}[1]{\boldsymbol{#1}}     % NB: overrides the LaTeX arrow accent


% -----------------------------------------------------------------------------
% 9. Key-rate numerics
% -----------------------------------------------------------------------------
\newcommand{\cG}{\mathcal{G}}               % protocol map
\newcommand{\cZ}{\mathcal{Z}}               % pinching map
\newcommand{\Amap}{\mathcal{A}_{\mathrm{ann}}} % announcement + measurement map
\newcommand{\Amapnb}{\mathcal{A}_{\mathrm{ann}|\mathrm{nb}}}
\newcommand{\simpl}{}                       % empty suffix for simplified block maps
\newcommand{\acSet}{\mathcal{S}_{\mathrm{acc}}} % acceptance set
\newcommand{\testE}{\mathcal{E}_{\mathrm{test}}} % test-statistics map
\newcommand{\ecLeak}{\mathrm{leak}^{\mathrm{EC}}_{\mathrm{obs}}}
\newcommand{\pgen}{p_{\mathrm{gen}}}
\newcommand{\askeyr}{R_{\infty}}             % asymptotic key rate
\newcommand{\ppass}{p_{\mathrm{pass}}}       % sifting pass probability


% -----------------------------------------------------------------------------
% 10. Finite-size analysis / Renyi entropies
% -----------------------------------------------------------------------------

% --- Round types: test vs. key generation (register T, alphabet \Fset) --------
\newcommand{\tst}{\mathrm{test}}            % test round
\newcommand{\gen}{\mathrm{gen}}             % key-generation round
\newcommand{\Fset}{\mathcal{T}}             % round-type alphabet {test, gen} (not \Tset)
\newcommand{\Chat}{\widehat{C}}             % acceptance-test register
\newcommand{\Chatset}{\widehat{\mathcal{C}}}% its alphabet

% --- Basis / key extraction functions ----------------------------------------
\newcommand{\bA}{b_{A}}                     % Alice's basis-extraction function
\newcommand{\bB}{b_{B}}                     % Bob's basis-extraction function
\newcommand{\kB}{k_{B}}                     % Bob's key-extraction function

% --- Maps --------------------------------------------------------------------
\newcommand{\Mprot}{\mathcal{M}}            % single-round protocol map
\newcommand{\Phimap}{\boldsymbol{\Phi}}     % test-statistics map
\newcommand{\Pdeph}{\mathcal{P}}            % dephasing (pinching) map on \Blreg

% --- Isometries --------------------------------------------------------------
\newcommand{\Vmsm}{V_{\mathrm{msm}}}        % measurement isometry
\newcommand{\Vann}{V_{\mathrm{ann}}}        % announcement isometry
\newcommand{\VI}{V_{I}}                     % Eve's copy of the announcements
\newcommand{\VChat}{V_{\Chat}}              % acceptance-data isometry
\newcommand{\VR}{V_{R}}                     % sifting + key map -> raw key R
\newcommand{\VZ}{V_{Z}}                     % raw key -> key register Z
\newcommand{\Wprot}{W}                      % full protocol isometry

% --- Renyi entropies ---------------------------------------------------------
% SECURITY-CRITICAL: see spec/notation.md §4. Never change an arrow to
% make an inequality work out; never convert silently between the sandwiched
% (tilde) and Petz (overline) families.
\newcommand{\Hsandup}{\widetilde{H}^{\uparrow}_{\alpha}}   % sandwiched, up
\newcommand{\Hsanddown}{\widetilde{H}^{\downarrow}_{\alpha}}% sandwiched, down
\newcommand{\Hpetzup}{\overline{H}^{\uparrow}_{\alpha}}    % Petz, up
\newcommand{\Hpetzdown}{\overline{H}^{\downarrow}_{\alpha}}% Petz, down
\newcommand{\hup}{h^{\uparrow}_{\alpha}}                   % single-round rate
\newcommand{\hqkd}{h^{\uparrow,\mathrm{QKD}}_{\alpha}}     % single-round QKD rate
\newcommand{\hlow}{\underline{h}_{\alpha}}                 % lower bound on \hqkd
% --- spec/notation.md 3.5 and 4.5 -------------------------------------------
% \hlb  is the first-order linearization lower bound on \hlow: the headline
%       key rate (the `first-order` tier of foundations/def:tiers). NOT \hlow
%       itself.
% \tnum is the instance's one declared numerical tolerance. Its
%       near-collision with the declared marginal \tau_A is disclosed in
%       spec/notation.md 3.5; bare \tau is reserved and denotes neither.
\newcommand{\hlb}{\underline{h}^{\mathrm{FO}}_{\alpha}}    % first-order lower bound on \hlow
\newcommand{\tnum}{\tau_{\mathrm{num}}}                    % declared numerical tolerance


% =============================================================================
%  PART B -- module-local macros, grouped by module. Deprecated aliases,
%  defined so that a text written with them compiles, used in no new contract: \Lop,
%  \testE, \Pdeph, \pass, \ecLeak (spec/notation.md section 3.9).
% =============================================================================

% --- meat-core -----------------------------------------------------------------
\newcommand{\SecSet}{\mathbb{S}}            % sector index set of Q
\newcommand{\Qsys}{Q}                       % declared sector-carrying measured system
\newcommand{\Vset}{\mathcal{V}}             % declared joint outcome alphabet of the Q-POVM
\newcommand{\CAset}{\mathcal{C}_{A}}        % Alice announcement alphabet
\newcommand{\CBset}{\mathcal{C}_{B}}        % Bob announcement alphabet
\newcommand{\Party}{\mathfrak{P}}           % declared party partition
\newcommand{\Oset}[1]{\mathcal{O}_{#1}}     % readout alphabet of party #1
\newcommand{\rd}[1]{\pi_{#1}}               % readout map of party #1
\newcommand{\Sadm}{\mathcal{S}_{A}}         % admissible Alice subspace
\newcommand{\nomarg}{\mathsf{none}}         % "no trusted marginal declared"
\newcommand{\Zset}{\mathcal{Z}_{\perp}}     % raw-key alphabet incl. \perp
\newcommand{\Zval}{\mathcal{Z}}             % key alphabet without \perp
\newcommand{\Etil}{\widetilde{E}}           % single-round purifying register
\newcommand{\Ehat}{\widehat{E}}             % Eve's n-round register
\newcommand{\OmAT}{\Omega_{\mathrm{AT}}}    % acceptance event
\newcommand{\epa}{\varepsilon_{\mathrm{PA}}}
\newcommand{\eev}{\varepsilon_{\mathrm{EV}}}
\newcommand{\lEC}{\lambda_{\mathrm{EC}}}
\newcommand{\Decl}{\mathfrak{D}}            % protocol declaration
\newcommand{\CList}{\mathfrak{L}}           % declared constraint list
\newcommand{\PinS}{\mathcal{P}_{\SecSet}}   % sector pinching on Q
\newcommand{\Feas}{\mathcal{W}}             % feasible set of block states
\newcommand{\Chkey}{\Chatset_{\mathrm{key}}}% key-live acceptance values
\newcommand{\Kann}{\mathcal{K}}             % announcement-class alphabet
\newcommand{\hold}{\mathsf{h}}              % declared key holder (A or B)
\newcommand{\sfD}{D}                        % classical KL divergence (base 2)
\newcommand{\Dtwo}{D}                       % quantum relative entropy (base 2)
\newcommand{\Dsand}{\widetilde{D}_{\alpha}} % sandwiched Renyi divergence
\newcommand{\idmap}{\mathrm{id}}            % identity channel
\newcommand{\Aev}{\mathcal{A}}              % Eve's per-round attack channel
\newcommand{\AttSet}{\mathfrak{A}}          % licensed attack-channel set
\newcommand{\Chacc}{\Chatset_{\mathrm{acc}}}% acceptance support of \acSet
\newcommand{\epsstar}{\varepsilon^{\star}}      % distance to an exact POVM
\newcommand{\dfH}{\Delta_{\mathrm{herm}}}       % Hermiticity defect
\newcommand{\dfP}{\Delta_{\mathrm{pos}}}        % negative part
\newcommand{\dfC}{\Delta_{\mathrm{cpl}}}        % completeness defect
\newcommand{\lmoq}{g}                       % acceptance-side gradient block
\newcommand{\lmoJ}{G}                       % feasible-side gradient block
\newcommand{\bacc}{\beta_{\mathrm{acc}}}    % certified acceptance-side minimum
\newcommand{\Itset}{\mathfrak{I}}           % finite index set of iterates (C22)

% --- reductions ----------------------------------------------------------------
\newcommand{\Bsys}{B}                        % physical (pre-squashing) system
\newcommand{\Tcl}[1]{\mathcal{T}(#1)}        % trace-class operators on H(#1)
\newcommand{\Ins}{\mathcal{I}}               % generalized-instrument branch map
\newcommand{\redNmsm}{\mathcal{N}}              % POVM measurement channel
\newcommand{\Eprot}{\mathcal{E}}             % n-round protocol channel
\newcommand{\Grnd}{\mathsf{G}}               % one round map of a protocol
\newcommand{\Kfin}{\mathsf{K}}               % end-of-protocol processing
\newcommand{\Crit}{\mathsf{sec}}             % value of a security criterion
\newcommand{\Pflag}{\Pi_{\Flag}}             % flag-sector projector on Q
\newcommand{\Pleq}{\Pi_{\leqN}}              % preserved-block projector
\newcommand{\Pgt}{\Pi_{\gtN}}                % above-cutoff projector
\newcommand{\Sph}{S^{\mathrm{ph}}}           % family of physical inputs
\newcommand{\Srch}{S^{\mathrm{rch}}}         % reachable squashed states
\newcommand{\Ssq}{S^{\sq}}                   % squashed admissible set
\newcommand{\Ffss}{F^{\mathrm{FSS}}}         % flag-state-squasher target
\newcommand{\Fwp}{F^{\mathrm{WP}}}           % multi-bound WPFSS target
\newcommand{\Mwp}{M^{\mathrm{WP}}}           % renormalized above-cutoff POVM
\newcommand{\Mdec}{{M'}^{\mathrm{ph}}}       % declared physical family on B
\newcommand{\Ppvm}{P^{\Flag}}                % flag PVM
\newcommand{\Lmp}{{\Lambda_{\mathrm{mp}}}}   % measure-and-prepare channel
\newcommand{\Lwp}{{\Lambda^{\mathrm{WP}}}}   % multi-bound WPFSS channel
\newcommand{\Lfss}{{\Lambda^{\mathrm{FSS}}}} % flag-state-squasher channel
\newcommand{\lmax}{\lambda_{\max}}           % maximum eigenvalue
\newcommand{\Lnoise}{\widetilde{\Lambda}}    % flag-noise channel
\newcommand{\cg}{\varpi}                     % declared coarse-graining map
\newcommand{\MCK}{\mathsf{MC}}               % multiclick raw-pattern set
\newcommand{\redChain}{\mathfrak{C}}            % reduction chain
\newcommand{\Lnk}{\ell}                      % a single typed link
\newcommand{\dsq}{\preceq}                      % the squashing relation
\newcommand{\opn}[1]{\left\lVert #1\right\rVert}           % operator (spectral) norm
\newcommand{\mxn}[1]{\left\lVert #1\right\rVert_{\max}}    % entrywise max modulus
\newcommand{\Fld}{\mathbb{F}}                   % exact ordered substrate field
\newcommand{\Fldi}{\mathbb{F}_{i}}              % its Gaussian extension F + iF
\newcommand{\lat}[1]{\big\lfloor #1\big\rfloor_{\kappa}} % round DOWN onto the lattice

% --- protocol-compiler ---------------------------------------------------------
\newcommand{\Spec}{\mathfrak{S}}             % protocol specification (compiler input)
\newcommand{\pcChain}{\mathfrak{R}}            % reduction chain (two typed segments)
\newcommand{\Chsrc}{\pcChain_{\mathrm{src}}}   % its Alice-side source-map segment
\newcommand{\Chrec}{\pcChain_{\mathrm{rec}}}   % its receiver-side squashing segment
\newcommand{\rvirt}{\rho^{\mathrm{v}}}       % virtual (PRE-source-map) signal state
\newcommand{\Fph}{F^{\mathrm{ph}}}           % physical (pre-reduction) POVM element
\newcommand{\Vph}{\mathcal{V}^{\mathrm{ph}}} % physical outcome alphabet
\newcommand{\Tpub}{\mathcal{T}_{\mathrm{pub}}}% public-choice alphabet
\newcommand{\Uset}{\mathcal{U}}              % residual outcome alphabet
\newcommand{\idxm}{\mathsf{idx}}             % label-to-index map
\newcommand{\Asig}{A'}                       % transmitted (signal) system -- the REAL one
\newcommand{\Avirt}{A''}                     % virtual transmitted system (source descriptor of Xi_1)
\newcommand{\Abar}{\bar{A}}                  % Alice's measured register
\newcommand{\Ahat}{\hat{A}}                  % Alice's shield register
\newcommand{\pcNmsm}{\mathcal{N}_{\mathsf A}}  % Alice's measurement channel (see the RULES block)
\newcommand{\Phip}{\Phi^{+}}                 % maximally entangled vector label
\newcommand{\Smap}{\mathfrak{s}}             % the emitted sector map (fraktur: a compile-time object)

% --- receiver ------------------------------------------------------------------
\newcommand{\Modes}{\mathbb{M}}              % declared ordered mode set
\newcommand{\Inp}{\mathbb{I}}                % declared input modes
\newcommand{\InpC}{\mathbb{I}_{\Chan}}       % channel input modes
\newcommand{\InpS}{\mathbb{I}_{\Prb}}        % probe (source) input modes
\newcommand{\Env}{\mathbb{E}}                % dilation (environment) modes
\newcommand{\Port}[1]{\mathbb{P}_{#1}}       % mode set of detector port #1
\newcommand{\Emd}[1]{E_{#1}}                 % coordinate projector on modes #1
\newcommand{\Gmat}{G}                        % instrument (transfer) matrix
\newcommand{\Fock}{\mathfrak{F}}             % Fock / second-quantization functor
\newcommand{\Symn}[1]{\mathrm{Sym}^{#1}}     % symmetric tensor power
\newcommand{\normord}[1]{{:}\mspace{2mu}#1\mspace{2mu}{:}}   % normal ordering
\newcommand{\rk}{\operatorname{rank}}        % rank (physics.sty owns \rank)
\newcommand{\SCK}{\mathsf{SC}}               % single-click pattern set
\newcommand{\Recv}{\mathfrak{R}}             % receiver instrument tuple
\newcommand{\Cat}{\mathfrak{C}}              % component catalogue
\newcommand{\Grph}{\mathfrak{E}}             % declared element graph
\newcommand{\defic}{\mathbb{L}}              % loss deficiency operator
\newcommand{\Pvac}[1]{\Pi_{\vac}(#1)}        % joint no-click element on ports #1
\newcommand{\pmc}{p_{\mc}}                   % uniform high-photon mc weight
\newcommand{\Bnd}{\beta}                     % the scalar bound function beta(n)
\newcommand{\Rot}{R}                         % planar rotation matrix
\newcommand{\Wplate}{W}                      % retarder (wave-plate) matrix
\newcommand{\Bsp}{B}                         % beamsplitter core matrix

% --- statistics ----------------------------------------------------------------
\newcommand{\honest}{p^{\mathrm{hon}}}          % honest per-round distribution
\newcommand{\ecomp}{\varepsilon_{\mathrm{comp}}}% completeness budget
\newcommand{\simplex}[1]{\Delta(#1)}            % probability simplex
\newcommand{\Chatcon}{\Chatset_{\mathrm{con}}}  % constrained outcomes
\newcommand{\Ireg}[3]{\mathcal{I}_{#1}(#2,#3)}  % regularized incomplete beta
\newcommand{\Ncount}[1]{\mathsf{N}_{#1}}        % per-outcome count
\newcommand{\qlo}{\mathbf{q}^{-}}               % box lower bounds
\newcommand{\qhi}{\mathbf{q}^{+}}               % box upper bounds
 for it, so the
%  file is read from spec/modules/<id>/.
%
%  RULES
%    * A macro that spec/notation.md catalogues is a verbatim copy of it
%      (part A). Changing one is a notation change: bump spec/notation.md.
%    * Part B holds the module-local macros. A macro defined by two modules
%      with the same expansion appears once, under the first module (in the
%      order below) that defines it. Two macros with the same name and
%      DIFFERENT expansions are both kept under a module prefix, and the
%      claims write the prefixed names:
%        \Chain  -> \redChain (reductions, \mathfrak{C}: the reduction chain)
%                   \pcChain  (protocol-compiler, \mathfrak{R}: the two-segment
%                             chain; \Chsrc and \Chrec are built on it)
%        \Nmsm   -> \redNmsm  (reductions, \mathcal{N}: a POVM measurement
%                             channel)
%                   \pcNmsm   (\mathcal{N}_{\mathsf A}: Alice's measurement
%                             channel, written by
%                             transmitter/thm:certificate-sufficiency,
%                             transmitter/rem:xreg and
%                             reductions/lem:source-transfer)
%    * Deliberately omitted: a \provenance{tag}{details} macro. Provenance
%      is a header field of every claim file (spec/claim-format.md section
%      3); a claim body that calls such a macro fails to compile, which is
%      the intended signal.
%
%  REQUIRED PACKAGES, loaded by the assembler's preamble before this file:
%    amsmath, amssymb, amsfonts, mathtools, dsfont, braket, physics, xcolor
%    (braket before physics, so that physics wins for \ket, \bra, \tr, ...).
%
%  This file contains only definitions: no \usepackage, no document body.
% =============================================================================


% =============================================================================
%  PART A -- notation macros, the ones spec/notation.md catalogues, each a
%  verbatim copy of its catalogue row; grouped under numbered headings.
% =============================================================================

% -----------------------------------------------------------------------------
% 1. State spaces and operator sets (a "system label" -> spaces of operators)
% -----------------------------------------------------------------------------
\newcommand{\Hs}[1]{\mathcal{H}(#1)}        % Hilbert space of a system label
\newcommand{\Lin}[1]{\mathcal{L}(#1)}       % bounded linear operators on H(#1)
\newcommand{\Linn}[2]{\mathcal{L}(#1,#2)}   % linear maps H(#1) -> H(#2)
\newcommand{\Pos}[1]{\mathrm{Pos}(#1)}      % positive semidefinite operators on H(#1)
\newcommand{\Density}{\mathrm{D}}           % density operators: \Density(A) = D(A)
\newcommand{\Seq}[1]{S_{=}(#1)}             % normalized states (trace = 1)
\newcommand{\Sleq}[1]{S_{\leq}(#1)}         % subnormalized states (0 <= trace <= 1)
\newcommand{\id}{\mathds{1}}                % identity operator
\newcommand{\idon}[1]{\mathds{1}_{#1}}      % identity on a named system, e.g. \idon{C}

% DEPRECATED: applies \mathcal to both arguments itself, so it takes bare
% letters, not system labels. Prefer \Linn; do not use in new contracts.
\providecommand{\Lop}[2]{\mathcal{L}(\mathcal{#1}, \mathcal{#2})}


% -----------------------------------------------------------------------------
% 2. Composite systems
% -----------------------------------------------------------------------------
\newcommand{\dsum}{\oplus}                  % direct sum of system labels
% Tensor product: use \otimes directly.
% n-fold direct sum of an operator space: \Lin{A}^{\oplus n}


% -----------------------------------------------------------------------------
% 3. Events, conditioning and classical registers (for STATES)
%    Convention:  \wedge  -> subnormalized "and-event" component, e.g. \rho_{A\wedge x}
%                 |        -> normalized conditional state,        e.g. \rho_{A|x}
% (Write these directly in subscripts; no extra macros needed.)
% -----------------------------------------------------------------------------


% -----------------------------------------------------------------------------
% 4. Photon-number subspaces / direct-sum blocks
%    Convention:  operators carry the restriction "on top", in brackets;
%                 states carry it as a \wedge subscript (see Sec. 3).
% -----------------------------------------------------------------------------
\newcommand{\Pnum}[1]{P_{#1}}               % photon-number projector (system from context)
\newcommand{\Pnumon}[2]{P_{#1}^{#2}}        % with explicit system superscript, e.g. P_n^C
\newcommand{\blk}[2]{#1^{(#2)}}             % operator block, e.g. \blk{M}{\leqN} = M^{(<=N)}
\newcommand{\leqN}{\leq N}                  % use as \blk{M}{\leqN}
\newcommand{\gtN}{> N}                      % use as \blk{M}{\gtN}
\newcommand{\eqn}{= n}                      % use as \blk{M}{\eqn}


% -----------------------------------------------------------------------------
% 5. System labels used in the protocol
% -----------------------------------------------------------------------------
\newcommand{\Chan}{C}                       % quantum-channel (Eve) input to Bob
\newcommand{\Prb}{P}                        % probe input
\newcommand{\Flag}{{\mathrm{Flg}}}          % flag system (EXISTS ONLY AFTER squashing)
\newcommand{\Loss}{{\mathrm{Lss}}}          % loss system
\newcommand{\Vc}{{\mathrm{Vcm}}}            % vacuum system
\newcommand{\Nout}{N_{\mathrm{out}}}
\newcommand{\Nin}{N_{\mathrm{in}}}


% -----------------------------------------------------------------------------
% 6. Alphabets
% -----------------------------------------------------------------------------
\newcommand{\Yset}{\mathcal{Y}}             % Bob's (coarse-grained) detection alphabet
\newcommand{\Yraw}{\mathcal{Y}_{\mathrm{raw}}}% raw click patterns {0,1}^4
\newcommand{\Yscset}{\mathcal{Y}_{\singlec}}% single-click outcomes
\newcommand{\Ymcset}{\mathcal{Y}_{\mc}}     % multi-click outcomes
\newcommand{\Tset}{\mathcal{T}_{\mathrm{set}}}% private setting set (the probe settings in the optical instance)
\newcommand{\Xset}{\mathcal{X}}             % Alice's alphabet
\newcommand{\Dset}{\{H,V,+,-\}}             % polarization / detector labels
\newcommand{\pass}{\top}                    % pass sifting


% -----------------------------------------------------------------------------
% 7. Mode / outcome abbreviations
% -----------------------------------------------------------------------------
\newcommand{\on}{{\mathrm{on}}}
\newcommand{\off}{{\mathrm{off}}}
\newcommand{\dc}{{\mathrm{dc}}}             % dark count
\newcommand{\sq}{\mathrm{sq}}               % squashed
\newcommand{\mc}{\mathrm{mc}}               % multi-click
\newcommand{\nc}{\mathrm{nc}}               % no-click
\newcommand{\singlec}{\mathrm{sc}}          % single-click
\newcommand{\vac}{\mathrm{vac}}             % vacuum
\newcommand{\kout}{{k_\mathrm{out}}}
\newcommand{\kin}{{k_\mathrm{in}}}


% -----------------------------------------------------------------------------
% 8. Maps, channels and miscellaneous operators
% -----------------------------------------------------------------------------
\newcommand{\Prob}{\mathrm{Prob}}           % probability (use \Prob[\cdot] or \Prob(\cdot))
\newcommand{\lmin}{\lambda_{\min}}          % minimum eigenvalue
\newcommand{\Wflag}{W_{\Flag}}              % flag-space weight
\newcommand{\transpose}{\mathsf{T}}
\newcommand{\innerp}[2]{\left\langle #1 \vert #2 \right\rangle}
\renewcommand{\vec}[1]{\boldsymbol{#1}}     % NB: overrides the LaTeX arrow accent


% -----------------------------------------------------------------------------
% 9. Key-rate numerics
% -----------------------------------------------------------------------------
\newcommand{\cG}{\mathcal{G}}               % protocol map
\newcommand{\cZ}{\mathcal{Z}}               % pinching map
\newcommand{\Amap}{\mathcal{A}_{\mathrm{ann}}} % announcement + measurement map
\newcommand{\Amapnb}{\mathcal{A}_{\mathrm{ann}|\mathrm{nb}}}
\newcommand{\simpl}{}                       % empty suffix for simplified block maps
\newcommand{\acSet}{\mathcal{S}_{\mathrm{acc}}} % acceptance set
\newcommand{\testE}{\mathcal{E}_{\mathrm{test}}} % test-statistics map
\newcommand{\ecLeak}{\mathrm{leak}^{\mathrm{EC}}_{\mathrm{obs}}}
\newcommand{\pgen}{p_{\mathrm{gen}}}
\newcommand{\askeyr}{R_{\infty}}             % asymptotic key rate
\newcommand{\ppass}{p_{\mathrm{pass}}}       % sifting pass probability


% -----------------------------------------------------------------------------
% 10. Finite-size analysis / Renyi entropies
% -----------------------------------------------------------------------------

% --- Round types: test vs. key generation (register T, alphabet \Fset) --------
\newcommand{\tst}{\mathrm{test}}            % test round
\newcommand{\gen}{\mathrm{gen}}             % key-generation round
\newcommand{\Fset}{\mathcal{T}}             % round-type alphabet {test, gen} (not \Tset)
\newcommand{\Chat}{\widehat{C}}             % acceptance-test register
\newcommand{\Chatset}{\widehat{\mathcal{C}}}% its alphabet

% --- Basis / key extraction functions ----------------------------------------
\newcommand{\bA}{b_{A}}                     % Alice's basis-extraction function
\newcommand{\bB}{b_{B}}                     % Bob's basis-extraction function
\newcommand{\kB}{k_{B}}                     % Bob's key-extraction function

% --- Maps --------------------------------------------------------------------
\newcommand{\Mprot}{\mathcal{M}}            % single-round protocol map
\newcommand{\Phimap}{\boldsymbol{\Phi}}     % test-statistics map
\newcommand{\Pdeph}{\mathcal{P}}            % dephasing (pinching) map on \Blreg

% --- Isometries --------------------------------------------------------------
\newcommand{\Vmsm}{V_{\mathrm{msm}}}        % measurement isometry
\newcommand{\Vann}{V_{\mathrm{ann}}}        % announcement isometry
\newcommand{\VI}{V_{I}}                     % Eve's copy of the announcements
\newcommand{\VChat}{V_{\Chat}}              % acceptance-data isometry
\newcommand{\VR}{V_{R}}                     % sifting + key map -> raw key R
\newcommand{\VZ}{V_{Z}}                     % raw key -> key register Z
\newcommand{\Wprot}{W}                      % full protocol isometry

% --- Renyi entropies ---------------------------------------------------------
% SECURITY-CRITICAL: see spec/notation.md §4. Never change an arrow to
% make an inequality work out; never convert silently between the sandwiched
% (tilde) and Petz (overline) families.
\newcommand{\Hsandup}{\widetilde{H}^{\uparrow}_{\alpha}}   % sandwiched, up
\newcommand{\Hsanddown}{\widetilde{H}^{\downarrow}_{\alpha}}% sandwiched, down
\newcommand{\Hpetzup}{\overline{H}^{\uparrow}_{\alpha}}    % Petz, up
\newcommand{\Hpetzdown}{\overline{H}^{\downarrow}_{\alpha}}% Petz, down
\newcommand{\hup}{h^{\uparrow}_{\alpha}}                   % single-round rate
\newcommand{\hqkd}{h^{\uparrow,\mathrm{QKD}}_{\alpha}}     % single-round QKD rate
\newcommand{\hlow}{\underline{h}_{\alpha}}                 % lower bound on \hqkd
% --- spec/notation.md 3.5 and 4.5 -------------------------------------------
% \hlb  is the first-order linearization lower bound on \hlow: the headline
%       key rate (the `first-order` tier of foundations/def:tiers). NOT \hlow
%       itself.
% \tnum is the instance's one declared numerical tolerance. Its
%       near-collision with the declared marginal \tau_A is disclosed in
%       spec/notation.md 3.5; bare \tau is reserved and denotes neither.
\newcommand{\hlb}{\underline{h}^{\mathrm{FO}}_{\alpha}}    % first-order lower bound on \hlow
\newcommand{\tnum}{\tau_{\mathrm{num}}}                    % declared numerical tolerance


% =============================================================================
%  PART B -- module-local macros, grouped by module. Deprecated aliases,
%  defined so that a text written with them compiles, used in no new contract: \Lop,
%  \testE, \Pdeph, \pass, \ecLeak (spec/notation.md section 3.9).
% =============================================================================

% --- meat-core -----------------------------------------------------------------
\newcommand{\SecSet}{\mathbb{S}}            % sector index set of Q
\newcommand{\Qsys}{Q}                       % declared sector-carrying measured system
\newcommand{\Vset}{\mathcal{V}}             % declared joint outcome alphabet of the Q-POVM
\newcommand{\CAset}{\mathcal{C}_{A}}        % Alice announcement alphabet
\newcommand{\CBset}{\mathcal{C}_{B}}        % Bob announcement alphabet
\newcommand{\Party}{\mathfrak{P}}           % declared party partition
\newcommand{\Oset}[1]{\mathcal{O}_{#1}}     % readout alphabet of party #1
\newcommand{\rd}[1]{\pi_{#1}}               % readout map of party #1
\newcommand{\Sadm}{\mathcal{S}_{A}}         % admissible Alice subspace
\newcommand{\nomarg}{\mathsf{none}}         % "no trusted marginal declared"
\newcommand{\Zset}{\mathcal{Z}_{\perp}}     % raw-key alphabet incl. \perp
\newcommand{\Zval}{\mathcal{Z}}             % key alphabet without \perp
\newcommand{\Etil}{\widetilde{E}}           % single-round purifying register
\newcommand{\Ehat}{\widehat{E}}             % Eve's n-round register
\newcommand{\OmAT}{\Omega_{\mathrm{AT}}}    % acceptance event
\newcommand{\epa}{\varepsilon_{\mathrm{PA}}}
\newcommand{\eev}{\varepsilon_{\mathrm{EV}}}
\newcommand{\lEC}{\lambda_{\mathrm{EC}}}
\newcommand{\Decl}{\mathfrak{D}}            % protocol declaration
\newcommand{\CList}{\mathfrak{L}}           % declared constraint list
\newcommand{\PinS}{\mathcal{P}_{\SecSet}}   % sector pinching on Q
\newcommand{\Feas}{\mathcal{W}}             % feasible set of block states
\newcommand{\Chkey}{\Chatset_{\mathrm{key}}}% key-live acceptance values
\newcommand{\Kann}{\mathcal{K}}             % announcement-class alphabet
\newcommand{\hold}{\mathsf{h}}              % declared key holder (A or B)
\newcommand{\sfD}{D}                        % classical KL divergence (base 2)
\newcommand{\Dtwo}{D}                       % quantum relative entropy (base 2)
\newcommand{\Dsand}{\widetilde{D}_{\alpha}} % sandwiched Renyi divergence
\newcommand{\idmap}{\mathrm{id}}            % identity channel
\newcommand{\Aev}{\mathcal{A}}              % Eve's per-round attack channel
\newcommand{\AttSet}{\mathfrak{A}}          % licensed attack-channel set
\newcommand{\Chacc}{\Chatset_{\mathrm{acc}}}% acceptance support of \acSet
\newcommand{\epsstar}{\varepsilon^{\star}}      % distance to an exact POVM
\newcommand{\dfH}{\Delta_{\mathrm{herm}}}       % Hermiticity defect
\newcommand{\dfP}{\Delta_{\mathrm{pos}}}        % negative part
\newcommand{\dfC}{\Delta_{\mathrm{cpl}}}        % completeness defect
\newcommand{\lmoq}{g}                       % acceptance-side gradient block
\newcommand{\lmoJ}{G}                       % feasible-side gradient block
\newcommand{\bacc}{\beta_{\mathrm{acc}}}    % certified acceptance-side minimum
\newcommand{\Itset}{\mathfrak{I}}           % finite index set of iterates (C22)

% --- reductions ----------------------------------------------------------------
\newcommand{\Bsys}{B}                        % physical (pre-squashing) system
\newcommand{\Tcl}[1]{\mathcal{T}(#1)}        % trace-class operators on H(#1)
\newcommand{\Ins}{\mathcal{I}}               % generalized-instrument branch map
\newcommand{\redNmsm}{\mathcal{N}}              % POVM measurement channel
\newcommand{\Eprot}{\mathcal{E}}             % n-round protocol channel
\newcommand{\Grnd}{\mathsf{G}}               % one round map of a protocol
\newcommand{\Kfin}{\mathsf{K}}               % end-of-protocol processing
\newcommand{\Crit}{\mathsf{sec}}             % value of a security criterion
\newcommand{\Pflag}{\Pi_{\Flag}}             % flag-sector projector on Q
\newcommand{\Pleq}{\Pi_{\leqN}}              % preserved-block projector
\newcommand{\Pgt}{\Pi_{\gtN}}                % above-cutoff projector
\newcommand{\Sph}{S^{\mathrm{ph}}}           % family of physical inputs
\newcommand{\Srch}{S^{\mathrm{rch}}}         % reachable squashed states
\newcommand{\Ssq}{S^{\sq}}                   % squashed admissible set
\newcommand{\Ffss}{F^{\mathrm{FSS}}}         % flag-state-squasher target
\newcommand{\Fwp}{F^{\mathrm{WP}}}           % multi-bound WPFSS target
\newcommand{\Mwp}{M^{\mathrm{WP}}}           % renormalized above-cutoff POVM
\newcommand{\Mdec}{{M'}^{\mathrm{ph}}}       % declared physical family on B
\newcommand{\Ppvm}{P^{\Flag}}                % flag PVM
\newcommand{\Lmp}{{\Lambda_{\mathrm{mp}}}}   % measure-and-prepare channel
\newcommand{\Lwp}{{\Lambda^{\mathrm{WP}}}}   % multi-bound WPFSS channel
\newcommand{\Lfss}{{\Lambda^{\mathrm{FSS}}}} % flag-state-squasher channel
\newcommand{\lmax}{\lambda_{\max}}           % maximum eigenvalue
\newcommand{\Lnoise}{\widetilde{\Lambda}}    % flag-noise channel
\newcommand{\cg}{\varpi}                     % declared coarse-graining map
\newcommand{\MCK}{\mathsf{MC}}               % multiclick raw-pattern set
\newcommand{\redChain}{\mathfrak{C}}            % reduction chain
\newcommand{\Lnk}{\ell}                      % a single typed link
\newcommand{\dsq}{\preceq}                      % the squashing relation
\newcommand{\opn}[1]{\left\lVert #1\right\rVert}           % operator (spectral) norm
\newcommand{\mxn}[1]{\left\lVert #1\right\rVert_{\max}}    % entrywise max modulus
\newcommand{\Fld}{\mathbb{F}}                   % exact ordered substrate field
\newcommand{\Fldi}{\mathbb{F}_{i}}              % its Gaussian extension F + iF
\newcommand{\lat}[1]{\big\lfloor #1\big\rfloor_{\kappa}} % round DOWN onto the lattice

% --- protocol-compiler ---------------------------------------------------------
\newcommand{\Spec}{\mathfrak{S}}             % protocol specification (compiler input)
\newcommand{\pcChain}{\mathfrak{R}}            % reduction chain (two typed segments)
\newcommand{\Chsrc}{\pcChain_{\mathrm{src}}}   % its Alice-side source-map segment
\newcommand{\Chrec}{\pcChain_{\mathrm{rec}}}   % its receiver-side squashing segment
\newcommand{\rvirt}{\rho^{\mathrm{v}}}       % virtual (PRE-source-map) signal state
\newcommand{\Fph}{F^{\mathrm{ph}}}           % physical (pre-reduction) POVM element
\newcommand{\Vph}{\mathcal{V}^{\mathrm{ph}}} % physical outcome alphabet
\newcommand{\Tpub}{\mathcal{T}_{\mathrm{pub}}}% public-choice alphabet
\newcommand{\Uset}{\mathcal{U}}              % residual outcome alphabet
\newcommand{\idxm}{\mathsf{idx}}             % label-to-index map
\newcommand{\Asig}{A'}                       % transmitted (signal) system -- the REAL one
\newcommand{\Avirt}{A''}                     % virtual transmitted system (source descriptor of Xi_1)
\newcommand{\Abar}{\bar{A}}                  % Alice's measured register
\newcommand{\Ahat}{\hat{A}}                  % Alice's shield register
\newcommand{\pcNmsm}{\mathcal{N}_{\mathsf A}}  % Alice's measurement channel (see the RULES block)
\newcommand{\Phip}{\Phi^{+}}                 % maximally entangled vector label
\newcommand{\Smap}{\mathfrak{s}}             % the emitted sector map (fraktur: a compile-time object)

% --- receiver ------------------------------------------------------------------
\newcommand{\Modes}{\mathbb{M}}              % declared ordered mode set
\newcommand{\Inp}{\mathbb{I}}                % declared input modes
\newcommand{\InpC}{\mathbb{I}_{\Chan}}       % channel input modes
\newcommand{\InpS}{\mathbb{I}_{\Prb}}        % probe (source) input modes
\newcommand{\Env}{\mathbb{E}}                % dilation (environment) modes
\newcommand{\Port}[1]{\mathbb{P}_{#1}}       % mode set of detector port #1
\newcommand{\Emd}[1]{E_{#1}}                 % coordinate projector on modes #1
\newcommand{\Gmat}{G}                        % instrument (transfer) matrix
\newcommand{\Fock}{\mathfrak{F}}             % Fock / second-quantization functor
\newcommand{\Symn}[1]{\mathrm{Sym}^{#1}}     % symmetric tensor power
\newcommand{\normord}[1]{{:}\mspace{2mu}#1\mspace{2mu}{:}}   % normal ordering
\newcommand{\rk}{\operatorname{rank}}        % rank (physics.sty owns \rank)
\newcommand{\SCK}{\mathsf{SC}}               % single-click pattern set
\newcommand{\Recv}{\mathfrak{R}}             % receiver instrument tuple
\newcommand{\Cat}{\mathfrak{C}}              % component catalogue
\newcommand{\Grph}{\mathfrak{E}}             % declared element graph
\newcommand{\defic}{\mathbb{L}}              % loss deficiency operator
\newcommand{\Pvac}[1]{\Pi_{\vac}(#1)}        % joint no-click element on ports #1
\newcommand{\pmc}{p_{\mc}}                   % uniform high-photon mc weight
\newcommand{\Bnd}{\beta}                     % the scalar bound function beta(n)
\newcommand{\Rot}{R}                         % planar rotation matrix
\newcommand{\Wplate}{W}                      % retarder (wave-plate) matrix
\newcommand{\Bsp}{B}                         % beamsplitter core matrix

% --- statistics ----------------------------------------------------------------
\newcommand{\honest}{p^{\mathrm{hon}}}          % honest per-round distribution
\newcommand{\ecomp}{\varepsilon_{\mathrm{comp}}}% completeness budget
\newcommand{\simplex}[1]{\Delta(#1)}            % probability simplex
\newcommand{\Chatcon}{\Chatset_{\mathrm{con}}}  % constrained outcomes
\newcommand{\Ireg}[3]{\mathcal{I}_{#1}(#2,#3)}  % regularized incomplete beta
\newcommand{\Ncount}[1]{\mathsf{N}_{#1}}        % per-outcome count
\newcommand{\qlo}{\mathbf{q}^{-}}               % box lower bounds
\newcommand{\qhi}{\mathbf{q}^{+}}               % box upper bounds

\makeatletter
\newcommand*{\ClaimHeadFont}{\bfseries}   % switched by prominence
\newcommand*{\ClaimNoteFont}{\mdseries}
\newcommand*{\ClaimShortId}{}
\newcommand*{\SetShortId}[1]{\renewcommand*{\ClaimShortId}{#1}}
\newcommand*{\RestatedName}{}
% Upright bodies throughout: italic is reserved for non-normative text.
% \@endparpenalty\@M keeps the meta line on the statement's page.
\newtheoremstyle{claim}{\topsep}{\topsep}{\normalfont\@endparpenalty\@M}{}%
  {\ClaimHeadFont}{.}{.5em}{\thmname{#1}\thmnumber{ #2}\thmnote{ {\ClaimNoteFont(#3)}}}
\newtheoremstyle{claimass}{\topsep}{\topsep}{\normalfont\@endparpenalty\@M}{}%
  {\ClaimHeadFont}{.}{.5em}{\thmname{#1}\thmnumber{ #2}%
  \ifx\ClaimShortId\@empty\else\ [\textsf{\ClaimShortId}]\fi%
  \thmnote{ {\ClaimNoteFont(#3)}}}
\newtheoremstyle{claimremark}{.5\topsep}{.5\topsep}{\small\@endparpenalty\@M}{}%
  {\itshape}{.}{.5em}{\thmname{#1}\thmnumber{ #2}\thmnote{ {\upshape(#3)}}}
\newtheoremstyle{claimreq}{\topsep}{\topsep}{\normalfont\@endparpenalty\@M}{}%
  {\color{reqAccent}\scshape}{.}{.5em}%
  {\thmname{#1}\thmnumber{ #2}\thmnote{ {\upshape(#3)}}}
\makeatother
\theoremstyle{claim}
\newtheorem{theorem}{Theorem}[section]
\newtheorem{lemma}[theorem]{Lemma}
\newtheorem{proposition}[theorem]{Proposition}
\newtheorem{corollary}[theorem]{Corollary}
\newtheorem{fact}[theorem]{Fact}
\newtheorem{definition}[theorem]{Definition}
\newtheorem*{restatedclaim}{\RestatedName}
\theoremstyle{claimass}
\newtheorem{assumption}[theorem]{Assumption}
\newtheorem*{restatedassumption}{\RestatedName}
\theoremstyle{claimremark}
\newtheorem{remark}[theorem]{Remark}
\newtheorem{example}[theorem]{Example}
\theoremstyle{claimreq}
\newtheorem{requirement}[theorem]{Requirement}
\newtheorem*{restatedrequirement}{\RestatedName}
%% --- 8. Prominence: main results a margin rule and a bold heading ----------
\newtcolorbox{mainclaim}{enhanced jigsaw, breakable, blanker,
  borderline west={1.5pt}{-8pt}{mainRule}, left=0pt, right=0pt, top=0pt, bottom=0pt,
  boxsep=0pt, before skip=\topsep, after skip=\topsep, parbox=false,
  before upper={\renewcommand*{\ClaimNoteFont}{\bfseries}}}
\newenvironment{technicalclaim}{\renewcommand*{\ClaimHeadFont}{\mdseries\color{techGrey}}%
  \renewcommand*{\ClaimNoteFont}{}}{}
%% --- 9. Non-normative text, the meta line, navigation ------------------------
\makeatletter
% Tie what follows to this paragraph: no page break, no list penalty before
% the next theorem head (the \@nobreak route a section heading uses).
\newcommand{\keepwithnext}{\par\nobreak\vspace{3pt}\@nobreaktrue
  \everypar{\if@nobreak\@nobreakfalse\clubpenalty\@M\else\clubpenalty\@clubpenalty\everypar{}\fi}}
\newcommand{\claimsummary}[2]{% #1 claim id (named destination), #2 summary
  \par\addvspace{\bigskipamount}\Needspace{6\baselineskip}%
  {\small\itshape\color{nnGrey}\noindent\hypertarget{claim-#1}{}#2\par}\keepwithnext}
\newcommand{\claimanchor}[1]{% a claim without a summary: the destination alone
  \par\addvspace{\bigskipamount}\Needspace{6\baselineskip}%
  \phantomsection\hypertarget{claim-#1}{}\keepwithnext}
\newcommand{\claimleadin}[1]{{\color{nnGrey}\noindent#1\par}\keepwithnext}
\newcommand{\claimmeta}[1]{\par\nobreak\vspace{2pt}%
  {\footnotesize\color{metaGrey}\raggedright\interlinepenalty\@M\noindent#1\par}%
  \addvspace{\medskipamount}}
% After an assumption: the sentence generated from its verification field; after
% any claim with one: its plain-text reading (user_statement). Both
% non-normative, both grey.
\newcommand{\claimverification}[1]{\par\nobreak\vspace{2pt}%
  {\footnotesize\color{metaGrey}\noindent#1\par}\nobreak}
\newcommand{\claimreading}[1]{\par\nobreak\vspace{2pt}%
  {\footnotesize\color{nnGrey}\noindent\emph{Reading:}\enspace#1\par}\nobreak}
\makeatother
\newcommand{\claimbookmark}[4]{% #1 id, #2 kind word, #3 label, #4 plain title
  \bookmark[dest=claim-#1,level=2]{#2 \getrefnumber{#3} #4}}
\newcommand{\metasep}{\unskip\enspace\textperiodcentered\enspace\ignorespaces}
\newcommand{\metalabel}[1]{\emph{#1}\enspace}
\newcommand{\sectionintro}[1]{{\color{nnGrey}\noindent#1\par}\smallskip}
\newcommand{\sectionimports}[1]{{\footnotesize\color{metaGrey}\raggedright\noindent\metalabel{Imports}#1\par}\smallskip}
\newenvironment{overview}{\par\color{nnGrey}\itshape\noindent\ignorespaces}{\par}
% References: the whole "Def. 2.1" is the link; ids are the fallback text of
% every link into another PDF (a relative URI to a named destination).
\newcommand{\mref}[2]{\hyperref[#2]{#1~\ref*{#2}}}
\newcommand{\aref}[2]{\hyperref[#2]{\textsf{#1}}}
\DeclareUrlCommand{\claimid}{\urlstyle{sf}\def\UrlBreaks{\do\/\do\:\do\-\do\.\do\_}%
  \def\UrlBigBreaks{}\def\UrlNoBreaks{}}
\newcommand{\extref}[2]{\hyperref{#1}{}{claim-#2}{\claimid{#2}}}
\newcommand{\extaref}[3]{\hyperref{#1}{}{claim-#2}{\textsf{#3}}}
\newcommand{\coderef}[2]{\href{#1}{\claimid{#2}}}
% Claim-body commands (spec/claim-format.md section 4): \xref{<full id>} prints
% the kind, title and id of a claim of another contract, linked to its stable
% destination in that contract's PDF (one \DeclareXref per cited claim below);
% \term[gloss]{phrase} marks a defined phrase for the index (the gloss feeds
% the Terms table); \tref[phrase]{text} is a use of a defined term. The
% assembler rewrites both in the body into \termdef (the definition's named
% destination term-<slug>) and \termref / \exttermref (a link to it, in this
% or another contract's PDF); the fallbacks below serve text it copies
% elsewhere, such as the digest. \iface is an interface key.
\makeatletter
\newcommand{\DeclareXref}[4]{% id, kind word, title, contract PDF
  \@namedef{xref@#1}{\hyperref{#4}{}{claim-#1}{#2 ``#3'' ({\small\claimid{#1}})}}}
\newcommand{\xref}[1]{\@ifundefined{xref@#1}%
  {\textbf{[unresolved \nolinkurl{#1}]}\@latex@warning{Unresolved xref #1}}%
  {\@nameuse{xref@#1}}}
\makeatother
\DeclareUrlCommand{\iface}{\urlstyle{tt}}
\newcommand{\term}[2][]{\emph{#2}\index{#2}}
\newcommand{\tref}[2][]{#2}
\newcommand{\termdef}[2]{\hypertarget{#1}{}\emph{#2}\index{#2}}
\newcommand{\termref}[2]{\hyperlink{#1}{#2}}
\newcommand{\exttermref}[3]{\hyperref{#1}{}{#2}{#3}}
\hypersetup{pdftitle={Protocol compiler (protocol-compiler) -- contract 0.24.0}}

% Claims of other contracts cited here (\xref, [[...]]).
\DeclareXref{foundations/def:binary64}{Definition}{Binary64 numbers, declared values and dyadic lattices}{../foundations/contract.pdf}
\DeclareXref{foundations/def:declared-data}{Definition}{Declared data and the parties to a module boundary}{../foundations/contract.pdf}
\DeclareXref{foundations/def:declared-order}{Definition}{Sums in a declared order and sums that round nowhere}{../foundations/contract.pdf}
\DeclareXref{foundations/def:exact-admissibility}{Definition}{Exact admissibility of a declared operator family}{../foundations/contract.pdf}
\DeclareXref{foundations/def:exact-island}{Definition}{The exact island}{../foundations/contract.pdf}
\DeclareXref{foundations/def:factor-enclosure}{Definition}{Factor enclosure of a block}{../foundations/contract.pdf}
\DeclareXref{foundations/def:round-homogeneity}{Definition}{Round homogeneity of a per-round family}{../foundations/contract.pdf}
\DeclareXref{foundations/def:sector-map}{Definition}{Sector maps and block-list operators}{../foundations/contract.pdf}
\DeclareXref{foundations/lem:exact-accumulation}{Lemma}{Exactness of a sum that rounds nowhere}{../foundations/contract.pdf}
\DeclareXref{foundations/rem:computed-completeness}{Remark}{A complete computed sum from an incomplete family}{../foundations/contract.pdf}
\DeclareXref{foundations/rem:summation-order}{Remark}{Dependence of a computed sum on its order}{../foundations/contract.pdf}
\DeclareXref{foundations/req:boundary-admission}{Requirement}{Exact admissibility at every module boundary}{../foundations/contract.pdf}
\DeclareXref{foundations/req:consumer-decides-exactly}{Requirement}{Exact decision by every consumer}{../foundations/contract.pdf}
\DeclareXref{foundations/req:declared-lattice}{Requirement}{The declared lattice condition}{../foundations/contract.pdf}
\DeclareXref{foundations/req:declared-order}{Requirement}{Summation order declared with the data}{../foundations/contract.pdf}
\DeclareXref{foundations/req:exact-emission}{Requirement}{Exact verification before emission}{../foundations/contract.pdf}
\DeclareXref{foundations/req:exact-island}{Requirement}{Exact arithmetic on the island, unbudgeted rounding elsewhere}{../foundations/contract.pdf}
\DeclareXref{foundations/req:tolerance-scope}{Requirement}{Scope of the declared tolerance}{../foundations/contract.pdf}
\DeclareXref{foundations/req:who-snaps}{Requirement}{Snapping by the declarer only}{../foundations/contract.pdf}
\DeclareXref{meat-core/ass:lab}{Assumption}{Each party's laboratory is sealed}{../meat-core/contract.pdf}
\DeclareXref{meat-core/cor:composed-security}{Corollary}{Security of the protocol and its net key}{../meat-core/contract.pdf}
\DeclareXref{meat-core/cor:reduction-constrained}{Corollary}{The reduction for arbitrary constraint lists}{../meat-core/contract.pdf}
\DeclareXref{meat-core/def:admissible-declaration}{Definition}{Admissible declaration}{../meat-core/contract.pdf}
\DeclareXref{meat-core/def:declaration}{Definition}{The protocol declaration as the MEAT core reads it}{../meat-core/contract.pdf}
\DeclareXref{meat-core/def:roundmodel}{Definition}{Round model and licensed attack sets}{../meat-core/contract.pdf}
\DeclareXref{meat-core/prop:a-neutral}{Proposition}{Value neutrality along Alice's sectors}{../meat-core/contract.pdf}
\DeclareXref{meat-core/prop:reference-companion}{Proposition}{The compiled reference companion is an admissible declaration}{../meat-core/contract.pdf}
\DeclareXref{meat-core/prop:trivial-alice-system}{Proposition}{The two source modes coincide at a trivial Alice system}{../meat-core/contract.pdf}
\DeclareXref{meat-core/rem:cgate}{Remark}{Degeneracy of the licensed-class half of the gate}{../meat-core/contract.pdf}
\DeclareXref{meat-core/req:companion-admission}{Requirement}{Admission of the reference companion beside its declaration}{../meat-core/contract.pdf}
\DeclareXref{meat-core/thm:reduction}{Theorem}{Reduction of the rate to the convex program}{../meat-core/contract.pdf}
\DeclareXref{receiver/def:clickalphabet}{Definition}{Click-pattern alphabet}{../receiver/contract.pdf}
\DeclareXref{receiver/def:coarse}{Definition}{Declared coarse-graining and the emitted POVM}{../receiver/contract.pdf}
\DeclareXref{receiver/thm:coarse-graining}{Theorem}{Coarse-graining of the click-pattern POVM}{../receiver/contract.pdf}
\DeclareXref{reductions/ass:meas}{Assumption}{Each measuring party's device as one POVM per round}{../reductions/contract.pdf}
\DeclareXref{reductions/ass:tsm}{Assumption}{Output-state test statistics}{../reductions/contract.pdf}
\DeclareXref{reductions/cor:entries}{Corollary}{The emitted constraint entries}{../reductions/contract.pdf}
\DeclareXref{reductions/def:attset}{Definition}{Admissible attack sets in the consumer's format}{../reductions/contract.pdf}
\DeclareXref{reductions/def:chain}{Definition}{Reduction chain}{../reductions/contract.pdf}
\DeclareXref{reductions/def:joint-family}{Definition}{Joint alphabet, classes and the joint family of a private setting}{../reductions/contract.pdf}
\DeclareXref{reductions/def:licensed}{Definition}{Licensed entries and the party-local form}{../reductions/contract.pdf}
\DeclareXref{reductions/def:link}{Definition}{Typed link}{../reductions/contract.pdf}
\DeclareXref{reductions/def:roundmodel}{Definition}{Round model and attack sequences, typed by source mode}{../reductions/contract.pdf}
\DeclareXref{reductions/def:squashing}{Definition}{Squashing relation}{../reductions/contract.pdf}
\DeclareXref{reductions/domination/def:domination-link}{Definition}{The domination link}{../reductions/contract.pdf}
\DeclareXref{reductions/ex:identity}{Example}{The identity link}{../reductions/contract.pdf}
\DeclareXref{reductions/lem:linear}{Lemma}{Linearity: post-processing, private settings and mixtures}{../reductions/contract.pdf}
\DeclareXref{reductions/lem:party-local}{Lemma}{Party-local transport}{../reductions/contract.pdf}
\DeclareXref{reductions/lem:precomposition}{Lemma}{Licensed admissibility survives precomposition with a source channel}{../reductions/contract.pdf}
\DeclareXref{reductions/lem:row-relaxation}{Lemma}{A rounded constraint entry with its constant relaxed}{../reductions/contract.pdf}
\DeclareXref{reductions/lem:source-transfer}{Lemma}{Exactness of the source map in trusted-source mode}{../reductions/contract.pdf}
\DeclareXref{reductions/rem:intertwining}{Remark}{Classical maps before a squasher, and relabelling after one}{../reductions/contract.pdf}
\DeclareXref{reductions/rem:ordering}{Remark}{Recorded stage order, immaterial to soundness for licensed lists}{../reductions/contract.pdf}
\DeclareXref{reductions/req:factorised-emission}{Requirement}{Emission of several measuring parties as factors}{../reductions/contract.pdf}
\DeclareXref{reductions/req:party-local-tag}{Requirement}{The party-local tag and its scope}{../reductions/contract.pdf}
\DeclareXref{reductions/req:reference-companion}{Requirement}{Emitting the reference companion beside the artefact}{../reductions/contract.pdf}
\DeclareXref{reductions/thm:chain}{Theorem}{Composition of a chain}{../reductions/contract.pdf}
\DeclareXref{reductions/thm:emit}{Theorem}{Exact emission of a reduction instance}{../reductions/contract.pdf}
\DeclareXref{reductions/thm:emitted-product}{Theorem}{The emission of several measuring parties}{../reductions/contract.pdf}
\DeclareXref{reductions/thm:licence}{Theorem}{The attack-set licence}{../reductions/contract.pdf}
\DeclareXref{reductions/thm:transfer}{Theorem}{Security transfer from a constrained squashed model}{../reductions/contract.pdf}
\DeclareXref{reductions/thm:two-sided}{Theorem}{The two-sided reduction in trusted-source mode}{../reductions/contract.pdf}
\DeclareXref{reductions/wpfss/prop:flat-instance}{Proposition}{The flat instance of the post-processed receiver}{../reductions/contract.pdf}
\DeclareXref{transmitter/ass:ens}{Assumption}{Emission of the specified ensemble}{../transmitter/contract.pdf}
\DeclareXref{transmitter/ass:seq}{Assumption}{Sequential round structure}{../transmitter/contract.pdf}
\DeclareXref{transmitter/ass:src}{Assumption}{Retention of the system A in trusted-source mode}{../transmitter/contract.pdf}
\DeclareXref{transmitter/decoy/cor:emitted-decoy}{Corollary}{The emitted decoy declaration}{../transmitter/contract.pdf}
\DeclareXref{transmitter/def:source-segment}{Definition}{Source segment and its interface conditions}{../transmitter/contract.pdf}
\DeclareXref{transmitter/def:source-spec}{Definition}{Source specification}{../transmitter/contract.pdf}
\DeclareXref{transmitter/def:srcert}{Definition}{Source-replacement certificate}{../transmitter/contract.pdf}
\DeclareXref{transmitter/fact:sr}{Fact}{Source replacement for prepare-and-measure protocols}{../transmitter/contract.pdf}
\DeclareXref{transmitter/lem:ensemble-presentation}{Lemma}{Certificate defined by an ensemble presentation}{../transmitter/contract.pdf}
\DeclareXref{transmitter/rem:certificate-forms}{Remark}{Purity and exact declarability of the certificate forms}{../transmitter/contract.pdf}
\DeclareXref{transmitter/rem:da1}{Remark}{The one-dimensional retained system}{../transmitter/contract.pdf}
\DeclareXref{transmitter/rem:segment-scope}{Remark}{Scope of the source-segment conditions}{../transmitter/contract.pdf}
\DeclareXref{transmitter/rem:source-mode}{Remark}{The two source modes}{../transmitter/contract.pdf}
\DeclareXref{transmitter/req:attestation-record}{Requirement}{Pass-through of the retention attestation}{../transmitter/contract.pdf}
\DeclareXref{transmitter/req:failing-certificate}{Requirement}{Refusal of a certificate that fails a decided condition}{../transmitter/contract.pdf}
\DeclareXref{transmitter/req:licence-scope}{Requirement}{Certificate and attestation demanded only at $d_A\ge2$}{../transmitter/contract.pdf}
\DeclareXref{transmitter/req:marginal-fill}{Requirement}{The unique marginal served at $d_A=1$}{../transmitter/contract.pdf}
\DeclareXref{transmitter/thm:certificate-sufficiency}{Theorem}{Checkability and sufficiency of the source-replacement certificate}{../transmitter/contract.pdf}

\begin{document}
\begin{center}
{\LARGE\bfseries Protocol compiler (protocol-compiler)}\\[4pt]
{\large Contract, version 0.24.0}
\end{center}
\begin{overview}
This module compiles a \termref{term-protocol-specification}{protocol specification} into the finite \termref{term-protocol-declaration}{protocol declaration} the key-rate engine consumes. The specification has five parts: Alice's source as the transmitter module emits it, the receiver's measurement, a reduction chain in two segments (source links that replace Alice's signal ensemble, then receiver links that replace Bob's measurement), the public announcements with sifting and key rule, and the constraint entries; the declaration comes with an \termref{term-index-map}{index map} fixing the order of every alphabet and of the sectors. Operators on the measured system are lists of sector blocks against a declared sector map, every numerical condition is decided exactly on the declared binary64 numbers at no tolerance, and the compiler passes declared data through unchanged, constructing only the raw-key map, the index map and the measurement families that carry the weights of public choices. Its main results are three. A public choice of known bias, such as a basis choice, enters as a weight inside the measurement operators and never as a separate fraction: the weighted family stays a measurement, a weight that depends on no private setting is reproduced exactly as the probability of its choice value wherever the weighted family is the one emitted, and a weight that depends on the private setting of a measuring station, of either party, is applied after that station's reduction chain and commutes with it (\hyperlink{claim-protocol-compiler/thm:choice-encoding}{Theorem~\ref*{thm:choice-encoding}}). Passing the validator's eight checks establishes every target condition of a \termref{term-well-formed-declaration}{well-formed declaration}, and every check except the feasibility check is complete, so a well-formed declaration is refused only when neither a default candidate nor a supplied witness state is verified to be feasible (\hyperlink{claim-protocol-compiler/thm:validator-soundness}{Theorem~\ref*{thm:validator-soundness}}). Gates at emission decide whether a well-formed declaration may back a key-length claim, and every gate failure has one fixed outcome: a refusal, or an emission flagged as not eligible to back such a claim (\hyperlink{claim-protocol-compiler/thm:gates}{Theorem~\ref*{thm:gates}}). The compiler decides what the record of the reduction chain and the transmitter's output make decidable and names the rest: the soundness of each reduction link, the licence that restricts the attacks considered, and the physical fidelity of the declared data are obligations it records and never establishes.
\end{overview}
\medskip
\noindent{\small\textcolor{nnGrey}{\itshape Grey text}, here and throughout, is non-normative: it explains, and nothing in it is assumed or proved.}

\smallskip
\noindent{\small The \href{https://github.com/fraunhofersg/certiqs-key/blob/v0.1.0/spec/modules/protocol-compiler/design.md}{design notes} (non-normative) explain why the module is built as it is and draw its processing flowchart. The dependencies between the sections are drawn in \href{graphs/overview.svg}{the section dependency matrix}, and \href{https://github.com/fraunhofersg/certiqs-key/blob/v0.1.0/spec/modules/protocol-compiler/graphs/README.md}{the graph index} opens every section's dependency graph. The \href{digest.pdf}{statements digest} holds every definition, assumption, result and requirement without the proofs. \href{https://github.com/fraunhofersg/certiqs-key/blob/v0.1.0/harness/contract-reading.md}{How to read a contract} explains this layout and its links.}

{\small\tableofcontents}

\section*{Main results}\addcontentsline{toc}{section}{Main results}
\sectionintro{The module's most important results, those its overview names as what it establishes, with every assumption each rests on, directly or through the claims it uses. Each links to its statement and proof; the statements digest holds them in full.}
\par\noindent\hyperref[thm:choice-encoding]{\textbf{Theorem~\ref*{thm:choice-encoding}}}\enspace Encoding of public choices in the measurement operators\par\nobreak
{\small\itshape\color{nnGrey}\noindent A public choice of known bias enters as weights inside the measurement operators: the family stays a POVM, fixed weights are reproduced exactly as choice probabilities where the weighted family is emitted, and setting-dependent weights commute with the chain. A fraction outside them is refused.\par}\nobreak
\claimmeta{\metalabel{Rests on}\aref{A-COIN}{ass:coin}, \aref{A-LINK}{ass:link}, \aref{A-SPEC}{ass:spec}}
\par\noindent\hyperref[thm:validator-soundness]{\textbf{Theorem~\ref*{thm:validator-soundness}}}\enspace Soundness, completeness and decidability of the validator\par\nobreak
{\small\itshape\color{nnGrey}\noindent Passing the validator's eight checks establishes every target condition of a well-formed declaration, and every check but the feasibility check is complete, so a well-formed declaration is refused only when no default or supplied candidate is verified feasible. Every decision is finite and exact.\par}\nobreak
\claimmeta{\metalabel{Rests on}\aref{A-SPEC}{ass:spec}}
\par\noindent\hyperref[thm:gates]{\textbf{Theorem~\ref*{thm:gates}}}\enspace The emission gates and their outcomes\par\nobreak
{\small\itshape\color{nnGrey}\noindent Each gate is decidable or a presence test, and each failure has one fixed outcome, a refusal or an emission flagged as backing no key length, never a default, dropped entry, changed source mode or licence the compiler writes. Passing asserts nothing about the apparatus or attack-set inclusion.\par}\nobreak
\claimmeta{\metalabel{Rests on}\aref{A-ATK}{ass:atk}, \aref{A-RHOM}{ass:rhom}, \extaref{../transmitter/contract.pdf}{transmitter/ass:seq}{A-SEQ}, \aref{A-SPEC}{ass:spec}, \extaref{../transmitter/contract.pdf}{transmitter/ass:src}{A-SRC}}

\section*{Assumptions at a glance}\addcontentsline{toc}{section}{Assumptions at a glance}
\sectionintro{Every assumption this contract states or its claims assume, with its category and the claims of this contract that assume it directly. An assumption of another contract links into that contract's PDF.}
{\footnotesize\begin{xltabular}{\linewidth}{@{}>{\raggedright\arraybackslash}p{1.6cm}>{\raggedright\arraybackslash}p{1.9cm}>{\raggedright\arraybackslash}X>{\raggedright\arraybackslash}p{2.3cm}>{\raggedright\arraybackslash}p{2.8cm}@{}}
\toprule
Assumption & Category & Title & Stated as & Assumed directly by\\
\midrule
\endhead
\aref{A-IDEAL}{ass:ideal} & receiver & Ideal qubit receiver on the Band-A route & \mref{Assumption}{ass:ideal} & \mref{Thm.}{thm:band-a}\\
\addlinespace
\extaref{../reductions/contract.pdf}{reductions/ass:meas}{A-MEAS} & receiver & Each measuring party's device as one POVM per round & \extref{../reductions/contract.pdf}{reductions/ass:meas} & \mref{Req.}{req:link-row-provenance}\\
\addlinespace
\aref{A-ALICE-MEAS}{ass:alice-meas} & protocol structure & Alice's declared measurement where no source certificate is demanded & \mref{Assumption}{ass:alice-meas} & ---\\
\addlinespace
\aref{A-RHOM}{ass:rhom} & protocol structure & Round homogeneity of the declaration & \mref{Assumption}{ass:rhom} & \mref{Thm.}{thm:gates}\\
\addlinespace
\aref{A-SPEC}{ass:spec} & protocol structure & Specification faithful to the apparatus & \mref{Assumption}{ass:spec} & \mref{Thm.}{thm:typing}, \mref{Thm.}{thm:measurements-and-statistics}, \mref{Thm.}{thm:sifting-and-key}\\
\addlinespace
\aref{A-COIN}{ass:coin} & local randomness & Public choices as private coins of declared bias & \mref{Assumption}{ass:coin} & \mref{Thm.}{thm:choice-encoding}, \mref{Rem.}{rem:exactness}, \mref{Rem.}{rem:relabel}\\
\addlinespace
\aref{A-ATK}{ass:atk} & supplied artefact & Attack-set licence & \mref{Assumption}{ass:atk} & \mref{Thm.}{thm:gates}\\
\addlinespace
\aref{A-LINK}{ass:link} & supplied artefact & Soundness of every reduction link & \mref{Assumption}{ass:link} & \mref{Thm.}{thm:choice-encoding}, \mref{Thm.}{thm:chain-consumption}\\
\bottomrule
\end{xltabular}}

\section*{Notation}\addcontentsline{toc}{section}{Notation}
\sectionintro{The symbols of the repository's notation catalogue (\href{../../notation.md}{spec/notation.md}) that this contract uses, in the catalogue's order, with their meaning.}
{\small\begin{xltabular}{\linewidth}{@{}>{\raggedright\arraybackslash}p{3.6cm}>{\raggedright\arraybackslash}X@{}}
\toprule
Symbol & Meaning\\
\midrule
\endhead
$\Hs{A}$ & Hilbert space of the system label $A$\\
$\Lin{A}$ & Bounded linear operators on $\mathcal{H}(A)$\\
$\Pos{A}$ & Positive semidefinite operators\\
$\Seq{A}$ & Normalized states\\
$\id$, $\idon{C}$ & Identity operator (contextual / named)\\
$\idmap$ & Identity channel\\
$\leqN$, $\gtN$, $\eqn$ & Block indices for \texttt{\textbackslash{}blk}\\
$\Decl$ & Protocol declaration $\big(A,\tau_A;\,Q,\mathbb{S};\,\{M_x\};\,\{F_v\};\,\mathfrak{P};\,a_A,a_B;\,\phi;\,r,\mathsf{h};\,\mathfrak{L}\big)$ \textemdash{} the tuple of \extref{../meat-core/contract.pdf}{meat-core/def:declaration}: the party partition $\mathfrak{P}$ is a slot, and the marginal slot $\tau_A$ takes a state or the no-marginal value (trusted- vs untrusted-source mode)\\
$\SecSet$ & Sector index set of $Q$\\
$\Qsys$ & The declared sector-carrying measured system\\
$\Party$ & The declared party partition: the readout maps $\pi_{\mathsf A}$, $\pi_{\mathsf B}$ onto the alphabets $\mathcal{O}_{\mathsf A}$, $\mathcal{O}_{\mathsf B}$, with a party label for $A$ and, for each sector $s$, a declared tensor factorisation of $\mathcal{H}(Q_s)$ with a party label per factor\\
$\Xset$ & Outcome alphabet of the measurement on $A$ \textemdash{} Alice's under the default partition\\
$\Vset$ & The measuring parties' declared joint outcome alphabet \textemdash{} Bob's alone under the default partition\\
$\CAset$, $\CBset$ & The alphabets of the first and second message of the transcript, which is the $L=2$ case\\
$\Kann$ & Transcript alphabet $\prod_{m=1}^{L}\mathcal{C}_m$, which at $L=2$ is $\mathcal{C}_A\times\mathcal{C}_B$\\
$\Chat$, $\Chatset$ & The acceptance-statistic register $\widehat{C}$ and its alphabet $\widehat{\mathcal{C}}$\\
$\Chkey$ & Key-live acceptance values\\
$\Zval$, $\Zset$ & Key alphabet $\{0,\dots,d_Z{-}1\}$ and $\mathcal{Z}\sqcup\{\perp\}$\\
$\hold$ & Declared key holder, $\mathsf{h}\in\{\mathsf A,\mathsf B\}$\\
$\CList$ & Declared constraint list $\{(\hat O_k,\bowtie_k,\Theta_k)\}$\\
$\Phimap$ & The test-statistics map $\mathcal{L}(AQ)\to\mathbb{R}^{\widehat{\mathcal{C}}}$\\
$\Feas$ & Feasible set of sector-block-diagonal marginals: $\operatorname{tr}\omega=1$ always; $\omega_A=\tau_A$ in trusted-source mode only; $\mathfrak{L}$ if declared (\extref{../meat-core/contract.pdf}{meat-core/def:feasible})\\
$\Chan$, $\Prb$ & Channel input; probe input\\
\bottomrule
\end{xltabular}}

\section*{Terms}\addcontentsline{toc}{section}{Terms}
\sectionintro{The phrases this contract defines, and the terms of other contracts it uses, each with a one-line gloss and the definition that fixes it. In the text, a linked term opens its definition.}
{\small\begin{xltabular}{\linewidth}{@{}>{\raggedright\arraybackslash}p{3.6cm}>{\raggedright\arraybackslash}X>{\raggedright\arraybackslash}p{2.6cm}@{}}
\toprule
Term & Meaning & Defined in\\
\midrule
\endhead
\termref{term-announcement-map}{announcement map} & A map from one party's readout alphabet to the alphabet of what that party announces. & \mref{Definition}{def:declaration}\\
\addlinespace[2pt]
\exttermref{../foundations/contract.pdf}{term-attack-sequence}{attack sequence} & one channel per round whose output is the per-round system, with Alice's system where the adversary supplies it, and Eve's new memory & \extref{../foundations/contract.pdf}{foundations/def:round-model}\\
\addlinespace[2pt]
\exttermref{../foundations/contract.pdf}{term-attack-set}{attack set} & the channels whose adjoint maps every listed operator into the stated relation with its constant & \extref{../foundations/contract.pdf}{foundations/def:round-model}\\
\addlinespace[2pt]
\termref{term-band-a}{Band A} & The direct route: a finite measurement declared on one sector, an empty source segment and one identity receiver link, given as data, with empty slots, that delivers the declared measurement. & \mref{Definition}{def:banda}\\
\addlinespace[2pt]
\termref{term-basic-format}{basic format} & The format with one Alice sector, one measured system, Alice's measurement as a POVM, two messages neither reading the other, and no party-local entry. & \mref{Definition}{def:declaration}\\
\addlinespace[2pt]
\exttermref{../foundations/contract.pdf}{term-binary64-numbers}{binary64 numbers} & the finite values of the IEEE 754 binary64 format, each a dyadic rational & \extref{../foundations/contract.pdf}{foundations/def:binary64}\\
\addlinespace[2pt]
\exttermref{../transmitter/contract.pdf}{term-block-form}{block form} & The certificate form declared as positive blocks, one per sector, on matched sector structures of Alice's system and the virtual system. & \extref{../transmitter/contract.pdf}{transmitter/def:srcert}\\
\addlinespace[2pt]
\exttermref{../foundations/contract.pdf}{term-block-list}{block list} & one square complex matrix per block of a sector map, of that block's dimension, identified with their direct sum & \extref{../foundations/contract.pdf}{foundations/def:sector-map}\\
\addlinespace[2pt]
\termref{term-choice-alphabet}{choice alphabet} & The finite alphabet, with a declared order, of a public choice whose weight depends on the private setting of the station that measures the family. & \mref{Definition}{def:weighted}\\
\addlinespace[2pt]
\termref{term-choice-weighted-family}{choice-weighted family} & The measurement whose elements are declared measurement elements multiplied by the declared weight of a public choice value. & \mref{Definition}{def:weighted}\\
\addlinespace[2pt]
\termref{term-compilation-map}{compilation map} & The map from a protocol specification to the protocol declaration it fills, slot by slot. & \mref{Definition}{def:compile}\\
\addlinespace[2pt]
\exttermref{../foundations/contract.pdf}{term-computed-sequential-sum}{computed sequential sum} & the left-to-right accumulation of a finite sequence from zero with every addition rounded by fl, entrywise, with no re-ordering and no re-association & \extref{../foundations/contract.pdf}{foundations/def:declared-order}\\
\addlinespace[2pt]
\exttermref{../foundations/contract.pdf}{term-consumer}{consumer} & a module other than a supplier that reads declared data & \extref{../foundations/contract.pdf}{foundations/def:declared-data}\\
\addlinespace[2pt]
\exttermref{../reductions/contract.pdf}{term-data-sector}{data sector} & A sector of a reduction instance's target whose blocks are the datum's matrices. & \extref{../reductions/contract.pdf}{reductions/def:reduction-instance}\\
\addlinespace[2pt]
\exttermref{../receiver/contract.pdf}{term-declared-coarse-graining}{declared coarse-graining} & A surjection from the raw click patterns onto the declared emitted alphabet, supplied as data by the protocol declaration. & \extref{../receiver/contract.pdf}{receiver/def:coarse}\\
\addlinespace[2pt]
\exttermref{../foundations/contract.pdf}{term-declared-data}{declared data} & the data an instance states and the data a producer emits as its own declaration, with the sets such data define & \extref{../foundations/contract.pdf}{foundations/def:declared-data}\\
\addlinespace[2pt]
\exttermref{../reductions/contract.pdf}{term-declared-datum}{declared datum} & The datum the emission normalizer outputs: exactly admissible on the declared lattice, and carried entry by entry in the emitted target. & \extref{../reductions/contract.pdf}{reductions/def:emitnorm}\\
\addlinespace[2pt]
\exttermref{../reductions/contract.pdf}{term-declared-model}{declared model} & The handed model with its datum replaced by the declared datum, as the instance's model map defines it. & \extref{../reductions/contract.pdf}{reductions/def:emitnorm}\\
\addlinespace[2pt]
\exttermref{../foundations/contract.pdf}{term-declared-order}{declared order} & a total order on a finite alphabet, fixed together with the alphabet & \extref{../foundations/contract.pdf}{foundations/def:declared-order}\\
\addlinespace[2pt]
\exttermref{../foundations/contract.pdf}{term-declared-value}{declared value} & a binary64 number stated as a datum, which denotes exactly that dyadic rational and approximates nothing & \extref{../foundations/contract.pdf}{foundations/def:binary64}\\
\addlinespace[2pt]
\exttermref{../foundations/contract.pdf}{term-declarer}{declarer} & the party whose declaration declared data are: the instance that states them or the producer that emits them & \extref{../foundations/contract.pdf}{foundations/def:declared-data}\\
\addlinespace[2pt]
\termref{term-default-party-partition}{default party partition} & The party partition in which Alice reads x and holds her system, and Bob reads v and holds every sector. & \mref{Definition}{def:target}\\
\addlinespace[2pt]
\exttermref{../foundations/contract.pdf}{term-defects}{defects} & the entrywise maximum moduli of the differences that the predicates of data name: the Hermiticity defect and the completeness, column, trace and sum residues & \extref{../foundations/contract.pdf}{foundations/def:tolerance}\\
\addlinespace[2pt]
\termref{term-ensemble-form}{ensemble form} & The trusted-source form of Alice's measurement slot that holds positive blocks summing to the marginal, and denotes a POVM-form declaration. & \mref{Definition}{def:declaration}\\
\addlinespace[2pt]
\exttermref{../transmitter/contract.pdf}{term-ensemble-presentation}{ensemble presentation} & The certificate form that declares only the weighted virtual ensemble in sector blocks and defines the certificate from it. & \extref{../transmitter/contract.pdf}{transmitter/def:srcert}\\
\addlinespace[2pt]
\exttermref{../foundations/contract.pdf}{term-entrywise-enclosure}{entrywise enclosure} & a Gaussian-rational Hermitian centre together with symmetric non-negative rational radii that bound, entry by entry, the modulus of the difference between a matrix and that centre & \extref{../foundations/contract.pdf}{foundations/def:exact-reference}\\
\addlinespace[2pt]
\exttermref{../foundations/contract.pdf}{term-exact-island}{exact island} & the computations of the engine, its instances' declarations included, by which a consumer of data or the declarer about to declare them decides whether they satisfy the conditions their declaration states, a declarer computes and decides the reference of a measurement family, a consumer's replacement of an emitted scalar bound is decided to lie on the bound's weakening side, a producer constructs the operator data it emits, the acceptance side of the optimisation is constructed and certified, or a lower bound on a linear minimum over the feasible side of the optimisation is certified & \extref{../foundations/contract.pdf}{foundations/def:exact-island}\\
\addlinespace[2pt]
\exttermref{../foundations/contract.pdf}{term-exactly-admissible}{exactly admissible} & said of a declared operator family whose sector map agrees with the given one, whose blocks are exactly Hermitian and positive semidefinite, and whose sum in the declared order equals the constant times the identity on every block and rounds nowhere; the constant is one where none is named & \extref{../foundations/contract.pdf}{foundations/def:exact-admissibility}\\
\addlinespace[2pt]
\termref{term-factor}{factor} & One entry of a factorised sector map: a system tag, a party label and a sector map of its own. & \mref{Definition}{def:declaration}\\
\addlinespace[2pt]
\exttermref{../foundations/contract.pdf}{term-factor-enclosure}{factor enclosure} & a matrix box of columns together with an exact positive diagonal core, such that the matrix is the product of a matrix L the box contains, the core and the adjoint of L & \extref{../foundations/contract.pdf}{foundations/def:factor-enclosure}\\
\addlinespace[2pt]
\termref{term-feasible-set}{feasible set} & The set of normalised block-list states on A and the measured system that meet the declared marginal and every constraint entry, each read exactly. & \mref{Definition}{def:feasible}\\
\addlinespace[2pt]
\termref{term-gate-outcome-table}{gate outcome table} & The table that assigns to each gate failure its outcome, a refusal or an emission that is not key-length-eligible, or none for the two certificate-absent rows, the refusal taking precedence where a declaration meets both kinds of row. & \mref{Definition}{def:gate-outcome-table}\\
\addlinespace[2pt]
\termref{term-general}{general} & Said of a constraint entry declaring a Hermitian matrix on Alice's system as its right-hand side. & \mref{Definition}{def:homog}\\
\addlinespace[2pt]
\termref{term-homogeneous}{homogeneous} & Said of a constraint entry, not declared party-local, whose right-hand side is a real scalar; in trusted-source mode it denotes that scalar times the marginal. & \mref{Definition}{def:homog}\\
\addlinespace[2pt]
\termref{term-index-map}{index map} & The emitted choice of a bijection from each alphabet and from the sector set onto an initial segment of the non-negative integers, through which every emitted list is indexed. & \mref{Definition}{def:compile}\\
\addlinespace[2pt]
\exttermref{../foundations/contract.pdf}{term-instance}{instance} & the party that fixes one protocol's device data, protocol data and run configuration & \extref{../foundations/contract.pdf}{foundations/def:declared-data}\\
\addlinespace[2pt]
\exttermref{../transmitter/contract.pdf}{term-interface-conditions}{interface conditions} & The conditions (S1) to (S4) a source segment satisfies: descriptor match, index preservation, the source identity, and the slots each link carries. & \extref{../transmitter/contract.pdf}{transmitter/def:source-segment}\\
\addlinespace[2pt]
\termref{term-interface-conditions-on-a-chain}{interface conditions on a chain} & The conditions a reduction chain presented to the compiler must satisfy, on each of its two segments and between them. & \mref{Definition}{def:chaincond}\\
\addlinespace[2pt]
\exttermref{../reductions/contract.pdf}{term-joint-alphabet}{joint alphabet} & The product of the outcome alphabet and the settings, ordered with the outcome varying fastest. & \extref{../reductions/contract.pdf}{reductions/def:joint-family}\\
\addlinespace[2pt]
\exttermref{../reductions/contract.pdf}{term-joint-family}{joint family} & The POVM over the joint alphabet whose element at an outcome and a setting is the class weight times the conditional element of that outcome given that setting. & \extref{../reductions/contract.pdf}{reductions/def:joint-family}\\
\addlinespace[2pt]
\termref{term-joint-sector}{joint sector} & A pair of an Alice sector and a sector of the measured system; the joint sectors are ordered with the Alice sector slowest. & \mref{Definition}{def:declaration}\\
\addlinespace[2pt]
\termref{term-key-length-eligible}{key-length-eligible} & An emitted declaration that may back a key-length claim: every check passes, the source-replacement licence record is present and re-decided where required, and a non-empty constraint list is licensed and carries an attack-set licence record. & \mref{Definition}{def:eligible}\\
\addlinespace[2pt]
\termref{term-key-live}{key-live} & Said of an acceptance value that some announcement class passing sifting maps to. & \mref{Definition}{def:declaration}\\
\addlinespace[2pt]
\exttermref{../foundations/contract.pdf}{term-lattice-exponent}{lattice exponent} & an integer k with 1 \ensuremath{\le} k \ensuremath{\le} 53 & \extref{../foundations/contract.pdf}{foundations/def:binary64}\\
\addlinespace[2pt]
\exttermref{../reductions/contract.pdf}{term-licence-record}{licence record} & The record of the attack-set licence a link carries exactly when its constraint list is non-empty. & \extref{../reductions/contract.pdf}{reductions/def:link}\\
\addlinespace[2pt]
\termref{term-licensed-for-the-declared-source-mode}{licensed for the declared source mode} & Said of a constraint list every entry of which is in the form the declared source mode admits, or is party-local. & \mref{Definition}{def:homog}\\
\addlinespace[2pt]
\termref{term-lifted-operator}{lifted operator} & The operator on the whole measured system that a scoped entry's operator denotes: the identity on every other factor. & \mref{Definition}{def:declaration}\\
\addlinespace[2pt]
\termref{term-marginal-slot}{marginal slot} & The slot of a declaration that holds Alice's marginal state, or the value marking untrusted-source mode. & \mref{Definition}{def:declaration}\\
\addlinespace[2pt]
\exttermref{../foundations/contract.pdf}{term-matrix-box}{matrix box} & a matrix of Gaussian-rational centres together with non-negative rational radii that bound, entry by entry, the modulus of the difference between a rectangular complex matrix and that centre & \extref{../foundations/contract.pdf}{foundations/def:factor-enclosure}\\
\addlinespace[2pt]
\termref{term-measured-factor}{measured factor} & A factor tagged Q: a system one party measures, with its own sector map, alphabet and family. & \mref{Definition}{def:declaration}\\
\addlinespace[2pt]
\exttermref{../reductions/contract.pdf}{term-measured-system}{measured system} & The quantum system the measuring parties receive in one round, and no other system; Bob's alone where he alone measures. & \extref{../reductions/contract.pdf}{reductions/def:protocol}\\
\addlinespace[2pt]
\exttermref{../reductions/contract.pdf}{term-measurement-channel}{measurement channel} & The channel that measures a POVM on a quantum input and writes the outcome into a classical register. & \extref{../reductions/contract.pdf}{reductions/def:msm}\\
\addlinespace[2pt]
\exttermref{../reductions/contract.pdf}{term-measuring-parties}{measuring parties} & The parties a declaration names as measuring; their measurements are the factors of the round's measurement channel. & \extref{../reductions/contract.pdf}{reductions/def:protocol}\\
\addlinespace[2pt]
\termref{term-message}{message} & One entry of a transcript protocol: a sender, a finite alphabet and a total table. & \mref{Definition}{def:declaration}\\
\addlinespace[2pt]
\exttermref{../foundations/contract.pdf}{term-module-boundary}{module boundary} & the passage of declared data from their declarer to a consumer & \extref{../foundations/contract.pdf}{foundations/def:declared-data}\\
\addlinespace[2pt]
\exttermref{../transmitter/contract.pdf}{term-obligation-record}{obligation record} & What a theorem-supplied source link carries in place of a decided identity: its kind, the soundness claim, the assumptions the obligation carries, and the family member its descriptors are. & \extref{../transmitter/contract.pdf}{transmitter/def:source-link}\\
\addlinespace[2pt]
\exttermref{../reductions/contract.pdf}{term-outcome-alphabet}{outcome alphabet} & The finite ordered set of device outcomes that a private setting is joined to. & \extref{../reductions/contract.pdf}{reductions/def:joint-family}\\
\addlinespace[2pt]
\termref{term-party-partition}{party partition} & The readout maps of the two parties, with a party label for Alice's system and for each tensor factor of every sector. & \mref{Definition}{def:declaration}\\
\addlinespace[2pt]
\termref{term-party-local}{party-local} & Said of a scoped constraint entry, tagged so, whose condition holds against the marginal of the state on the systems other than its factor. & \mref{Definition}{def:homog}\\
\addlinespace[2pt]
\exttermref{../reductions/contract.pdf}{term-party-local-form}{party-local form} & The form of a constraint entry scoped to one factor of a product measured system: tagged party-local and read as an operator inequality against the marginal actually present. & \extref{../reductions/contract.pdf}{reductions/def:licensed}\\
\addlinespace[2pt]
\termref{term-povm-form}{POVM form} & The form of Alice's measurement slot that holds the POVM itself. & \mref{Definition}{def:declaration}\\
\addlinespace[2pt]
\exttermref{../foundations/contract.pdf}{term-producer}{producer} & a module other than a supplier that emits numerical data, constructed from declared data it read, to another module as its own declaration & \extref{../foundations/contract.pdf}{foundations/def:declared-data}\\
\addlinespace[2pt]
\termref{term-protocol-declaration}{protocol declaration} & The finite object the compiler emits and the key-rate engine consumes: systems, sector map, marginal slot, measurement families, public maps, raw-key map and constraint list. & \mref{Definition}{def:declaration}\\
\addlinespace[2pt]
\termref{term-protocol-specification}{protocol specification} & The compiler's input: a source part, a receiver part, a reduction chain, a public part and a constraint part. & \mref{Definition}{def:spec}\\
\addlinespace[2pt]
\termref{term-public-choice-alphabet}{public-choice alphabet} & The finite alphabet of values of a per-round public choice of known bias, each value with a declared positive weight. & \mref{Definition}{def:weighted}\\
\addlinespace[2pt]
\exttermref{../transmitter/contract.pdf}{term-pure-form}{pure form} & The certificate form declared as a rank-one projector, on the one-sector structure. & \extref{../transmitter/contract.pdf}{transmitter/def:srcert}\\
\addlinespace[2pt]
\exttermref{../reductions/contract.pdf}{term-reachable-set-statement}{reachable-set statement} & The set of reachable squashed states of a link, for its channel witness and the declared family of inputs. & \extref{../reductions/contract.pdf}{reductions/def:link}\\
\addlinespace[2pt]
\termref{term-readout-map}{readout map} & A total map from the outcome pairs to one party's declared finite alphabet, saying what that party reads. & \mref{Definition}{def:declaration}\\
\addlinespace[2pt]
\exttermref{../reductions/contract.pdf}{term-receiver-segment}{receiver segment} & The list of receiver-stage typed links of a reduction chain. & \extref{../reductions/contract.pdf}{reductions/def:chain}\\
\addlinespace[2pt]
\exttermref{../reductions/contract.pdf}{term-reduction-chain}{reduction chain} & A pair of a source segment and a receiver segment of typed links. & \extref{../reductions/contract.pdf}{reductions/def:chain}\\
\addlinespace[2pt]
\exttermref{../foundations/contract.pdf}{term-round-model}{round model} & n rounds of a per-round system S, with the attack sequences and attack sets below & \extref{../foundations/contract.pdf}{foundations/def:round-model}\\
\addlinespace[2pt]
\exttermref{../foundations/contract.pdf}{term-round-homogeneous}{round-homogeneous} & one fixed list of objects describes every round, and every round behaves as they describe whatever the other rounds do & \extref{../foundations/contract.pdf}{foundations/def:round-homogeneity}\\
\addlinespace[2pt]
\exttermref{../foundations/contract.pdf}{term-rounds-nowhere}{rounds nowhere} & said of an accumulation every one of whose binary64 additions is exact & \extref{../foundations/contract.pdf}{foundations/def:declared-order}\\
\addlinespace[2pt]
\termref{term-scope}{scope} & The measured factor that an entry of the constraint list is declared on. & \mref{Definition}{def:declaration}\\
\addlinespace[2pt]
\exttermref{../foundations/contract.pdf}{term-sector-map}{sector map} & a block kind together with a non-empty, ordered list of block identifiers, each with a dimension & \extref{../foundations/contract.pdf}{foundations/def:sector-map}\\
\addlinespace[2pt]
\exttermref{../reductions/contract.pdf}{term-sector-set}{sector set} & The finite ordered list of a reduction instance's target sectors, data and scalar, each with a declared dimension. & \extref{../reductions/contract.pdf}{reductions/def:reduction-instance}\\
\addlinespace[2pt]
\exttermref{../transmitter/contract.pdf}{term-sector-structure}{sector structure} & A decomposition of Alice's retained system into orthogonal sectors, each spanned by declared basis vectors. & \extref{../transmitter/contract.pdf}{transmitter/def:source-spec}\\
\addlinespace[2pt]
\termref{term-setting-alphabet}{setting alphabet} & The finite alphabet of the private per-round setting of the station that measures the family, onto which the declared setting map sends each outcome. & \mref{Definition}{def:weighted}\\
\addlinespace[2pt]
\termref{term-setting-component}{setting component} & The value of the measuring station's private setting that an outcome records, read by the declared setting map. & \mref{Definition}{def:weighted}\\
\addlinespace[2pt]
\termref{term-setting-dependent-form}{setting-dependent form} & The choice-weighted family of a measured system after its chain: each delivered element times a stochastic weight read through the measuring station's setting. & \mref{Definition}{def:weighted}\\
\addlinespace[2pt]
\termref{term-setting-independent-form}{setting-independent form} & The choice-weighted family whose weights are one declared number per choice value, the same in every round. & \mref{Definition}{def:weighted}\\
\addlinespace[2pt]
\termref{term-soundness-condition-of-a-party-local-entry}{soundness condition of a party-local entry} & The condition a party-local entry is sound under: every attack reaches its scoped factor through a channel on that factor alone whose adjoint, applied to the entry's operator, stands in the entry's relation to its scalar times the identity. & \mref{Definition}{def:homog}\\
\addlinespace[2pt]
\exttermref{../transmitter/contract.pdf}{term-soundness-registry}{soundness registry} & The closed list, emitted with every artefact, of admitted source-link kinds with their soundness claims and carried assumptions. & \extref{../transmitter/contract.pdf}{transmitter/def:source-segment}\\
\addlinespace[2pt]
\exttermref{../transmitter/contract.pdf}{term-source-descriptor}{source descriptor} & The transmitted system and signal ensemble a source link starts from. & \extref{../transmitter/contract.pdf}{transmitter/def:source-link}\\
\addlinespace[2pt]
\exttermref{../reductions/contract.pdf}{term-source-family}{source family} & The system and POVM a receiver-stage link starts from. & \extref{../reductions/contract.pdf}{reductions/def:link}\\
\addlinespace[2pt]
\exttermref{../transmitter/contract.pdf}{term-source-link}{source link} & A link of a reduction chain's source stage that maps one Alice-side ensemble to another with the same labels and weights. & \extref{../transmitter/contract.pdf}{transmitter/def:source-link}\\
\addlinespace[2pt]
\exttermref{../transmitter/contract.pdf}{term-source-model}{source model} & A submodule's named definition of the signal ensemble as mathematical objects, possibly infinite dimensional, rather than as matrices. & \extref{../transmitter/contract.pdf}{transmitter/def:source-spec}\\
\addlinespace[2pt]
\exttermref{../transmitter/contract.pdf}{term-source-segment}{source segment} & The ordered, possibly empty list of source links leading from the virtual ensemble to the specified one. & \extref{../transmitter/contract.pdf}{transmitter/def:source-segment}\\
\addlinespace[2pt]
\exttermref{../transmitter/contract.pdf}{term-source-specification}{source specification} & The declared description of Alice's side of a protocol: alphabet, retained system with sectors, measurement, signal ensemble and source mode. & \extref{../transmitter/contract.pdf}{transmitter/def:source-spec}\\
\addlinespace[2pt]
\exttermref{../transmitter/contract.pdf}{term-source-replacement-certificate}{source-replacement certificate} & An operator on Alice's retained system and the virtual transmitted system whose marginal on her side is her declared marginal and which, measured on her side by her POVM, leaves the weighted virtual members, with retention attested. & \extref{../transmitter/contract.pdf}{transmitter/def:srcert}\\
\addlinespace[2pt]
\exttermref{../transmitter/contract.pdf}{term-specified-signal-ensemble}{specified signal ensemble} & The weights and signal states, one per letter, that Alice's source is specified to emit. & \extref{../transmitter/contract.pdf}{transmitter/def:source-spec}\\
\addlinespace[2pt]
\exttermref{../reductions/contract.pdf}{term-squashed-sequence}{squashed sequence} & The attack sequence followed, in every round, by the squashing channel on the measured system and the identity on everything else. & \extref{../reductions/contract.pdf}{reductions/def:roundmodel}\\
\addlinespace[2pt]
\exttermref{../reductions/contract.pdf}{term-squashing-relation}{squashing relation} & The relation in which the adjoint of a channel takes every target POVM element to the physical one and every carried operator to its prescribed physical operator. & \extref{../reductions/contract.pdf}{reductions/def:squashing}\\
\addlinespace[2pt]
\exttermref{../foundations/contract.pdf}{term-sum-in-the-declared-order}{sum in the declared order} & the sequential sum of a family indexed by an alphabet, taken in the alphabet's declared order & \extref{../foundations/contract.pdf}{foundations/def:declared-order}\\
\addlinespace[2pt]
\termref{term-support-of-the-weighting}{support of the weighting} & The pairs of an outcome and a choice value whose weight under the declared stochastic matrix is positive. & \mref{Definition}{def:weighted}\\
\addlinespace[2pt]
\termref{term-system-tag}{system tag} & The label A or Q of a factor, saying whether it is Alice's system or a measured system. & \mref{Definition}{def:declaration}\\
\addlinespace[2pt]
\exttermref{../transmitter/contract.pdf}{term-target-descriptor}{target descriptor} & The transmitted system and signal ensemble a source link produces. & \extref{../transmitter/contract.pdf}{transmitter/def:source-link}\\
\addlinespace[2pt]
\exttermref{../reductions/contract.pdf}{term-target-family}{target family} & The target POVM of a reduction instance at a datum, assigned one sector block at a time from the datum and the scalar values. & \extref{../reductions/contract.pdf}{reductions/def:reduction-instance}\\
\addlinespace[2pt]
\exttermref{../reductions/contract.pdf}{term-target-povm}{target POVM} & The POVM on the squashed system that a squashing relation pulls back to the physical POVM. & \extref{../reductions/contract.pdf}{reductions/def:squashing}\\
\addlinespace[2pt]
\exttermref{../transmitter/contract.pdf}{term-theorem-constructor}{theorem constructor} & A named claim of a submodule that supplies a source link's channel for each member of a declared family of descriptors. & \extref{../transmitter/contract.pdf}{transmitter/def:source-link}\\
\addlinespace[2pt]
\exttermref{../reductions/contract.pdf}{term-total-weight}{total weight} & The sum of the certified weights over all outcomes. & \extref{../reductions/contract.pdf}{reductions/wpfss/def:targets}\\
\addlinespace[2pt]
\termref{term-transcript-protocol}{transcript protocol} & The announcement slot of the general format: an ordered list of messages, each a function of its sender's readout and of the earlier messages. & \mref{Definition}{def:declaration}\\
\addlinespace[2pt]
\exttermref{../transmitter/contract.pdf}{term-transmitted-system}{transmitted system} & The system Alice sends into the channel, carrying the specified signal ensemble. & \extref{../transmitter/contract.pdf}{transmitter/def:source-spec}\\
\addlinespace[2pt]
\exttermref{../transmitter/contract.pdf}{term-trusted-source-mode}{trusted-source mode} & The source mode in which Alice's retained system never reaches the adversary, with a declared marginal and, above dimension one, a certificate. & \extref{../transmitter/contract.pdf}{transmitter/def:source-spec}\\
\addlinespace[2pt]
\exttermref{../reductions/contract.pdf}{term-typed-link}{typed link} & A link of a reduction chain, at the source stage or the receiver stage, carrying what is needed to re-check it. & \extref{../reductions/contract.pdf}{reductions/def:link}\\
\addlinespace[2pt]
\exttermref{../transmitter/contract.pdf}{term-untrusted-source-mode}{untrusted-source mode} & The source mode with no certificate and no marginal, in which the adversary may have prepared Alice's retained system. & \extref{../transmitter/contract.pdf}{transmitter/def:source-spec}\\
\addlinespace[2pt]
\exttermref{../transmitter/contract.pdf}{term-virtual-ensemble}{virtual ensemble} & The finite-dimensional ensemble the source segment starts from and the certificate is written against. & \extref{../transmitter/contract.pdf}{transmitter/def:source-segment}\\
\addlinespace[2pt]
\exttermref{../transmitter/contract.pdf}{term-virtual-transmitted-system}{virtual transmitted system} & The transmitted system of the virtual ensemble, where the source segment starts. & \extref{../transmitter/contract.pdf}{transmitter/def:source-segment}\\
\addlinespace[2pt]
\exttermref{../transmitter/contract.pdf}{term-weighted-members}{weighted members} & The positive operators, weight times state, one per letter, through which an ensemble is declared and read. & \extref{../transmitter/contract.pdf}{transmitter/def:source-spec}\\
\addlinespace[2pt]
\termref{term-well-formed-declaration}{well-formed declaration} & A declaration, together with its index map, that meets the eight target conditions D(i) to D(viii). & \mref{Definition}{def:target}\\
\addlinespace[2pt]
\termref{term-well-formed-specification}{well-formed specification} & A protocol specification with finite alphabets and systems, total tables, exactly admissible declared families, the presented physical one included, normalised weights and Hermitian, tagged constraint blocks. & \mref{Definition}{def:well-formed-spec}\\
\bottomrule
\end{xltabular}}

\clearpage
\part{Mathematics}

\section{The protocol declaration}\label{group-protocol-compiler/def:declaration}
\sectionintro{The object the compiler emits is a \termref{term-protocol-declaration}{protocol declaration}: finite systems, a sector map, Alice's \termref{term-marginal-slot}{marginal slot}, the two measurement families, the public maps, the raw-key map and a typed constraint list, every operator on the measured system being given as the list of its sector blocks. In its general format the sector map is an ordered list of factors, Alice's system may carry sectors, each \termref{term-measured-factor}{measured factor} declares its own family and no entry of their product is declared, Alice's measurement may be declared in \termref{term-ensemble-form}{ensemble form}, as positive blocks summing to her marginal, and the public discussion is a transcript of messages. Declared operator data are exactly admissible when five conditions hold of the declared numbers, one of them that every binary64 addition forming a completeness sum is exact, and the eight target conditions state when a declaration, with the \termref{term-index-map}{index map} that orders its alphabets and sectors, is well-formed. The consuming engine decides the same five conditions, and the declarations admitted here are a subset of those it admits (\hyperlink{claim-protocol-compiler/rem:engine-admission}{Remark~\ref*{rem:engine-admission}}).}
\sectionimports{\extref{../foundations/contract.pdf}{foundations/def:binary64}, \extref{../foundations/contract.pdf}{foundations/def:exact-admissibility}, \extref{../foundations/contract.pdf}{foundations/lem:exact-accumulation}}

% claim protocol-compiler/def:declaration
\claimsummary{protocol-compiler/def:declaration}{A protocol declaration fixes the systems, sector map, marginal slot, both measurements, public maps, raw-key map and a typed constraint list, operators on the measured system being block lists. The general format adds factors, Alice sectors, an ensemble form, a transcript and scoped entries.}
\claimbookmark{protocol-compiler/def:declaration}{Definition}{def:declaration}{Protocol declaration}
\claimleadin{%
The key-rate engine consumes a protocol declaration. We state it here as this
contract's own definition, with its conditions in
Definitions~\ref{def:exact} and~\ref{def:target}, so that every theorem of
this contract is proved from text inside it. Whether restricting a state
to the \exttermref{../foundations/contract.pdf}{term-block-list}{block lists} of the general format is sound is the subject of
Remark~\ref{rem:block-restriction}. How an entry's form reaches the emitted
declaration is Requirement~\ref{req:pass-through}.
}
\begin{definition}[Protocol declaration]\label{def:declaration}
A \termdef{term-protocol-declaration}{protocol declaration}
is a tuple
\begin{equation}\label{eq:declaration}
  \Decl=\Big(A,\ \tau_A;\ \Qsys,\ \SecSet;\ \{M_x\}_{x\in\Xset};\
  \{F_v\}_{v\in\Vset};\ \Party;\ a_A,\ a_B;\ \phi;\ r,\ \hold;\
  \CList\Big),
\end{equation}
with components as follows. $A$ and $\Qsys$ are finite-dimensional systems,
$\Hs{\Qsys}=\bigoplus_{s\in\SecSet}\Hs{\Qsys_s}$ with a finite index set
$\SecSet$ of sectors, sector projectors $\Pi_s$ and dimensions
$d_s=\dim\Hs{\Qsys_s}$, $d_{\Qsys}=\sum_sd_s$, $d_A=\dim\Hs{A}$. The
\exttermref{../foundations/contract.pdf}{term-sector-map}{sector map} $\Smap$ is the ordered list
$\big((s,d_s)\big)_{s\in\SecSet}$ of block identifiers with their dimensions,
together with its block kind, a label naming what the identifiers mean
(\xref{foundations/def:sector-map}). It is one object, carried with the
declaration beside the slot $\Qsys,\SecSet$ of \eqref{eq:declaration}, and it
is what fixes $\SecSet$, its total order and $(d_s)_s$.
The
\termdef{term-marginal-slot}{marginal slot}
holds either a state $\tau_A\in\Seq{A}$, in \exttermref{../transmitter/contract.pdf}{term-trusted-source-mode}{trusted-source mode}, or
the value $\nomarg$, in \exttermref{../transmitter/contract.pdf}{term-untrusted-source-mode}{untrusted-source mode}.
$\{M_x\}_{x\in\Xset}$ is a POVM on $A$, given as one matrix per $x$, since $A$
carries no sector decomposition. The words ``state'' and ``POVM'' in the two
sentences above name the intended objects, and they are also exactly the
conditions this contract decides: what a declaration is required to satisfy
is D(i)--D(viii) of Definition~\ref{def:target}, every operator predicate of
which is decided \emph{exactly} on the declared \exttermref{../foundations/contract.pdf}{term-binary64-numbers}{binary64 numbers},
together with the summation condition of Definition~\ref{def:exact}, without
which a computed sum would not be the sum it is read as.

\emph{Operators on $\Qsys$ are \exttermref{../foundations/contract.pdf}{term-block-list}{block lists}.} An operator
on $\Qsys$ is given as one array $F^{(s)}\in\Lin{\Qsys_s}$ per
$s\in\SecSet$, against the sector map $\Smap$, and that list is the
operator: it is a block list against $\Smap$, the operator it denotes is
$\bigoplus_{s}F^{(s)}$, and the
declaration holds no array larger than a single sector block. Where
$|\SecSet|=1$, say $\SecSet=\{s_0\}$, that one block \emph{is} the whole of
$\Qsys$ and is held and emitted as such, $d_{s_0}=d_{\Qsys}$, as on the
\termref{term-band-a}{Band-A} route (Theorem~\ref{thm:band-a}\,(i)). Sector block-diagonality is therefore a
property of the \emph{type}: there is no off-sector entry an operator of this
form could carry, and correspondingly nothing to verify. So
$\{F_v\}_{v\in\Vset}$ is a family of such block lists over a finite alphabet,
one list per $v$; $F_v^{(s)}$ denotes its $s$-block. A state $\omega$ on
$A\Qsys$ is likewise a list, one array $\omega^{(s)}\in\Lin{A\Qsys_s}$ per
$s$, and $\omega=\bigoplus_s\omega^{(s)}$ in the same sense. $\Party$ is a
\termdef{term-party-partition}{party partition}:
two total maps $\rd{\mathsf A}:\Xset\times\Vset\to\Oset{\mathsf A}$ and
$\rd{\mathsf B}:\Xset\times\Vset\to\Oset{\mathsf B}$ onto declared finite
alphabets, each a
\termdef{term-readout-map}{readout map},
together with a party label for $A$ and, for each $s$, a declared tensor
factorization of $\Hs{\Qsys_s}$ with a party label per factor.
$a_A:\Oset{\mathsf A}\to\CAset$ and $a_B:\Oset{\mathsf B}\to\CBset$, each an
\termdef{term-announcement-map}{announcement map},
are written jointly
\begin{equation}\label{eq:jointann}
  a(x,v)\coloneqq\big(a_A(\rd{\mathsf A}(x,v)),\,a_B(\rd{\mathsf B}(x,v))\big)
  \in\Kann\coloneqq\CAset\times\CBset ,
\end{equation}
$\phi:\Kann\to\Chatset$ is the acceptance-statistic map,
$r:\Xset\times\Vset\to\Zset$ the raw-key map with
$\Zset=\Zval\sqcup\{\perp\}$ and $\Zval=\{0,\dots,d_Z-1\}$,
$\hold\in\{\mathsf A,\mathsf B\}$ the declared key holder, and
$\CList=\{(\hat O_k,\bowtie_k,\mathrm{rhs}_k)\}_{k=1}^{m}$, $m\ge0$, a finite
constraint list, each entry an operator $\hat O_k$, a relation symbol
$\bowtie_k$ and a right-hand side $\mathrm{rhs}_k$ (which symbols a
well-formed declaration admits is D(vi) of Definition~\ref{def:target}),
with every $\hat O_k$ a block list against the same $\Smap$, Hermitian block
by block.

\emph{The right-hand side of an entry carries a declared tag, and the tag is
part of the datum rather than a property measured of it.} In trusted-source
mode an entry is declared either \termref{term-homogeneous}{homogeneous}, carrying a real scalar
$c_k$, from which the right-hand side it denotes is \emph{constructed} as
$\Theta_k\coloneqq c_k\tau_A$, or \termref{term-general}{general}, carrying a Hermitian matrix
$\Theta_k$ on $A$. In untrusted-source mode only the scalar form exists and
the right-hand side is a real number $\Theta_k=c_k$. So an entry is written
$(\hat O_k,\bowtie_k,c_k)$ exactly when it is declared homogeneous or is
untrusted-mode, and $(\hat O_k,\bowtie_k,\Theta_k)$ when it is declared
general (Definition~\ref{def:homog}). For these two forms the tag
is the type of the right-hand-side datum, a real scalar or a Hermitian
matrix, which is why Definition~\ref{def:homog} calls them the untagged
forms; the third form, party-local, exists only in the general format below
and is declared by its tag (DX6). We write $\kappa=(c_A,c_B)\in\Kann$, a
value of the joint announcement \eqref{eq:jointann}, for an announcement
class and
\begin{equation}\label{eq:chkey}
  \Chkey\coloneqq\big\{\hat c\in\Chatset:\ \exists\,\kappa\in\phi^{-1}(\hat c)
  \text{ with }\varsigma(\kappa)=\mathrm{pass}\big\}
\end{equation}
for the
\termdef{term-key-live}{key-live}
acceptance values. Here $\varsigma$ is the sifting predicate of the
specification's public part (Definition~\ref{def:spec}\,(d)), read as
$\varsigma(\kappa)=\mathrm{fail}$ on every class
$\kappa\notin a(\Xset\times\Vset)$, which no outcome pair realises; on the
realised classes it is, under the hypotheses of
Theorem~\ref{thm:sifting-and-key}, a witness of D(v) of
Definition~\ref{def:target}
(Theorem~\ref{thm:sifting-and-key}\,(ii)). It is \emph{not} a slot of
\eqref{eq:declaration}, so \eqref{eq:chkey} is the set of acceptance values
to which some realised class passing sifting maps, fixed by the
specification and not by the declaration; the key-live set of the consuming
engine, \xref{meat-core/def:declaration}, takes the same convention on the
classes no outcome pair realises.

\emph{The general format.} Everything above is written for the format in
which $A$ carries one sector, one system is measured, Alice's measurement is
declared as the POVM $\{M_x\}$, the public discussion is two messages neither
of which reads the other, and no entry is declared \termref{term-party-local}{party-local}; we
call it the
\termdef{term-basic-format}{basic format}.
It is the case $|\mathbb{T}|=1$, $|\mathbb{F}|=1$, POVM form, $L=2$ of the
following one, \emph{entry for entry} on every component a derivation reads:
at those values every such object defined in (DX1)--(DX6) is the object of
the same name above, and nothing else is declared. On the party labels of the
tensor factors of each sector block, which no derivation of this module reads
and which the consuming engine's laboratory premise reads as naming the
parties' laboratories, the general format is narrower than the basic one. In general a declaration is the
tuple above with its slots read as follows.
\begin{enumerate}[label=(DX\arabic*),leftmargin=3.2em]
\item \emph{(factorised sector map)} The sector map is an ordered list, each
entry of which is a
\termdef{term-factor}{factor}.
Each factor carries a
\termdef{term-system-tag}{system tag},
$A$ or $\Qsys$, a party label, and a sector map of its own, an ordered list of
block identifiers with their dimensions; all factors carry one block kind. At
most one factor is tagged $A$, and where there is one it is the first of the
list; at least one factor is tagged $\Qsys$. The factor tagged $A$ fixes a
finite ordered index set $\mathbb{T}$ and a decomposition
$\Hs{A}=\bigoplus_{t\in\mathbb{T}}\Hs{A_t}$ into subspaces spanned by vectors
of the declared basis of $A$, with projectors $\Pi^{A}_t$ and dimensions
$d_{A_t}$, $d_A=\sum_td_{A_t}$; where no factor is tagged $A$, the system
$A$ is still present and is one Alice sector, $\mathbb{T}=\{t_0\}$ and
$A_{t_0}=A$. A factor tagged $\Qsys$ is a
\termdef{term-measured-factor}{measured factor},
and the measured factors form the ordered set $\mathbb{F}$;
factor $f$ fixes a system $\Qsys_f$ with
$\Hs{\Qsys_f}=\bigoplus_{s_f\in\SecSet^{f}}\Hs{\Qsys_{f,s_f}}$, projectors
$\Pi^{f}_{s_f}$ and dimensions $d_{f,s_f}$, and for $s=(s_f)_{f}$
\begin{equation}\label{eq:productmap}
  \Qsys=\bigotimes_{f\in\mathbb{F}}\Qsys_f,\qquad
  \SecSet=\prod_{f\in\mathbb{F}}\SecSet^{f},\qquad
  \Hs{\Qsys_s}=\bigotimes_{f}\Hs{\Qsys_{f,s_f}},\qquad
  \Pi_s=\bigotimes_f\Pi^{f}_{s_f},\qquad d_s=\prod_fd_{f,s_f},
\end{equation}
all products in the order of $\mathbb{F}$, and $\SecSet$ ordered with the
last factor fastest, a product identifier being the factors' identifiers
joined in that order. Each pair $(t,s)\in\mathbb{T}\times\SecSet$ is a
\termdef{term-joint-sector}{joint sector};
the joint sectors are ordered with $t$ slowest, and
$\Hs{A\Qsys}=\bigoplus_{(t,s)}\Hs{A_t}\otimes\Hs{\Qsys_s}$, the $A$ index
first inside a block. The party partition labels $A$ with the party of the
factor tagged $A$, where there is one, and otherwise with the party label it
declares for $A$, as in the basic format; and it labels every tensor factor of
$\Hs{\Qsys_{f,s_f}}$ with the party of $f$.
\item \emph{(Alice's data and $\omega$ are block lists)} $\tau_A$, every
operator of Alice's measurement slot (DX4) and the right-hand side $\Theta_k$
of a general entry are lists of one array on $\Hs{A_t}$ per
$t\in\mathbb{T}$, and the list is the operator:
$\tau_A=\bigoplus_t\tau_A^{(t)}$, $M_x=\bigoplus_tM_x^{(t)}$,
$\Theta_k=\bigoplus_t\Theta_k^{(t)}$. A state $\omega$ on $A\Qsys$ is a list
of one array $\omega^{(t,s)}\in\Lin{A_t\Qsys_s}$ per joint sector,
$\omega=\bigoplus_{(t,s)}\omega^{(t,s)}$. Block diagonality along
$\mathbb{T}$ is a property of the type: no off-sector entry can be written.
By itself this typing licenses nothing; whether restricting $\omega$ to such
lists is sound is not settled by this definition.
\item \emph{(factorised family)} The \exttermref{../reductions/contract.pdf}{term-outcome-alphabet}{outcome alphabet} is the product
$\Vset=\prod_{f\in\mathbb{F}}\Vset^{f}$ of the factors' ordered alphabets, in
the product order, the last factor fastest. What is declared is one family of
block lists $\{F^{f}_{v_f}\}_{v_f\in\Vset^{f}}$ per measured factor, against
that factor's own sector map, together with the rule
\begin{equation}\label{eq:productfamily}
  F_v\coloneqq\bigotimes_{f\in\mathbb{F}}F^{f}_{v_f},\qquad
  F^{(s)}_v=\bigotimes_{f\in\mathbb{F}}F^{f,(s_f)}_{v_f}
  \qquad\big(v=(v_f)_f,\ s=(s_f)_f\big).
\end{equation}
No entry of a product is declared. That $\{F_v\}_{v\in\Vset}$ is then a
POVM on $\Qsys$, block diagonal for the product sector map with the blocks of
\eqref{eq:productfamily}, follows for the exact numbers the factors declare:
a tensor product of Hermitian positive semidefinite operators is Hermitian
and positive semidefinite,
$\sum_{v}\bigotimes_fF^{f}_{v_f}=\bigotimes_f\sum_{v_f}F^{f}_{v_f}
=\idon{\Qsys}$ by multilinearity, and the tensor product distributes over the
direct sums $F^{f}_{v_f}=\bigoplus_{s_f}F^{f,(s_f)}_{v_f}$.
\item \emph{(the two forms of Alice's measurement slot)} In
\termdef{term-povm-form}{POVM form}
the slot holds $\{M_x\}_{x\in\Xset}$. In
\termdef{term-ensemble-form}{ensemble form},
which exists in trusted-source mode only, it holds a family
$\{N_x\}_{x\in\Xset}$ of block lists along $\mathbb{T}$, every block
Hermitian and positive semidefinite, with
\begin{equation}\label{eq:ensembleform}
  \sum_{x\in\Xset}N^{(t)}_x=\tau^{(t)}_A
  \quad\text{and}\quad\tau^{(t)}_A\ne0
  \qquad\text{for every }t\in\mathbb{T},
\end{equation}
the sum formed entrywise in the \exttermref{../foundations/contract.pdf}{term-declared-order}{declared order} of $\Xset$ and
rounding nowhere, and the second condition a comparison of declared entries with the
datum $0$. An ensemble-form declaration \emph{denotes} the POVM-form
declaration with the same remaining slots and
\begin{equation}\label{eq:ensemblepovm}
  M^{(t)}_x\coloneqq\big(\tau^{(t)}_A\big)^{-1/2}N^{(t)}_x
  \big(\tau^{(t)}_A\big)^{-1/2}
  +\frac{\tr N^{(t)}_x}{\tr\tau^{(t)}_A}\,\big(\idon{A_t}-P_t\big),
\end{equation}
where $P_t$ is the projector onto the support of $\tau^{(t)}_A$ and the
inverse square root is taken on that support and is zero off it. The
operators \eqref{eq:ensemblepovm} are defined and in general irrational; every statement made of an
ensemble-form declaration is made of the declaration it denotes. For a
POVM-form declaration in trusted-source mode put
$N_x\coloneqq\tau_A^{1/2}M_x\tau_A^{1/2}$; for the declaration an ensemble
form denotes, this returns the declared $N_x$.
\item \emph{(transcript protocol)} The announcement slot holds a
\termdef{term-transcript-protocol}{transcript protocol},
an ordered list of $L\ge1$ entries, each a
\termdef{term-message}{message},
message $m$ carrying a sender $\pi_m\in\{\mathsf A,\mathsf B\}$, a finite
alphabet $\mathcal{C}_m$ and a total table
$a_m:\Oset{\pi_m}\times\mathcal{C}_1\times\dots\times\mathcal{C}_{m-1}
\to\mathcal{C}_m$. The joint announcement function is \emph{derived} from the
tables by composition,
\begin{equation}\label{eq:transcript}
  \kappa_m(x,v)\coloneqq a_m\big(\rd{\pi_m}(x,v),\,\kappa_1(x,v),\dots,
  \kappa_{m-1}(x,v)\big),\qquad
  a(x,v)\coloneqq\big(\kappa_m(x,v)\big)_{m=1}^{L}\in
  \Kann\coloneqq\prod_{m=1}^{L}\mathcal{C}_m ,
\end{equation}
and the acceptance-statistic map $\phi$, the sifting predicate $\varsigma$ of
\eqref{eq:chkey} and the holder's key map $r_{\hold}$ of D(viii) of
Definition~\ref{def:target} are maps on this $\Kann$. The two
maps $a_A$, $a_B$ of the basic format are the transcript protocol with $L=2$,
$\pi=(\mathsf A,\mathsf B)$, $\mathcal{C}_1=\CAset$, $\mathcal{C}_2=\CBset$,
$a_1=a_A$ and $a_2(o,\kappa_1)=a_B(o)$. A message is a function of its
sender's readout and of the earlier messages of the same round, and of
nothing else. No measurement of a round depends on a message of that round,
the measurements being the fixed families of (DX3) and (DX4); processing
across rounds is not a round message.
\item \emph{(party-scoped entries)} Each entry of $\CList$ carries the label
of one measured factor, its
\termdef{term-scope}{scope}
$f_k$, and its operator $\hat O_k$ is a block list against that factor's own
sector map, Hermitian block by block. Unless the entry is declared
\termref{term-party-local}{party-local} it denotes the
\termdef{term-lifted-operator}{lifted operator}
$\idon{}\otimes\hat O_k\otimes\idon{}$ on $\Qsys$, with $\hat O_k$ on the
tensor factor $\Qsys_{f_k}$, whose $s$-block is
$\idon{}\otimes\hat O_k^{(s_{f_k})}\otimes\idon{}$, and it is read in the forms stated
above for that lifted operator. With one measured factor the scope is that
factor and the lifted operator is $\hat O_k$ itself. An entry may instead be
declared party-local, in either source mode, carrying a real $c_k$; its
reading is the third form of Definition~\ref{def:homog}.
\end{enumerate}

\emph{What this module decides of the general format} is
Definition~\ref{def:exact} for the operator data and D(i)--D(viii) of
Definition~\ref{def:target} for the rest; that the family denoted by an
ensemble form is a POVM is Theorem~\ref{thm:sifting-and-key}\,(v).
\end{definition}
\claimmeta{\metalabel{Used by}\mref{Rem.}{rem:party-label-restriction}, \mref{Def.}{def:exact}, \mref{Def.}{def:target}, \mref{Def.}{def:feasible}, \mref{Lem.}{lem:hermitian-constraint-values}, \mref{Rem.}{rem:wellposed}, \mref{Def.}{def:homog}, \mref{Def.}{def:spec}, \mref{Def.}{def:compile}, \mref{Thm.}{thm:measurements-and-statistics}, \mref{Def.}{def:weighted}, \mref{Rem.}{rem:alphabet}, \mref{Thm.}{thm:sifting-and-key}, \mref{Thm.}{thm:feasibility}, \mref{Def.}{def:validator}, \mref{Thm.}{thm:validator-soundness}, \mref{Def.}{def:chaincond}, \mref{Thm.}{thm:chain-consumption}, \mref{Def.}{def:eligible}, \mref{Thm.}{thm:gates}, \mref{Ass.}{ass:coin}, \mref{Req.}{req:declaration-keys}, \mref{Req.}{req:operator-form}, \mref{Req.}{req:derived-objects}, \mref{Req.}{req:da1-source-identity}, \mref{Req.}{req:reference-companion-compiled}, \extref{../meat-core/contract.pdf}{meat-core/def:declaration}
\metasep newly derived
\metasep implemented
\metasep \metalabel{Code}\coderef{https://github.com/fraunhofersg/certiqs-key/blob/v0.1.0/src/certiqs_key/protocol_compiler/declaration.py}{certiqs_key.protocol_compiler.declaration:CandidateDeclaration}}

% claim protocol-compiler/rem:party-label-restriction
\claimsummary{protocol-compiler/rem:party-label-restriction}{Deriving the party labels from the factor list is narrower than the basic format, which may label the tensor factors of one sector differently. Only the premises placing the laboratories read the labels; two parties inside one measured system need two measured factors.}
\claimbookmark{protocol-compiler/rem:party-label-restriction}{Remark}{rem:party-label-restriction}{Party labels in the general format}
\begin{remark}[Party labels in the general format]\label{rem:party-label-restriction}
The party labels that (DX1) of Definition~\ref{def:declaration} derives from
the factor list are narrower than the party partition of the basic format at
one measured factor: there a declared tensor factorization of $\Hs{\Qsys_s}$
may carry different party labels on different tensor factors, while in the
general format every tensor factor of one measured factor's block carries
that factor's party. The narrowing is admitted and not repaired. No
derivation of the compiler or of the consuming engine reads those labels ---
what the derivations consume is the readout maps $\rd{\mathsf A}$,
$\rd{\mathsf B}$ --- and correspondingly no check compares a declared label
with a derived one. The labels are read by the premises that place the
parties' laboratories: the consuming engine's laboratory premise
(\xref{meat-core/ass:lab}) takes them, with the party label of $A$, as naming
the parties' laboratories, and the source licence
(\xref{transmitter/ass:src}) reads the party label of $A$; so the narrowing restricts which laboratory boundaries a
general-format declaration can state, and a declaration that wants two parties inside one
\exttermref{../reductions/contract.pdf}{term-measured-system}{measured system} declares two measured factors instead. So the general format
coincides with the basic format entry for entry on every component a
derivation reads, and on the party labels it is a restriction.
\end{remark}
\claimmeta{\metalabel{Uses}\mref{Def.}{def:declaration}
\metasep newly derived
\metasep accepted}

% claim protocol-compiler/def:exact
\claimsummary{protocol-compiler/def:exact}{Declared binary64 operator data are admitted by sector-map agreement, Hermiticity, positivity, completeness and a completeness sum that rounds nowhere, each decided exactly on the declared numbers. No tolerance is read and no ideal operator is inferred.}
\claimbookmark{protocol-compiler/def:exact}{Definition}{def:exact}{Exact admissibility of declared operator data}
\claimleadin{%
Of the five conditions of the definition, four are conditions on an operator
family and one, (S), is a condition on the binary64 additions that form a
completeness sum. Binary64 addition is not associative, so each such sum comes
with the order in which it is accumulated
(\xref{foundations/rem:summation-order}).

Deciding every admission predicate on \exttermref{../foundations/contract.pdf}{term-declared-data}{declared data} exactly, on the
\exttermref{../foundations/contract.pdf}{term-declared-value}{declared values}, is what \xref{foundations/req:consumer-decides-exactly} requires of
every consumer.
The conditions the definition states on Alice's data, on $\{M_x\}$ and on
$\tau_A$, are producer obligations of this module and not only consumer
checks, since the consuming engine decides them exactly and refuses
otherwise, so a declaration this module emitted without deciding them would
be refused downstream on data this module had passed. For $\{M_x\}$ and
for the trace sum of $\tau_A$ the obligation is
\xref{foundations/req:exact-emission}\,(ii)--(v), and for the Hermiticity,
positivity and unit trace of $\tau_A$ it is
\xref{foundations/req:boundary-admission}\,(ii). Each is decided here, before
emission, and refused with a named reason otherwise
(Requirement~\ref{req:validator-refusals}).
}
\begin{definition}[Exact admissibility of declared operator data]
\label{def:exact}
Declared operator data are \exttermref{../foundations/contract.pdf}{term-binary64-numbers}{binary64 numbers}, each the dyadic rational
it denotes and not an approximation of anything
(\xref{foundations/def:binary64}\,(iii)), and every predicate this module
decides on them is decided \emph{exactly}, on those numbers. This definition lists
those predicates, together with the one condition of the list that is a
condition on the arithmetic that forms a sum rather than a predicate on an
operator.

\emph{The declared summation order.} The \termref{term-compilation-map}{compilation map} emits a
total order on every declared alphabet and on $\SecSet$: the
\termref{term-index-map}{label-to-index map} $\idxm$, which numbers the letters of
each declared alphabet, and the sectors, consecutively from $0$
(Definition~\ref{def:compile}, Theorem~\ref{thm:typing}); the
\exttermref{../foundations/contract.pdf}{term-declared-order}{declared order} of an alphabet is its $\idxm$ order. For a finite
family $\{a_j\}_{j\in J}$ of $n\times n$ matrices indexed by a declared
alphabet $J$ with $\idxm$ order $j_1\prec\dots\prec j_{|J|}$, the
\exttermref{../foundations/contract.pdf}{term-sum-in-the-declared-order}{sum in the declared order} is accumulated as the
\exttermref{../foundations/contract.pdf}{term-computed-sequential-sum}{computed sequential sum} (\xref{foundations/def:declared-order})
\begin{equation}\label{eq:declorder}
  \Sigma_0\coloneqq0,\qquad
  \Sigma_i\coloneqq\Sigma_{i-1}+a_{j_i}\ \ (1\le i\le|J|),\qquad
  \Sigma\coloneqq\Sigma_{|J|},
\end{equation}
formed left to right, entrywise, in binary64, with no re-ordering and no
re-association; the trace of a matrix is likewise accumulated over its row
index in increasing order, which is the order of the emitted array.

\emph{The five conditions.} Let $\{a_j\}_{j\in J}$ be a finite family of
block lists against a sector map $\Smap$, one array $a_j^{(s)}$ on
$\Hs{\Qsys_s}$ per $s\in\SecSet$. The family is \exttermref{../foundations/contract.pdf}{term-exactly-admissible}{exactly admissible}
for $\Smap$ when
\begin{itemize}[leftmargin=2.6em]
\item[(T)] \emph{(sector map)} the sector map the supplied family carries is
  $\Smap$, the map this module holds: the same block kind, the same block
  identifiers in the same order, the same dimensions, so that
  $a_j=\bigoplus_sa_j^{(s)}$ holds \emph{by construction}. This is a type
  condition: it is decided, never measured, and an off-block entry is not
  small but inexpressible;
\item[(S)] \emph{(the sum \exttermref{../foundations/contract.pdf}{term-rounds-nowhere}{rounds nowhere})} writing $\Sigma^{(s)}$ for
  the total $\Sigma$ that \eqref{eq:declorder} produces on the block family
  $\{a_j^{(s)}\}_j$: for each $s$, every one of the $|J|$ entrywise binary64
  additions forming it is exact, so that the computed $\Sigma^{(s)}$ \emph{is}
  the real sum $\sum_ja_j^{(s)}$;
\item[(H)] \emph{(Hermiticity)} $a_j^{(s)}=a_j^{(s)\dagger}$ for every $(j,s)$;
\item[(P)] \emph{(positivity)} $a_j^{(s)}\succeq0$ for every $(j,s)$; the
  matrix being Hermitian by (H), this says that its spectrum is non-negative;
\item[(C)] \emph{(completeness)} $\Sigma^{(s)}=\idon{\Qsys_s}$ for every $s$.
\end{itemize}
Each of (H), (P), (C) is an \emph{identity or an order relation between
exactly represented numbers} and is decided as such: (H) and (C) entry by
entry, (P) by exact symmetric elimination on the block with the zero-pivot
row test (a negative pivot fails, and a zero pivot passes only when the rest
of its row is zero), cubic in $d_s$. Together the five conditions are exact
admissibility for $\Smap$ with constant $1$, that is with every completeness
sum equal to the identity, in the sense of
\xref{foundations/def:exact-admissibility}, the declared order of $J$ being
its $\idxm$ order: (T), (H), (P) and (S) are the conditions of that
definition bearing the same names, and (C), which reads the computed sum, is
its completeness identity, because under (S) the computed sum $\Sigma^{(s)}$
is the exact sum $\sum_ja_j^{(s)}$
(\xref{foundations/lem:exact-accumulation}).

\emph{The same conditions on Alice's data, where there is one block and no
sector map.} $A$ carries no sector decomposition, so a family on $A$ is one
array per index and (T) is empty for it. $\{M_x\}_{x\in\Xset}$ is
\emph{exactly admissible} when (S), (H), (P) and (C) hold of it on $A$ with
$\idon{A}$ on the right of (C); the declared marginal $\tau_A$ satisfies its
own three conditions when $\tau_A=\tau_A^{\dagger}$, $\tau_A\succeq0$ and
$\tr\tau_A=1$ with the trace accumulated in the row order \emph{without
rounding}, the trace being a sum, read below as an identity of real numbers
exactly as a completeness sum is.

\emph{What is not on this list.} A strict inequality on a declared scalar,
such as the condition of Definition~\ref{def:weighted} that every
public-choice weight be positive, is a domain condition on a declared real
number, decided by its sign and not by any of the five. A relation symbol, a
declared tag and the type of a right-hand side are likewise decided and never
measured (D(vi)). And nothing in this definition is a \emph{norm}: no
quantity here is compared against a tolerance window, and no
\exttermref{../foundations/contract.pdf}{term-defects}{defect}, the entrywise size of the difference a condition
names, is computed.

\emph{The general format.} For a declaration in the general format of
Definition~\ref{def:declaration}\,(DX1)--(DX6) the same conditions are read
as follows. \emph{Per measured factor:} for every $f\in\mathbb{F}$ the
family $\{F^{f}_{v_f}\}_{v_f\in\Vset^{f}}$ is exactly admissible for the
sector map of $f$ --- (T) against that factor's map, and (S), (H), (P), (C)
on that factor's blocks in that factor's declared order. No condition is
decided on a product block \eqref{eq:productfamily}, none being a datum; that
the product family then satisfies (H), (P) and (C) on $\Qsys$ as identities
of real numbers is shown in (DX3). An entry's operator $\hat O_k$ is decided
(T) and (H) against the map of its scope. \emph{Alice's data, block by
block:} with a factor tagged $A$, condition (T) is not empty for Alice's
data: $\tau_A$, the operators of her measurement slot and a general
$\Theta_k$ are block lists against that factor's map. In POVM form
$\{M^{(t)}_x\}_x$ satisfies (S), (H), (P) and (C) on every $A_t$ with
$\idon{A_t}$ on the right of (C); $\tau_A$ satisfies (H) and (P) block by
block and $\sum_t\tr\tau^{(t)}_A=1$, the traces accumulated in the row order
and then over the blocks in their declared order, rounding nowhere.
\emph{Ensemble form:} the family $\{N^{(t)}_x\}_x$ satisfies (S), (H) and
(P) on every $A_t$, and in place of (C) the identity
$\Sigma^{(t)}=\tau^{(t)}_A$ of \eqref{eq:ensembleform}, entry by entry,
between the declared-order sum and the declared block of $\tau_A$; and
$\tau^{(t)}_A\ne0$ for every $t$, a comparison of declared entries with the
datum $0$. These are all the conditions of the form: nothing is decided of
the denoted $\{M_x\}$ of \eqref{eq:ensemblepovm}, its being a POVM a theorem
(Theorem~\ref{thm:sifting-and-key}\,(v)). At one Alice sector and one measured factor in
POVM form these readings are the paragraphs above.
\end{definition}
\claimmeta{\metalabel{Uses}\mref{Def.}{def:declaration}, \extref{../foundations/contract.pdf}{foundations/def:binary64}, \extref{../foundations/contract.pdf}{foundations/def:exact-admissibility}, \extref{../foundations/contract.pdf}{foundations/lem:exact-accumulation}
\metasep \metalabel{Used by}\mref{Rem.}{rem:summation-condition}, \mref{Def.}{def:target}, \mref{Rem.}{rem:engine-admission}, \mref{Def.}{def:feasible}, \mref{Def.}{def:well-formed-spec}, \mref{Def.}{def:compile}, \mref{Thm.}{thm:measurements-and-statistics}, \mref{Def.}{def:weighted}, \mref{Rem.}{rem:condfamily}, \mref{Ex.}{ex:condfamily}, \mref{Thm.}{thm:sifting-and-key}, \mref{Thm.}{thm:feasibility}, \mref{Def.}{def:validator}, \mref{Thm.}{thm:validator-soundness}, \mref{Def.}{def:chaincond}, \mref{Thm.}{thm:chain-consumption}, \mref{Thm.}{thm:band-a}, \mref{Def.}{def:eligible}, \mref{Thm.}{thm:gates}, \mref{Prop.}{prop:conservative}, \mref{Req.}{req:derived-objects}, \mref{Req.}{req:no-window}, \mref{Req.}{req:physical-family-decision}
\metasep newly derived
\metasep implemented (effective: accepted)
\metasep \metalabel{Code}\coderef{https://github.com/fraunhofersg/certiqs-key/blob/v0.1.0/src/certiqs_key/foundations/exact.py}{certiqs_key.foundations.exact:sum_in_declared_order}}

% claim protocol-compiler/rem:summation-condition
\claimsummary{protocol-compiler/rem:summation-condition}{Completeness alone says that the computed floating-point sum is the identity; the summation condition is what makes it an identity of the declared numbers. A family of exactly representable entries may therefore still be refused, since the condition is on the additions.}
\claimbookmark{protocol-compiler/rem:summation-condition}{Remark}{rem:summation-condition}{The summation condition beside completeness}
\begin{remark}[The summation condition beside completeness]\label{rem:summation-condition}
The total $\Sigma^{(s)}$ of Definition~\ref{def:exact} is by
\eqref{eq:declorder} a \emph{floating-point} sum. So condition (C) alone says
that the \emph{computed} sum is the identity, and \emph{not} that
$\sum_ja_j=\idon{\Qsys}$ holds of the declared numbers, which is what every
derivation of this contract that consumes completeness reads it as, and what
a consumer of the emitted declaration reads it as. Condition (S) is exactly
what makes the two the same statement
(\xref{foundations/lem:exact-accumulation}), and it is the reason a family
whose entries are all exactly representable may still be refused: the
condition is on the \emph{additions}, not on the operands.
\xref{foundations/rem:computed-completeness} exhibits a family satisfying
(H), (P) and completeness of the computed sum whose true sum is not the
identity.
\end{remark}
\claimmeta{\metalabel{Uses}\mref{Def.}{def:exact}
\metasep newly derived
\metasep accepted}

% claim protocol-compiler/def:target
\claimsummary{protocol-compiler/def:target}{A declaration with its index map is well-formed when D(i) to D(viii) hold: agreeing finite data, one sector order, an admissible marginal slot, exact measurement families, raw-key values in the key alphabet or none, announcement-determined sifting, typed entries, feasibility and party computability.}
\claimbookmark{protocol-compiler/def:target}{Definition}{def:target}{Target conditions of a well-formed declaration}
\claimleadin{%
The sector order a consumer compares against its own sector map and the
summation order it reads off the emitted \termref{term-index-map}{index map} must be one order,
so D(i) states their agreement as a condition on the declaration together
with its index map.
}
\begin{definition}[Target conditions of a well-formed declaration]\label{def:target}
A declaration \eqref{eq:declaration}, \emph{together with the label-to-index
map $\idxm$ emitted alongside it} (Definition~\ref{def:compile}), is a
\termdef{term-well-formed-declaration}{well-formed declaration}
when the following hold, every operator predicate below being decided
\emph{exactly} in the sense of Definition~\ref{def:exact} and every
combinatorial condition exactly as well:
\begin{enumerate}[label=D(\roman*),leftmargin=3.2em]
\item all systems are finite-dimensional, all alphabets are finite and
  non-empty, all declared maps are total, and \emph{the sector map agrees}:
  there is one sector map $\Smap$; every declared $\Qsys$-operator is a
  block list against it, with the same block kind, the same
  identifiers in the same order and the same dimensions; and
  $d_{\Qsys}=\sum_{s\in\SecSet}d_s$. That much is condition~(T) of
  Definition~\ref{def:exact}, stated once for the declaration. \emph{And the
  order declared on the sectors is one order and not two}: the restriction of
  $\idxm$ to $\SecSet$ is the total order $\Smap$ fixes. This last clause is a
  condition on the pair $(\Decl,\idxm)$ and not on a slot of
  \eqref{eq:declaration};
\item the marginal slot holds $\tau_A$ with $\tau_A=\tau_A^{\dagger}$,
  $\tau_A\succeq0$ and $\tr\tau_A=1$, the trace summed in the row order and
  \exttermref{../foundations/contract.pdf}{term-rounds-nowhere}{rounding nowhere} --- so that $\tau_A\in\Seq{A}$
  holds of the declared numbers, and a slot holding the zero matrix is
  inadmissible, its trace being $0$ --- or the value $\nomarg$;
\item \emph{(typing)} $\{F_v\}$ and every $\hat O_k$ are block lists against
  $\Smap$ in the sense of Definition~\ref{def:declaration}, so that
  $\Pi_sF_v\Pi_{s'}=0$ and $\Pi_s\hat O_k\Pi_{s'}=0$ for $s\ne s'$ hold
  \emph{by construction} and are not conditions on data; and
  \emph{(five predicates)} $M_x\succeq0$ for each $x$ and $\{M_x\}$ complete on
  $A$; $F_v^{(s)}\succeq0$ for each $(v,s)$ and $\{F_v^{(s)}\}_v$ complete on
  $\Qsys_s$ for each $s$; and $\hat O_k^{(s)}$ Hermitian for each $(k,s)$,
  each decided exactly, each positivity predicate being the two decisions
  (H) and (P) of Definition~\ref{def:exact} on the same block, and each
  completeness predicate carrying with it the condition~(S) that its
  declared-order sum round nowhere. This is where Hermiticity of Alice's and
  Bob's emitted families enters the target conditions, and it is the only
  place it does;
\item $r$ maps into $\Zval\sqcup\{\perp\}$;
\item \emph{sifting is announcement-determined}: there is
  $\varsigma:\Kann\to\{\mathrm{pass},\mathrm{fail}\}$ with
  $r(x,v)=\perp\iff\varsigma(a(x,v))=\mathrm{fail}$;
\item in trusted-source mode $\bowtie_k\in\{\preceq,=,\succeq\}$ and
  each entry carries one of the two declared forms of
  Definition~\ref{def:homog}, \termref{term-homogeneous}{homogeneous} with $c_k\in\mathbb{R}$ or
  \termref{term-general}{general} with $\Theta_k$ Hermitian on $A$; in
  untrusted-source mode $\bowtie_k\in\{\le,=,\ge\}$ with
  $\Theta_k=c_k\in\mathbb{R}$. The relation symbol and the type of the
  right-hand side, which is what fixes the tag (Definition~\ref{def:homog}),
  are combinatorial and are read exactly; the Hermiticity of a general
  $\Theta_k$ is the exact predicate (H) on one matrix;
\item the \termref{term-feasible-set}{feasible set} of Definition~\ref{def:feasible} is non-empty,
  $\Feas\ne\varnothing$;
\item \emph{party computability}: each party announces a function of its own
  readout ($a_A$ is defined on $\Oset{\mathsf A}$ and $a_B$ on
  $\Oset{\mathsf B}$, jointly \eqref{eq:jointann}) and the declared holder
  computes the raw key from its own readout and the announcements: there is a
  total $r_{\hold}:\Oset{\hold}\times\Kann\to\Zset$ with
  \begin{equation}\label{eq:partycomp}
    r(x,v)=r_{\hold}\big(\rd{\hold}(x,v),\,a(x,v)\big)
    \qquad\text{for all }(x,v)\in\Xset\times\Vset .
  \end{equation}
\end{enumerate}
The
\termdef{term-default-party-partition}{default party partition}
is $\rd{\mathsf A}(x,v)=x$, $\rd{\mathsf B}(x,v)=v$, with $A$ labelled
$\mathsf A$ and every $\Qsys_s$ labelled $\mathsf B$.

\emph{The general format.} For a declaration in the general format of
Definition~\ref{def:declaration}\,(DX1)--(DX6) the conditions read as
follows, and in the \termref{term-basic-format}{basic format} --- one Alice sector, one measured
factor, POVM form, two messages neither reading the other and no party-local
entry --- they are the eight above.
D(i): the sector map is one factorised map (at most one factor tagged $A$,
and then the first, at least one tagged $\Qsys$, one block kind), every
declared family on a measured factor and every entry's operator is a block
list against the map of its factor, Alice's data against the factor tagged
$A$, $d_{\Qsys}=\prod_f\sum_{s_f}d_{f,s_f}$, and the restriction of $\idxm$
to each factor's sectors is the order that factor's map fixes, the product
order being derived and not declared. D(ii): $\tau_A$ is a block list along
$\mathbb{T}$ with every block Hermitian and positive semidefinite and total
trace $1$, summed without rounding; in ensemble form additionally every
block is non-zero. D(iii): the typing and the predicates of
Definition~\ref{def:exact} in their general-format reading, per measured
factor and per Alice sector, the ensemble-form conditions in place of those
on $\{M_x\}$ where that form is declared; ensemble form is refused in
untrusted-source mode. D(vi): an entry carries its scope, which names a
measured factor of the map, and one of the forms of
Definition~\ref{def:homog}, the \termref{term-party-local}{party-local} form included, in either
mode, with a real $c_k$; a general $\Theta_k$ is a block list along
$\mathbb{T}$, Hermitian block by block. \emph{And the relation symbol is
read by form and not by mode alone}: an \emph{untagged} entry carries
$\bowtie_k\in\{\preceq,=,\succeq\}$ in trusted-source mode and
$\bowtie_k\in\{\le,=,\ge\}$ in untrusted-source mode, exactly as stated for
D(vi), while an entry declared party-local carries
$\bowtie_k\in\{\preceq,=,\succeq\}$ \emph{in either mode}, its reading being
the operator relation of the last line of \eqref{eq:feas-general}
(Definition~\ref{def:homog}, the relation-symbol clause of the \exttermref{../reductions/contract.pdf}{term-party-local-form}{party-local
form}); where every row of that entry is $1\times1$ (every Alice block, and
every sector block of every measured factor other than the entry's scope, of
dimension one, $d_A=1$ at one measured factor included) that relation is
between $1\times1$ matrices and the three operator symbols are the three
scalar symbols, which is the degenerate case. Where no entry is party-local,
as in every declaration of the \termref{term-basic-format}{basic format}, this clause is D(vi) as
stated. D(viii): the announcement slot is a \termref{term-transcript-protocol}{transcript protocol} (DX5),
every table $a_m$ being total on
$\Oset{\pi_m}\times\mathcal{C}_1\times\dots\times\mathcal{C}_{m-1}$, so that
a message reads its sender's readout and earlier messages \emph{by its
signature}; the joint function $a$ is the composition \eqref{eq:transcript},
and \eqref{eq:partycomp} holds with it. D(iv), D(v) and D(vii) keep their
wording, with $a$ the derived function and $\Feas$ the set of
\eqref{eq:feas-general}.
\end{definition}
\claimmeta{\metalabel{Uses}\mref{Def.}{def:declaration}, \mref{Def.}{def:exact}
\metasep \metalabel{Used by}\mref{Rem.}{rem:engine-admission}, \mref{Lem.}{lem:hermitian-constraint-values}, \mref{Def.}{def:compile}, \mref{Thm.}{thm:typing}, \mref{Thm.}{thm:measurements-and-statistics}, \mref{Thm.}{thm:sifting-and-key}, \mref{Rem.}{rem:vacuous}, \mref{Def.}{def:validator}, \mref{Thm.}{thm:validator-soundness}, \mref{Thm.}{thm:chain-consumption}
\metasep newly derived
\metasep accepted}

% claim protocol-compiler/rem:engine-admission
\claimsummary{protocol-compiler/rem:engine-admission}{The consuming engine decides the same five exact conditions as this module, and decides them alike. The declarations admitted here are a subset of those the engine admits: D(i) and D(iii) are stronger here and check V7 is one-sided.}
\claimbookmark{protocol-compiler/rem:engine-admission}{Remark}{rem:engine-admission}{Relation to the consuming engine's admission conditions}
\begin{remark}[Relation to the consuming engine's admission conditions]\label{rem:engine-admission}
The five conditions of Definition~\ref{def:exact} are also the consuming
engine's admission conditions, which that engine decides exactly, allowing
no nonzero \exttermref{../foundations/contract.pdf}{term-defects}{defect}
(\xref{foundations/req:consumer-decides-exactly}), so the two modules
decide \emph{those five conditions} alike. \emph{The agreement is on the five
conditions and not on the declarations they are conditions of}: this module
additionally imposes D(i)'s clause on $\idxm|_{\SecSet}$
(Definition~\ref{def:target}), for which that engine's definition of an
admissible declaration (\xref{meat-core/def:admissible-declaration}) has no
counterpart, and D(iii)'s Hermiticity clause, which that definition leaves to
the typing of the tuple;
and its feasibility check V7 (Definition~\ref{def:validator}) is one-sided:
it may refuse a declaration whose feasible set is non-empty, for want of a
verified witness. So the declarations admitted
here are a \emph{subset} of those that engine admits, and where the two
differ this module refuses declarations that engine would admit, never the
reverse.
\end{remark}
\claimmeta{\metalabel{Uses}\mref{Def.}{def:exact}, \mref{Def.}{def:target}
\metasep newly derived
\metasep accepted}

\section{Feasible states and licensed constraints}\label{group-protocol-compiler/def:feasible}
\sectionintro{The \termref{term-feasible-set}{feasible set} holds the normalised positive block-list states that meet every constraint entry and, in trusted-source mode, the declared marginal, each condition read exactly; Hermiticity of the constraint operators and of the state, with that of the declared right-hand sides, is what makes the two sides of each constraint condition comparable (\hyperlink{claim-protocol-compiler/lem:hermitian-constraint-values}{Lemma~\ref*{lem:hermitian-constraint-values}}). Restricting the states to sector-block-diagonal ones is a restriction in the unsafe direction that the operator type carries and the consuming engine proves sound (\hyperlink{claim-protocol-compiler/rem:block-restriction}{Remark~\ref*{rem:block-restriction}}). A constraint list is licensed for its source mode when every entry is declared in the form the mode admits or carries the party-local tag, and licensing reads the declared tag and compares no numbers. A party-local entry, whose condition holds against the state's marginal on the systems other than its own factor, is sound only under a certificate owed by the producer of the reduction chain, which the attack-set licence of the list attests (\hyperlink{claim-protocol-compiler/rem:party-local-soundness}{Remark~\ref*{rem:party-local-soundness}}). The example writes out a declaration for an untrusted source serving two measuring parties, with a party-local entry.}

% claim protocol-compiler/def:feasible
\claimsummary{protocol-compiler/def:feasible}{The feasible set holds the normalised positive block-list states that meet the declared marginal and every constraint entry, each read exactly. The source mode decides whether an entry is an operator or a scalar condition, and the set carries no parameter and admits no window.}
\claimbookmark{protocol-compiler/def:feasible}{Definition}{def:feasible}{Feasible set}
\begin{definition}[Feasible set]\label{def:feasible}
For a declaration $\Decl$, the
\termdef{term-feasible-set}{feasible set}
$\Feas$ is the set of $\omega$ presented as a
\exttermref{../foundations/contract.pdf}{term-block-list}{sector-block list} on $A\Qsys$ --- $\omega=\bigoplus_s\omega^{(s)}$ with $\omega^{(s)}\in\Lin{A\Qsys_s}$,
so that $(\idon{A}\otimes\Pi_s)\omega(\idon{A}\otimes\Pi_{s'})=0$ for
$s\ne s'$ holds by the presentation and is not a condition --- such that
\begin{equation}\label{eq:feas}
  \omega^{(s)}\succeq0\ \ \forall s,\qquad \tr\omega=1,\qquad
  \omega\models\Decl,\qquad \omega\models\CList ,
\end{equation}
each read \emph{exactly}, with $\tr\omega=\sum_s\tr\omega^{(s)}$ and with
$\omega^{(s)}\succeq0$ carrying, as in Definition~\ref{def:exact},
Hermiticity of the block. Here $\omega\models\Decl$ is the marginal condition,
$\omega_A=\tau_A$ in \exttermref{../transmitter/contract.pdf}{term-trusted-source-mode}{trusted-source mode}, with
$\omega_A=\sum_s\tr_{\Qsys_s}\omega^{(s)}$, and no condition beyond
$\tr\omega=1$ in \exttermref{../transmitter/contract.pdf}{term-untrusted-source-mode}{untrusted-source mode}; and $\omega\models\CList$ reads
\begin{equation}\label{eq:feas-modes}
\begin{aligned}
  \text{trusted:}\quad&
  \tr_{\Qsys}\!\big[(\idon{A}\otimes\hat O_k)\,\omega\big]
  =\sum_s\tr_{\Qsys_s}\!\big[(\idon{A}\otimes\hat O^{(s)}_k)\,\omega^{(s)}\big]
  \ \bowtie_k\ \Theta_k
  \quad\forall k;\\
  \text{untrusted:}\quad&
  \tr\!\big[(\idon{A}\otimes\hat O_k)\,\omega\big]
  =\sum_s\tr\!\big[(\idon{A}\otimes\hat O^{(s)}_k)\,\omega^{(s)}\big]
  \ \bowtie_k\ c_k
  \quad\forall k ,
\end{aligned}
\end{equation}
where in the trusted line $\Theta_k$ is the \emph{constructed} right-hand side
$c_k\tau_A$ for an entry declared \termref{term-homogeneous}{homogeneous} and the \emph{declared}
Hermitian matrix for one declared \termref{term-general}{general}
(Definition~\ref{def:homog}). The set carries no parameter, admits no
tolerance window and relaxes no conjunct; in particular positivity is never
relaxed.

\emph{The general format.} For a declaration in the general format of
Definition~\ref{def:declaration}\,(DX1)--(DX6), $\Feas$ is the set of
$\omega$ presented as a list over the
\termref{term-joint-sector}{joint sectors},
$\omega=\bigoplus_{(t,s)}\omega^{(t,s)}$ with
$\omega^{(t,s)}\in\Lin{A_t\Qsys_s}$ and every block positive semidefinite, so
that block diagonality along $\mathbb{T}$ and along $\SecSet$ both hold by
the presentation, subject to the same conditions $\tr\omega=1$,
$\omega\models\Decl$ and $\omega\models\CList$. Write
$\hat O^{\uparrow(s)}_k\coloneqq\idon{}\otimes\hat O^{(s_{f_k})}_k\otimes\idon{}$
for the $s$-block of the \termref{term-lifted-operator}{lifted operator} of an entry scoped to $f_k$
(DX6), and for a \termref{term-party-local}{party-local} entry scoped to $f$ split
$s=(s_{f^c},s_f)$, with $\Qsys_{f^c,s_{f^c}}$ the product of the other
factors' blocks. The conditions then read, block by block,
\begin{equation}\label{eq:feas-general}
\begin{aligned}
  &\textstyle\sum_{(t,s)}\tr\omega^{(t,s)}=1
  &&\text{(normalization)};\\
  &\textstyle\sum_{s}\tr_{\Qsys_s}\omega^{(t,s)}=\tau^{(t)}_A
  \quad\forall t
  &&\text{(marginal, trusted-source mode)};\\
  &\textstyle\sum_{s}\tr_{\Qsys_s}\!\big[(\idon{A_t}\otimes
  \hat O^{\uparrow(s)}_k)\,\omega^{(t,s)}\big]\bowtie_k\Theta^{(t)}_k
  \quad\forall t
  &&\text{(trusted-source mode)};\\
  &\textstyle\sum_{(t,s)}\tr\!\big[(\idon{A_t}\otimes
  \hat O^{\uparrow(s)}_k)\,\omega^{(t,s)}\big]\bowtie_kc_k
  &&\text{(untrusted-source mode)};\\
  &\textstyle\sum_{s_f\in\SecSet^{f}}\tr_{\Qsys_{f,s_f}}\!\Big[\big(\idon{}
  \otimes(\hat O^{(s_f)}_k-c_k\idon{})\big)\,
  \omega^{(t,(s_{f^c},s_f))}\Big]\bowtie_k0
  \quad\forall(t,s_{f^c})
  &&\text{(party-local, either mode)},
\end{aligned}
\end{equation}
the third line with $\Theta^{(t)}_k=c_k\tau^{(t)}_A$ for an entry declared
homogeneous, the fourth for a scalar entry, and the last as an operator on
$\Hs{A_t}\otimes\Hs{\Qsys_{f^c,s_{f^c}}}$. Each
line is the corresponding operator condition on $A\Qsys$: the marginal
$\omega_A$, the entry value $\tr_{\Qsys}[(\idon{A}\otimes\hat O_k)\omega]$,
$\tau_A$ and every $\Theta_k$ are block diagonal along $\mathbb{T}$, the
party-local value
$\tr_{\Qsys_f}[(\idon{}\otimes(\hat O_k-c_k\idon{\Qsys_f}))\,\omega]$ is block
diagonal along the pairs $(t,s_{f^c})$, and an equality or a semidefinite
inequality between operators block diagonal for one decomposition holds
exactly when it holds block by block. With $|\mathbb{T}|=1$, one
\termref{term-measured-factor}{measured factor} and no party-local entry, \eqref{eq:feas-general} is
\eqref{eq:feas} with \eqref{eq:feas-modes}, line for line. With one measured
factor the party-local line reads
$\sum_s\tr_{\Qsys_s}[(\idon{A_t}\otimes(\hat O^{(s)}_k-c_k\idon{}))
\omega^{(t,s)}]\bowtie_k0$ on $\Hs{A_t}$, that is
$\tr_{\Qsys}[(\idon{A}\otimes\hat O_k)\omega]\bowtie_kc_k\,\omega_A$. The
set is convex and compact: every line is an affine equality or an affine
semidefinite or scalar inequality, and $\omega\succeq0$ with $\tr\omega=1$.
\end{definition}
\claimmeta{\metalabel{Uses}\mref{Def.}{def:declaration}, \mref{Def.}{def:exact}
\metasep \metalabel{Used by}\mref{Lem.}{lem:hermitian-constraint-values}, \mref{Rem.}{rem:wellposed}, \mref{Rem.}{rem:block-restriction}, \mref{Def.}{def:homog}, \mref{Def.}{def:phi}, \mref{Thm.}{thm:measurements-and-statistics}, \mref{Thm.}{thm:feasibility}, \mref{Thm.}{thm:validator-soundness}, \mref{Prop.}{prop:conservative}, \extref{../meat-core/contract.pdf}{meat-core/def:feasible}
\metasep newly derived
\metasep implemented (effective: accepted)
\metasep \metalabel{Code}\coderef{https://github.com/fraunhofersg/certiqs-key/blob/v0.1.0/src/certiqs_key/protocol_compiler/feasibility.py}{certiqs_key.protocol_compiler.feasibility:in_feasible_set}
\metasep \metalabel{Oracles}\coderef{https://github.com/fraunhofersg/certiqs-key/blob/v0.1.0/tests/oracles/registry_protocol_compiler.py}{registry_protocol_compiler.py}}

% claim protocol-compiler/lem:hermitian-constraint-values
\claimsummary{protocol-compiler/lem:hermitian-constraint-values}{For a Hermitian constraint operator and a Hermitian \ensuremath{\omega}, the partial trace of their product over the measured system is Hermitian and its trace real; under D(ii) and D(vi) every trusted right-hand side is Hermitian. So both sides of each constraint line can be compared.}
\claimbookmark{protocol-compiler/lem:hermitian-constraint-values}{Lemma}{lem:hermitian-constraint-values}{Hermiticity of the constraint values}
\claimleadin{%
The trusted line of \eqref{eq:feas-modes} compares two operators on $A$ in
the L\"owner order, which is defined on Hermitian operators, and the untrusted
line compares a trace with a real number. That the left-hand side is
Hermitian is \emph{not} automatic from $\omega=\omega^{\dagger}$ alone, a
product of two Hermitian operators need not be Hermitian; it follows from the
Hermiticity of $\hat O_k$ that D(iii) requires.
}
\begin{technicalclaim}
\begin{lemma}[Hermiticity of the constraint values]\label{lem:hermitian-constraint-values}
Let $\hat O_k$ be an operator on $\Qsys$ with
$\hat O_k=\hat O_k^{\dagger}$, as D(iii) of Definition~\ref{def:target}
requires of every constraint operator, and let $\omega$ be an operator on
$A\Qsys$ with
$\omega=\omega^{\dagger}$. Then:
\begin{enumerate}[label=(\roman*),leftmargin=2.4em]
\item $\tr_{\Qsys}[(\idon{A}\otimes\hat O_k)\,\omega]$ is Hermitian;
\item $\tr[(\idon{A}\otimes\hat O_k)\,\omega]$ is real;
\item where D(ii) and D(vi) hold, the right-hand side of the trusted line of
  \eqref{eq:feas-modes} is Hermitian: $c_k\tau_A$ is, for $c_k$ real and
  $\tau_A$ Hermitian by D(ii), and a general $\Theta_k$ is, by D(vi).
\end{enumerate}
\end{lemma}
\end{technicalclaim}
\claimmeta{\metalabel{Uses}\mref{Def.}{def:declaration}, \mref{Def.}{def:target}, \mref{Def.}{def:feasible}
\metasep \metalabel{Used by}\mref{Rem.}{rem:wellposed}, \mref{Thm.}{thm:measurements-and-statistics}, \mref{Prop.}{prop:conservative}
\metasep newly derived
\metasep accepted}

\begin{proof}
Write $X\coloneqq\idon{A}\otimes\hat O_k$, which is Hermitian because
$\hat O_k$ is. For any operator $Y$ on $A\Qsys$,
\begin{equation}\label{eq:ptrcyc}
  \tr_{\Qsys}\!\big[(\idon{A}\otimes\hat O_k)\,Y\big]
  =\tr_{\Qsys}\!\big[Y\,(\idon{A}\otimes\hat O_k)\big] ,
\end{equation}
an identity of \emph{linear} maps, holding for every $Y$ whatsoever: in any
product bases $\{\ket{i}\}$ of $\Hs{A}$ and $\{\ket{q}\}$ of $\Hs{\Qsys}$,
the $(i,j)$ entry of either side equals
$\sum_{q,q'}\bra{q}\hat O_k\ket{q'}\,\bra{i\,q'}Y\ket{j\,q}$, for the
left-hand side directly, and for the right-hand side after exchanging the
names of the two summation indices.

(i) Using $\omega=\omega^{\dagger}$, $X=X^{\dagger}$ and \eqref{eq:ptrcyc},
$\big(\tr_{\Qsys}[X\omega]\big)^{\dagger}=\tr_{\Qsys}[(X\omega)^{\dagger}]
=\tr_{\Qsys}[\omega X]=\tr_{\Qsys}[X\omega]$.

(ii) $\tr[X\omega]$ is the trace of the Hermitian operator
$\tr_{\Qsys}[X\omega]$ of (i), and so is real.

(iii) $c_k\tau_A$ is Hermitian, $c_k$ being real and $\tau_A$ Hermitian by
D(ii); a general $\Theta_k$ is Hermitian by D(vi).
\end{proof}

% claim protocol-compiler/rem:wellposed
\claimsummary{protocol-compiler/rem:wellposed}{The left-hand sides of both constraint lines are Hermitian, respectively real, because the constraint operators and \ensuremath{\omega} are Hermitian (\hyperlink{claim-protocol-compiler/lem:hermitian-constraint-values}{Lemma~\ref*{lem:hermitian-constraint-values}}). Without Hermiticity the lines lose their meaning, which is why it is a condition of the declaration, decided exactly.}
\claimbookmark{protocol-compiler/rem:wellposed}{Remark}{rem:wellposed}{Hermiticity as a declared condition}
\begin{remark}[Hermiticity as a declared condition]
\label{rem:wellposed}
By Lemma~\ref{lem:hermitian-constraint-values} both sides of the trusted line
of \eqref{eq:feas-modes} are Hermitian, so that their L\"owner comparison is
defined, and the left-hand side of the untrusted line is real, so that it can
be compared with $c_k\in\mathbb{R}$. Drop Hermiticity and both lines lose
their meaning rather than merely their soundness, which is why it is a
condition of Definition~\ref{def:declaration} and of D(iii), and not only a
check.

\emph{Where each hypothesis of that lemma comes from.} $X=X^{\dagger}$, for
$X=\idon{A}\otimes\hat O_k$, is D(iii)'s Hermiticity clause on $\hat O_k$,
decided by check V3 of the validator, which decides the operators block by
block (Definition~\ref{def:validator}). $\omega=\omega^{\dagger}$ is the
Hermiticity half of the positivity condition of \eqref{eq:feas}, decided by
the feasibility check V7 on the candidate it verifies: a condition on
$\omega$ and not an assumption about it. Both are decided exactly.
\end{remark}
\claimmeta{\metalabel{Uses}\mref{Def.}{def:declaration}, \mref{Def.}{def:feasible}, \mref{Lem.}{lem:hermitian-constraint-values}
\metasep newly derived
\metasep accepted}

% claim protocol-compiler/rem:block-restriction
\claimsummary{protocol-compiler/rem:block-restriction}{The feasible set ranges over sector-block-diagonal states alone, a restriction in the unsafe direction that the operator type carries. The consuming engine's master reduction proves it sound by sector dephasing.}
\claimbookmark{protocol-compiler/rem:block-restriction}{Remark}{rem:block-restriction}{Restriction of the feasible set to block-diagonal states}
\begin{remark}[Restriction of the feasible set to block-diagonal states]\label{rem:block-restriction}
\emph{The restriction.} The consuming engine's argument is sound when its
minimization ranges over the \emph{sector-block-diagonal} $\omega$ of
\eqref{eq:feas} rather than over all of $\Lin{A\Qsys}$: every state that
argument ranges over is carried into sector-block-diagonal form at no increase
of the bound. This module states the restriction and does not discharge it.
Its direction is the \emph{unsafe} one --- a smaller domain gives a larger
infimum, so an unlicensed restriction overstates a key length. And it is
carried by the operator \emph{type}: an $\omega$ \emph{is} a list of blocks
and no expressible object has an off-sector entry
(Definition~\ref{def:feasible}), so nothing verifies it and nothing can. The
argument that the sectors are a symmetry of every state the security proof
must admit belongs to the consuming engine and to the reductions module; this
module could not make it, since it never sees the reachable set.
Definition~\ref{def:feasible} \emph{defines} the \termref{term-feasible-set}{feasible set}, and it
is not unsound if the restriction is unlicensed.

\emph{Where it is proved.} The consuming engine's master reduction
(\xref{meat-core/thm:reduction}, step~1 of its proof, the sector dephasing,
and \xref{meat-core/cor:reduction-constrained} for a non-empty list) proves
the restriction: an arbitrary feasible state is dephased, feasibility and
every constraint value survive because every $F_v$ and every $\hat O_k$ is
sector-block-diagonal, and the rest of the proof bounds the entropy by the
objective at the dephased state.

\emph{In the general format} (Definition~\ref{def:feasible}, the set of
\eqref{eq:feas-general}) the restriction is along the \emph{joint} sectors:
along the sectors of the \exttermref{../reductions/contract.pdf}{term-measured-system}{measured system} as before and along Alice's
declared sectors. It is proved by the same step of the same theorem, which
dephases by the joint pinching, under preconditions that hold by
\emph{type}: Alice's measurement slot, $\tau_A$ and every right-hand side are
\exttermref{../foundations/contract.pdf}{term-block-list}{block lists} against a \emph{declared} factor tagged $A$,
as $\{F_v\}$ and every $\hat O_k$ are against the
\termref{term-measured-factor}{measured factors} (\xref{meat-core/thm:reduction},
Step~1, with \xref{meat-core/prop:a-neutral} for the statement that the
restriction along Alice's sectors changes no value). The direction is
unchanged and so is the warning: the licence is that theorem under a declared
factor. A producer that emits a block-diagonal $\tau_A$ or $\{M_x\}$ without
a factor tagged $A$ has declared one Alice sector, and nothing downstream may
read a restriction off the shape of its data.
\end{remark}
\claimmeta{\metalabel{Uses}\mref{Def.}{def:feasible}
\metasep newly derived
\metasep accepted}

% claim protocol-compiler/def:homog
\claimsummary{protocol-compiler/def:homog}{A constraint list is licensed when every entry is declared in the form its source mode admits, homogeneous scalar data in trusted mode and scalars in untrusted mode, or carries the party-local tag. Licensing reads the declared tag and compares no numbers.}
\claimbookmark{protocol-compiler/def:homog}{Definition}{def:homog}{Licensed constraint list}
\claimleadin{%
Reading the tag of an entry off the type of its datum is a decision on the
data handed over, and infers nothing from a property of an object the entry
was built from (\xref{foundations/req:consumer-decides-exactly}). The
reductions module, which defines the \exttermref{../reductions/contract.pdf}{term-reduction-chain}{reduction chain} this module consumes,
declares every entry's right-hand side as a real constant
(\xref{reductions/def:licensed}, \xref{reductions/cor:entries}). On a
one-factor emission of that module the form beside it is the tag value
homogeneous in trusted-source mode, so every entry is homogeneous, the
licensed form in trusted-source mode, and an arriving tag field agrees with
the type read off the datum; where the \exttermref{../reductions/contract.pdf}{term-measured-system}{measured system} is a declared product
of party-scoped factors, an entry may instead arrive tagged party-local, with
a scope (\xref{reductions/req:party-local-tag}). That licensing is
decided exactly and at no window is part of Requirement~\ref{req:no-window}.
}
\begin{definition}[Licensed constraint list]\label{def:homog}
Each entry of $\CList$ carries its right-hand side in a declared form, and
the form is part of the datum rather than a property measured of it
(Definition~\ref{def:declaration}). There are two \emph{untagged} forms,
called so because their tag is the type of their datum and needs no field of
its own: a
\termdef{term-homogeneous}{homogeneous}
entry declares a real scalar $c_k$, a
\termdef{term-general}{general}
entry declares a Hermitian $\Theta_k$ on $A$; in untrusted-source mode only
the scalar form exists. For these two forms
the tag \emph{is} the type of the right-hand side an entry carries, and it is
read off that type: an entry whose right-hand-side slot holds a real scalar
is homogeneous, and one whose slot holds a Hermitian matrix on $A$ is
general. \emph{That
reading is a type decision}: it inspects which of two objects a slot holds,
compares no numbers, reads no matrix entry and measures nothing. The third
form, party-local, is declared by its tag and is not read off the type of the
datum. $\CList$ is
\termdef{term-licensed-for-the-declared-source-mode}{licensed for the declared source mode}
when every entry is in the form that mode admits:
\begin{enumerate}[label=(\alph*),leftmargin=2.2em]
\item \emph{trusted-source mode:} every entry is declared
  homogeneous, carrying $c_k\in\mathbb{R}$; the right-hand side it denotes is
  then \emph{constructed} as $\Theta_k\coloneqq c_k\tau_A$ from the declared
  $c_k$ and the declared $\tau_A$, and the entry of \eqref{eq:feas-modes}
  reads $\tr_{\Qsys}[(\idon{A}\otimes\hat O_k)\,\omega]\bowtie_kc_k\tau_A$;
\item \emph{untrusted-source mode:} every right-hand side is a real number
  $\Theta_k=c_k$ and the entry reads
  $\tr[(\idon{A}\otimes\hat O_k)\,\omega]\bowtie_kc_k$.
\end{enumerate}
Licensing is therefore a \emph{type} check on the tag, and it compares no
numbers: what is inspected is which of the two forms an entry declares; no
operator identity is tested against any number, and nothing here is decided
at a window. \emph{Only in cases (a) and
(b) are the entries written $(\hat O_k,\bowtie_k,c_k)$}; a
general entry is written $(\hat O_k,\bowtie_k,\Theta_k)$ and is the form the
notation table records under $\Theta_k$. A general entry is admissible
input, every structural claim of this contract being proved for it, but it is
not licensed, so a list containing one carries no key-length claim
(Theorem~\ref{thm:gates}).

\emph{The \exttermref{../reductions/contract.pdf}{term-party-local-form}{party-local form}, a third licensed form.} In the general format of
Definition~\ref{def:declaration}\,(DX6) an entry may be declared
\termdef{term-party-local}{party-local}:
it carries its scope $f$, a block list $\hat O_k$ against the
\exttermref{../foundations/contract.pdf}{term-sector-map}{sector map} of the \termref{term-measured-factor}{measured factor} $f$ and a real $c_k$, it is
written
$(\hat O_k,\bowtie_k,c_k;f)$, and its state-level reading, in either source
mode, is
\begin{equation}\label{eq:partylocal-entry}
  \tr_{\Qsys_f}\!\Big[\big(\idon{}\otimes
  (\hat O_k-c_k\,\idon{\Qsys_f})\big)\,\omega\Big]\;\bowtie_k\;0
  \qquad\text{as an operator on the systems of }\omega
  \text{ other than }\Qsys_f ,
\end{equation}
those systems being $A$ and the measured factors $\Qsys_{f'}$, $f'\ne f$;
equivalently $\tr_{\Qsys_f}[(\idon{}\otimes\hat O_k)\,\omega]\bowtie_k
c_k\tr_{\Qsys_f}\omega$, the entry being homogeneous with respect to the
marginal that is actually there. Block by block it is the last line of
\eqref{eq:feas-general}. The list is then \emph{licensed} when every entry
is either in the form the source mode admits, (a) or (b), read for its
\termref{term-lifted-operator}{lifted operator}, or is declared party-local. This is still a type
check and compares no numbers: what is inspected is the declared tag. A
party-local entry's right-hand-side slot holds a real scalar, as a
homogeneous entry's does, so the type of the datum does not separate the
two: an entry is party-local exactly when it carries that tag.

\emph{The soundness condition of a party-local entry.} A party-local entry
$(\hat O_k,\bowtie_k,c_k;f)$ meets the
\termdef{term-soundness-condition-of-a-party-local-entry}{soundness condition of a party-local entry}
when every attack reaches $\Qsys_f$ through a channel $\Lambda_f$ acting on
that factor alone which carries the spectral certificate
$\Lambda_f^{\dagger}(\hat O_k)\bowtie_kc_k\idon{}$, that is, the adjoint of
the channel applied to the entry's operator stands in the entry's relation to
$c_k$ times the identity.

\emph{What the tag changes, case by case.} With one measured factor and
trusted-source mode, \eqref{eq:partylocal-entry} reads
$\tr_{\Qsys}[(\idon{A}\otimes\hat O_k)\,\omega]\bowtie_kc_k\,\omega_A$, which
under the marginal condition $\omega_A=\tau_A$ is the form (a) of the same
$(\hat O_k,\bowtie_k,c_k)$. With one measured factor and $d_A=1$, in either
mode, it is the comparison of real numbers
$\tr[\hat O_k\omega]\bowtie_kc_k\tr\omega=c_k$, the form (b). \emph{With one
measured factor, untrusted-source mode and $d_A>1$ it is neither}: it is an
operator inequality on $A$ against the marginal $\omega_A$ that is actually
there, and the scalar entry (b) of the same $(\hat O_k,\bowtie_k,c_k)$ is its
trace, implied by it and not equivalent to it. With more than one measured
factor it implies, in either mode, the untagged reading of the same entry,
which is obtained from it by the partial trace over the other measured
factors (and over $A$ in untrusted-source mode), and it is in general not
implied by it. So the tag never weakens an entry, and of the four one-factor
cases, by source mode and by whether $d_A=1$, it changes something only in
untrusted-source mode at $d_A>1$.

\emph{Which relation symbols a tagged entry carries, in either mode.} An entry
declared party-local carries $\bowtie_k\in\{\preceq,=,\succeq\}$ \emph{in
either source mode}, because \eqref{eq:partylocal-entry} is a relation
between \emph{operators} in either mode: its left-hand side is an operator on
the systems of $\omega$ other than $\Qsys_{f_k}$, which are $A$ and the other
measured factors. The scalar case is the degenerate one: where
$d_{A_t}\prod_{f\ne f_k}d_{f,s_f}=1$ for every row --- every Alice block, and
every sector block of every measured factor other than $f_k$, of dimension
one, which with one measured factor means every Alice block of dimension one,
whether $d_A=1$ or $d_A=|\mathbb{T}|>1$ --- each row is a comparison of
$1\times1$ matrices, that is of real numbers, and the three operator symbols
are then the three scalar symbols $\le$, $=$, $\ge$. With one measured factor
and $d_A=1$ there is one row, and the reading is the form (b) of the same
$(\hat O_k,\bowtie_k,c_k)$, as the case-by-case paragraph states; otherwise
there is one such comparison per row, and the reading relates to the untagged
forms as that paragraph states for its case. \emph{Nothing about the
untagged forms changes}: an untagged entry carries the symbols its mode
admits, and the operator symbols beside a tag are a property of the third
form and of no other. No entry of a declaration in the \termref{term-basic-format}{basic format} is
tagged.
\end{definition}
\claimmeta{\metalabel{Uses}\mref{Def.}{def:declaration}, \mref{Def.}{def:feasible}
\metasep \metalabel{Used by}\mref{Rem.}{rem:party-local-soundness}, \mref{Def.}{def:spec}, \mref{Def.}{def:compile}, \mref{Thm.}{thm:typing}, \mref{Thm.}{thm:feasibility}, \mref{Def.}{def:validator}, \mref{Thm.}{thm:validator-soundness}, \mref{Def.}{def:eligible}, \mref{Thm.}{thm:gates}, \mref{Prop.}{prop:conservative}, \mref{Ass.}{ass:atk}, \mref{Req.}{req:declaration-keys}, \mref{Req.}{req:tag-reading}, \extref{../meat-core/contract.pdf}{meat-core/def:homog}, \extref{../meat-core/contract.pdf}{meat-core/rem:tagged-entries-and-accumulation}
\metasep newly derived
\metasep implemented (effective: accepted)
\metasep \metalabel{Code}\coderef{https://github.com/fraunhofersg/certiqs-key/blob/v0.1.0/src/certiqs_key/protocol_compiler/gates.py}{certiqs_key.protocol_compiler.gates:licensed_for_the_emitted_mode}}

% claim protocol-compiler/rem:party-local-soundness
\claimsummary{protocol-compiler/rem:party-local-soundness}{A party-local entry shrinks the feasible set more than its untagged reading and is sound only if every attack reaches its factor through a channel on that factor alone with a spectral certificate. The reduction chain's producer owes it and the attack-set licence attests it; nothing here decides it.}
\claimbookmark{protocol-compiler/rem:party-local-soundness}{Remark}{rem:party-local-soundness}{Soundness of a party-local entry}
\begin{remark}[Soundness of a party-local entry]\label{rem:party-local-soundness}
A party-local entry of Definition~\ref{def:homog} shrinks the feasible set
further than its untagged reading does, and it is sound only where it meets
the soundness condition of a party-local entry stated in that definition, the
spectral certificate on a channel acting on $\Qsys_f$ alone through which
every attack reaches that factor. That is a certificate of the module that produces the chain,
stated by the reductions module (\xref{reductions/thm:licence}\,(v), with
\xref{reductions/lem:party-local} and the \exttermref{../foundations/contract.pdf}{term-attack-set}{attack set} of
\xref{reductions/def:attset}); it is what the attack-set licence of a list
carrying the tag attests, and nothing in this module decides it. A general
entry whose right-hand side is $\Theta_k=c_k\idon{A}$ stays unlicensed in
every case: the right-hand side of \eqref{eq:partylocal-entry} is $c_k$
times a marginal of $\omega$, never $c_k$ times an identity.
\end{remark}
\claimmeta{\metalabel{Uses}\mref{Def.}{def:homog}
\metasep newly derived
\metasep accepted}

% claim protocol-compiler/ex:two-party-shape
\claimsummary{protocol-compiler/ex:two-party-shape}{An untrusted source serves two measuring parties; the measured system is a product of two factors with their own sector maps and families, and no product matrix is formed. A party-local entry bounds Bob's flag weight by 1/8 conditionally on every state of Alice's measured factor.}
\claimbookmark{protocol-compiler/ex:two-party-shape}{Example}{ex:two-party-shape}{Declaration of a two-party protocol with an untrusted source}
\claimleadin{%
The example writes out the general format of
Definition~\ref{def:declaration} for a protocol in which an untrusted source
serves two \exttermref{../reductions/contract.pdf}{term-measuring-parties}{measuring parties}. It illustrates the format and asserts nothing
about any physical protocol.
}
\begin{example}[Declaration of a two-party protocol with an untrusted source]\label{ex:two-party-shape}
\emph{Systems and sector map.}
\exttermref{../transmitter/contract.pdf}{term-untrusted-source-mode}{Untrusted-source mode}, $d_A=1$: the
\termref{term-marginal-slot}{marginal slot} holds $\nomarg$, $\Xset=\{x_0\}$ and $M_{x_0}=[1]$
(\xref{transmitter/rem:da1}). The sector map has no factor tagged $A$ and two
factors tagged $\Qsys$, in this order: $f=\mathsf a$ with party label
$\mathsf A$ and $f=\mathsf b$ with party label $\mathsf B$
(Definition~\ref{def:declaration}\,(DX1)). Each factor has a \exttermref{../reductions/contract.pdf}{term-data-sector}{data sector} of
dimension $2$ and a flag sector of dimension $1$,
$\SecSet^{f}=(\mathrm{d},\mathrm{fl})$, $d_{f,\mathrm d}=2$,
$d_{f,\mathrm{fl}}=1$. The product map has the four sectors
$(\mathrm d,\mathrm d),(\mathrm d,\mathrm{fl}),(\mathrm{fl},\mathrm d),
(\mathrm{fl},\mathrm{fl})$ in that order, the last factor fastest, of
dimensions $4,2,2,1$; $|\mathbb{T}|=1$, so these are also the
\termref{term-joint-sector}{joint sectors}, and a state $\omega$ is a list of four blocks of those sizes.

\emph{Family.} Each factor declares its own family over
$\Vset^{f}=(0,1,\mathrm{fl})$, in this order:
$F^{f,(\mathrm d)}_0=\mathrm{diag}(1,0)$,
$F^{f,(\mathrm d)}_1=\mathrm{diag}(0,1)$, $F^{f,(\mathrm d)}_{\mathrm{fl}}=0$,
and $F^{f,(\mathrm{fl})}_0=F^{f,(\mathrm{fl})}_1=[0]$,
$F^{f,(\mathrm{fl})}_{\mathrm{fl}}=[1]$. Each is \exttermref{../foundations/contract.pdf}{term-exactly-admissible}{exactly admissible}
for its factor's map, its two completeness sums rounding nowhere
(Definition~\ref{def:exact}, per measured factor). The \exttermref{../reductions/contract.pdf}{term-outcome-alphabet}{outcome alphabet} is
$\Vset=\Vset^{\mathsf a}\times\Vset^{\mathsf b}$ with nine letters, and
$F_{(v_{\mathsf a},v_{\mathsf b})}=F^{\mathsf a}_{v_{\mathsf a}}\otimes
F^{\mathsf b}_{v_{\mathsf b}}$ by the rule \eqref{eq:productfamily}; its block
on $(\mathrm d,\mathrm{fl})$, for instance, is
$F^{\mathsf a,(\mathrm d)}_{v_{\mathsf a}}\otimes
F^{\mathsf b,(\mathrm{fl})}_{v_{\mathsf b}}$, a $2\times2$ matrix that is
written down by nobody. Nine $9\times9$ arrays are never formed.

\emph{\termref{term-party-partition}{Party partition} and announcements.}
$\rd{\mathsf A}(x_0,v)=v_{\mathsf a}$ and $\rd{\mathsf B}(x_0,v)=v_{\mathsf b}$,
which is \emph{not} the \termref{term-default-party-partition}{default party partition}: under the default one
both parties' data would sit in $\rd{\mathsf B}$ and party computability would
be vacuous. A transcript with $L=2$: $\pi=(\mathsf A,\mathsf B)$, each party
announcing whether its outcome is a flag,
$a_1(v_{\mathsf a})=a_2(v_{\mathsf b},\kappa_1)=\mathrm{fl}$ on the flag
letter and $\mathrm{ok}$ otherwise. Sifting passes on
$\kappa=(\mathrm{ok},\mathrm{ok})$; $\hold=\mathsf A$, $d_Z=2$, and the key
rule $r^{0}_{\mathsf A}$ of the specification's public part
(Definition~\ref{def:spec}\,(d)), into $\Zval=\{0,1\}$, is
$r^{0}_{\mathsf A}(v_{\mathsf a},\kappa)=v_{\mathsf a}$ for
$v_{\mathsf a}\in\{0,1\}$ and $r^{0}_{\mathsf A}(\mathrm{fl},\kappa)=0$, a
value no passing round reads, since $v_{\mathsf a}=\mathrm{fl}$ announces
$\kappa_1=\mathrm{fl}$. Party computability, \eqref{eq:partycomp}, requires
the raw key to be a function of Alice's own readout and the transcript, as it
is.

\emph{A party-local entry.} Let $\hat O$ be the block list on factor
$\mathsf b$ with $\hat O^{(\mathrm d)}=0$ and $\hat O^{(\mathrm{fl})}=[1]$,
the projector onto Bob's flag sector, and declare
$(\hat O,\preceq,\tfrac18;\mathsf b)$ with the tag party-local. By the last
line of \eqref{eq:feas-general}, with $s_{f^c}=s_{\mathsf a}$, it reads, for
each of Alice's two sectors $s_{\mathsf a}$,
\begin{equation*}
  \tr_{\Qsys_{\mathsf b,\mathrm{fl}}}\omega^{(s_{\mathsf a},\mathrm{fl})}
  -\tfrac18\Big(\tr_{\Qsys_{\mathsf b,\mathrm d}}
  \omega^{(s_{\mathsf a},\mathrm d)}
  +\tr_{\Qsys_{\mathsf b,\mathrm{fl}}}\omega^{(s_{\mathsf a},\mathrm{fl})}\Big)
  \;\preceq\;0
  \qquad\text{on }\Hs{\Qsys_{\mathsf a,s_{\mathsf a}}} ,
\end{equation*}
a $2\times2$ and a $1\times1$ condition: Bob's flag weight is at most
$\tfrac18$ \emph{conditionally on every state of Alice's measured factor
$\mathsf a$}. The same
entry without the tag reads
$\tr[(\idon{}\otimes\hat O)\omega]\le\tfrac18$, one real inequality, which is
the trace of the sum of the two conditions above and does not imply them: it
allows all of Bob's flag weight to sit where Alice's measured factor is in
one fixed state. Which reading the \exttermref{../foundations/contract.pdf}{term-producer}{producer} of the entries may emit is decided by
its certificate (Definition~\ref{def:homog},
Remark~\ref{rem:party-local-soundness}), not here.
\end{example}
\claimmeta{newly derived
\metasep accepted}

\section{Specification, compilation map and typing}\label{group-protocol-compiler/def:spec}
\sectionintro{A specification supplies the pieces, and the \termref{term-compilation-map}{compilation map} assembles the declaration from them, passing declared data through unchanged and constructing only the raw-key map, the \termref{term-index-map}{index map} and the weighted families. A specification is well-formed when its alphabets and systems are finite, its tables total, every declared measurement family, the presented physical one included, exactly admissible, its weights normalised and its constraint blocks Hermitian and tagged. Every declaration the map emits from a \termref{term-well-formed-specification}{well-formed specification} meets the typing conditions, and its index map is a bijection on every alphabet it indexes (\hyperlink{claim-protocol-compiler/thm:typing}{Theorem~\ref*{thm:typing}}). Its declared measurement families are POVMs block by block as identities of the declared numbers, per factor and, in \termref{term-povm-form}{POVM form}, per Alice sector in the general format, where the product family and an \termref{term-ensemble-form}{ensemble form}'s denoted family are POVMs by theorem; the statistics functional is a probability distribution on every normalised positive block-list state (\hyperlink{claim-protocol-compiler/thm:measurements-and-statistics}{Theorem~\ref*{thm:measurements-and-statistics}}).}
\sectionimports{\extref{../transmitter/contract.pdf}{transmitter/def:source-segment}, \extref{../transmitter/contract.pdf}{transmitter/def:source-spec}, \extref{../transmitter/contract.pdf}{transmitter/def:srcert}}

% claim protocol-compiler/def:spec
\claimsummary{protocol-compiler/def:spec}{A protocol specification is the compiler's input: the transmitter artefact, the receiver part, a two-segment reduction chain, the public part and the constraint entries. In the general format the receiver part and the chain come per measured factor, with a transcript and scoped entries.}
\claimbookmark{protocol-compiler/def:spec}{Definition}{def:spec}{Protocol specification}
\claimleadin{%
The compiler's input comes in five parts. Parts (a)--(e) describe the
specification of a declaration in the \termref{term-basic-format}{basic format}; the paragraphs on
the general format read the same five parts more widely. The descriptors a
link carries are, for a \exttermref{../transmitter/contract.pdf}{term-source-link}{source link}, the transmitter's
\exttermref{../transmitter/contract.pdf}{term-source-descriptor}{source descriptor} and \exttermref{../transmitter/contract.pdf}{term-target-descriptor}{target descriptor}; for a receiver link
they are receiver-side. How the compiler treats the slots it reads is
Requirement~\ref{req:pass-through}. Which objects check V1 compares, and why
a supplied witness is compared at check V7 instead, is
Remark~\ref{rem:chain-and-validator}.

In the general format, how (DX1) restricts the \termref{term-party-partition}{party partition} is
Remark~\ref{rem:party-label-restriction}; that the product rule
\eqref{eq:productfamily} of (b$'$) is formed by no one is
Requirement~\ref{req:derived-objects}; and the ordered
list of the factors' artefacts in (c$'$) is the one
\xref{reductions/req:factorised-emission} emits. A certificate declared in
\exttermref{../transmitter/contract.pdf}{term-block-form}{block form} is refused at gate (G2) of Definition~\ref{def:eligible}.
The ground is not that this module demands purity of every certificate: at several Alice sectors global purity
fails for a presentation too. What separates the two forms is stated in
Remark~\ref{rem:gate-by-certificate-form}, and the operative ground for the
refusal is the scope rule: admitting a block form would widen an
unsafe-direction gate at the degenerate values.
}
\begin{definition}[Protocol specification]\label{def:spec}
A \termdef{term-protocol-specification}{protocol specification} is a tuple
\begin{equation*}
  \Spec=\big(\Spec_{\mathrm{src}},\ \Spec_{\mathrm{rec}},\ \pcChain,\
  \Spec_{\mathrm{pub}},\ \Spec_{\mathrm{con}}\big)
\end{equation*}
whose parts are:

\emph{(a) Source part $\Spec_{\mathrm{src}}$.} A \exttermref{../transmitter/contract.pdf}{term-source-specification}{source specification}
in the sense of \xref{transmitter/def:source-spec}, presented to this module as
the \emph{transmitter artefact} in its declared mode: a finite non-empty
alphabet $\Xset$ with its \exttermref{../foundations/contract.pdf}{term-declared-order}{declared order}, a finite-dimensional system
$A$ with declared computational basis, a POVM $\{M_x\}_{x\in\Xset}$ on $A$
--- itself possibly presented as a \termref{term-choice-weighted-family}{choice-weighted family}
(Definition~\ref{def:weighted}) --- and the \exttermref{../transmitter/contract.pdf}{term-specified-signal-ensemble}{specified signal ensemble}
$\{p_x,\rho_x\}_{x\in\Xset}$ on the \exttermref{../transmitter/contract.pdf}{term-transmitted-system}{transmitted system} $\Asig$
(\xref{transmitter/ass:ens}). The mode is read off the artefact's
\termref{term-marginal-slot}{marginal slot} and is never a second datum. In
\exttermref{../transmitter/contract.pdf}{term-trusted-source-mode}{trusted-source mode} the slot holds $\tau_A\in\Seq{A}$ --- at $d_A=1$
the slot the transmitter filled with $[1]$ included --- and at $d_A\ge2$ the
artefact additionally carries a \exttermref{../transmitter/contract.pdf}{term-source-replacement-certificate}{source-replacement certificate}
$(\tau_A,\Psi,\{p_x,\rvirt_x\}_{x\in\Xset})$ on $A\Avirt$ in the sense of
\xref{transmitter/def:srcert}, written against the \exttermref{../transmitter/contract.pdf}{term-virtual-ensemble}{virtual ensemble}
the \exttermref{../transmitter/contract.pdf}{term-source-segment}{source segment} \emph{starts from}, which lives on the
\exttermref{../transmitter/contract.pdf}{term-virtual-transmitted-system}{virtual transmitted system} $\Avirt$, the \exttermref{../transmitter/contract.pdf}{term-source-descriptor}{source descriptor} of
$\Xi_1$ (\xref{transmitter/def:source-segment}\,(S1), (S3)). When
$L_{\mathrm s}=0$ that system and that ensemble are the specified ones:
$\Avirt=\Asig$ and $\rvirt_x=\rho_x$. In \exttermref{../transmitter/contract.pdf}{term-untrusted-source-mode}{untrusted-source mode} the slot
holds $\nomarg$ and there is no certificate.

\emph{(b) Receiver part $\Spec_{\mathrm{rec}}$.} A finite-dimensional physical
receiver system with a POVM $\{\Fph_w\}_{w\in\Vph}$, possibly presented as a
choice-weighted family; and, optionally, a public choice whose weight
depends on the private setting of the measuring apparatus --- a
stochastic matrix $P(f\,|\,t)$ with the setting map
$t:\Vset\to\mathcal{T}$ it is read through --- to be applied by
this module \emph{after} the \exttermref{../reductions/contract.pdf}{term-reduction-chain}{reduction chain}, in the
\termref{term-setting-dependent-form}{setting-dependent form} of Definition~\ref{def:weighted}.

\emph{(c) Reduction chain $\pcChain=(\Chsrc,\Chrec)$.} A reduction chain in
the sense of \xref{reductions/def:chain}, presented to this module as a
record: an ordered \emph{pair} of ordered lists of \exttermref{../reductions/contract.pdf}{term-typed-link}{typed links}, the
\emph{source segment} $\Chsrc=(\Xi_1,\dots,\Xi_{L_{\mathrm s}})$ with
$L_{\mathrm s}\ge0$, which is a source segment in the sense of
\xref{transmitter/def:source-segment} --- its links act on the transmitted
system $\Asig$ and replace Alice's signal ensemble --- and the \emph{\exttermref{../reductions/contract.pdf}{term-receiver-segment}{receiver
segment}} $\Chrec=(\Lambda_1,\dots,\Lambda_L)$ with $L\ge1$, whose links act on
the receiver system and replace Bob's measurement. What this module reads of a
link, stated in its own words, is this: each link of either segment carries a
type tag, source and target system descriptors, a status tag, its channel as
data where the link type has one, the constraint entries it transports, and a
licence slot. A receiver-stage link in the sense of
\xref{reductions/def:link} carries more --- its \exttermref{../reductions/contract.pdf}{term-reachable-set-statement}{reachable-set statement} and
its carried operator pair --- and this module reads none of it. The last link
of $\Chrec$ has target descriptor $(\Qsys,\SecSet,(d_s)_s)$ with the
sector projectors $\Pi_s$, and \exttermref{../reductions/contract.pdf}{term-target-povm}{target POVM} $\{F_v\}_{v\in\Vset}$. The two
segments are separate lists because they have different types; the pair is
composed source segment first. The \termref{term-interface-conditions-on-a-chain}{interface conditions on a chain} that
this module decides on the record are those of
Definition~\ref{def:chaincond}.

\emph{(d) Public part $\Spec_{\mathrm{pub}}$.} A \termref{term-party-partition}{party partition}
$\Party$ (readouts $\rd{\mathsf A},\rd{\mathsf B}$ onto $\Oset{\mathsf A}$,
$\Oset{\mathsf B}$, and the party labels), \termref{term-announcement-map}{announcement maps}
$a_A:\Oset{\mathsf A}\to\CAset$ and $a_B:\Oset{\mathsf B}\to\CBset$, a sifting
predicate $\varsigma:\Kann\to\{\mathrm{pass},\mathrm{fail}\}$, a key rule
$r_{\hold}^{0}:\Oset{\hold}\times\Kann\to\Zval$, a holder
$\hold\in\{\mathsf A,\mathsf B\}$, a finite acceptance alphabet $\Chatset$ and
a map $\phi:\Kann\to\Chatset$. All maps are total.

\emph{(e) Constraint part $\Spec_{\mathrm{con}}$.} A finite, possibly empty,
list of entries. Each $\hat O_k$ is presented as a \exttermref{../foundations/contract.pdf}{term-block-list}{block list} on
$\Qsys$ together with the \exttermref{../foundations/contract.pdf}{term-sector-map}{sector map} it is written against, which is
carried with it as for every other supplied $\Qsys$-operator family, and is
Hermitian block by block. Each entry carries a right-hand-side datum of one of
two types, and the type fixes the entry's tag: \termref{term-homogeneous}{homogeneous} with a real
$c_k$, or \termref{term-general}{general} with a Hermitian $\Theta_k$ on $A$, in trusted-source
mode; a real $c_k$ in untrusted-source mode (Definition~\ref{def:homog}). The
constraint part also carries the optional attack-set \exttermref{../reductions/contract.pdf}{term-licence-record}{licence record} and an
optional feasibility witness $\omega^{\star}$, likewise a block list with its
own map. The map each $\hat O_k$ carries is compared with the emitted $\Smap$
by check V1, an entry declared here being on the same footing as one
transported by a link. The witness's map is not compared there: $\omega^{\star}$
is a block list on $A\Qsys$ and not a $\Qsys$-operator, it is no slot of
\eqref{eq:declaration}, and branch~(c) of check V7 is where it is required to
be a block list against $\Smap$ (Definition~\ref{def:validator}).

\emph{The general format.} A specification to be compiled into the general
format of Definition~\ref{def:declaration}\,(DX1)--(DX6) has the same five
parts, read as follows. At one Alice sector, one \termref{term-measured-factor}{measured factor},
\termref{term-povm-form}{POVM form}, two messages neither of which reads the other and no
party-local entry these readings are (a)--(e) above, part for part on every
component a derivation reads: at those values every such object named below
is the object of the same name there and nothing else is presented. On the
party labels of the tensor factors of each $\Hs{\Qsys_s}$, which no
derivation of this module reads and which the consuming engine's laboratory
premise reads as naming the parties' laboratories, (DX1) restricts the party
partition of (d). The party set is unchanged: it is
$\{\mathsf A,\mathsf B\}$, with the two \termref{term-readout-map}{readout maps} of $\Party$. What
is new is that several measured factors may carry a party label and that the
factor tagged $A$ carries one too (DX1).

\emph{(a$'$) Source part.} The artefact's marginal slot may hold a block list
$(\tau^{(t)}_A)_{t\in\mathbb{T}}$ along the Alice sectors the artefact
declares, and Alice's measurement arrives in one of the two forms of (DX4):
the POVM slot, or, in trusted-source mode only, the ensemble slot holding the
block lists $\{N_x\}_{x\in\Xset}$ of \eqref{eq:ensembleform}. The
source-replacement certificate may accordingly be declared in
\exttermref{../transmitter/contract.pdf}{term-pure-form}{pure form}, in \exttermref{../transmitter/contract.pdf}{term-block-form}{block form}, or as an
\exttermref{../transmitter/contract.pdf}{term-ensemble-presentation}{ensemble presentation}, all three in the sense of
\xref{transmitter/def:srcert}. Which form a certificate is in is read off the
artefact and is never a second datum, exactly as the mode is. That a form may
be \emph{presented} is the transmitter's rule and is stated here. Which forms
this module \emph{admits}, and what gate (G2) decides of each form, is
Definition~\ref{def:eligible}'s and is stated once there.

\emph{(b$'$) Receiver part.} $\Spec_{\mathrm{rec}}$ is an ordered list, one
entry per measured factor $f\in\mathbb{F}$ and in the order of $\mathbb{F}$: a
party label; a finite-dimensional physical receiver system with a POVM
$\{F^{\mathrm{ph},f}_{w_f}\}_{w_f\in\mathcal{V}^{\mathrm{ph},f}}$, possibly
presented as a choice-weighted family; and, optionally, a
setting-dependent public choice at that factor --- a stochastic matrix
$P_f(\,\cdot\mid\cdot\,)$ with the setting map on that factor's alphabet it is
read through --- to be applied after that factor's segment, in the
setting-dependent form of Definition~\ref{def:weighted}. No entry of the list
carries an operator on another factor's system, and no product of physical
families is presented: the product is the rule \eqref{eq:productfamily},
formed by no one.

\emph{(c$'$) Reduction chain.} $\pcChain=(\Chsrc,\Chrec)$ with $\Chsrc$ as
in (c) and $\Chrec=(\pcChain^{f}_{\mathrm{rec}})_{f\in\mathbb{F}}$ one
receiver segment per measured factor, each
$\pcChain^{f}_{\mathrm{rec}}=(\Lambda^{f}_1,\dots,\Lambda^{f}_{L_f})$,
$L_f\ge1$, a list of receiver-stage links whose source descriptor is factor
$f$'s physical receiver system and whose last target descriptor is
$(\Qsys_f,\SecSet^{f},(d_{f,s_f})_{s_f})$ with target POVM
$\{F^{f}_{v_f}\}_{v_f\in\Vset^{f}}$. What this module reads of a link is what
(c) lists and nothing more, of each segment separately: the segments are
separate lists because their links act on different systems, and no link of
one is composable with a link of another. Where the factors' data are the
ordered list of the factors' artefacts that the reductions module emits,
factor $f$'s receiver segment and \exttermref{../reductions/contract.pdf}{term-target-family}{target family} are that factor's
own, and the product rule of \xref{reductions/thm:emitted-product} is
emitted in place of a \exttermref{../reductions/contract.pdf}{term-joint-family}{joint family}; that list
is what the factor list of (DX1) is built from, in its order. The interface
conditions on a chain that this module decides on the record are those of
Definition~\ref{def:chaincond}, per factor for the receiver segments.

\emph{(d$'$) Public part.} In place of $a_A$ and $a_B$ the public part
declares the \termref{term-transcript-protocol}{transcript} of (DX5): the number $L\ge1$ of messages and,
for each $m$, its sender $\pi_m$, its finite alphabet $\mathcal{C}_m$ and
its total table
$a_m:\Oset{\pi_m}\times\mathcal{C}_1\times\dots\times\mathcal{C}_{m-1}
\to\mathcal{C}_m$. The sifting predicate $\varsigma$, the key rule
$r^{0}_{\hold}$, the holder $\hold$, the acceptance alphabet $\Chatset$ and
$\phi$ are declared on the derived $\Kann=\prod_{m}\mathcal{C}_m$ of
\eqref{eq:transcript}, which is a set determined by the tables and declared by
no one. All tables are total. The party partition $\Party$ is as in (d), with
a party label per measured factor and, where there is one, for the factor
tagged $A$.

\emph{(e$'$) Constraint part.} Each entry presents, in addition to what (e)
lists, its \termref{term-scope}{scope}: the label of one measured factor of the chain; and
$\hat O_k$ is a block list on $\Qsys_{f_k}$ together with the sector map of
that factor, Hermitian block by block. Its right-hand-side datum is of one of
the three forms of Definition~\ref{def:homog}: a real $c_k$ tagged
homogeneous, a Hermitian $\Theta_k$ tagged general and then itself a block
list along $\mathbb{T}$, or a real $c_k$ tagged \termref{term-party-local}{party-local}. The
arriving tag is compared with the type of that datum at check V6, which also
decides that the scope names a measured factor of the emitted map. An optional
witness $\omega^{\star}$ is a block list over the \termref{term-joint-sector}{joint sectors} with
its own map, compared at branch~(c) of V7 and nowhere else, for the reason
(e) gives.
\end{definition}
\claimmeta{\metalabel{Uses}\mref{Def.}{def:declaration}, \mref{Def.}{def:homog}, \extref{../transmitter/contract.pdf}{transmitter/def:source-segment}, \extref{../transmitter/contract.pdf}{transmitter/def:source-spec}, \extref{../transmitter/contract.pdf}{transmitter/def:srcert}
\metasep \metalabel{Used by}\mref{Def.}{def:well-formed-spec}, \mref{Thm.}{thm:typing}, \mref{Thm.}{thm:choice-encoding}, \mref{Rem.}{rem:alphabet}, \mref{Thm.}{thm:sifting-and-key}, \mref{Rem.}{rem:sidecond}, \mref{Def.}{def:validator}, \mref{Thm.}{thm:chain-consumption}, \mref{Thm.}{thm:gates}, \mref{Ass.}{ass:spec}, \mref{Ass.}{ass:alice-meas}, \mref{Req.}{req:da1-source-identity}, \mref{Req.}{req:choice-weighting-record}, \mref{Req.}{req:band-a-ideal-receiver}
\metasep newly derived
\metasep implemented (effective: accepted)
\metasep \metalabel{Code}\coderef{https://github.com/fraunhofersg/certiqs-key/blob/v0.1.0/src/certiqs_key/protocol_compiler/spec.py}{certiqs_key.protocol_compiler.spec:ProtocolSpecification}}

% claim protocol-compiler/def:well-formed-spec
\claimsummary{protocol-compiler/def:well-formed-spec}{A specification is well-formed when its declared data are as declared: finite alphabets and systems, total tables, exactly admissible measurement families, the presented physical one included, normalised weights and Hermitian, tagged constraint blocks. The compiler decides every clause.}
\claimbookmark{protocol-compiler/def:well-formed-spec}{Definition}{def:well-formed-spec}{Well-formed specification}
\claimleadin{%
Every clause below is a condition the compiler decides on the declared
numbers and refuses when it fails, so a specification that is not well-formed
yields a refusal and not a declaration. Clauses (i)--(iii) are decided by
check V1 of Definition~\ref{def:validator} on the emitted data and by the
input types on the consumed tables, and the sector-map half of (vii) by check
V1; the
measurement families of (iv) by check V3 and condition~(C3) of
Definition~\ref{def:chaincond}, the conditional families of a
choice-weighted presentation by Requirement~\ref{req:conditional-family-decision},
and the physical receiver family as presented by
Requirement~\ref{req:physical-family-decision};
(v) by Requirement~\ref{req:weight-list}; (vi) by
Requirement~\ref{req:stochastic-weights}; the Hermiticity in (vii) by check
V3 and the tag by check V6 (Requirement~\ref{req:tag-reading}); a failing
check is refused (Requirement~\ref{req:validator-refusals}). The physical
receiver space carries no sector map this module holds, so the physical
family is read on it as one block, as Alice's data are read on $A$.
}
\begin{definition}[Well-formed specification]\label{def:well-formed-spec}
A protocol specification $\Spec$ (Definition~\ref{def:spec}) is a
\termdef{term-well-formed-specification}{well-formed specification}
when its \exttermref{../foundations/contract.pdf}{term-declared-data}{declared data} are as declared:
(i) all alphabets are finite and non-empty;
(ii) all systems are finite-dimensional with the declared dimensions;
(iii) all tabulated maps are total on their stated domains;
(iv) all declared POVM families are \exttermref{../foundations/contract.pdf}{term-exactly-admissible}{exactly admissible}: every family
other than the physical receiver family $\{\Fph_w\}_{w\in\Vph}$ as presented,
block by block in the sense of Definition~\ref{def:exact}; and the physical
receiver family as presented (in the general format, each
\termref{term-measured-factor}{measured factor}'s, on that factor's physical receiver system), on its
physical receiver space read as one block, condition~(T) being empty for it
as Definition~\ref{def:exact} reads Alice's data: where it is presented
directly, (S), (H), (P) and (C) of that definition hold of it in the
\exttermref{../foundations/contract.pdf}{term-declared-order}{declared order} of its alphabet; where it is presented as a
\termref{term-choice-weighted-family}{choice-weighted family} (Definition~\ref{def:weighted}), its
conditional families are exactly admissible, block by block against the map
they carry, and condition~(S) holds of its own declared-order sum. A
choice-weighted family whose conditional families carry a sector map is a
block list against that map, and reading it as that block list or as one
block decides the same (S), (H), (P) and (C), its off-block entries being
exact zeros;
(v) all public-choice weights are strictly positive --- a domain condition on
a declared scalar (Definition~\ref{def:weighted}) --- and sum to one;
(vi) every setting-dependent weighting, at whichever \exttermref{../reductions/contract.pdf}{term-measured-system}{measured system} or
measured factor it is declared, is a stochastic matrix, its
entries non-negative and its rows summing to one exactly (the
\termref{term-setting-dependent-form}{setting-dependent form} of Definition~\ref{def:weighted});
(vii) every constraint operator is a \exttermref{../foundations/contract.pdf}{term-block-list}{block list} against the upstream
module's \exttermref{../foundations/contract.pdf}{term-sector-map}{sector map} and Hermitian block by block, and every entry
carries its declared right-hand-side tag.
\end{definition}
\claimmeta{\metalabel{Uses}\mref{Def.}{def:exact}, \mref{Def.}{def:spec}, \mref{Def.}{def:weighted}
\metasep \metalabel{Used by}\mref{Thm.}{thm:typing}, \mref{Thm.}{thm:measurements-and-statistics}, \mref{Thm.}{thm:sifting-and-key}, \mref{Req.}{req:physical-family-decision}
\metasep newly derived
\metasep accepted}

% claim protocol-compiler/def:compile
\claimsummary{protocol-compiler/def:compile}{The compilation map fills the declaration from the specification, passing declared data through and constructing only the raw-key map, the index map and the weighted families. In the general format it also assembles the factor list.}
\claimbookmark{protocol-compiler/def:compile}{Definition}{def:compile}{The compilation map}
\claimleadin{%
The compilation map assembles the declaration from the specification. It
passes the \exttermref{../foundations/contract.pdf}{term-declared-data}{declared data} through, weights a measurement family where a public
choice is presented, and constructs only the raw-key map, the index map and
the weighted families. The list below fills the \termref{term-basic-format}{basic format}; the
paragraphs on the general format extend it slot for slot. That the engine
passes the declared data through unchanged, constructs no right-hand side, and
forms no product of two factors' families, no \termref{term-lifted-operator}{lifted operator} and no
other derived object of the general format as data is
Requirements~\ref{req:pass-through} and~\ref{req:derived-objects}; the
refusals the validator issues on the factor list and on Alice's slots are
Requirement~\ref{req:validator-refusals}'s.

At $d_A=1$ a trusted-source request may leave the marginal slot empty; the
transmitter then fills it with $[1]$ and records the fill
(\xref{transmitter/req:marginal-fill}); why the two
source modes impose the same condition there is \xref{transmitter/rem:da1}. A
certificate that is present and fails any of the five conditions of gate (G2)
is refused at every $d_A$, row~3 of the \termref{term-gate-outcome-table}{gate outcome table} of
Definition~\ref{def:gate-outcome-table}. The order in which the binding rule below has
condition~(S) decided on a supplied family is the order that family was
declared in and also the order the consuming engine re-accumulates in (\xref{foundations/req:declared-order}): one order
at all three places.
}
\begin{definition}[The compilation map]\label{def:compile}
The \termdef{term-compilation-map}{compilation map} is the map $\Spec\mapsto\Decl$ that fills
\eqref{eq:declaration} as follows.
\begin{itemize}[leftmargin=1.6em]
\item \emph{Alice's system and measurement.} $A$, $d_A$ and $\{M_x\}$ from
  $\Spec_{\mathrm{src}}$, weighted by Definition~\ref{def:weighted} where a
  choice family is presented.
\item \emph{Alice's marginal.} $\tau_A$ is what the transmitter artefact's
  marginal slot holds: in trusted-source mode the matrix in the slot --- at
  $d_A=1$ the value $[1]$ the transmitter filled the slot with, the unique
  element of $\Seq{A}$ --- and
  $\tau_A=\nomarg$ where the slot is empty. The mode is read off that slot and
  off nothing else. A certificate the artefact carries is consumed by gate
  (G2) of Definition~\ref{def:eligible}.
\item \emph{The \exttermref{../reductions/contract.pdf}{term-measured-system}{measured system} and its family.} $(\Qsys,\SecSet,(d_s)_s)$
  and $\{F_v\}$ from the last link of $\Chrec$. Where $\Spec_{\mathrm{rec}}$
  declares a setting-dependent public choice, the family is weighted
  after the chain by the setting-dependent form of
  Definition~\ref{def:weighted}: the emitted family is then
  $N_{v,f}=P(f\,|\,t(v))\,F_v$ over the support of $P$, the emitted \exttermref{../reductions/contract.pdf}{term-outcome-alphabet}{outcome
  alphabet} is that support, and the setting map $t$ through which $P$ is read
  is the class map of the \exttermref{../reductions/contract.pdf}{term-joint-alphabet}{joint alphabet} of the chain's \exttermref{../foundations/contract.pdf}{term-producer}{producer} where
  the chain was emitted by the reductions module
  (\xref{reductions/def:joint-family}\,(a)), and is declared with the family
  otherwise.
\item \emph{The public maps.} $\Party$, $a_A$, $a_B$, $\phi$, $\hold$
  verbatim from $\Spec_{\mathrm{pub}}$, declared on the emitted alphabets.
\item \emph{The constraint list.} $\CList$ is the concatenation of
  $\Spec_{\mathrm{con}}$'s entries with the entries transported by the links
  of $\Chsrc$ and then those of $\Chrec$, each segment in its own order, each
  entry with the datum its right-hand-side slot holds and with the tag that
  datum's type determines (Definition~\ref{def:homog}).
\item \emph{The raw-key map.}
\begin{equation}\label{eq:rdef}
  r(x,v)\coloneqq
  \begin{cases}
    r^{0}_{\hold}\big(\rd{\hold}(x,v),\,a(x,v)\big),
      & \varsigma(a(x,v))=\mathrm{pass},\\[0.3ex]
    \perp, & \varsigma(a(x,v))=\mathrm{fail},
  \end{cases}
  \qquad\text{with }d_Z\coloneqq|\Zval| .
\end{equation}
\end{itemize}

\emph{The index map.} Alongside $\Decl$ the compiler emits the
\termdef{term-index-map}{index map}
$\idxm$, the label-to-index map: a choice, for each of $\Xset$, $\Vset$,
$\SecSet$, $\Oset{\mathsf A}$, $\Oset{\mathsf B}$, $\CAset$, $\CBset$,
$\Chatset$ and $\Zset$, of a bijection onto an initial segment of
$\mathbb{Z}_{\ge0}$; every emitted list, matrix and function table is indexed
through $\idxm$, and every semantic label is carried as metadata only. Its
component on $\SecSet$ is the total order the sector map $\Smap$ already
fixes, and not a second choice: one order on the sectors, required by the last
clause of D(i) (Definition~\ref{def:target}) and decided by V1. And $\idxm$ is
what fixes the declared summation order of Definition~\ref{def:exact} on every
alphabet it indexes; it indexes no basis of a matrix, a trace being
accumulated in the row order of the emitted array.

\emph{The binding rule.} On the remaining alphabets the choice is the
compiler's, under one binding rule. Where a supplied family arrives with a
declared order on the alphabet that indexes it --- as the target POVM of the
last link of $\Chrec$ does, its \exttermref{../foundations/contract.pdf}{term-declarer}{declarer} having established its own
completeness sum in that order --- the compiler takes \emph{that} order as
$\idxm$'s component there. So the order in which condition~(S) of
Definition~\ref{def:exact} is decided on a supplied family here is the order
that family was declared in, and no family is refused here for an ordering
its declarer never chose. Where no order arrives with an alphabet the choice is free, and (S) is
decided in the order the compiler chose and emits. This rule prescribes
nothing to the declarer of a family: it constrains only which bijection this
module picks.

Where a choice-weighted presentation is resolved, this same rule settles the
order the decision on each conditional family is read in, and adds nothing to
it. The family the weighting produces is indexed by one alphabet --- $\Xset$
where the presentation is Alice's, the physical receiver alphabet where it is
a measuring party's, which in \termref{term-band-a}{Band~A} is the emitted $\Vset$ --- and where a declared order
arrives with that alphabet the rule above takes \emph{it}, here as everywhere:
$\Xset$ arrives ordered with the source part
(Definition~\ref{def:spec}\,(a)), and the physical alphabet with the order the
fibre sum of Definition~\ref{def:chaincond}\,(C4) is accumulated in.
Condition~(S) on the conditional family of a choice value is then read in the
\emph{restriction of that order} to the elements that choice value
contributes, which is a total order on $\Uset_t$ however the record listed
it; the two conditions are independent, only the order restricting. Where no order arrives with that alphabet the
free choice is made as Definition~\ref{def:weighted} says, by the record's own
listing, and it is emitted. No declarer's order is overridden and no
presentation is refused for an order. Where the setting-dependent weighting is applied,
the emitted alphabet is the support of $P$ and its order is the product order
of Definition~\ref{def:weighted}: $v$ in its declarer's declared order and,
within each $v$, $f$ in the declared order of the \termref{term-choice-alphabet}{choice alphabet}, the
pairs off the support omitted; that is the order (S) is decided in on the
emitted family and the order $\idxm$ emits on it.

\emph{The general format.} For a specification presented as
Definition~\ref{def:spec}'s (a$'$)--(e$'$) the map fills the general format
of Definition~\ref{def:declaration}\,(DX1)--(DX6) as follows. In the
\termref{term-basic-format}{basic format} --- one Alice sector, one measured factor, POVM form, two
messages neither reading the other and no party-local entry --- every clause
below is the clause of the same subject above, slot for slot, and the map
emits what those clauses emit.

\emph{The factor list} (DX1) is assembled and not declared as a whole: the
factor tagged $A$ is the Alice sector map the transmitter artefact declares,
where it declares one, and it is the first of the list; where the artefact
declares none, $\mathbb{T}=\{t_0\}$ and no $A$ factor is emitted. The factors
tagged $\Qsys$ are the entries of $\Spec_{\mathrm{rec}}$ in their declared
order, each carrying that entry's party label and the sector map of the last
target descriptor of that entry's receiver segment, which
Definition~\ref{def:chaincond}\,(C1) decides per factor. The one block kind
is the kind those maps carry, and the factors' agreement in it is decided by
the same four-component comparison V1 makes, taken on the factor list
(Definition~\ref{def:validator}, V1).
The product sector map, its order, the product alphabet and the
\termref{term-joint-sector}{joint sectors} are derived from that list by (DX1) and (DX3) and are
declared, formed and emitted by no one: what is emitted is the factor list and
the rules.

\emph{Alice's slots.} $\tau_A$ is what the artefact's marginal slot holds,
the block list along $\mathbb{T}$ included; Alice's measurement slot is the
artefact's, in the form the artefact declares --- the POVM slot or the
ensemble slot of (DX4) --- weighted by Definition~\ref{def:weighted} in that
form where a choice family is presented. The mode is read off the
marginal slot as above.

\emph{Per measured factor.} $(\Qsys_f,\SecSet^{f},(d_{f,s_f})_{s_f})$ and
$\{F^{f}_{v_f}\}_{v_f\in\Vset^{f}}$ from the last link of that factor's
receiver segment and, where that factor's entry of $\Spec_{\mathrm{rec}}$
declares a setting-dependent public choice, weighted after that segment by the
setting-dependent form of Definition~\ref{def:weighted}, so that factor $f$'s
emitted family is $P_f(\,\cdot\mid\cdot\,)$ times that factor's delivered
family over the support of $P_f$, and factor $f$'s emitted alphabet
$\Vset^{f}$ is that support in its product order. Each factor's weighting is
applied to its own family and to nothing else. No product of two factors'
families is formed by this map, whether emitted or handed on; so the subject of the
representability decision of Definition~\ref{def:weighted} is, as in the basic
format, the products \emph{that definition's} two constructions form, and no
product block \eqref{eq:productfamily} is a product this map forms.

\emph{The public part and the raw key.} $\Party$, $\phi$, $\hold$ verbatim
from $\Spec_{\mathrm{pub}}$ and the \termref{term-transcript-protocol}{transcript} verbatim: the sender,
the alphabet and the table of each message, declared on the emitted
alphabets. The joint announcement function $a$ is not a slot: it is the
composition \eqref{eq:transcript} of those tables, and $r$ is \eqref{eq:rdef}
with \emph{that} $a$, which is what Theorem~\ref{thm:sifting-and-key}\,(iv), (vi) is
stated of. $\CList$ is the concatenation above, each entry carrying in
addition its \termref{term-scope}{scope}, and each with the datum its right-hand-side slot
holds and the tag that datum's type determines; for a party-local entry that
is the tag field it arrives with beside a real scalar
(Definition~\ref{def:homog}).

\emph{The index map in the general format.} $\idxm$ has one component per
factor alphabet $\Vset^{f}$ and one per factor \exttermref{../reductions/contract.pdf}{term-sector-set}{sector set} $\SecSet^{f}$, and,
where an $A$ factor is emitted, one on $\mathbb{T}$, which is the total order
that factor's own map fixes. Its components on $\Vset$ and on $\SecSet$ are
the product orders (DX1) and (DX3) derive from the factors' orders and are not
further choices, as its component on $\SecSet$ above is not; the joint sectors
are ordered with $t$ slowest by (DX1). The binding rule above is applied per
factor: where a factor's target family arrives with a declared order on its
alphabet, the compiler takes that order there. So condition~(S) of
Definition~\ref{def:exact} is decided per factor, in that factor's arriving
order, and per Alice sector on Alice's data; no sum over a product alphabet is
accumulated by this map or by its consumer.
\end{definition}
\claimmeta{\metalabel{Uses}\mref{Def.}{def:declaration}, \mref{Def.}{def:exact}, \mref{Def.}{def:target}, \mref{Def.}{def:homog}, \extref{../transmitter/contract.pdf}{transmitter/def:source-spec}
\metasep \metalabel{Used by}\mref{Thm.}{thm:typing}, \mref{Rem.}{rem:rbuilt}, \mref{Thm.}{thm:measurements-and-statistics}, \mref{Def.}{def:weighted}, \mref{Thm.}{thm:choice-encoding}, \mref{Rem.}{rem:condfamily}, \mref{Rem.}{rem:relabel}, \mref{Thm.}{thm:sifting-and-key}, \mref{Thm.}{thm:feasibility}, \mref{Def.}{def:validator}, \mref{Rem.}{rem:chain-and-validator}, \mref{Thm.}{thm:chain-consumption}, \mref{Thm.}{thm:band-a}, \mref{Thm.}{thm:gates}, \mref{Req.}{req:declaration-keys}, \mref{Req.}{req:pass-through}, \mref{Req.}{req:da1-source-identity}, \mref{Req.}{req:conditional-family-decision}, \mref{Req.}{req:reference-companion-compiled}, \extref{../meat-core/contract.pdf}{meat-core/cor:reference-companion-factors}, \extref{../meat-core/contract.pdf}{meat-core/prop:reference-companion}
\metasep newly derived
\metasep implemented (effective: accepted)
\metasep \metalabel{Code}\coderef{https://github.com/fraunhofersg/certiqs-key/blob/v0.1.0/src/certiqs_key/protocol_compiler/compile.py}{certiqs_key.protocol_compiler.compile:compile_specification}
\metasep \metalabel{Oracles}\coderef{https://github.com/fraunhofersg/certiqs-key/blob/v0.1.0/tests/oracles/registry_protocol_compiler.py}{registry_protocol_compiler.py}}

% claim protocol-compiler/thm:typing
\claimsummary{protocol-compiler/thm:typing}{A declaration the compilation map emits satisfies the typing conditions D(i), D(ii), D(iv) and D(vi), and its index map is a bijection on every alphabet it indexes. Each follows from the construction or from a check decided on the emitted data, never from a source certificate.}
\claimbookmark{protocol-compiler/thm:typing}{Theorem}{thm:typing}{Typing of an emitted declaration}
\claimleadin{%
A declaration is emitted only when no check refused it
(Requirement~\ref{req:gate-outcomes}, row~1). The typing conditions then
follow either from how the compilation map builds the declaration or from a
check that decides them on the emitted data.
}
\begin{theorem}[Typing of an emitted declaration]\label{thm:typing}
Let $\Spec$ be a \termref{term-well-formed-specification}{well-formed specification}
(Definition~\ref{def:well-formed-spec}) satisfying
Assumption~\ref{ass:spec}, and
let $\Decl$ be a declaration that the compilation map of Definition~\ref{def:compile}
\emph{emits}, that is, one for which no check of
Definition~\ref{def:validator} refused. Then $\Decl$, together with the index
map $\idxm$ emitted alongside it, satisfies the target conditions D(i), D(ii),
D(iv) and D(vi) of Definition~\ref{def:target}, and $\idxm$ is a bijection on
each alphabet it indexes, so that every emitted operator list is totally
ordered and its index is recoverable from its semantic label.
\end{theorem}
\claimmeta{\metalabel{Uses}\mref{Def.}{def:target}, \mref{Def.}{def:homog}, \mref{Def.}{def:spec}, \mref{Def.}{def:well-formed-spec}, \mref{Def.}{def:compile}, \mref{Def.}{def:validator}
\metasep \metalabel{Assumes}\aref{A-SPEC}{ass:spec}
\metasep \metalabel{Used by}\mref{Thm.}{thm:measurements-and-statistics}
\metasep newly derived
\metasep accepted
\metasep \metalabel{Code}\coderef{https://github.com/fraunhofersg/certiqs-key/blob/v0.1.0/src/certiqs_key/protocol_compiler/compile.py}{certiqs_key.protocol_compiler.compile:compile_specification}
\metasep \metalabel{Oracles}\coderef{https://github.com/fraunhofersg/certiqs-key/blob/v0.1.0/tests/oracles/registry_protocol_compiler.py}{registry_protocol_compiler.py}}

\begin{proof}
D(i): each system of $\Spec$ is declared finite-dimensional and each alphabet
finite and non-empty (Definition~\ref{def:well-formed-spec}\,(i), (ii)); $\Qsys$ is finite-dimensional
because $d_{\Qsys}=\sum_{s\in\SecSet}d_s$ is a finite sum of finite terms. The
emitted maps are $a_A$, $a_B$, $\phi$, $\rd{\mathsf A}$, $\rd{\mathsf B}$ ---
total by Definition~\ref{def:spec}\,(d) --- and $r$, which is total because
\eqref{eq:rdef} is defined by a case distinction on the value of the total map
$\varsigma\circ a$, both branches of which are defined on all of
$\Xset\times\Vset$. The sector-map clause: the emitted map $\Smap$ is the
target descriptor of the last link of $\Chrec$
(Definition~\ref{def:compile}), and it is emitted once; every emitted
$\Qsys$-operator is a block list against \emph{that} object, because
Definition~\ref{def:compile} fills $\{F_v\}$ from that same link's target POVM
and $\CList$ from lists each carried against the map its \exttermref{../foundations/contract.pdf}{term-declarer}{declarer} wrote
it against. What is not automatic, and is where the hypothesis of the theorem
is used, is that each declarer's map \emph{is} $\Smap$: a family presented
against a different block kind, a different identifier list, a different order
or a different dimension is a well-typed object of the wrong type, and check
V1 rejects exactly that, so a declaration that is emitted at all is one for
which V1 did not refuse. $d_{\Qsys}=\sum_sd_s$ is then the identity that
defines $d_{\Qsys}$ from the emitted map. The last clause of D(i), that the
restriction of $\idxm$ to $\SecSet$ is the order $\Smap$ fixes, holds
\emph{by construction}: Definition~\ref{def:compile} takes that component of
$\idxm$ to be the sector map's own order and not a second order. Check V1 of
Definition~\ref{def:validator} re-decides it on the emitted data, a condition
on an emitted object being decided on that object and never inferred from how
it was built, and a declaration that is emitted at all is one for which V1 did
not refuse.

D(ii): the marginal slot holds either a matrix or $\nomarg$, and no third value
is emitted (Definition~\ref{def:compile}). The three numerical conjuncts are
decided on the emitted matrix and on nothing else: check V2 of
Definition~\ref{def:validator} decides $\tau_A=\tau_A^{\dagger}$,
$\tau_A\succeq0$ and $\tr\tau_A=1$, each exactly, and decides that the trace
sum \exttermref{../foundations/contract.pdf}{term-rounds-nowhere}{rounds nowhere}; and a declaration that is emitted at all is one for
which V2 did not refuse, so D(ii) holds of the declared numbers. This is the
only route, a condition on the emitted matrix being decided exactly on that
matrix and never inferred from the certificate
(\xref{foundations/req:boundary-admission}\,(ii); for the trace sum
\xref{foundations/req:exact-emission}\,(v)): an emitted declaration
satisfies D(ii) even where a certificate is malformed. The certificate
supplies a \emph{modelling} statement about the same matrix, not a substitute
for that decision: in trusted-source mode $\tau_A$ is the certificate's first
component, and (SR1), decided on the emitted data by gate (G2), relates it to
$\tr_{\Avirt}\Psi$, which is positive semidefinite with trace $\tr\Psi=1$.
That is what makes the emitted marginal \emph{the certificate's} marginal
rather than an independently declared number; it is not what makes D(ii)
hold.

D(iv): the range of \eqref{eq:rdef} is contained in
$\operatorname{ran}r^0_{\hold}\cup\{\perp\}\subseteq\Zval\sqcup\{\perp\}$,
because $r^0_{\hold}$ maps into $\Zval$ by
Definition~\ref{def:spec}\,(d).

D(vi): the emitted list is the concatenation of the specification's entries
with the entries transported by the links of both segments; each entry is
emitted with its relation symbol and the datum its right-hand-side slot holds
unchanged, and with the tag that datum's type fixes
(Definition~\ref{def:homog}). Nothing in the hypotheses on $\Spec$ prevents a
link from transporting an entry whose declared form does not match the
emitted source mode --- an entry carrying a matrix where the mode admits only
a scalar, or one whose datum is of neither admitted type --- and that is a
property of the declared data, not of the compilation map; so this clause is where the hypothesis of
the theorem is used: check V6 of Definition~\ref{def:validator} rejects
exactly such an entry, and a declaration that is emitted at all is one for
which V6 did not refuse.

$\idxm$: each component is chosen as a bijection onto an initial segment of
$\mathbb{Z}_{\ge0}$, which exists because the alphabet is finite and is
injective and surjective by construction; the inverse table is emitted with it,
so the semantic label of an index is recoverable.

\emph{The general format.} For a declaration the general-format clauses of
Definition~\ref{def:compile} emit, the four conditions are read as in the
general-format paragraph of Definition~\ref{def:target}, and the argument is
the one above with the checks read as in the general-format paragraph of
Definition~\ref{def:validator}; a declaration that is emitted at all is one
for which none of them refused. D(i): finiteness and totality as above, every
table $a_m$ of the transcript included; the factor list with at most one factor
tagged $A$ and then first, at least one tagged $\Qsys$ and one block kind, the
dimension identities, the agreement of the map carried by each per-factor
family, by each entry's operator and by Alice's declared data with the map of
its factor, and the restriction of $\idxm$ to each factor's sectors and to
$\mathbb{T}$ are exactly the clauses V1 decides in its general-format
reading, the product orders being derived by (DX1) and (DX3) and not further
conditions. D(ii): V2 in its general-format reading decides every Alice block
of $\tau_A$ Hermitian and positive semidefinite and their total trace equal to
$1$ with every addition rounding nowhere, and, where Alice's slot is in
\termref{term-ensemble-form}{ensemble form}, every block non-zero. D(iv): as above, $r^{0}_{\hold}$
being typed into $\Zval$ by Definition~\ref{def:spec}\,(d), read on the
derived $\Kann$ by (d$'$). D(vi): each entry is emitted with its scope, its
relation symbol and its datum, and with the tag that datum's type fixes or,
for the \termref{term-party-local}{party-local} form, the tag field it arrives with
(Definition~\ref{def:homog}); V6 in its general-format reading decides the
relation symbol by form, the datum against its tag, a general $\Theta_k$
Hermitian block by block along $\mathbb{T}$, and that the scope names a
measured factor of the map, and refuses an entry failing any of them. The
index map: each component --- one per factor alphabet, one per factor
\exttermref{../reductions/contract.pdf}{term-sector-set}{sector set} and, where an $A$ factor is emitted, one on $\mathbb{T}$ ---
is chosen as above, and its components on $\Vset$ and on $\SecSet$ are the
product orders, bijections because each factor's component is.
\end{proof}

% claim protocol-compiler/rem:rbuilt
\claimsummary{protocol-compiler/rem:rbuilt}{Built from the sifting predicate and the holder's key rule, the raw-key map meets D(v) and D(viii) by construction. The only residual obligation is that the key rule map into the key values, a typing condition on one function table that \hyperlink{claim-protocol-compiler/ass:spec}{A-SPEC} carries.}
\claimbookmark{protocol-compiler/rem:rbuilt}{Remark}{rem:rbuilt}{The raw-key map built from sifting and the key rule}
\begin{remark}[The raw-key map built from sifting and the key rule]\label{rem:rbuilt}
A specification that supplied $r$ directly would have to be \emph{checked}
for D(v) and D(viii). Built as \eqref{eq:rdef}, $r$ satisfies both by
construction (Theorem~\ref{thm:sifting-and-key}), and the only residual obligation is the
side condition that $r^{0}_{\hold}$ maps into $\Zval$ rather than into
$\Zset$, which is a typing condition on the specification, decidable by
inspection of one function table and carried as a clause of
Assumption~\ref{ass:spec} (Remark~\ref{rem:sidecond}).
\end{remark}
\claimmeta{\metalabel{Uses}\mref{Def.}{def:compile}
\metasep newly derived
\metasep accepted}

% claim protocol-compiler/def:phi
\claimsummary{protocol-compiler/def:phi}{The joint outcome distribution of a declaration is given by measurement traces, and the acceptance statistic groups it through the public maps. Both functionals are affine on all operators and are probability distributions on every feasible state.}
\claimbookmark{protocol-compiler/def:phi}{Definition}{def:phi}{Declared statistics functional}
\claimleadin{%
The two functionals below are defined on every operator on $A\Qsys$, and not
only on positive ones, because affinity is what
Theorem~\ref{thm:measurements-and-statistics}\,(iv) and Lemma~\ref{lem:weighted}\,(iii) use and
neither needs positivity. The consuming engine minimises over the same set:
its own definition of the \termref{term-feasible-set}{feasible set} uses
Definition~\ref{def:feasible}.
}
\begin{definition}[Declared statistics functional]\label{def:phi}
For a declaration $\Decl$ and $\omega\in\Lin{A\Qsys}$ put
\begin{equation}\label{eq:jointdist}
  P_{\Decl}[\omega](x,v)\coloneqq\tr\!\big[(M_x\otimes F_v)\,\omega\big],
  \qquad
  \Phimap[\omega](\hat c)\coloneqq
  \sum_{(x,v):\,\phi(a(x,v))=\hat c}P_{\Decl}[\omega](x,v).
\end{equation}
Both are defined by trace formulas and are therefore defined, and affine, on
all of $\Lin{A\Qsys}$. The \termref{term-feasible-set}{feasible set} $\Feas$ of
Definition~\ref{def:feasible} is contained in $\Pos{A\Qsys}$, and on it the
values are a probability distribution
(Theorem~\ref{thm:measurements-and-statistics}\,(iv), and (iv$'$) in the
general format).
\end{definition}
\claimmeta{\metalabel{Uses}\mref{Def.}{def:feasible}
\metasep \metalabel{Used by}\mref{Thm.}{thm:measurements-and-statistics}, \mref{Thm.}{thm:choice-encoding}, \mref{Rem.}{rem:alphabet}
\metasep newly derived
\metasep implemented (effective: accepted)
\metasep \metalabel{Code}\coderef{https://github.com/fraunhofersg/certiqs-key/blob/v0.1.0/src/certiqs_key/protocol_compiler/statistics.py}{certiqs_key.protocol_compiler.statistics:phi_distribution}
\metasep \metalabel{Oracles}\coderef{https://github.com/fraunhofersg/certiqs-key/blob/v0.1.0/tests/oracles/registry_protocol_compiler.py}{registry_protocol_compiler.py}}

% claim protocol-compiler/thm:measurements-and-statistics
\claimsummary{protocol-compiler/thm:measurements-and-statistics}{Declared measurement families are exact POVMs block by block, per factor and, in POVM form, per Alice sector in the general format; the product and denoted ensemble families are POVMs by theorem. So D(iii) holds, and the statistics are a distribution on every normalised positive block-list state.}
\claimbookmark{protocol-compiler/thm:measurements-and-statistics}{Theorem}{thm:measurements-and-statistics}{Exact measurement families and the statistics functional}
\begin{theorem}[Exact measurement families and the statistics functional]\label{thm:measurements-and-statistics}
Let $\Spec$ be a \termref{term-well-formed-specification}{well-formed specification}
(Definition~\ref{def:well-formed-spec}) satisfying
Assumption~\ref{ass:spec}, and
let $\Decl$ be emitted by Definition~\ref{def:compile} in the
\termref{term-basic-format}{basic format}. Then:
\begin{enumerate}[label=(\roman*),leftmargin=2.2em]
\item $\{M_x\}_{x\in\Xset}$ is a POVM on $A$ and, for each $s\in\SecSet$, the
block family $\{F_v^{(s)}\}_{v\in\Vset}$ is a POVM on $\Qsys_s$ --- each
predicate holding \emph{of the declared numbers}, decided exactly --- and each
of the completeness sums $\sum_xM_x$ and $\sum_vF_v^{(s)}$, formed in the
declared order, rounds nowhere, so that each computed sum \emph{is} the sum it
is written as (Definition~\ref{def:exact});
\item every $F_v$ and every $\hat O_k$ is sector-block-diagonal
\emph{by construction}: it is emitted as the list
$\big(F_v^{(s)}\big)_{s}$, respectively $\big(\hat O_k^{(s)}\big)_{s}$, and
the operator it denotes is $\bigoplus_sF_v^{(s)}$. The two formulations of
block-diagonality remain equivalent --- for an operator $F$ on $\Qsys$,
$\Pi_sF\Pi_{s'}=0$ for all $s\ne s'$ if and only if
$F=\bigoplus_s\Pi_sF\Pi_s$ --- and that equivalence is what licenses the
\emph{type}: a list of blocks denotes exactly the operators the vanishing
condition picks out, and no other. It is not checked, because the condition
has no data to be checked on;
\item hence D(iii) holds: its typing half by (ii) and check V1; its first
four predicates --- those on $\{M_x\}$ and on each block family
$\{F_v^{(s)}\}_v$ --- by (i); and its fifth, Hermiticity of
$\hat O_k^{(s)}$ for each $(k,s)$, by the last clause of check V3
(Definition~\ref{def:validator}), which decides it exactly on the emitted
constraint entries and which (i) does not state; and
\item $P_{\Decl}[\omega]$ of \eqref{eq:jointdist} and $\Phimap[\omega]$ are
affine in $\omega$, $\Phimap[\omega]$ depends on $\omega$ only through
$P_{\Decl}[\omega]$, and the fibres of $\phi\circ a$ partition
$\Xset\times\Vset$, so $\Phimap[\omega]$ is the pushforward of
$P_{\Decl}[\omega]$ and inherits its total mass. For every block list $\omega$
with $\omega^{(s)}\succeq0$ for each $s$ and $\tr\omega=1$ --- in particular
for every $\omega\in\Feas$ --- $P_{\Decl}[\omega]$ is a probability
distribution on $\Xset\times\Vset$ and $\Phimap[\omega]$ one on $\Chatset$.
\end{enumerate}

\emph{The general format.} Let $\Decl$ instead be emitted by the
general-format clauses of Definition~\ref{def:compile}, in the general format
of Definition~\ref{def:declaration}\,(DX1)--(DX6). Then (i)--(iv) hold with
the following readings, of which the four above are the case of the basic
format.
\begin{enumerate}[label=(\roman*$'$),leftmargin=2.6em]
\item for every \termref{term-measured-factor}{measured factor} $f$ and every $s_f\in\SecSet^{f}$ the
block family $\{F^{f,(s_f)}_{v_f}\}_{v_f\in\Vset^{f}}$ is a POVM on
$\Qsys_{f,s_f}$; and for every Alice sector $t$, in \termref{term-povm-form}{POVM form}
$\{M^{(t)}_x\}_{x\in\Xset}$ is a POVM on $A_t$, and in \termref{term-ensemble-form}{ensemble form}
every $N^{(t)}_x$ is Hermitian and positive semidefinite, with
$\sum_xN^{(t)}_x=\tau^{(t)}_A$ and $\tau^{(t)}_A\ne0$ --- each predicate
holding of the declared numbers, decided exactly, and each of these sums,
formed in its declared order, rounding nowhere. The product family
\eqref{eq:productfamily} is a POVM on $\Qsys$ and, in ensemble form, the
family \eqref{eq:ensemblepovm} is a POVM on $A$, both as identities of real
numbers, and neither is formed, decided or summed;
\item every $F^{f}_{v_f}$ is a block list against factor $f$'s sector map and
every $\hat O_k$ one against the sector map of its scope, and $\tau_A$, every
operator of Alice's measurement slot and every general $\Theta_k$ is a block
list along $\mathbb{T}$, each block diagonal by construction as in (ii); and
the product family is block diagonal for the product sector map,
$\Pi_sF_v\Pi_{s'}=0$ for $s\ne s'$, with the $s$-block of
\eqref{eq:productfamily};
\item hence D(iii) holds in its general-format reading
(Definition~\ref{def:target}); and
\item (iv) holds for every list $\omega$ over the
\termref{term-joint-sector}{joint sectors} with every block positive semidefinite and
$\tr\omega=1$ --- in particular for every $\omega\in\Feas$ --- with $F_v$ the
product \eqref{eq:productfamily} and, in ensemble form, $M_x$ the operator
\eqref{eq:ensemblepovm} of the declaration $\Decl$ denotes: $P_{\Decl}[\omega]$
is a probability distribution on $\Xset\times\Vset$ and $\Phimap[\omega]$ one
on $\Chatset$.
\end{enumerate}
\end{theorem}
\claimmeta{\metalabel{Uses}\mref{Def.}{def:declaration}, \mref{Def.}{def:exact}, \mref{Def.}{def:target}, \mref{Def.}{def:feasible}, \mref{Lem.}{lem:hermitian-constraint-values}, \mref{Def.}{def:well-formed-spec}, \mref{Def.}{def:compile}, \mref{Thm.}{thm:typing}, \mref{Def.}{def:phi}, \mref{Def.}{def:weighted}, \mref{Thm.}{thm:sifting-and-key}, \mref{Def.}{def:validator}
\metasep \metalabel{Assumes}\aref{A-SPEC}{ass:spec}
\metasep \metalabel{Used by}\mref{Thm.}{thm:choice-encoding}, \mref{Rem.}{rem:exactness}, \mref{Thm.}{thm:validator-soundness}, \mref{Thm.}{thm:chain-consumption}
\metasep newly derived
\metasep accepted
\metasep \metalabel{Code}\coderef{https://github.com/fraunhofersg/certiqs-key/blob/v0.1.0/src/certiqs_key/protocol_compiler/validation.py}{certiqs_key.protocol_compiler.validation:check_v3}
\metasep \metalabel{Oracles}\coderef{https://github.com/fraunhofersg/certiqs-key/blob/v0.1.0/tests/oracles/registry_protocol_compiler.py}{registry_protocol_compiler.py}}

\begin{proof}
(i) The predicates are decided on the emitted objects, not inferred
(\xref{foundations/req:exact-emission}): check V3 of
Definition~\ref{def:validator} decides, exactly, $M_x=M_x^{\dagger}$ and
$M_x\succeq0$ for each $x$, $\sum_xM_x=\idon{A}$ with the sum rounding
nowhere, and, per block, $F_v^{(s)}=F_v^{(s)\dagger}$, $F_v^{(s)}\succeq0$ and
$\sum_vF_v^{(s)}=\idon{\Qsys_s}$ with that sum rounding nowhere; a declaration
that is emitted at all is one for which V3 did not refuse. Two constructions in
this contract would have licensed an inference instead, and neither is used as
one. Where the specification presents a family as choice-weighted,
Lemma~\ref{lem:weighted}\,(i) shows the weighted family is a POVM whenever the
conditional families are; the emitted family's predicates are nonetheless
decided, because the emitted matrices are the products $q_tN_{u|t}$ as
represented, which the compiler forms exactly and then decides to be binary64
data before emitting them at all (Definition~\ref{def:weighted}), and because
the lemma says nothing about condition~(S) on the emitted family's own sum.
Where $\{F_v\}$ arrives as the target POVM of the last link of $\pcChain$,
completeness is likewise decided by V3 on the emitted blocks and never inferred
from the physical family; Theorem~\ref{thm:chain-consumption} records what the chain is and
is not asked to guarantee. No floating-point budget is therefore needed
anywhere below: nothing here depends on how much rounding a construction
introduced, because the result is decided.

(ii) By Definition~\ref{def:compile} the emitted $\{F_v\}$ and $\{\hat O_k\}$
are block lists against $\Smap$, so the first sentence is the emission rule.
For the equivalence: any $F\in\Lin{\Qsys}$ satisfies
$F=\sum_{s,s'}\Pi_sF\Pi_{s'}$, since $\sum_s\Pi_s=\idon{\Qsys}$ and the $\Pi_s$
are mutually orthogonal projectors; hence $\Pi_sF\Pi_{s'}=0$ for all $s\ne s'$
holds if and only if $F=\sum_s\Pi_sF\Pi_s=\bigoplus_s\Pi_sF\Pi_s$, i.e.\ if and
only if $F$ is the direct sum of its diagonal blocks. So the block lists denote
exactly the block-diagonal operators, which is what makes the type faithful:
nothing expressible as a list violates the condition, and nothing satisfying
the condition fails to be expressible as a list. This is a statement about the
representation, not a check on data. Where the specification presents
$\{F_v\}$ as a choice-weighted family, the weighting is applied blockwise,
$N_{u,t}^{(s)}=q_tN_{u|t}^{(s)}$, so the emitted object is a block list
whatever the presentation (Lemma~\ref{lem:weighted}\,(ii)).

(iii) is (i) and (ii) together with D(iii)'s statement. Its typing half is
(ii), and the sector-map agreement that half rests on is check V1
(Theorem~\ref{thm:typing}, D(i)). Of D(iii)'s five predicates, the four on
$\{M_x\}$ and on the block families $\{F_v^{(s)}\}_v$ are (i). The fifth is
not among them and is cited from where it is decided: D(iii) requires in
addition that $\hat O_k^{(s)}$ be Hermitian for each $(k,s)$, which (i) does
not state, and which the last clause of check V3
(Definition~\ref{def:validator}) decides exactly on the emitted constraint
entries --- a declaration that is emitted at all being one for which V3 did
not refuse, which is the route (i)'s own four take
(Theorem~\ref{thm:typing}). It is among the decisions \eqref{eq:predcount}
counts, as that equation's $m\,|\SecSet|$ term, and it is the hypothesis of
Lemma~\ref{lem:hermitian-constraint-values}, which makes
\eqref{eq:feas-modes} well posed.

(iv) Affinity of both maps is immediate from linearity of the trace in
$\omega$; that $\Phimap[\omega]$ depends on $\omega$ only through
$P_{\Decl}[\omega]$ is the definition \eqref{eq:jointdist}. Since $\phi\circ a$
is a total map $\Xset\times\Vset\to\Chatset$ (Theorem~\ref{thm:typing}, D(i)), the
fibres $\{(x,v):\phi(a(x,v))=\hat c\}$ partition $\Xset\times\Vset$, so
$\Phimap[\omega]$ is the pushforward of $P_{\Decl}[\omega]$ and has the same
total mass.

For the distribution statement, let $\omega$ be a block list with
$\omega^{(s)}\succeq0$ for each $s$ and $\tr\omega=1$. By (i) each $M_x$ and
each $F_v^{(s)}$ is Hermitian with non-negative spectrum, hence
$M_x\succeq0$ and $F_v=\bigoplus_sF_v^{(s)}\succeq0$, positivity of a
block-diagonal operator being positivity of each of its blocks; and
$\omega=\bigoplus_s\omega^{(s)}\succeq0$ for the same reason. Then
non-negativity and realness together, since $M_x\otimes F_v\succeq0$ and
$\omega\succeq0$ give $\tr[(M_x\otimes F_v)\omega]\ge0$ --- the trace of a
product of two positive semidefinite operators is a non-negative \emph{real}
number. And normalization:
\begin{equation*}
  \sum_{x,v}P_{\Decl}[\omega](x,v)
  =\tr\!\Big[\Big(\sum_xM_x\otimes\sum_vF_v\Big)\omega\Big]
  =\tr[(\idon{A}\otimes\idon{\Qsys})\,\omega]=\tr\omega=1 ,
\end{equation*}
by bilinearity of the tensor product and the two completeness identities of
(i). This step is where condition~(S) is used: what (i) delivers is
$\sum_xM_x=\idon{A}$ and $\sum_vF_v^{(s)}=\idon{\Qsys_s}$ as identities of the
declared numbers, and it delivers them only because the sums forming them round
nowhere --- a computed sum equal to the identity would not license this
display. Hence $P_{\Decl}[\omega]$ is a probability distribution, and so is its
pushforward $\Phimap[\omega]$.

\emph{The general format.} (i$'$) Check V3 in its general-format reading
(Definition~\ref{def:validator}) decides, exactly: per measured factor,
$F^{f,(s_f)}_{v_f}$ Hermitian and positive semidefinite for every
$(v_f,s_f)$ and, for each $s_f$,
$\sum_{v_f}F^{f,(s_f)}_{v_f}=\idon{\Qsys_{f,s_f}}$ with its declared-order sum
rounding nowhere; on Alice's data, block by block along $\mathbb{T}$, in POVM
form $M^{(t)}_x$ Hermitian and positive semidefinite for every $(x,t)$ and
$\sum_xM^{(t)}_x=\idon{A_t}$ for every $t$ with that sum rounding nowhere,
and in ensemble form $N^{(t)}_x$ Hermitian and positive semidefinite for
every $(x,t)$ and $\sum_xN^{(t)}_x=\tau^{(t)}_A$ entry by entry with that sum
rounding nowhere; and check V2 in its general-format reading decides
$\tau^{(t)}_A\ne0$ for every $t$ where the slot is in ensemble form. A
declaration that is emitted at all is one for which neither check refused, so
each of these holds of the declared numbers. The product family: every
$F^{f}_{v_f}=\bigoplus_{s_f}F^{f,(s_f)}_{v_f}$ is then Hermitian and positive
semidefinite and $\sum_{v_f}F^{f}_{v_f}=\idon{\Qsys_f}$, an identity of the
declared numbers because each of its blocks' sums rounds nowhere; a tensor
product of Hermitian positive semidefinite operators is Hermitian and
positive semidefinite, and
$\sum_{v}F_v=\sum_{v}\bigotimes_fF^{f}_{v_f}=\bigotimes_f\sum_{v_f}F^{f}_{v_f}
=\bigotimes_f\idon{\Qsys_f}=\idon{\Qsys}$ by multilinearity, the sum running
over $\Vset=\prod_f\Vset^{f}$ (Definition~\ref{def:declaration}\,(DX3)). The
denoted family: Theorem~\ref{thm:sifting-and-key}\,(v), whose hypothesis is the
conditions Definition~\ref{def:exact} states for the ensemble form, which are
those just decided. Neither family is a datum
(Requirement~\ref{req:derived-objects}): nothing of either is formed, decided
or summed, so no summation condition arises for either, and both statements
are identities of the real numbers the declared data denote.

(ii$'$) By the general-format clauses of Definition~\ref{def:compile} each
$F^{f}_{v_f}$ is emitted as a block list against factor $f$'s map and each
$\hat O_k$ against the map of its scope, and by
Definition~\ref{def:declaration}\,(DX2) Alice's data are block lists along
$\mathbb{T}$; for each such list the argument of (ii), with the projectors
$\Pi^{f}_{s_f}$, respectively $\Pi^{A}_t$, in place of $\Pi_s$, gives the
equivalence. For the product family, $\Pi_s=\bigotimes_f\Pi^{f}_{s_f}$ by
\eqref{eq:productmap}, so
$\Pi_sF_v\Pi_{s'}=\bigotimes_f\Pi^{f}_{s_f}F^{f}_{v_f}\Pi^{f}_{s'_f}$; for
$s\ne s'$ some factor has $s_f\ne s'_f$, and its tensor factor vanishes because
$F^{f}_{v_f}$ is block diagonal for factor $f$'s map; and
$\Pi_sF_v\Pi_s=\bigotimes_fF^{f,(s_f)}_{v_f}$ is the block of
\eqref{eq:productfamily}.

(iii$'$) D(iii) in its general-format reading (Definition~\ref{def:target})
asks for the typing and for the predicates of Definition~\ref{def:exact} in
their general-format reading, per measured factor and per Alice sector, with
the ensemble-form conditions in place of those on $\{M_x\}$ where that form is
declared, and it refuses ensemble form in untrusted-source mode. The typing
is (ii$'$), the agreement of each list's map with the map of its factor being
decided by check V1 in its general-format reading (Theorem~\ref{thm:typing},
D(i)); the predicates and summation conditions on the per-factor families and
on Alice's data are (i$'$); Hermiticity of $\hat O^{(s_{f_k})}_k$ for every
entry and every block of its scope is the last clause of check V3 in its
general-format reading; and an ensemble-form slot beside a marginal slot
holding $\nomarg$ is refused by that check's clause on which slot is filled.
A declaration that is emitted at all is one for which V1 and V3 did not
refuse.

(iv$'$) Affinity and the pushforward are as in (iv), with $a$ the function
\eqref{eq:transcript} derives, total into $\Kann$
(Theorem~\ref{thm:sifting-and-key}\,(iv)), and $\phi$ total on $\Kann$
(Theorem~\ref{thm:typing}, D(i)). For the distribution statement let $\omega$
be such a list. $M_x\succeq0$: in POVM form $M_x=\bigoplus_tM^{(t)}_x$ with
every block positive semidefinite by (i$'$), and in ensemble form by
Theorem~\ref{thm:sifting-and-key}\,(v); $F_v\succeq0$ by (i$'$); and
$\omega=\bigoplus_{(t,s)}\omega^{(t,s)}\succeq0$, positivity of a
block-diagonal operator being positivity of each block. So
$\tr[(M_x\otimes F_v)\omega]$ is a non-negative real number. And
$\sum_xM_x=\bigoplus_t\sum_xM^{(t)}_x=\bigoplus_t\idon{A_t}=\idon{A}$, in POVM
form an identity of the declared numbers by (i$'$) and in ensemble form an
identity of real numbers by Theorem~\ref{thm:sifting-and-key}\,(v), while
$\sum_vF_v=\idon{\Qsys}$ by (i$'$); so the normalization display of (iv) holds
with these operators and gives $\sum_{x,v}P_{\Decl}[\omega](x,v)=\tr\omega=1$.
Condition~(S) is used there per measured factor and, in POVM form, per Alice
sector, which is what (i$'$) delivers; no sum over the product alphabet and no
sum of the denoted family is computed, so none needs a summation condition.
\end{proof}

% claim protocol-compiler/rem:phitotal
\claimsummary{protocol-compiler/rem:phitotal}{Were the acceptance-statistic map partial, its fibres would not partition the outcome pairs and the acceptance statistic would lose mass. Totality is therefore decided by check V1, not assumed.}
\claimbookmark{protocol-compiler/rem:phitotal}{Remark}{rem:phitotal}{Totality of the acceptance-statistic map}
\begin{remark}[Totality of the acceptance-statistic map]\label{rem:phitotal}
If $\phi$ were partial, the fibres would not partition $\Xset\times\Vset$: the
total mass of $\Phimap[\omega]$ would fall short of that of
$P_{\Decl}[\omega]$ by the mass of the unhandled classes, for some $\omega$ ---
a lost probability, and an acceptance statistic that does not sum to one.
Totality is therefore decided by check V1 of Definition~\ref{def:validator},
not assumed, and a specification with an unhandled announcement class is
refused (Requirement~\ref{req:validator-refusals}).
\end{remark}
\claimmeta{newly derived
\metasep accepted}

\section{Public choices as measurement weights}\label{group-protocol-compiler/def:weighted}
\sectionintro{A public choice of known bias enters the declaration as a weight multiplying a complete conditional measurement, either independently of every setting or, at a \termref{term-measured-factor}{measured factor} of either party after its chain, through a stochastic matrix read off the private setting of the station that measures it. The main result, \hyperlink{claim-protocol-compiler/thm:choice-encoding}{Theorem~\ref*{thm:choice-encoding}}, shows that the encoding keeps the family a measurement, that the \termref{term-setting-independent-form}{setting-independent form} reproduces the declared weights exactly wherever it is the emitted family, that the setting-dependent relabelling commutes with the reduction chain, and that no separate fixed fraction is ever emitted. Each declared conditional family is decided before any product is formed; the remarks and the example say why that decision cannot be left to the validator, where the relabelling is applied, and why the outcome alphabet is chosen upstream. A factor enclosure of a delivered block survives the setting-dependent weighting, the weight going into its core (\hyperlink{claim-protocol-compiler/lem:weighted-factor}{Lemma~\ref*{lem:weighted-factor}}).}
\sectionimports{\extref{../foundations/contract.pdf}{foundations/def:binary64}, \extref{../foundations/contract.pdf}{foundations/def:factor-enclosure}}

% claim protocol-compiler/def:weighted
\claimsummary{protocol-compiler/def:weighted}{A public choice enters as a weight on a complete conditional measurement: one declared weight per choice value, or, at a measured factor after its chain, a stochastic weight read through its station's setting. Weights and families are binary64 numbers read exactly; products are formed exactly.}
\claimbookmark{protocol-compiler/def:weighted}{Definition}{def:weighted}{Choice-weighted families}
\claimleadin{%
A per-round public choice of known bias --- a basis choice, a choice between
test and generation rounds, a probe setting --- is encoded as a weight inside
the measurement operators and never as a separate fraction. There are two
constructions: one declared weight per choice value, independent of every
setting, and, at a \termref{term-measured-factor}{measured factor} after the \exttermref{../reductions/contract.pdf}{term-reduction-chain}{reduction chain}, a weight
that depends on the private setting of the station that measures it.

Assumption~\ref{ass:coin} and Remark~\ref{rem:relabel} are why a
setting-dependent weight is read through the station's setting and through
nothing else, and Remark~\ref{rem:relabel} and Theorem~\ref{thm:choice-encoding} why
that form is sound after the chain and is not a link of it. Which conditions
of the weighted family do and do not follow from the decision on the
conditional families, and why that decision is taken here rather than among
the target conditions, is Remark~\ref{rem:condfamily}; that the weighted
family can satisfy all five while none of the conditional families it was
built from is a POVM is Example~\ref{ex:condfamily}. The order in which the
weights are summed is declared with the public-choice record, as
\xref{foundations/req:declared-order} requires of every summed predicate, and
a presentation indexed by a product alphabet is refused under
Requirement~\ref{req:derived-objects}.
}
\begin{definition}[Choice-weighted families]\label{def:weighted}
Let $S$ be a finite-dimensional system, $\Tpub$ a finite non-empty
\termdef{term-public-choice-alphabet}{public-choice alphabet}
with declared weights $q_t>0$, $\sum_{t\in\Tpub}q_t=1$, and for each
$t\in\Tpub$ a POVM $\{N_{u|t}\}_{u\in\Uset_t}$ on $S$. The
\termdef{term-choice-weighted-family}{choice-weighted family}
is indexed by the disjoint union
$\Vset\coloneqq\bigsqcup_{t\in\Tpub}\Uset_t\times\{t\}$ and defined by
\begin{equation}\label{eq:weighted}
  N_{u,t}\coloneqq q_t\,N_{u|t},\qquad (u,t)\in\Vset .
\end{equation}
Several independent choices compose by iterating \eqref{eq:weighted}; the
resulting weight is the product of the individual weights. \emph{What is
declared, and therefore what is decided, is the innermost family and each
weight list}: an iterated presentation declares $\{N_{u|t}\}_u$ for each value
of each choice together with one weight list per choice, the decision on the
conditional families is taken on those, and the intermediate
families are this module's own products, which inherit admissibility by
Lemma~\ref{lem:weighted}\,(i) at each step and carry no second declared
condition; the representability decision is taken on every product the map
forms, the intermediate ones included. This is the
\termdef{term-setting-independent-form}{setting-independent form}:
the weight of a choice value is one declared number, the same in every round
and whatever else either party's apparatus did.

\emph{A weighting that depends on a station's setting.} Let
$\{F_v\}_{v\in\Vset}$ be a POVM on $S$ over a finite alphabet $\Vset$ with a
\exttermref{../foundations/contract.pdf}{term-declared-order}{declared order} --- the \exttermref{../reductions/contract.pdf}{term-target-family}{target family} of the last link of $\Chrec$, as
Definition~\ref{def:compile} takes it --- and let
$t:\Vset\to\mathcal{T}$ be a declared total map onto a finite non-empty
\termdef{term-setting-alphabet}{setting alphabet},
the
\termdef{term-setting-component}{setting component}
of the outcome: the value of the private per-round setting of the measuring
apparatus that the outcome $v$ records. Where the chain was emitted by the reductions module, $t$ is the
class map of that module's \exttermref{../reductions/contract.pdf}{term-joint-alphabet}{joint alphabet}
(\xref{reductions/def:joint-family}\,(a): $\Vset=\mathcal{Y}\times\mathcal{T}$
with classes $\Vset_t$); where the family is declared directly, $t$ is
declared with it, the singleton $\mathcal{T}=\{\ast\}$ being the case of no
setting. Let $\mathcal{F}$ be a finite non-empty
\termdef{term-choice-alphabet}{choice alphabet}
with a declared order and let $P(f\,|\,t)$,
$(t,f)\in\mathcal{T}\times\mathcal{F}$, be a \emph{stochastic matrix}:
$P(f\,|\,t)\ge0$ and $\sum_{f\in\mathcal{F}}P(f\,|\,t)=1$ for every
$t\in\mathcal{T}$. The
\termdef{term-support-of-the-weighting}{support of the weighting}
is
\begin{equation}\label{eq:support}
  \Vset_P\coloneqq\big\{(v,f)\in\Vset\times\mathcal{F}:\ P(f\,|\,t(v))>0\big\},
\end{equation}
ordered by the product order --- $v$ in the declared order of $\Vset$ and,
within each $v$, $f$ in the declared order of $\mathcal{F}$, the pairs off
the support omitted --- and the choice-weighted family in its
\termdef{term-setting-dependent-form}{setting-dependent form}
is
\begin{equation}\label{eq:weighted-t}
  N_{v,f}\coloneqq P(f\,|\,t(v))\,F_v,\qquad (v,f)\in\Vset_P .
\end{equation}
The setting-independent form is the case $P(f\,|\,t)=q_f$ for every $t$, in
which \eqref{eq:weighted-t} is \eqref{eq:weighted} with
$\Tpub=\mathcal{F}$ and $N_{v|f}=F_v$ for every $f$ --- one conditional
family, weighted by a choice that does not see the setting; conversely,
\eqref{eq:weighted} with conditional families that differ with the choice is
not of the form \eqref{eq:weighted-t}, and the two forms are two constructions
overlapping in that case. \emph{In the overlap the two constructions do not
agree on an order, and the order is a property of the presentation and not of
the family}: the alphabet the first form produces is
$\bigsqcup_f\Uset_f\times\{f\}$, listed choice-major, while the support
\eqref{eq:support} of the second is listed outcome-major, so the two are
transposes of one another and one family declared through the two records has
its condition~(S) decided in two different orders. (S) is order-dependent by
construction (Definition~\ref{def:exact}), and the two orders are not
reconciled: what fixes the order is the record that was declared, and a
\exttermref{../foundations/contract.pdf}{term-declarer}{declarer} that declares the same family both ways must expect the
decision its own record's order gives. \emph{$P$ is read through the setting
and through nothing else.} It is declared on $\mathcal{T}\times\mathcal{F}$
and reaches an outcome only through $t(v)$. Zero entries of $P$ are admitted, and the support
is what they are for: a pair $(v,f)$ with $P(f\,|\,t(v))=0$ is not an outcome,
is not emitted and indexes nothing, so every emitted element of
\eqref{eq:weighted-t} carries a strictly positive weight, exactly as $q_t>0$
requires of every element of \eqref{eq:weighted}. Which choice values a given
setting admits is therefore declared by the zero pattern of $P$ and not by a
zero-weight outcome. The
form is applied by the compiler \emph{after} the reduction chain, to the
family the last link delivered (Definition~\ref{def:compile}).

\emph{The reading of each declared condition.} The weights are
\exttermref{../foundations/contract.pdf}{term-binary64-numbers}{binary64 numbers} and the conditional families binary64 matrices. So
$\sum_tq_t=1$ is decided
\emph{exactly}, with the sum accumulated in the declared order of $\Tpub$ ---
the order in which the declared public-choice record lists $\Tpub$, fixed
there because $\idxm$ does not index $\Tpub$
(Definition~\ref{def:compile}) --- and
required to \exttermref{../foundations/contract.pdf}{term-rounds-nowhere}{round nowhere}; and ``$\{N_{u|t}\}_u$ is a
POVM on $S$'' is exact admissibility in the sense of
Definition~\ref{def:exact}, per block where $S$ carries sectors: (H), (P) and
(C) per block; (S) on the completeness sum $\sum_{u\in\Uset_t}N_{u|t}$,
accumulated in the order stated below and required to round nowhere; and (T)
as the agreement --- in block kind, in block identifiers, in their order and in
their dimensions --- of the maps \emph{the conditional families of one
presentation themselves carry}: there is one such map for the presentation,
and the family the weighting produces is a \exttermref{../foundations/contract.pdf}{term-block-list}{block list} against
\emph{it}. On a system carrying no sector decomposition, Alice's $A$, (T) is
empty. This condition is decided of each $t\in\Tpub$ separately, on the
declared numbers.

\emph{In which order (S) is read on a conditional family.} The family the
weighting produces is indexed by one alphabet, and that alphabet carries one
order: $\Xset$ where the presentation is Alice's, which arrives with its
declared order (Definition~\ref{def:spec}\,(a)); the physical receiver
alphabet where it is a measuring party's, which is the order the fibre sum of
Definition~\ref{def:chaincond}\,(C4) is accumulated in and, in \termref{term-band-a}{Band A}, is
the emitted alphabet's own. Condition~(S) on the conditional family
$\{N_{u|t}\}_u$ is decided in the \emph{restriction of that order} to the
elements the choice value $t$ contributes --- a total order on $\Uset_t$
whatever the presentation listed, so no compatibility of the record's listing
with it is required and none is decided. Where a declared order arrives with
that alphabet the compiler takes it, by the binding rule of
Definition~\ref{def:compile}, to which this definition adds nothing; where
none arrives, the choice the compiler makes and emits is the record's own
listing, $t$ in the order it lists $\Tpub$ in and, within each $t$, $u$ in the
order it lists $\Uset_t$ in. \emph{No presentation is refused for an order}:
what is fixed here is which order a decision is read in, not what a declarer
may declare.

\emph{The remaining declared conditions, at the same reading.} The strict
inequality $q_t>0$ is a domain condition on a declared real scalar, decided by
its sign, and Definition~\ref{def:exact} carries no
strict-inequality row because none is wanted --- $q_t\ne0$ is what
Lemma~\ref{lem:weighted}\,(ii) needs, and a negative $q_t$ would reverse the
sign of an element of the emitted family, which no public choice of known bias
can mean. For the setting-dependent form the same reading applies to $P$: each
entry is a binary64 number, $P(f\,|\,t)\ge0$ is decided by sign, each row sum
$\sum_fP(f\,|\,t)=1$ is decided exactly in the declared order of
$\mathcal{F}$ and required to round nowhere, the support \eqref{eq:support}
is decided by the sign of each entry, and ``$\{F_v\}$ is a POVM on $S$'' is
the exact admissibility the chain's (C3) already decided of the family the
last link delivered. \emph{This form has no second conditional family to
decide}: the choice value is drawn after the measurement, $\{F_v\}$ is the
whole of what the decision on the conditional families has as its subject
here, and (C3) is where it was taken.

\emph{What this module constructs, and the two emission decisions.} The
products $q_tN_{u|t}$ of \eqref{eq:weighted} and $P(f\,|\,t(v))F_v$ of
\eqref{eq:weighted-t} are formed \emph{exactly} on the declared dyadic
numbers. A product of two binary64 numbers need not be one, so the compiler
decides, before emitting anything, that every entry of every $N_{u,t}$,
respectively of every $N_{v,f}$, is a binary64 number --- on every product
this map forms, whether the product is emitted or handed on --- and that the
emitted family satisfies Definition~\ref{def:exact} --- (H), (P), (C) and, on
its own declared-order sum over $\Vset$, respectively over $\Vset_P$ in its
product order, (S). These are the two \emph{emission decisions}.

\emph{The lattice of the products.} Let the arriving entries be integer
multiples of $2^{-k}$ of modulus at most $1$, and let each $q_t$ be an
integer multiple of $2^{-k_q}$; then every entry of $q_tN_{u|t}$ is an integer
multiple of $2^{-(k+k_q)}$ of modulus at most $1$, so the constructed family
lies on a lattice $k_q$ steps finer than the one it was handed and is exactly
binary64 whenever $k+k_q\le53$. Where a $q_t$ is not a power of two the finer
lattice can leave binary64 altogether: the decimal literal $0.9$ is an integer
multiple of $2^{-53}$ and of no coarser one, so $k+k_q>53$ for every $k\ge1$
and the exact product of a declared entry by it need not be a binary64
number. The same arithmetic governs $P(f\,|\,t(v))F_v$ with the exponent
$k_P$ of the entries of $P$ in place of $k_q$: a $P$ whose non-zero entries
are dyadic with small exponent --- $\tfrac12$, $\tfrac14$, $1$ --- composes
with any family the chain delivers at $k\le53-k_P$, and a decimal literal in
$P$ can leave binary64 exactly as a decimal $q_t$ can.

\emph{The general format.} In the general format of
Definition~\ref{def:declaration}\,(DX1)--(DX6) both forms above are read
\emph{on the system of the party that declares them}, and nothing else of
this definition changes. At one \termref{term-measured-factor}{measured factor}, in \termref{term-povm-form}{POVM form}
and with no factor tagged $A$, the readings below are the text above,
sentence for sentence.

\emph{Which system $S$ is.} For a presentation of Alice's it is $A$, and (T)
is \emph{not empty where a factor is tagged $A$}: the conditional families
are then block lists along $\mathbb{T}$ and (T) is the four-component
agreement of the maps they carry with that factor's map, as on any other
sectored system. For a presentation of a measuring party's it is that party's
\emph{own} measured factor: one presentation per factor, decided against that
factor's own map; \emph{a presentation indexed by a product alphabet has no
declared system to be a family on} --- the product is the rule
\eqref{eq:productfamily} and is no datum (Definition~\ref{def:compile}, the
general format). So the decision ``$\{N_{u|t}\}_u$ is
a POVM on $S$'' and the products this module forms are per factor, and the
representability decision keeps exactly the subject it has: the products
\emph{these two constructions} form, and no product block of
\eqref{eq:productfamily}.

\emph{The \termref{term-ensemble-form}{ensemble form} of Alice's slot.} Where Alice's measurement slot
is declared in ensemble form (DX4), a choice-weighted presentation of it
declares, for each $t\in\Tpub$, a family $\{N_{x|t}\}_{x\in\Xset}$ of block
lists along $\mathbb{T}$ with every block Hermitian and positive
semidefinite, and \emph{in place of the POVM condition} the ensemble
condition of \eqref{eq:ensembleform}: for every Alice sector $t'$,
$\sum_{x}N^{(t')}_{x|t}=\tau^{(t')}_A$, the sum accumulated in the order
stated above and rounding nowhere. That condition is decided \emph{of each
$t$ separately and on the declared numbers}, with (H), (P) per block, that
identity in place of (C), and (T) as above. \emph{Why the
construction is the same construction.} With $N_{x,t}\coloneqq q_tN_{x|t}$ of
\eqref{eq:weighted}, the exact real numbers the data denote satisfy
\begin{equation}\label{eq:weighted-ensemble}
  \sum_{(x,t)\in\Vset}N^{(t')}_{x,t}
  =\sum_{t\in\Tpub}q_t\sum_{x\in\Xset}N^{(t')}_{x|t}
  =\Big(\sum_{t\in\Tpub}q_t\Big)\tau^{(t')}_A=\tau^{(t')}_A
\end{equation}
for every $t'$, by the decision just taken and by $\sum_tq_t=1$; so a
choice-weighted presentation in ensemble form whose conditional ensembles are
admitted denotes a weighted family with the ensemble identity of
\eqref{eq:ensembleform}. \emph{That is a consistency statement and not a
substitute for the decision on the emitted family}: the emitted family's own
(S), (H), (P) and its identity with the declared $\tau_A$ are decided at check
V3, on the emitted blocks and in the emitted order, and neither condition
implies the other: \eqref{eq:weighted-ensemble} is an identity of real
numbers and the emitted decision is about the additions that compute it. In
\exttermref{../transmitter/contract.pdf}{term-untrusted-source-mode}{untrusted-source mode} there is no ensemble form and this paragraph has
no subject.

\emph{The setting-dependent form, per measured factor.} It is applied by the
compiler after \emph{that factor's} \exttermref{../reductions/contract.pdf}{term-receiver-segment}{receiver segment}, to the family that
factor's last link delivered, with that factor's own setting map, support
\eqref{eq:support} and product order; the setting component read is a
component of \emph{that factor's} outcome, the private per-round setting of
the station that measures that factor, whichever party's label the factor
carries. This form has no second
conditional family here either, per factor: (C3) of
Definition~\ref{def:chaincond}, read for that factor, is where the delivered
family was decided.
\end{definition}
\claimmeta{\metalabel{Uses}\mref{Def.}{def:declaration}, \mref{Def.}{def:exact}, \mref{Def.}{def:compile}
\metasep \metalabel{Used by}\mref{Def.}{def:well-formed-spec}, \mref{Thm.}{thm:measurements-and-statistics}, \mref{Lem.}{lem:weighted}, \mref{Lem.}{lem:weighted-factor}, \mref{Thm.}{thm:choice-encoding}, \mref{Rem.}{rem:exactness}, \mref{Rem.}{rem:condfamily}, \mref{Ex.}{ex:condfamily}, \mref{Rem.}{rem:relabel}, \mref{Rem.}{rem:alphabet}, \mref{Thm.}{thm:band-a}, \mref{Ass.}{ass:coin}, \mref{Rem.}{rem:coin-scope}, \mref{Req.}{req:choice-weighting-record}, \mref{Req.}{req:conditional-family-decision}, \mref{Req.}{req:binary64-weights}, \mref{Req.}{req:weight-preparation}, \mref{Req.}{req:setting-only-weighting}, \mref{Req.}{req:weight-list}, \mref{Req.}{req:stochastic-weights}, \mref{Req.}{req:weighted-family-emission}, \mref{Req.}{req:no-lattice-exponent}, \mref{Req.}{req:physical-family-decision}, \mref{Req.}{req:band-a-ideal-receiver}, \mref{Req.}{req:band-a-identity-link}, \mref{Req.}{req:reference-companion-compiled}, \extref{../meat-core/contract.pdf}{meat-core/cor:reference-companion-factors}, \extref{../meat-core/contract.pdf}{meat-core/prop:reference-companion}
\metasep newly derived
\metasep implemented (effective: accepted)
\metasep \metalabel{Code}\coderef{https://github.com/fraunhofersg/certiqs-key/blob/v0.1.0/src/certiqs_key/protocol_compiler/weights.py}{certiqs_key.protocol_compiler.weights:decide_conditional_families}
\metasep \metalabel{Oracles}\coderef{https://github.com/fraunhofersg/certiqs-key/blob/v0.1.0/tests/oracles/registry_protocol_compiler.py}{registry_protocol_compiler.py}}

% claim protocol-compiler/lem:weighted
\claimsummary{protocol-compiler/lem:weighted}{Weighting keeps a measurement a measurement and acts block by block. The independent form has the declared state-independent choice marginal; the setting-dependent form returns the delivered family when the choice is forgotten, and its choice marginal may depend on the state through the setting.}
\claimbookmark{protocol-compiler/lem:weighted}{Lemma}{lem:weighted}{Properties of the choice-weighted families}
\claimleadin{%
Both constructions multiply declared measurement elements by positive numbers
(a zero entry of $P$ indexes no element), so what survives is what a positive
scalar preserves; the choice marginal is where the two forms part. Part~(i) is
about the mathematical products; what is decided of the emitted binary64
family, and where, is Remark~\ref{rem:exactness}. No identity of the form
\eqref{eq:choicemass} is wanted for the setting-dependent form: that the
marginal of a public choice may depend on the state through the setting is
the output-state test-statistics model (\xref{reductions/ass:tsm}), which
this lemma neither uses nor contradicts (Theorem~\ref{thm:choice-encoding}).
That the emitted declaration carries that model is
Requirement~\ref{req:no-fixed-fraction}.
}
\begin{lemma}[Properties of the choice-weighted families]\label{lem:weighted}
With the notation of Definition~\ref{def:weighted}, and reading its conditions
exactly:
\begin{enumerate}[label=(\roman*),leftmargin=2.2em]
\item $\{N_{u,t}\}_{(u,t)\in\Vset}$ is a POVM on $S$; and so is
$\{N_{v,f}\}_{(v,f)\in\Vset_P}$ of the setting-dependent form, over the
support of $P$;
\item if $S=\Qsys$ carries the sector decomposition and every $N_{u|t}$ is
sector-block-diagonal, then so is every $N_{u,t}$, and conversely (for
$q_t>0$); equivalently, on block lists the weighting acts blockwise,
$N_{u,t}^{(s)}=q_t\,N_{u|t}^{(s)}$, so a weighted family is a block list
against the same \exttermref{../foundations/contract.pdf}{term-sector-map}{sector map} as the conditional families it was built from.
The same holds of the setting-dependent form: every $N_{v,f}$ is a scalar
multiple of the $F_v$ it was built from, $N_{v,f}^{(s)}=P(f\,|\,t(v))\,F_v^{(s)}$
block by block, and $N_{v,f}$ is sector-block-diagonal iff $F_v$ is, the
weight being strictly positive on the support;
\item \emph{(the choice marginal, for the setting-independent form)} for every
$\omega\in\Lin{S}$ with $\tr\omega=1$ and every $t\in\Tpub$,
\begin{equation}\label{eq:choicemass}
  \sum_{u\in\Uset_t}\tr[N_{u,t}\,\omega]=q_t ,
\end{equation}
independently of $\omega$; and for any complementary system $R$ and any
$\omega\in\Lin{RS}$ of unit trace likewise
$\sum_{u}\tr[(\idon{R}\otimes N_{u,t})\,\omega]=q_t$, the operator being
written with its factors in the declared order of the tensor factors. For a
choice declared on Bob's side, $S=\Qsys$ and $R=A$, this reads
$\sum_u\tr[(\idon{A}\otimes N_{u,t})\omega]=q_t$; for a choice declared on
Alice's side, $S=A$ and $R=\Qsys$, it reads
$\sum_u\tr[(N_{u,t}\otimes\idon{\Qsys})\omega]=q_t$. \emph{This part is
asserted of the setting-independent form and of it alone};
\item \emph{(what is asserted of the setting-dependent form instead)}
forgetting the choice value returns the family the chain delivered,
\begin{equation}\label{eq:forgetf}
  \sum_{f:\,(v,f)\in\Vset_P}N_{v,f}=F_v\qquad\text{for every }v\in\Vset ,
\end{equation}
an exact operator identity on $S$; and the marginal of the choice value is
the declared matrix applied to the \emph{setting marginal} of the outcome
statistics,
\begin{equation}\label{eq:fmarginal}
  \sum_{v:\,(v,f)\in\Vset_P}\tr[N_{v,f}\,\omega]
  =\sum_{t\in\mathcal{T}}P(f\,|\,t)\sum_{v:\,t(v)=t}\tr[F_v\,\omega]
  \qquad\text{for every }\omega\in\Lin{S},
\end{equation}
which depends on $\omega$ only through the \emph{setting marginal}
$m_t(\omega)\coloneqq\tr[\Pi_t\omega]$, $\Pi_t\coloneqq\sum_{v:\,t(v)=t}F_v$,
with $\Pi_t\succeq0$ and $\sum_t\Pi_t=\idon{S}$. If $P(f\,|\,t)=q_f$ for
every $t$ it equals $q_f$ for every unit-trace $\omega$ --- the
setting-independent case, and \eqref{eq:choicemass} in the form the first
construction gives it. If every class sum $\Pi_t$ is a multiple
$\lambda_t\idon{S}$ of the identity --- as for the identity link, where the
emitted family is the physical \exttermref{../reductions/contract.pdf}{term-joint-family}{joint family} --- it is state-independent as
well, with value $\sum_tP(f\,|\,t)\lambda_t$. Otherwise it is, in general, a
function of the state, and \emph{no identity of the form
\eqref{eq:choicemass} is asserted for this form}.
\end{enumerate}
\end{lemma}
\claimmeta{\metalabel{Uses}\mref{Def.}{def:weighted}
\metasep \metalabel{Used by}\mref{Lem.}{lem:weighted-factor}, \mref{Thm.}{thm:choice-encoding}, \mref{Rem.}{rem:exactness}, \mref{Rem.}{rem:condfamily}, \mref{Ex.}{ex:condfamily}, \mref{Req.}{req:reference-companion-compiled}, \extref{../meat-core/contract.pdf}{meat-core/cor:reference-companion-factors}, \extref{../meat-core/contract.pdf}{meat-core/prop:reference-companion}
\metasep newly derived
\metasep implemented (effective: accepted)
\metasep \metalabel{Code}\coderef{https://github.com/fraunhofersg/certiqs-key/blob/v0.1.0/src/certiqs_key/protocol_compiler/weights.py}{certiqs_key.protocol_compiler.weights:choice_mass}
\metasep \metalabel{Oracles}\coderef{https://github.com/fraunhofersg/certiqs-key/blob/v0.1.0/tests/oracles/registry_protocol_compiler.py}{registry_protocol_compiler.py}}

\begin{proof}
(i) Each $N_{u,t}=q_tN_{u|t}\succeq0$ since $q_t>0$ and $N_{u|t}\succeq0$, and
$\sum_{(u,t)\in\Vset}N_{u,t}=\sum_{t}q_t\sum_{u\in\Uset_t}N_{u|t}
=\sum_tq_t\,\idon{S}=\idon{S}$, using completeness of each conditional POVM ---
an identity of the declared numbers by condition~(S) on each conditional family
--- and $\sum_tq_t=1$, likewise an identity by (S) on the weights. Hermiticity
passes through the real positive multiple. For the setting-dependent form,
each $N_{v,f}=P(f\,|\,t(v))F_v$ is Hermitian and positive semidefinite as a
positive multiple of $F_v$, and
$\sum_{(v,f)\in\Vset_P}N_{v,f}
=\sum_{v\in\Vset}\Big(\sum_{f:\,P(f|t(v))>0}P(f\,|\,t(v))\Big)F_v
=\sum_{v\in\Vset}F_v=\idon{S}$, because the omitted terms of each inner sum
are exactly the zero entries of the row $P(\,\cdot\,|\,t(v))$, so the inner
sum is the row sum, which is $1$ --- an identity of the declared numbers by
the exact row-sum decision of Definition~\ref{def:weighted} --- and
$\sum_vF_v=\idon{S}$ is a hypothesis on the family the lemma is given (it is
the completeness that condition (C3) of the
\termref{term-interface-conditions-on-a-chain}{interface conditions on a chain} (Definition~\ref{def:chaincond})
decides).

(ii) $\Pi_sN_{u,t}\Pi_{s'}=q_t\Pi_sN_{u|t}\Pi_{s'}$, which
vanishes iff $\Pi_sN_{u|t}\Pi_{s'}$ does, because $q_t\ne0$ --- and $q_t\ne0$
is needed for the converse direction only. Likewise
$\Pi_sN_{v,f}\Pi_{s'}=P(f\,|\,t(v))\Pi_sF_v\Pi_{s'}$ with
$P(f\,|\,t(v))>0$ on the support, and on block lists the scalar multiplies
each block.

(iii) $\sum_{u}\tr[q_tN_{u|t}\omega]=q_t\tr[(\sum_uN_{u|t})\omega]
=q_t\tr[\idon{S}\omega]=q_t\tr\omega=q_t$; the bipartite statement is the same
computation with the identity on $R$ adjoined and $\tr\omega=1$, and is
insensitive to which tensor slot $S$ occupies, since the trace is.

(iv) \eqref{eq:forgetf} is the inner sum of (i): for fixed $v$,
$\sum_{f:\,(v,f)\in\Vset_P}P(f\,|\,t(v))F_v=\big(\sum_fP(f\,|\,t(v))\big)F_v
=F_v$, the omitted terms being zero. For \eqref{eq:fmarginal}, group the
outcomes by their setting: $\sum_{v:\,(v,f)\in\Vset_P}\tr[N_{v,f}\omega]
=\sum_{t}\sum_{v:\,t(v)=t,\ P(f|t)>0}P(f\,|\,t)\tr[F_v\omega]
=\sum_tP(f\,|\,t)\sum_{v:\,t(v)=t}\tr[F_v\omega]$, the restriction
$P(f\,|\,t)>0$ being immaterial since the excluded terms vanish; the inner
sum is $\tr[\Pi_t\omega]=m_t(\omega)$ by linearity of the trace, $\Pi_t$ is a
sum of positive semidefinite operators, and $\sum_t\Pi_t=\sum_vF_v=\idon{S}$
since the classes partition $\Vset$. If
$P(f\,|\,t)=q_f$ for every $t$, the right-hand side
is $q_f\sum_tm_t(\omega)=q_f\tr\omega=q_f$ for every unit-trace $\omega$. If
$\Pi_t=\lambda_t\idon{S}$ for every $t$, then $m_t(\omega)=\lambda_t\tr\omega
=\lambda_t$ for every \emph{unit-trace} $\omega$ and the right-hand side is
$\sum_tP(f\,|\,t)\lambda_t$, whatever that state is. Outside these two cases
the right-hand side varies with the state for some families and does not for
others, and ``in general'' is asserted in that weak sense and in no stronger
one. At $|\mathcal{T}|=2$ it is exactly decided: writing
$m_{t_2}=\tr\omega-m_{t_1}$, the right-hand side is
$P(f\,|\,t_2)+\big(P(f\,|\,t_1)-P(f\,|\,t_2)\big)m_{t_1}(\omega)$ on
unit-trace states, which varies with $\omega$ precisely when $\Pi_{t_1}$ is
not a multiple of the identity --- so that $m_{t_1}$ takes at least two
values --- \emph{and} $P(f\,|\,t_1)\ne P(f\,|\,t_2)$; this witnesses that the
case occurs. Beyond two settings those two hypotheses are no longer
enough: on $S=\mathbb{C}^2$ take $\Vset=\{v_1,v_2,v_3\}$ with
$F_{v_1}=\tfrac12\dyad{0}$, $F_{v_2}=\tfrac12\dyad{1}$,
$F_{v_3}=\tfrac12\idon{S}$, a POVM since the three sum to $\idon{S}$, and
$t(v_i)=t_i$, so that $\Pi_{t_i}=F_{v_i}$ and $\Pi_{t_1}$ is not a multiple of
the identity; with $P(f\,|\,t_1)=P(f\,|\,t_2)=x\ne y=P(f\,|\,t_3)$ the
right-hand side is $x\,\tr[(\Pi_{t_1}+\Pi_{t_2})\omega]+y\,m_{t_3}(\omega)
=\tfrac12(x+y)$ for every unit-trace $\omega$, because
$\Pi_{t_1}+\Pi_{t_2}=\tfrac12\idon{S}=\Pi_{t_3}$. What makes the marginal
state-independent is a condition on the \emph{sums} $\Pi_t$ weighted by the
column $P(f\,|\,\cdot\,)$ --- that $\sum_tP(f\,|\,t)\Pi_t$ be a multiple of
the identity --- and this lemma asserts no such condition, because nothing
downstream needs one.
\end{proof}

% claim protocol-compiler/lem:weighted-factor
\claimsummary{protocol-compiler/lem:weighted-factor}{Weighting a delivered element by its positive stochastic weight keeps a factor enclosure of each of its blocks: the same box, the core multiplied by the weight. A pair outside the support of the weighting is no outcome and carries no block.}
\claimbookmark{protocol-compiler/lem:weighted-factor}{Lemma}{lem:weighted-factor}{A factor enclosure through a setting-dependent weighting}
\claimleadin{%
The compiler weights the family a chain delivers by a stochastic weight read
through the setting (Definition~\ref{def:weighted}). Where a block of a
delivered element is known through a \exttermref{../foundations/contract.pdf}{term-factor-enclosure}{factor enclosure}, as on a reference
companion (Requirement~\ref{req:reference-companion-compiled}), the weight is
a positive rational and goes into the core, and the columns are untouched.
}
\begin{technicalclaim}
\begin{lemma}[A factor enclosure through a setting-dependent weighting]\label{lem:weighted-factor}
With the notation of the setting-dependent form of
Definition~\ref{def:weighted}, let $(v,f)\in\Vset_P$, let $s$ be a sector of
$S$, and let $(L_c,\varrho,W)$ be a \exttermref{../foundations/contract.pdf}{term-factor-enclosure}{factor enclosure}
(\xref{foundations/def:factor-enclosure}) of the block $F^{(s)}_v$. Then
$(L_c,\varrho,P(f\,|\,t(v))\,W)$ is a factor enclosure of the block
$N^{(s)}_{v,f}=P(f\,|\,t(v))\,F^{(s)}_v$, with the same
\exttermref{../foundations/contract.pdf}{term-matrix-box}{matrix box} and the same number of columns. A pair $(v,f)\in\Vset\times\mathcal{F}$ outside
$\Vset_P$ indexes no element of the weighted family, and so no block.
\end{lemma}
\end{technicalclaim}
\claimmeta{\metalabel{Uses}\mref{Def.}{def:weighted}, \mref{Lem.}{lem:weighted}, \extref{../foundations/contract.pdf}{foundations/def:binary64}, \extref{../foundations/contract.pdf}{foundations/def:factor-enclosure}
\metasep \metalabel{Used by}\mref{Req.}{req:reference-companion-compiled}, \extref{../meat-core/contract.pdf}{meat-core/cor:reference-companion-factors}
\metasep newly derived
\metasep implemented (effective: accepted)
\metasep \metalabel{Code}\coderef{https://github.com/fraunhofersg/certiqs-key/blob/v0.1.0/src/certiqs_key/protocol_compiler/companion.py}{certiqs_key.protocol_compiler.companion:_weighted_factor}}

\begin{proof}
A pair $(v,f)$ lies in $\Vset_P$ exactly when $P(f\,|\,t(v))>0$
\eqref{eq:support}, and $P(f\,|\,t(v))$ is a binary64 number, so a positive
dyadic rational (\xref{foundations/def:binary64}\,(i)); hence
$P(f\,|\,t(v))\,W$ is diagonal with positive rational diagonal entries. By
Lemma~\ref{lem:weighted}\,(ii), $N^{(s)}_{v,f}=P(f\,|\,t(v))\,F^{(s)}_v$. The
triple $(L_c,\varrho,W)$ being a factor enclosure of $F^{(s)}_v$, there is an
$L$ that the box $(L_c,\varrho)$ contains with $F^{(s)}_v=LWL^{\dagger}$, so
$N^{(s)}_{v,f}=L\big(P(f\,|\,t(v))\,W\big)L^{\dagger}$ with the same $L$: the
box and the new core make a factor enclosure of $N^{(s)}_{v,f}$, with the
columns of the old one. The last sentence is the reading of the support in
Definition~\ref{def:weighted}: a pair with $P(f\,|\,t(v))=0$ is not an
outcome and indexes nothing.
\end{proof}

% claim protocol-compiler/thm:choice-encoding
\claimsummary{protocol-compiler/thm:choice-encoding}{A public choice of known bias enters as weights inside the measurement operators: the family stays a POVM, fixed weights are reproduced exactly as choice probabilities where the weighted family is emitted, and setting-dependent weights commute with the chain. A fraction outside them is refused.}
\claimbookmark{protocol-compiler/thm:choice-encoding}{Theorem}{thm:choice-encoding}{Encoding of public choices in the measurement operators}
\claimleadin{%
The theorem has two halves: (i)--(iv) for a choice of known bias declared on
either party's side, (v)--(viii) for a choice whose weight depends on the
private setting of the apparatus that measures a system, whichever party
draws it, applied after the chain. By (ii) the weight inside
the operators reproduces the declared choice statistics exactly wherever the
weighted family is the one emitted. A measuring party's presentation in front
of a receiver segment that squashes is the other case: there the weighted
family is the physical one handed to the chain, the emitted family is the
chain's, and the first half says what still holds. The
assumption of the output-state test-statistics model that (vii) declares is
\xref{reductions/ass:tsm}; it is named here, and the theorem does not use it.
Remark~\ref{rem:relabel} explains
why the relabelling of (vi) is not a link of the chain; it is the relabelling
of \xref{reductions/rem:intertwining}, which leaves it to the consumer to
apply.
}
\begin{mainclaim}
\begin{theorem}[Encoding of public choices in the measurement operators]
\label{thm:choice-encoding}
Let a per-round public choice of known bias --- a basis choice, a choice
between test and generation rounds, a probe setting --- be declared on either
party's side with alphabet $\Tpub$, weights $(q_t)$ and conditional POVMs
$\{N_{u|t}\}$, under Assumption~\ref{ass:coin}. Then the compiler forms the
choice-weighted family \eqref{eq:weighted}, and (iv) below holds. Where that
family is the emitted one --- the presentation is Alice's, or it is a
measuring party's and the last link of that party's receiver segment
delivers the weighted family itself, on its own alphabet, as its
\exttermref{../reductions/contract.pdf}{term-target-povm}{target POVM}, as on the \termref{term-band-a}{Band A} route
(Definition~\ref{def:banda}\,(b)) --- the compiler emits it, and (i)--(iii)
hold as well:
\begin{enumerate}[label=(\roman*),leftmargin=2.2em]
\item the emitted family is a POVM and preserves sector block-diagonality
(Lemma~\ref{lem:weighted}\,(i),(ii)), so the encoding cannot break D(iii) ---
and it is emitted only after the two decisions of
Definition~\ref{def:weighted}, so that what is emitted is binary64 data whose
own predicates D(iii) is then decided on
(Theorem~\ref{thm:measurements-and-statistics}\,(i), and (i$'$) in the
general format);
\item the encoding is \emph{exact}: for every state, the marginal probability
of the choice value $t$ computed from the emitted declaration equals the
declared weight $q_t$ (Lemma~\ref{lem:weighted}\,(iii));
\item the choice value is a component of the joint outcome, hence enters every
downstream quantity only through $P_{\Decl}[\omega]$ and
$\Phimap[\omega]$ of Definition~\ref{def:phi}; no factor $q_t$ is emitted
anywhere outside a POVM element; and
\item the alternative encoding --- declaring the conditional families
$\{N_{u|t}\}$ together with a separate state-independent fraction $q_t$ to be
applied outside the operators --- is never emitted, and a specification
requesting it is refused.
\end{enumerate}
Where the presentation is a measuring party's and the last link of that
party's receiver segment delivers any other family, the weighted family is
the physical receiver family handed to the chain
(Definition~\ref{def:spec}\,(b)) and is not emitted: it is a POVM on the
physical receiver space, and for every unit-trace state there the marginal
probability of $t$ computed from it is $q_t$
(Lemma~\ref{lem:weighted}\,(i), (iii)); no factor $q_t$ is emitted anywhere
outside a POVM element, and (iv) holds. The emitted family is that link's
target POVM, whose own conditions are (C3)'s and check V3's
(Definition~\ref{def:chaincond}), and
nothing here asserts (ii) or (iii) of it; the declared test-statistics model
is the output-state model, as (vii) says of the setting-dependent form.
Let, further, a public choice $f\in\mathcal{F}$ be declared whose weight
depends on the private setting of the apparatus that measures the
\exttermref{../reductions/contract.pdf}{term-measured-system}{measured system}, as a
stochastic matrix $P(f\,|\,t)$ read through the setting map
$t:\Vset\to\mathcal{T}$ of the family $\{F_v\}_{v\in\Vset}$ the last link of
$\Chrec$ delivered --- the setting-dependent form of
Definition~\ref{def:weighted}; in the general format, at a measured factor,
read through that factor's setting map on the family the last link of that
factor's receiver segment delivered, with $S$ that factor's system --- under
Assumption~\ref{ass:coin}. Then the compiler applies it \emph{after} the
reduction chain, emits the family \eqref{eq:weighted-t},
$N_{v,f}=P(f\,|\,t(v))F_v$ over the support $\Vset_P$, and:
\begin{enumerate}[label=(\roman*),leftmargin=2.2em,resume]
\item the emitted family is a POVM and preserves sector block-diagonality
(Lemma~\ref{lem:weighted}\,(i),(ii)), every element being a scalar multiple
of an element the chain delivered, so the encoding cannot break D(iii); and it
is emitted only after the two decisions of Definition~\ref{def:weighted} on
the products $P(f\,|\,t(v))F_v$, so that what is emitted is binary64 data
whose own predicates D(iii) is then decided on
(Theorem~\ref{thm:measurements-and-statistics}\,(i), and (i$'$) in the
general format);
\item \emph{(soundness: the relabelling commutes with the chain)} for every
channel $\Lambda$ into $S$ and every family $\{G_v\}_{v\in\Vset}$ on its
input space with $\Lambda^{\dagger}(F_v)=G_v$ for every $v$,
\begin{equation}\label{eq:relabel-sound}
  \Lambda^{\dagger}\big(N_{v,f}\big)=P(f\,|\,t(v))\,G_v
  \qquad\text{for every }(v,f)\in\Vset_P :
\end{equation}
the weighted target and the identically weighted input family stand in the
same relation under the \emph{same} channel, with the same target space, the
same sector map, the same constraint entries transported along the chain and
every pull-back unchanged. In this module's setting $G$ is the
$\gamma$-coarse-grained physical family of
Definition~\ref{def:chaincond}\,(C4), $G_v=\sum_{w\in\gamma^{-1}(v)}\Fph_w$,
so the emitted family and the physical family weighted by
$P(f\,|\,t(\gamma(w)))$ satisfy the pullback identity \eqref{eq:pullback}
with the outcome correspondence $(w,f)\mapsto(\gamma(w),f)$ restricted to the
support, wherever (C4) is decidable, and are related as (C4)'s families are
wherever it is not. This is \xref{reductions/lem:linear}\,(i) for the
stochastic matrix from $\Vset$ to $\Vset\times\mathcal{F}$ with entries
$\delta_{vv'}P(f\,|\,t(v'))$;
\item \emph{(what is not asserted)} the exactness clause (ii) is asserted for
the setting-independent form and for it alone. For this form the marginal of
$f$ is the quantity \eqref{eq:fmarginal} of
Lemma~\ref{lem:weighted}\,(iv), a function of the state through the setting
marginal, equal to the declared $q_f$ for every state when $P(f\,|\,t)=q_f$
and state-dependent in general otherwise; nothing here asserts otherwise, no
state-independent fraction is emitted, and the declared test-statistics model
is the output-state model.
What (iv) says of the fixed-fraction encoding holds of this form unchanged;
\item \emph{(refusals)} a weighting presented through anything but the
setting map --- indexed by $\Vset$, by the device outcome, or by any
component of $v$ other than $t(v)$ --- is refused; a $P$ with a negative
entry or a row whose declared-order sum is not exactly $1$ is refused; and
the two emission decisions of Definition~\ref{def:weighted} refuse, for a
product that is not representable or with the predicate that failed. Each
refusal emits no declaration.
\end{enumerate}
\end{theorem}
\end{mainclaim}
\claimmeta{\metalabel{Uses}\mref{Def.}{def:spec}, \mref{Def.}{def:compile}, \mref{Def.}{def:phi}, \mref{Thm.}{thm:measurements-and-statistics}, \mref{Def.}{def:weighted}, \mref{Lem.}{lem:weighted}, \mref{Def.}{def:chaincond}, \mref{Def.}{def:banda}
\metasep \metalabel{Assumes}\aref{A-COIN}{ass:coin}, \aref{A-LINK}{ass:link}
\metasep \metalabel{Used by}\mref{Rem.}{rem:relabel}, \mref{Req.}{req:setting-only-weighting}, \mref{Req.}{req:stochastic-weights}, \mref{Req.}{req:weighted-family-emission}, \mref{Req.}{req:no-fixed-fraction}
\metasep newly derived
\metasep implemented (effective: accepted)
\metasep \metalabel{Code}\coderef{https://github.com/fraunhofersg/certiqs-key/blob/v0.1.0/src/certiqs_key/protocol_compiler/weights.py}{certiqs_key.protocol_compiler.weights:resolve_povm}
\metasep \metalabel{Oracles}\coderef{https://github.com/fraunhofersg/certiqs-key/blob/v0.1.0/tests/oracles/registry_protocol_compiler.py}{registry_protocol_compiler.py}}

\begin{proof}
Under the hypothesis of (i)--(iii) the emitted family is the weighted family
itself, on its own alphabet: for Alice's presentation by
Definition~\ref{def:compile}, which emits her measurement slot from the
presentation; for a measuring party's because the emitted family is the
target POVM of the last link (Definition~\ref{def:compile}), which by
hypothesis is the weighted family on $\Vset=\bigsqcup_t\Uset_t\times\{t\}$
--- on the Band A route by Definition~\ref{def:banda}\,(b), with $\gamma$ the
identity on $\Vset=\Vph$ and $F_v=\Fph_v$ entry for entry. So (i) and (ii)
are Lemma~\ref{lem:weighted}, read of the emitted family and against D(iii)
and the emission decisions of Definition~\ref{def:weighted}. For (iii): the
emitted alphabet is $\bigsqcup_t\Uset_t\times\{t\}$ and the emitted matrices
are the $N_{u,t}$; the choice value $t$ is recoverable from the outcome index
alone, and no other emitted slot carries $q_t$. Since $P_{\Decl}[\omega]$ and
$\Phimap[\omega]$ are functions of $\omega$ and the emitted operators only
(Definition~\ref{def:phi}), the claim follows. (iv) is a rule of the
\termref{term-compilation-map}{compilation map}, not a derived statement, and is stated here because it
is the rule the encoding exists to enforce: the emitted declaration carries
the output-state model, and no path of the compiler produces an artefact
carrying a fixed-fraction model. Requirement~\ref{req:no-fixed-fraction}
binds its refusal.

For a measuring party's presentation whose last link delivers any other
family, Definition~\ref{def:spec}\,(b) makes the weighted family the physical
receiver family, the object (C2) and (C4) of Definition~\ref{def:chaincond}
read, and Definition~\ref{def:compile} fills the emitted family from the last
link's target POVM, so the weighted family is not emitted. Lemma~\ref{lem:weighted}\,(i)
and (iii) are stated on any finite-dimensional $S$ and apply with $S$ the
physical receiver space: the weighted family is a POVM there and its choice
marginal is $q_t$ for every unit-trace state there. No slot of the
declaration is filled from the weights: the emitted family is the target
POVM and the remaining slots are the other parts' (Definition~\ref{def:compile}),
so no factor $q_t$ is emitted outside a POVM element, and (iv) is the rule
of the compilation map as before. Nothing is asserted of the emitted family
beyond (C3) and check V3, and nothing more could be in general: (C2) admits
any tabulated total surjection $\gamma$, so $t$ need not be recoverable from
an emitted outcome; and where $\gamma$ keeps the choice value, every emitted
outcome carrying the $t$ of each physical outcome it receives, so that (C4)
summed over the emitted outcomes carrying $t$ reads
$\Lambda^{\dagger}(\Pi_t)=q_t\idon{}$ on the physical receiver space with
$\Pi_t$ the sum of the emitted elements carrying $t$, that fixes $\Pi_t$
only through $\Lambda^{\dagger}$ and does not make it $q_t\idon{\Qsys}$ off
the image of the channel. The
sentence on the test-statistics model is, as in (vii), a statement of what
the emitted declaration carries.

(v) is Lemma~\ref{lem:weighted}\,(i), (ii) for the setting-dependent form,
read against D(iii) and the emission decisions exactly as (i) is.

(vi) By linearity of $\Lambda^{\dagger}$,
$\Lambda^{\dagger}(P(f\,|\,t(v))F_v)=P(f\,|\,t(v))\Lambda^{\dagger}(F_v)
=P(f\,|\,t(v))G_v$, which is \eqref{eq:relabel-sound}; nothing else is used,
and in particular nothing about $\Lambda$ beyond linearity of its adjoint.
That the two weighted families are POVMs is (v) and its mirror on the input
side (the same computation with $G$ for $F$, $\sum_vG_v=\Lambda^{\dagger}
(\sum_vF_v)$ being the identity where $\Lambda^{\dagger}$ is unital, which is
the trace-preservation clause of the obligation of the chain's producer and
not a step of this proof). The identity \eqref{eq:relabel-sound} is the
instance of the post-processing part~(i) of \xref{reductions/lem:linear} for
the relation $G\dsq_{\Lambda}F$ and the stochastic matrix
$\Pi\big((v',f)\,|\,v\big)\coloneqq\delta_{vv'}P(f\,|\,t(v'))$ from $\Vset$ to
$\Vset\times\mathcal{F}$, every column of which sums to
$\sum_fP(f\,|\,t(v))=1$: $(\Pi F)_{(v,f)}=P(f\,|\,t(v))F_v$ and
$(\Pi G)_{(v,f)}=P(f\,|\,t(v))G_v$, and the computation above is that
lemma's argument for $\Pi$, written out. The target space, its
sector map and the transported entries are untouched because a relabelling
changes the index set of the two families and nothing else: every constraint
operator $\hat O_k$ is an operator on $\Qsys$, not a function of $v$, and its
pull-back $\Lambda^{\dagger}(\hat O_k)$ is computed through the same
$\Lambda$. For the compiler's instance of it, substitute (C4)'s
$G_v=\sum_{w\in\gamma^{-1}(v)}\Fph_w$: then
$\Lambda^{\dagger}(N_{v,f})=\sum_{w\in\gamma^{-1}(v)}P(f\,|\,t(\gamma(w)))\Fph_w$,
which is \eqref{eq:pullback} for the family
$\{P(f\,|\,t(\gamma(w)))\Fph_w\}_{(w,f)}$ under the correspondence
$(w,f)\mapsto(\gamma(w),f)$, a total surjection onto $\Vset_P$ because
$\gamma$ is one onto $\Vset$ and the weight depends on $w$ only through
$\gamma(w)$; where (C4) is not decidable from the data,
\eqref{eq:relabel-sound} is still an identity of linear maps, and it says of
the relation the chain's producer asserts and nobody here decides
(Definition~\ref{def:chaincond}\,(C4), second sentence) exactly what it says
of a decided one.

(vii) is Lemma~\ref{lem:weighted}\,(iv): \eqref{eq:fmarginal} is the marginal
of $f$ computed from the emitted declaration, it equals $q_f$ for every state
in the setting-independent case, and it is in general a function of the state
otherwise. The theorem asserts no value for it in that case; the sentence
about the \exttermref{../reductions/contract.pdf}{term-declared-model}{declared model} is a statement of what the emitted declaration
carries, as (iv) is.

(viii) is the rule of the compilation map that Definition~\ref{def:weighted}
states and Requirements~\ref{req:setting-only-weighting},
\ref{req:stochastic-weights} and~\ref{req:weighted-family-emission} bind,
stated beside (iv) because it is the rule the encoding exists to enforce: a
weighting that reaches an outcome through anything but its setting would make
$f$ a function of the device outcome, which is the inference
Assumption~\ref{ass:coin} excludes; a row of $P$ with non-negative entries
whose sum is not $1$ would make \eqref{eq:forgetf} false at every $v$ of that
setting with $F_v\ne0$, the inner sum over the support \eqref{eq:support}
being that row sum, and a negative entry, which the support never keeps and
which therefore reaches no emitted element, is refused because
Assumption~\ref{ass:coin} reads $P(\,\cdot\,|\,t)$ as the distribution the
choice $f$ is drawn from, and a negative number is no probability; and
the two emission decisions are those of the setting-independent form, taken on
the products of this one.
\end{proof}

% claim protocol-compiler/rem:exactness
\claimsummary{protocol-compiler/rem:exactness}{The exact choice marginal fixes the sum over outcomes, not the individual probabilities, so no state-independent weight can be factored out of a security argument. Nor does it establish D(iii) of the emitted binary64 family, which check V3 decides after the emission decisions.}
\claimbookmark{protocol-compiler/rem:exactness}{Remark}{rem:exactness}{Scope of the exact encoding}
\begin{remark}[Scope of the exact encoding]\label{rem:exactness}
Equation \eqref{eq:choicemass} says the choice \emph{marginal} reproduces the
declared weight for every state. What it fixes is that marginal alone: it
does \emph{not} license treating $q_t$ as a fraction that may be factored out
of a security argument, since $q_t$ appears in \eqref{eq:choicemass} only after
summing $\tr[N_{u,t}\omega]$ over $u$, and the terms $\tr[N_{u,t}\omega]$
that carry it depend on $\omega$. That is what separates the emitted encoding
from a fixed fraction applied outside the operators. The exactness statement
matters because it removes the only respectable motive for that unsound
encoding: the weight carried inside the operators already reproduces the
declared choice statistics exactly. Assumption~\ref{ass:coin} is
where the physical content sits --- a choice correlated with the adversary is
not a choice of known bias, and no encoding repairs that.

\emph{Nor does it say that the emitted family passes V3.}
Lemma~\ref{lem:weighted}\,(i) is a statement about the mathematical products
$q_tN_{u|t}$, respectively $P(f\,|\,t(v))F_v$; what is emitted is those
products as represented, as binary64 data, and two things can go wrong between
the two. A product of binary64 numbers need not be binary64, and a family
whose entries are all binary64 may still have a declared-order sum that
rounds. Both are decided before emission, by the emission decisions of
Definition~\ref{def:weighted}, and D(iii) is decided on the emitted blocks
afterwards, by check V3 (Theorem~\ref{thm:measurements-and-statistics}\,(i),
and (i$'$) in the general format), and never inferred from
the lemma, no part of which rests on that theorem. A weighting that fails
either decision is refused, which is the direction that costs a good
declaration rather than admitting a bad one. Both paragraphs are about the products of
Definition~\ref{def:weighted} in either of its forms; what differs between
the forms is only the first: for the setting-dependent form the exactness
statement, Theorem~\ref{thm:choice-encoding}\,(ii), is \emph{not asserted at all}, the
marginal of the choice value being state-dependent through the setting
marginal of the outcome statistics (Lemma~\ref{lem:weighted}\,(iv)), and no
one is tempted to factor it out.
\end{remark}
\claimmeta{\metalabel{Uses}\mref{Thm.}{thm:measurements-and-statistics}, \mref{Def.}{def:weighted}, \mref{Lem.}{lem:weighted}
\metasep \metalabel{Assumes}\aref{A-COIN}{ass:coin}
\metasep newly derived
\metasep accepted}

% claim protocol-compiler/rem:condfamily
\claimsummary{protocol-compiler/rem:condfamily}{Each conditional family is decided by the compilation map, not the validator, since the emitted family's conditions do not imply it; the converse gives (T), (H), (P) and, with (S), (C). For a choice-weighted physical family, only its own (S) is left to decide on the family itself.}
\claimbookmark{protocol-compiler/rem:condfamily}{Remark}{rem:condfamily}{Deciding the conditional families}
\begin{remark}[Deciding the conditional families]\label{rem:condfamily}
A choice-weighted presentation is two declared things and one constructed one:
the weights, the conditional families $\{N_{u|t}\}_{u\in\Uset_t}$, and the
products the compiler forms from them --- the \emph{weighted family}, which is
the emitted one where the presentation is Alice's or where the receiver segment
is the identity, and is the \emph{physical} family handed to the chain where it
is Bob's at a non-trivial one. The compiler decides all three
(Definition~\ref{def:weighted},
Requirement~\ref{req:conditional-family-decision}) --- the weights by sign and
by an exact declared-order sum, each conditional family for exact
admissibility in the sense of Definition~\ref{def:exact}, and every product it
forms for representability, with the five conditions decided once more on the
weighted family wherever that family is the emitted one. This remark records
why the middle decision cannot be dropped in favour of the last, what it buys
where the weighted family is \emph{not} emitted, and what it does not disturb.

\emph{It is a decision of the compilation map and not a check of the
validator.} The conditional families are declared input: they are no slot of
the emitted declaration, their index $(u,t)$ is not an emitted index, and the
family the consumer receives is the one product family. So the decision is
taken where the presentation is read, before any product is formed
(Requirement~\ref{req:conditional-family-decision}), beside the decisions on the weights and not
among V1--V8. Nothing in Definition~\ref{def:validator} gains or loses a
conjunct by it, and Theorem~\ref{thm:validator-soundness}\,(i) and~(ii) pair the same checks
with the same target conditions whether or not the decision is taken: the
target conditions D(i)--D(viii) quantify over the emitted declaration, which
this decision does not change, and a candidate that meets them is refused at
no check because of it. What the decision affects is which
\emph{specifications} reach a candidate at all, which is the compilation map's
business and not the validator's.

\emph{The emitted family's conditions do not imply the conditional families'.}
Example~\ref{ex:condfamily} is two weights and two one-element conditional
families, neither of them a POVM, whose weighted family satisfies (T), (H),
(P), (C) and (S) and passes V1's map comparison and check V3 with every product
exactly representable.
The reason it can happen is visible in the arithmetic: (C) on the emitted
family is one identity, $\sum_{t}q_t\sum_uN_{u|t}=\idon{S}$, and the inner sums
enter it only through that one combination, so a deficit under one choice value
is paid for by a surplus under another. Deciding the outer identity therefore
says nothing about the inner ones, and \eqref{eq:condfam-bad} is what ``nothing''
looks like.

\emph{The converse does hold, in four of the five, and it is stated of the
weighted family --- whether or not that family is the emitted one.} Suppose the
conditional families of one presentation all carry one map and are each
\exttermref{../foundations/contract.pdf}{term-exactly-admissible}{exactly admissible} for it, the weights satisfy $q_t>0$ and
$\sum_tq_t=1$ with that sum rounding nowhere, and the representability
decision has passed on the products, so that every entry of the weighted
family \emph{is} the real product $q_tN_{u|t}$. Then, of that family: (T)
holds \emph{against that same map} --- the map the conditional families
carry, which for a physical presentation is that of the module presenting it
and not $\Smap$ --- the weighting acting blockwise
(Lemma~\ref{lem:weighted}\,(ii)); (H) holds, a real multiple of a Hermitian
block being Hermitian; (P) holds, $q_t>0$ and $N_{u|t}^{(s)}\succeq0$ giving
$q_tN_{u|t}^{(s)}\succeq0$, the product being exact; and the \emph{real} sum
over the alphabet that indexes it is the
identity on each block, which is the computation of
Lemma~\ref{lem:weighted}\,(i) read on the declared numbers --- it is available
as an identity of those numbers precisely because (S) and (C) were decided of
each conditional family and of the weights. What does \emph{not} follow is (S)
on the weighted family's own declared-order sum: that condition is on binary64
additions --- of weighted operands, accumulated across the boundaries between
choice values --- and no hypothesis on the operands controls them
(Remark~\ref{rem:summation-condition}). \emph{Only the order restricts, never
the condition}: the order the conditional family's (S) is read in is the
restriction of the order on the alphabet that indexes the weighted family
(Definition~\ref{def:weighted}), and the two conditions are independent in
both directions. Given (S), (C) follows, the computed sum then being the real
one. So (S) is the one of the five that still needs its own decision, and the
emission decision of Definition~\ref{def:weighted} is not made redundant by
this one.

\emph{And here the two cases part company.} Where the weighted family
\emph{is} the emitted one --- a presentation on Alice's side, or one on Bob's
where the receiver segment is the identity --- (H), (P), (C) and (S) are decided
on it again by check V3 and its map (T) again by check V1
(Definition~\ref{def:validator}), because those are the conditions the
consumer's admission rule reads and this module decides what it emits rather
than inferring it (\xref{foundations/req:exact-emission}); what the converse
buys there is not a check saved but the knowledge that the surviving refusals
are (S)'s and representability's (Remark~\ref{rem:exactness}). Where it is
\emph{not} --- a physical receiver family presented choice-weighted at a
non-trivial receiver segment, which lives on the physical receiver space and
is no $\Qsys$-operator family --- the converse is what the decision on it
rests on (Requirement~\ref{req:physical-family-decision}): (T), (H) and (P)
hold of it by the decision on its conditional families, (S) on its own
declared-order sum is the one condition decided of the family itself, and (C)
of it follows from that (S) together with the conditional (C), (S) and the
weight sum. The emitted
family at such a segment is a different object, the \exttermref{../reductions/contract.pdf}{term-target-povm}{target POVM} of the last
link, and its own five are (C3)'s and then V1's and V3's
(Definition~\ref{def:chaincond}).

\emph{The setting-dependent form has no second subject.} There the family that
is weighted is the single family $\{F_v\}_{v\in\Vset}$ the last link of
$\Chrec$ delivered, and the choice value is drawn \emph{after} the
measurement, with probability $P(f\,|\,t(v))$; there is no family of
measurements indexed by $f$, and $\{N_{v,f}\}_v$ at fixed $f$ is not one and
is not meant to be (Lemma~\ref{lem:weighted}\,(iv),
Theorem~\ref{thm:choice-encoding}\,(vii)). Read through the overlap of the two forms ---
$P(f\,|\,t)=q_f$ for every $t$, where the one conditional family is $\{F_v\}$
itself --- the decision above applies with that family as its subject, and it
was taken already: condition~(C3) of Definition~\ref{def:chaincond} decides
the exact admissibility of the family the last link delivered, per block. So
no conditional family of the setting-dependent form goes undecided and no
second decision is added there; what the compiler decides of $P$ is what
Definition~\ref{def:weighted} states of it.

\emph{What stays assumed.} A specification whose
conditional families are not measurements is refused: what is lost is a
declaration whose per-choice reading is false, and Example~\ref{ex:condfamily}
shows that such a declaration implements a probability of the choice value
other than the one the declared public-choice record states and
Assumption~\ref{ass:coin} reads. Assumption~\ref{ass:spec}
keeps what no function of the \exttermref{../foundations/contract.pdf}{term-declared-data}{declared data} decides --- that the declared
dimensions, alphabets and tables describe the apparatus and the
declared measurements it names are the ones performed. The physical
receiver family as presented is decided on every route
(Requirement~\ref{req:physical-family-decision}); where it is presented in
choice-weighted form, its conditional families are read here like any other
and their five conditions are decided, so what is decided of the presented
family itself shrinks to condition~(S) on its own declared-order sum ---
from which its (C) then follows, by the converse above.
\end{remark}
\claimmeta{\metalabel{Uses}\mref{Def.}{def:exact}, \mref{Def.}{def:compile}, \mref{Def.}{def:weighted}, \mref{Lem.}{lem:weighted}, \mref{Def.}{def:validator}, \mref{Thm.}{thm:validator-soundness}
\metasep newly derived
\metasep accepted}

% claim protocol-compiler/ex:condfamily
\claimsummary{protocol-compiler/ex:condfamily}{Two weights of one half and two one-element conditional families that are not POVMs give a weighted family passing checks V1 and V3 exactly. The emitted conditions therefore do not imply the per-choice ones, and the declared weight is then not the probability the declaration implements.}
\claimbookmark{protocol-compiler/ex:condfamily}{Example}{ex:condfamily}{An exact weighted family built from incomplete conditional families}
\begin{example}[An exact weighted family built from incomplete conditional families]\label{ex:condfamily}
Let $S$ be two-dimensional carrying the single sector $\SecSet=\{s_0\}$ with
$d_{s_0}=2$, so that a block list against that map is one $2\times2$ array and
condition~(T) of Definition~\ref{def:exact} is the agreement of one identifier,
one dimension and one block kind. Let $\Tpub=\{t_0,t_1\}$ with
$q_{t_0}=q_{t_1}=\tfrac12$, let $\Uset_{t_0}=\Uset_{t_1}=\{u_0\}$ --- one
outcome per choice value --- and let the conditional families be
\begin{equation}\label{eq:condfam-bad}
  N_{u_0|t_0}=\operatorname{diag}\big(\tfrac12,\tfrac32\big),\qquad
  N_{u_0|t_1}=\operatorname{diag}\big(\tfrac32,\tfrac12\big).
\end{equation}
Every entry is an integer multiple of $2^{-1}$ and each weight is $\tfrac12$,
so every product of \eqref{eq:weighted} is exactly a binary64 number and the
representability decision of Definition~\ref{def:weighted} passes. The emitted
family is
\begin{equation*}
  N_{u_0,t_0}=\operatorname{diag}\big(\tfrac14,\tfrac34\big),\qquad
  N_{u_0,t_1}=\operatorname{diag}\big(\tfrac34,\tfrac14\big).
\end{equation*}

\emph{All five conditions hold of it.} (T): there is one block and one map.
(H) and (P): both arrays are real, diagonal and entrywise non-negative. Its
declared-order sum over $\Vset=\{(u_0,t_0),(u_0,t_1)\}$ is two binary64
additions per entry, $0+\tfrac14$ then $\tfrac14+\tfrac34$ and $0+\tfrac34$
then $\tfrac34+\tfrac14$, each exact, so
(S) holds; and its value is $\operatorname{diag}(1,1)=\idon{S}$, so (C) holds.
So check V3 passes on the emitted family --- it is V3 that decides (H), (P),
(C) and (S) --- and check V1's four-component comparison passes on its map,
which is where (T) is decided (Definition~\ref{def:validator}); and so does
every decision on the weights:
$q_t>0$ by sign, and $\tfrac12+\tfrac12=1$ exactly in the declared order of
$\Tpub$.

\emph{Neither conditional family is a POVM.} Their completeness sums are
\begin{equation*}
  \sum_{u\in\Uset_{t_0}}N_{u|t_0}
    =\operatorname{diag}\big(\tfrac12,\tfrac32\big),\qquad
  \sum_{u\in\Uset_{t_1}}N_{u|t_1}
    =\operatorname{diag}\big(\tfrac32,\tfrac12\big),
\end{equation*}
each a one-term sum, whose single addition $0+N_{u_0|t}$ of
\eqref{eq:declorder} is exact, and neither is $\idon{S}$: condition~(C) fails
for both, of the
declared numbers and through no rounding, while (T), (H), (P) and (S) hold of
both. So there is no experiment ``draw $t$ with probability $q_t$, then measure
$\{N_{u|t}\}_u$'' for these numbers, and the hypothesis of
Definition~\ref{def:weighted} --- that each $\{N_{u|t}\}_u$ be a POVM on $S$
--- is violated by a specification whose weighted family passes every decision
taken on the emitted data.

\emph{And the declaration then implements a choice probability the record does
not declare.} At $\omega=\dyad{0}$ the choice marginal computed from the
emitted declaration is $\sum_{u}\tr[N_{u,t_0}\,\omega]=\tfrac14$, while the
record declares $q_{t_0}=\tfrac12$ and Assumption~\ref{ass:coin} reads that
number as the bias of the coin Bob drew: \eqref{eq:choicemass} of
Lemma~\ref{lem:weighted}\,(iii) is false of the pair, its hypothesis having
failed, and with it the exactness Theorem~\ref{thm:choice-encoding}\,(ii) asserts.
\emph{No slot of the declaration carries the wrong number} --- no
public-choice weight is emitted outside a POVM element
(Theorem~\ref{thm:choice-encoding}\,(iii)), and a consumer of the emitted operators sees a
consistent POVM implementing $\tfrac14$. The failure is therefore between the
specification and the declaration the module built from it, which is where a
decision on the specification is the only thing that can catch it, and is why
the per-choice reading is needed.

\emph{The presentation is refused}, naming (C) at $t_0$ --- the first
failure, the decision stopping there
(Requirement~\ref{req:conditional-family-decision}); nothing here is
rescaled, renormalised or repaired (Remark~\ref{rem:condfamily}).
\end{example}
\claimmeta{\metalabel{Uses}\mref{Def.}{def:exact}, \mref{Def.}{def:weighted}, \mref{Lem.}{lem:weighted}
\metasep newly derived
\metasep accepted}

% claim protocol-compiler/rem:relabel
\claimsummary{protocol-compiler/rem:relabel}{A choice whose bias depends on a station's setting is applied after the chain, to both sides of the last link's relation, by one stochastic matrix; by linearity the relation keeps the same channel. It is not a link, and it is read through the setting alone.}
\claimbookmark{protocol-compiler/rem:relabel}{Remark}{rem:relabel}{The setting-dependent relabelling after the chain}
\begin{remark}[The setting-dependent relabelling after the chain]\label{rem:relabel}
A public choice whose bias depends on the private setting of the apparatus
that measures a system, whichever party draws it --- a test or
generation round drawn against a probe setting, a probe light that fires
only in test rounds --- is the case the setting-dependent form of
Definition~\ref{def:weighted} exists for, and this remark records how the
compiler handles it. In the general format everything below holds per
measured factor, with that factor's receiver segment, family and setting map
in place of $\Chrec$, $\{F_v\}$ and $t$ (Definition~\ref{def:compile}).

\emph{What is applied, and where.} The compiler applies the stochastic matrix
$P(f\,|\,t)$, read through the setting map $t:\Vset\to\mathcal{T}$, to the
family $\{F_v\}$ the last link of $\Chrec$ delivered, \emph{after} the
reduction chain has been consumed and its
\termref{term-interface-conditions-on-a-chain}{interface conditions} decided
(Definition~\ref{def:chaincond}) and \emph{before} the validator and the
gates read the emitted declaration: the emitted family is
$N_{v,f}=P(f\,|\,t(v))F_v$ over the support $\Vset_P$, the emitted alphabet is
$\Vset_P$ in its product order, and the public part is declared on it
(Definition~\ref{def:compile}). The same weighting is, in the same stroke,
applied to the physical side of the last link's relation: the family the
emitted one is a reduction of is
$\{P(f\,|\,t(\gamma(w)))\Fph_w\}_{(w,f)}$ under the correspondence
$(w,f)\mapsto(\gamma(w),f)$ (Theorem~\ref{thm:choice-encoding}\,(vi)). Both sides are
relabelled by one stochastic matrix, which is what makes the next statement
true.

\emph{Why it is sound.} The soundness is
Theorem~\ref{thm:choice-encoding}\,(vi): by linearity of the adjoint, the
relation the last link asserts survives the stochastic matrix
$\delta_{vv'}P(f\,|\,t(v'))$ from $\Vset$ to $\Vset\times\mathcal{F}$, whose
columns sum to one, with the \emph{same} channel:
$\Qsys$, its sector map and its dimension are unchanged, every element
$P(f\,|\,t(v))F_v$ is a scalar multiple of an existing one so that block
diagonality holds by construction (Lemma~\ref{lem:weighted}\,(ii)), every
transported constraint operator has the same pull-back, and every entry's
certificate is untouched. That theorem is this module's instance, with the
one-line linearity argument written out, of the post-processing part~(i) of
\xref{reductions/lem:linear} and of the relabelling that
\xref{reductions/rem:intertwining} states and leaves to the consumer. What
that theorem does \emph{not}
assert is the choice marginal: the probability of $f$ is state-dependent
through the setting marginal of the outcome statistics
(Lemma~\ref{lem:weighted}\,(iv)), which is the output-state test-statistics
model (\xref{reductions/ass:tsm}, Requirement~\ref{req:no-fixed-fraction}),
and no fraction is factored out.

\emph{Why it is not a link.} A link in the sense of
\xref{reductions/def:link} composes a \emph{new} channel with the chain and
carries a \exttermref{../reductions/contract.pdf}{term-reachable-set-statement}{reachable-set statement} for it; a relabelling composes no channel
--- it changes the alphabet at both ends of an existing relation and leaves
the channel where it was. It therefore has no kind, no register and no
position in a chain (\xref{reductions/def:chain}), and that module declares
none: the chain record this module consumes and re-emits is unchanged by the
weighting, the setting map it is read through is the class map of the joint
alphabet of \xref{reductions/def:joint-family}, and the weighting is recorded
beside the chain as this module's own construction
(Requirement~\ref{req:choice-weighting-record}). Putting $f$ into the alphabet
\emph{before} the chain instead is sound as well but strictly weaker where
the squasher is flag-shaped: every value of $f$ then opens flag sectors of its
own, and the squashed model gains adversarial freedom no physical state
produces; applied after the chain, no sector is indexed by $f$ and
$N_{v,f}=0$ as an operator wherever $P(f\,|\,t(v))=0$, so a zero-probability
pair is unreachable for every state rather than merely of zero weight. That
comparison is not a claim of this module; the claim is the soundness of the
relabelling, Theorem~\ref{thm:choice-encoding}\,(vi).

\emph{Why the weighting is read through the setting only.} $P$ is declared on
$\mathcal{T}\times\mathcal{F}$ and reaches an outcome through $t(v)$ alone.
A weighting that reached it through the device outcome would make $f$ an
inference from that outcome, not a choice, and the relabelling would then be
applied to a variable the adversary can influence through the state; that is
what Assumption~\ref{ass:coin} excludes and what
Requirement~\ref{req:setting-only-weighting} enforces by type. The setting
$t$ itself is a variable the squasher's channel depends on --- the conditional
families differ with it --- which is why it is in the alphabet before the
first link and $f$ is not.

\emph{With no setting-dependent choice declared}, the emitted family is the
chain's $\{F_v\}$ unchanged. What the emitted alphabet may be is the upstream
modules' question, not this module's (Remark~\ref{rem:alphabet}).
\end{remark}
\claimmeta{\metalabel{Uses}\mref{Def.}{def:compile}, \mref{Def.}{def:weighted}, \mref{Thm.}{thm:choice-encoding}, \mref{Def.}{def:chaincond}
\metasep \metalabel{Assumes}\aref{A-COIN}{ass:coin}
\metasep \metalabel{Used by}\mref{Req.}{req:choice-weighting-record}
\metasep newly derived
\metasep accepted}

% claim protocol-compiler/rem:alphabet
\claimsummary{protocol-compiler/rem:alphabet}{Bob's emitted outcome alphabet is chosen upstream, the receiver's declared coarse-graining deciding how much of the click pattern survives. Every object of this module indexed by it is stated over any finite alphabet, so a fine-grained instance changes declared data and no claim here.}
\claimbookmark{protocol-compiler/rem:alphabet}{Remark}{rem:alphabet}{The outcome alphabet is chosen upstream}
\begin{remark}[The outcome alphabet is chosen upstream]\label{rem:alphabet}
Nothing in this module fixes what Bob's \exttermref{../reductions/contract.pdf}{term-outcome-alphabet}{outcome alphabet} is. The emitted
$\Vset$ is the alphabet of the family the last link of $\Chrec$ delivered
(Definition~\ref{def:spec}\,(c)), and that family's alphabet is the upstream
modules': the raw outcome of a threshold-detector receiver is the
\emph{click pattern}, the set $S\subseteq[K]$ of ports that fired
(\xref{receiver/def:clickalphabet}), and it is the receiver's
\exttermref{../receiver/contract.pdf}{term-declared-coarse-graining}{declared coarse-graining} --- a surjection from $2^{[K]}$ onto the
emitted alphabet (\xref{receiver/def:coarse},
\xref{receiver/thm:coarse-graining}) --- that decides how much of the pattern
survives, \emph{before} the joint alphabet with Bob's setting is formed
(\xref{reductions/def:joint-family}) and before any reduction runs. A
coarse-graining that keeps every pattern, one that keeps only click, no-click
and multi-click, and one that keeps the single clicks and merges the rest are
three declarations of that map, and this module receives the result of
whichever was declared.

Every object of this module that is indexed by $\Vset$ is stated over an
arbitrary finite non-empty alphabet and admits any finite disjoint union.
The family slot $\{F_v\}_{v\in\Vset}$ of \eqref{eq:declaration} is a family
of block lists over a finite alphabet, and D(iii) is decided element by
element whatever $|\Vset|$ is. The outcome correspondence $\gamma$ of
Definition~\ref{def:chaincond}\,(C2) is a total surjection from any finite
physical alphabet onto any finite emitted one, the direction being the only
thing (C2) fixes (Remark~\ref{rem:coarse}), so a finer emitted alphabet is a
$\gamma$ with smaller fibres and nothing else. The readouts of the
\termref{term-party-partition}{party partition} and Bob's \termref{term-announcement-map}{announcement map} $a_B$
(Definition~\ref{def:spec}\,(d))
are declared on the alphabets the specification names, and $\phi$
(Definition~\ref{def:phi}) on the announcement classes; a finer $\Vset$ is a
larger table for each and no new condition. And both forms of
Definition~\ref{def:weighted} are constructions over a disjoint union ---
$\bigsqcup_t\Uset_t\times\{t\}$, and the support $\Vset_P\subseteq
\Vset\times\mathcal{F}$ --- so a disjoint union over settings and choice
values is what they produce, not something they have to be extended to.

A \emph{fine-grained instance} --- one whose receiver keeps more of the click
pattern than the coarsest sound choice, or one whose outcome carries a
setting and a choice value beside the detection result --- therefore changes
the receiver's declared coarse-graining and the instance's public part, and
changes no claim of this module: the compiler's claims are stated for the
alphabet it is handed. What it does change downstream is not this module's
either: the size of the declaration, the number of exact decisions V3 takes,
and which coarse outcomes carry a bound
(\xref{receiver/thm:coarse-graining}\,(iii): a bound on a coarse outcome does
not descend to its fibres), which is the receiver's trade-off and the reason
its coarse-graining is declared data rather than a convention. Which
alphabets that coarse-graining can reach is the receiver's question; where a
setting-dependent choice value enters the alphabet is
Remark~\ref{rem:relabel}.
\end{remark}
\claimmeta{\metalabel{Uses}\mref{Def.}{def:declaration}, \mref{Def.}{def:spec}, \mref{Def.}{def:phi}, \mref{Def.}{def:weighted}, \mref{Def.}{def:chaincond}
\metasep newly derived
\metasep accepted}

\section{Sifting, key computability and feasibility}\label{group-protocol-compiler/thm:sifting-and-key}
\sectionintro{Each announcement is a function of its sender's readout and of the announcements before it, sifting is determined by the announcements, and the holder computes the raw key locally, for a transcript of messages as for two \termref{term-announcement-map}{announcement maps} (\hyperlink{claim-protocol-compiler/thm:sifting-and-key}{Theorem~\ref*{thm:sifting-and-key}}). Feasibility is the one target condition that is not a typing or a finite syntactic check: an empty constraint list has an explicit product state, and otherwise the compiler verifies a default candidate and, if it fails, a supplied witness, and refuses when neither verifies (\hyperlink{claim-protocol-compiler/thm:feasibility}{Theorem~\ref*{thm:feasibility}}).}

% claim protocol-compiler/thm:sifting-and-key
\claimsummary{protocol-compiler/thm:sifting-and-key}{Each announcement is a function of its sender's readout and of the announcements before it, sifting is determined by the announcements, and the holder computes the raw key locally. An ensemble-form measurement slot of Alice's denotes a POVM.}
\claimbookmark{protocol-compiler/thm:sifting-and-key}{Theorem}{thm:sifting-and-key}{Announcement locality, sifting and party computability}
\begin{theorem}[Announcement locality, sifting and party computability]\label{thm:sifting-and-key}
Let $\Spec$ be a \termref{term-well-formed-specification}{well-formed specification}
(Definition~\ref{def:well-formed-spec}) satisfying
Assumption~\ref{ass:spec}, so
that in particular $r^{0}_{\hold}$ maps into $\Zval$
(Assumption~\ref{ass:spec}\,(ii)), and let $\Decl$ be emitted by
Definition~\ref{def:compile}. Then:
\begin{enumerate}[label=(\roman*),leftmargin=2.2em]
\item \emph{(announcement locality)} $a_A$ is a function of
$\rd{\mathsf A}(x,v)$ alone and $a_B$ of $\rd{\mathsf B}(x,v)$ alone, and the
joint announcement is \eqref{eq:jointann};
\item \emph{(D(v))} the specification's $\varsigma$ --- consumed, and not
itself a slot of \eqref{eq:declaration} --- witnesses announcement-determined
sifting: $r(x,v)=\perp\iff\varsigma(a(x,v))=\mathrm{fail}$;
\item \emph{(D(viii))} $r$ is party-computable: the total map
\begin{equation}\label{eq:rhold}
  r_{\hold}(o,\kappa)\coloneqq
  \begin{cases}
    r^{0}_{\hold}(o,\kappa), & \varsigma(\kappa)=\mathrm{pass},\\
    \perp, & \varsigma(\kappa)=\mathrm{fail},
  \end{cases}
  \qquad (o,\kappa)\in\Oset{\hold}\times\Kann,
\end{equation}
satisfies \eqref{eq:partycomp}.
\item \emph{(\termref{term-transcript-protocol}{transcript protocols})} if the
announcement slot holds a transcript protocol
(Definition~\ref{def:declaration}\,(DX5)), then for
every $m$ the $m$-th component of $a(x,v)$ is a function of
$\rd{\pi_m}(x,v)$ and of the components before it, and of nothing else; $a$ is
total; and (ii) and (iii) hold with this $a$. For $L=2$,
$\pi=(\mathsf A,\mathsf B)$ and $a_2$ independent of $\kappa_1$ this is (i);
\item \emph{(\termref{term-ensemble-form}{ensemble form})} if Alice's measurement slot is declared
in ensemble form and the conditions of Definition~\ref{def:exact} for that
form hold, then the operators \eqref{eq:ensemblepovm} satisfy
$M^{(t)}_x\succeq0$ and $\sum_xM^{(t)}_x=\idon{A_t}$ for every $t$, and
$(\tau^{(t)}_A)^{1/2}M^{(t)}_x(\tau^{(t)}_A)^{1/2}=N^{(t)}_x$: the denoted
declaration satisfies the POVM clause of D(iii) on $A$ as an identity of real
numbers, and $\tr[M_x\tau_A]=\sum_t\tr N^{(t)}_x$.
\item \emph{(the compiled general format)} if $\Spec$ is presented as
Definition~\ref{def:spec}\,(a$'$)--(e$'$) and $\Decl$ is emitted by the
general-format clauses of Definition~\ref{def:compile}, then: the
announcement slot of $\Decl$ holds the transcript $\Spec_{\mathrm{pub}}$
declares, message for message; the joint announcement function of $\Decl$ is
\eqref{eq:transcript} evaluated at \emph{those} tables and is a slot of
nothing; and Alice's measurement slot holds the form the transmitter artefact
declares. Consequently (iv) applies to $\Decl$, hence (ii) and (iii) hold of
it, and (v) applies to it exactly where that form is the ensemble form.
\end{enumerate}
Moreover, under the \termref{term-default-party-partition}{default party partition} these read:
$a_A:\Xset\to\CAset$, $a_B:\Vset\to\CBset$, and $r(x,v)=r_{\mathsf A}(x,a(x,v))$
for $\hold=\mathsf A$, respectively $r(x,v)=r_{\mathsf B}(v,a(x,v))$ for
$\hold=\mathsf B$.
\end{theorem}
\claimmeta{\metalabel{Uses}\mref{Def.}{def:declaration}, \mref{Def.}{def:exact}, \mref{Def.}{def:target}, \mref{Def.}{def:spec}, \mref{Def.}{def:well-formed-spec}, \mref{Def.}{def:compile}
\metasep \metalabel{Assumes}\aref{A-SPEC}{ass:spec}
\metasep \metalabel{Used by}\mref{Thm.}{thm:measurements-and-statistics}
\metasep newly derived
\metasep accepted
\metasep \metalabel{Code}\coderef{https://github.com/fraunhofersg/certiqs-key/blob/v0.1.0/src/certiqs_key/protocol_compiler/spec.py}{certiqs_key.protocol_compiler.spec:derived_holder_key_rule}
\metasep \metalabel{Oracles}\coderef{https://github.com/fraunhofersg/certiqs-key/blob/v0.1.0/tests/oracles/registry_protocol_compiler.py}{registry_protocol_compiler.py}}

\begin{proof}
(i) is Definition~\ref{def:spec}\,(d) together with \eqref{eq:jointann}: the
\termref{term-announcement-map}{announcement maps} are declared on the readout alphabets, so no other
signature is expressible.

(ii) Fix $(x,v)$ and write $\kappa=a(x,v)$. If
$\varsigma(\kappa)=\mathrm{fail}$ then \eqref{eq:rdef} gives $r(x,v)=\perp$. If
$\varsigma(\kappa)=\mathrm{pass}$ then
$r(x,v)=r^{0}_{\hold}(\rd{\hold}(x,v),\kappa)\in\Zval$, and $\perp\notin\Zval$
by $\Zset=\Zval\sqcup\{\perp\}$, so $r(x,v)\ne\perp$. The two implications are
the two directions of the equivalence.

(iii) $r_{\hold}$ of \eqref{eq:rhold} is total on $\Oset{\hold}\times\Kann$
because $\varsigma$ is total and both branches are defined everywhere on their
respective preimages. Fix $(x,v)$ and put $o=\rd{\hold}(x,v)$,
$\kappa=a(x,v)$. If $\varsigma(\kappa)=\mathrm{pass}$, both sides of
\eqref{eq:partycomp} equal $r^{0}_{\hold}(o,\kappa)$ by \eqref{eq:rdef} and
\eqref{eq:rhold}; if $\varsigma(\kappa)=\mathrm{fail}$, both equal $\perp$.

(iv) By \eqref{eq:transcript} the $m$-th component is
$a_m(\rd{\pi_m}(x,v),\kappa_1(x,v),\dots,\kappa_{m-1}(x,v))$, and $a_m$ has
no other argument, its domain being declared as
$\Oset{\pi_m}\times\mathcal{C}_1\times\dots\times\mathcal{C}_{m-1}$; totality
follows by induction on $m$ from totality of the tables and of the
\termref{term-readout-map}{readout maps}. The proofs of (ii) and (iii) use of $a$ only that it is a
total map into $\Kann$ on which $\varsigma$ and $r^{0}_{\hold}$ are defined, so
they apply unchanged. With $L=2$ and $a_2(o,\kappa_1)=a_B(o)$,
\eqref{eq:transcript} is \eqref{eq:jointann}.

(v) Fix $t$ and drop it. $\tau_A\ne0$ is positive semidefinite, so
$\tr\tau_A>0$ and $w_x\coloneqq\tr N_x/\tr\tau_A\ge0$ with $\sum_xw_x=1$ by
$\sum_xN_x=\tau_A$. With $S\coloneqq\tau_A^{-1/2}$ on the support, zero off
it, and $P$ the support projector: $SN_xS\succeq0$ and
$w_x(\idon{}-P)\succeq0$, so $M_x\succeq0$, and
$\sum_xM_x=S\tau_AS+(\idon{}-P)=\idon{}$. From
$0\preceq N_x\preceq\tau_A$, every vector in the kernel of $\tau_A$ is in
the kernel of $N_x$, so $PN_xP=N_x$; then
$\tau_A^{1/2}M_x\tau_A^{1/2}=PN_xP+w_x\tau_A^{1/2}(\idon{}-P)\tau_A^{1/2}
=N_x$, using $\tau_A^{1/2}S=P$ and $\tau_A^{1/2}(\idon{}-P)=0$. The trace
statement follows by taking traces and summing over $t$.

(vi) The three clauses are what the general-format clauses of
Definition~\ref{def:compile} do, and they are read off that definition and
not inferred: the transcript is passed through verbatim, \emph{$a$ is not a
slot} but the composition \eqref{eq:transcript} of the tables passed through,
and Alice's measurement slot is passed through in the form the artefact
declares, that definition converting neither form into the other. For the
consequences: Assumption~\ref{ass:spec}\,(ii) gives $r^{0}_{\hold}$ into
$\Zval$ and Definition~\ref{def:well-formed-spec}\,(iii) the totality of the
declared public tables, which is all (iv) consumes
besides the signatures, and the signatures are the ones (d$'$) of
Definition~\ref{def:spec} declares; so (iv) applies, and (ii), (iii) with it,
by (iv)'s last two sentences. Where the slot is the ensemble slot, the
conditions (v) is stated under are exactly the general-format conditions
Definition~\ref{def:exact} states for that form --- (S), (H), (P) of
$\{N^{(t)}_x\}_x$ per Alice sector, the identity of its declared-order sum
with $\tau^{(t)}_A$, and $\tau^{(t)}_A\ne0$ --- so (v) applies to it.
\emph{Nothing of (i) is claimed of such a declaration}: (i) is the two-map
case, which is $L=2$ with $a_2$ independent of $\kappa_1$, and (iv)'s last
sentence is where the two meet.

The default-partition statements follow by substituting
$\rd{\mathsf A}(x,v)=x$ and $\rd{\mathsf B}(x,v)=v$.
\end{proof}

% claim protocol-compiler/rem:sidecond
\claimsummary{protocol-compiler/rem:sidecond}{Without the condition that the key rule map into the key values, the specification's sifting predicate stops witnessing announcement-determined sifting. Check V5 catches only a key rule that withholds a key on part of a passing class, so the condition is a clause of \hyperlink{claim-protocol-compiler/ass:spec}{A-SPEC}.}
\claimbookmark{protocol-compiler/rem:sidecond}{Remark}{rem:sidecond}{Necessity of the key-rule side condition}
\begin{remark}[Necessity of the key-rule side condition]\label{rem:sidecond}
Drop the requirement that $r^{0}_{\hold}$ map into $\Zval$, and
Theorem~\ref{thm:sifting-and-key}\,(ii) fails: a key rule that returns $\perp$
on a \emph{passing} class makes $r(x,v)=\perp$ for pairs whose announcement
class $\varsigma$ calls $\mathrm{pass}$, so the specification's $\varsigma$ no
longer witnesses announcement-determined sifting. The two ways this can happen
differ in what is still decided. Where the key rule returns $\perp$ on part of
a passing class, no predicate witnesses D(v): whether a round is sifted is not
determined by the announcements, a sifting event becomes invisible to the
public record, and check V5 of Definition~\ref{def:validator} refuses the
candidate. Where it returns $\perp$ on the whole of a passing class, D(v) still
holds, witnessed by a predicate that calls that class $\mathrm{fail}$, and V5
passes; what is lost is that the specification's $\varsigma$ is that witness,
and \eqref{eq:chkey} of Definition~\ref{def:declaration}, which reads the
specification's $\varsigma$, may then count as key-live an acceptance value
that no key-producing class reaches. This is why the condition is stated in
the specification's type (Definition~\ref{def:spec}\,(d)) and carried as a
clause of Assumption~\ref{ass:spec}, V5 catching only its first failure.
\end{remark}
\claimmeta{\metalabel{Uses}\mref{Def.}{def:spec}
\metasep newly derived
\metasep accepted}

% claim protocol-compiler/rem:vacuous
\claimsummary{protocol-compiler/rem:vacuous}{Read against the joint outcome pair, party computability would hold for every raw-key map; read against the holder's readout it is a genuine restriction. It matters most where both parties' data share one outcome symbol, as when no source is trusted.}
\claimbookmark{protocol-compiler/rem:vacuous}{Remark}{rem:vacuous}{Party computability read against the readouts}
\begin{remark}[Party computability read against the readouts]
\label{rem:vacuous}
Read against the joint pair $(x,v)$ instead of against the readouts,
\eqref{eq:partycomp} would be satisfied by \emph{every} $r$ whatsoever, since
$r$ is by definition a function of $(x,v)$; the condition would be vacuous
exactly for the declaration shapes that need it most --- those which place both
parties' data in one outcome symbol, as the $d_A=1$ embedding of
\xref{transmitter/rem:da1} does. Checked against $\rd{\hold}$, it is a genuine
restriction: it fails as soon as two pairs with the same holder readout and the
same announcements carry different key values (check V8 of
Definition~\ref{def:validator}). That embedding --- Alice's measured
components folded into the receiver-side outcome, the \termref{term-party-partition}{party partition}
naming which components she holds --- is the shape in which a protocol with no
trusted source is declared, and the argument above is why party computability
is checked against the partition in that shape.
\end{remark}
\claimmeta{\metalabel{Uses}\mref{Def.}{def:target}
\metasep newly derived
\metasep accepted}

% claim protocol-compiler/ex:transcript
\claimsummary{protocol-compiler/ex:transcript}{Alice announces the kind of round, Bob discloses his bit only on test rounds, and Alice then discloses hers only on detected test rounds. Party computability holds by the signatures of the tables, and two independent announcement maps cannot express the protocol.}
\claimbookmark{protocol-compiler/ex:transcript}{Example}{ex:transcript}{A three-message transcript}
\begin{example}[A three-message transcript]\label{ex:transcript}
Default readouts, $\rd{\mathsf A}(x,v)=x$ and $\rd{\mathsf B}(x,v)=v$. Alice's
letter is $x=(z,b)$, with $z\in\{\mathrm{test},\mathrm{gen}\}$ saying whether
the round is a test or a generation round and a bit $b\in\{0,1\}$; Bob's
letter is $v\in\{0,1,\varnothing\}$, a bit or no detection. The transcript
(Definition~\ref{def:declaration}\,(DX5)) has $L=3$ messages:
\begin{enumerate}[label=(m\arabic*),leftmargin=3.0em]
\item $\pi_1=\mathsf A$, $\mathcal{C}_1=\{\mathrm{test},\mathrm{gen}\}$,
  $a_1(z,b)=z$: Alice announces the kind of round;
\item $\pi_2=\mathsf B$,
  $\mathcal{C}_2=\{0,1,\mathrm{click},\varnothing\}$, and
  $a_2(v,\kappa_1)=v$ if $\kappa_1=\mathrm{test}$ or $v=\varnothing$, and
  $a_2(v,\mathrm{gen})=\mathrm{click}$ for $v\in\{0,1\}$: Bob discloses his
  bit on test rounds and only the detection on generation rounds;
\item $\pi_3=\mathsf A$, $\mathcal{C}_3=\{0,1,\ast\}$, and
  $a_3((z,b),\kappa_1,\kappa_2)=b$ if $\kappa_1=\mathrm{test}$ and
  $\kappa_2\ne\varnothing$, and $\ast$ otherwise: Alice discloses her bit on
  detected test rounds.
\end{enumerate}
By \eqref{eq:transcript}, $a((\mathrm{test},b),v)=(\mathrm{test},v,b)$ for
$v\in\{0,1\}$, $a((\mathrm{test},b),\varnothing)
=(\mathrm{test},\varnothing,\ast)$,
$a((\mathrm{gen},b),v)=(\mathrm{gen},\mathrm{click},\ast)$ for $v\in\{0,1\}$
and $a((\mathrm{gen},b),\varnothing)=(\mathrm{gen},\varnothing,\ast)$: seven
classes occur out of the twenty-four of
$\Kann=\mathcal{C}_1\times\mathcal{C}_2\times\mathcal{C}_3$, and the tables
are total on all of their domains as D(viii) asks, their values on
combinations that never occur being immaterial. Sifting passes on
$(\mathrm{gen},\mathrm{click},\ast)$ alone; with $\hold=\mathsf A$ the key
rule is $r^{0}_{\mathsf A}((z,b),\kappa)=b$, a function of Alice's readout.
The acceptance map may send each test class $(\mathrm{test},v,b)$ to the
pair $(v,b)$, so that the test statistics see both bits and the generation
rounds disclose neither.

\emph{What two \termref{term-announcement-map}{announcement maps} cannot express.} In
the \termref{term-basic-format}{basic format}, with two maps $a_A$ on $\Oset{\mathsf A}$ and $a_B$
on $\Oset{\mathsf B}$, Bob's announcement is a function of $v$ alone and cannot
depend on the kind of round, which he learns only from (m1); the only way to
write this protocol there is to put $z$ into Bob's readout, which misstates
what Bob measures. Here party computability holds \emph{by the signatures}:
$a_2$ has no argument through which $x$ could enter except the message
$\kappa_1$ that Alice has made public (Theorem~\ref{thm:sifting-and-key}\,(iv)).

\emph{What it is read under.} Every message of a round follows every
measurement of that round (\xref{transmitter/ass:seq}): (m1) is sent after Bob
has measured, so his measurement cannot depend on it, and the declaration has
no slot in which such a dependence could be written, $\{F_v\}$ being one fixed
family. A protocol in which Bob chose his measurement after hearing $z$ is
outside this format and outside the fixed-measurement \exttermref{../foundations/contract.pdf}{term-round-model}{round model}.
\end{example}
\claimmeta{newly derived
\metasep accepted}

% claim protocol-compiler/thm:feasibility
\claimsummary{protocol-compiler/thm:feasibility}{An empty constraint list has an explicit product state in the feasible set. Otherwise the compiler verifies, exactly and block by block, the default candidate and, if it fails, a supplied witness, and refuses when neither verifies.}
\claimbookmark{protocol-compiler/thm:feasibility}{Theorem}{thm:feasibility}{Feasibility by an explicit or a verified candidate}
\claimleadin{%
D(vii) is the one target condition that is not a typing or a finite syntactic
check: whether $\Feas\ne\varnothing$ depends on the interaction of the marginal
with the constraint entries. It is therefore treated in two cases, and the
second one ends in a refusal rule rather than a theorem; the engine's order of
the two candidates and its refusal are Requirement~\ref{req:feasibility-witness}.
}
\begin{theorem}[Feasibility by an explicit or a verified candidate]\label{thm:feasibility}
Let $\Decl$ be a candidate declaration satisfying D(i)--D(iii), and let
$\Feas$ be as in Definition~\ref{def:feasible}. (The hypothesis
is on the conditions alone: nothing below uses that $\Decl$ was emitted by
Definition~\ref{def:compile}, and Theorem~\ref{thm:validator-soundness} applies this theorem to
candidates.)
\begin{enumerate}[label=(\roman*),leftmargin=2.2em]
\item If $\CList=\varnothing$ then $\Feas\ne\varnothing$. Explicitly, the
\exttermref{../foundations/contract.pdf}{term-block-list}{block list}
\begin{equation}\label{eq:witness-empty}
  \omega^{\star(s)}=
  \begin{cases}
    \tau_A\otimes\dfrac{\idon{\Qsys_s}}{d_{\Qsys}},
      & \text{trusted-source mode},\\[1.2ex]
    \dfrac{\idon{A}}{d_A}\otimes\dfrac{\idon{\Qsys_s}}{d_{\Qsys}},
      & \text{untrusted-source mode},
  \end{cases}
  \qquad s\in\SecSet ,
\end{equation}
which denotes $\tau_A\otimes\idon{\Qsys}/d_{\Qsys}$, respectively
$\idon{A}/d_A\otimes\idon{\Qsys}/d_{\Qsys}$, lies in $\Feas$.
\item If $\CList\ne\varnothing$ and $\Decl$ satisfies D(vi) as well, which
makes each constraint leg of \eqref{eq:feas} a well-posed comparison, then
$\Feas\ne\varnothing$ holds if and only
if some block list $\omega$ satisfies the conditions of \eqref{eq:feas}. The
compiler decides those conditions \emph{exactly}, by finitely many arithmetic
operations and per block, on at most two candidates and in a fixed order:
first the \emph{default candidate} \eqref{eq:witness-empty}, which satisfies
every conjunct of \eqref{eq:feas} except possibly the constraint legs
(by the proof of (i)) and which the compiler forms in order to test and never
emits; and, if that candidate fails a constraint leg, a \emph{supplied}
witness $\omega^{\star}$, consumed and verified and never emitted. It refuses
exactly when the default candidate fails and either no witness is supplied or
the supplied one fails.
\end{enumerate}
\emph{The general format.} For a candidate declaration in the general format
of Definition~\ref{def:declaration}\,(DX1)--(DX6) satisfying D(i)--D(iii),
and D(vi) for part (ii), with $\Feas$ the set of \eqref{eq:feas-general},
both parts hold with the
\termref{term-joint-sector}{joint sectors} in place of the sectors. In (i) the default candidate is
the block list over the joint sectors
\begin{equation}\label{eq:witness-general}
  \omega^{\star(t,s)}=
  \begin{cases}
    \tau^{(t)}_A\otimes\dfrac{\idon{\Qsys_s}}{d_{\Qsys}},
      & \text{trusted-source mode},\\[1.2ex]
    \dfrac{\idon{A_t}}{d_A}\otimes\dfrac{\idon{\Qsys_s}}{d_{\Qsys}},
      & \text{untrusted-source mode},
  \end{cases}
  \qquad (t,s)\in\mathbb{T}\times\SecSet ,
\end{equation}
and in (ii) the verification is the per-joint-block one listed in the proof,
which forms no array larger than one joint block. At $|\mathbb{T}|=1$ and one
\termref{term-measured-factor}{measured factor} \eqref{eq:witness-general} is \eqref{eq:witness-empty}
and the two lists of decisions coincide item for item.
\end{theorem}
\claimmeta{\metalabel{Uses}\mref{Def.}{def:declaration}, \mref{Def.}{def:exact}, \mref{Def.}{def:feasible}, \mref{Def.}{def:homog}, \mref{Def.}{def:compile}
\metasep \metalabel{Used by}\mref{Def.}{def:validator}, \mref{Thm.}{thm:validator-soundness}, \mref{Rem.}{rem:completeness}, \mref{Rem.}{rem:chain-and-validator}, \mref{Req.}{req:feasibility-witness}
\metasep newly derived
\metasep implemented (effective: accepted)
\metasep \metalabel{Code}\coderef{https://github.com/fraunhofersg/certiqs-key/blob/v0.1.0/src/certiqs_key/protocol_compiler/feasibility.py}{certiqs_key.protocol_compiler.feasibility:check_feasibility}
\metasep \metalabel{Oracles}\coderef{https://github.com/fraunhofersg/certiqs-key/blob/v0.1.0/tests/oracles/registry_protocol_compiler.py}{registry_protocol_compiler.py}}

\begin{proof}
\emph{Idea.} The default candidate is $\tau_A\otimes\idon{\Qsys}/d_{\Qsys}$,
or $\idon{A}/d_A\otimes\idon{\Qsys}/d_{\Qsys}$ where there is no marginal,
presented block by block; every conjunct of the \termref{term-feasible-set}{feasible set} except
the constraint legs holds of it by direct computation, and at an empty list
there are no constraint legs. It is formed and decided in
exact rational arithmetic and never emitted, so its entries need not be
binary64. Verifying any candidate is a finite list of exact decisions, each on
one block or on one $d_A\times d_A$ matrix. The general format repeats both
arguments over the joint sectors.

\emph{Step 1 (the default candidate lies in $\Feas$ at an empty list, (i)).}
Write $\chi_s\coloneqq\idon{\Qsys_s}/d_{\Qsys}$, so that
$\sum_s\tr\chi_s=\sum_sd_s/d_{\Qsys}=1$. The sector decomposition is not a
condition here: $\omega^{\star}$ is presented as a list, one array on
$A\Qsys_s$ per $s$, so the off-block conjunct holds by the presentation and
there is nothing to verify.

\emph{Positivity, per block.} In \exttermref{../transmitter/contract.pdf}{term-trusted-source-mode}{trusted-source mode}
$\omega^{\star(s)}=\tau_A\otimes\chi_s$ is Hermitian, $\tau_A$ being Hermitian
by D(ii) and $\chi_s$ a positive real multiple of an identity; and it is
positive semidefinite, being the tensor product of the positive semidefinite
$\tau_A$ (D(ii)) with the positive semidefinite $\chi_s$. In
\exttermref{../transmitter/contract.pdf}{term-untrusted-source-mode}{untrusted-source mode} $\omega^{\star(s)}=\idon{A}/d_A\otimes\chi_s\succ0$ outright.
\emph{Trace.} $\tr\omega^{\star}=\sum_s\tr\omega^{\star(s)}
=\tr\tau_A\sum_s\tr\chi_s=\tr\tau_A=1$ in trusted-source mode, by D(ii); in
untrusted-source mode $\tr\omega^{\star}=1$ directly.
\emph{Marginal.} In trusted-source mode
$\omega^{\star}_A=\sum_s\tr_{\Qsys_s}\omega^{\star(s)}
=\tau_A\sum_s\tr\chi_s=\tau_A$; in untrusted-source mode there is no marginal
condition beyond the trace, already checked. \emph{Constraints.} With
$\CList=\varnothing$ the conjunction of \eqref{eq:feas-modes} is empty and
holds vacuously. Hence $\omega^{\star}\in\Feas$.

\emph{Step 2 (the number model of Step 1).} Two points of that argument
concern the number model. First, every step used D(ii) \emph{as an identity of
the declared numbers}, which is what Definition~\ref{def:exact} delivers ---
$\tr\tau_A=1$ with its sum rounding nowhere, in particular. Second,
\emph{the entries of $\omega^{\star}$ need not be \exttermref{../foundations/contract.pdf}{term-binary64-numbers}{binary64 numbers}, and
nothing requires them to be}: $1/d_{\Qsys}$ is a rational, dyadic only when
$d_{\Qsys}$ is a power of two, and $\omega^{\star}$ is \emph{never emitted}.
It is formed and decided in exact rational arithmetic on the declared dyadic
entries --- the exact arithmetic in which the \exttermref{../foundations/contract.pdf}{term-exact-island}{exact island}
(\xref{foundations/def:exact-island}\,(i)) takes its decisions on
\exttermref{../foundations/contract.pdf}{term-declared-data}{declared data}, applied here to a candidate the compiler forms itself
and never emits --- so every predicate above is decided, not approximated,
and no representability question arises for an
object no consumer ever sees. A \emph{supplied} witness is a different
matter: it is \exttermref{../foundations/contract.pdf}{term-declared-data}{declared data}, hence dyadic, and its predicates are
decided exactly on those dyadic numbers.

\emph{Step 3 (the verification of a candidate, (ii)).} The stated equivalence
is Definition~\ref{def:feasible} read literally. The verification of a
candidate is finite and \emph{per block}. Writing $n_s\coloneqq d_Ad_s$ for
the size of the $s$-block, no array larger than one $n_s\times n_s$ block is
formed at any point; where $|\SecSet|=1$ that one block is $d_Ad_{\Qsys}$
square, and is a block and not an exception (Theorem~\ref{thm:band-a}\,(i)). It
consists of
\begin{itemize}[leftmargin=1.6em]
\item $|\SecSet|$ positivity decisions, each on an $n_s\times n_s$ matrix ---
  the Hermiticity decision $\omega^{(s)}=\omega^{(s)\dagger}$, entry by entry,
  and the semidefiniteness of the resulting Hermitian matrix by exact
  symmetric elimination, of which the largest is on $\max_sn_s$ and not on
  $d_Ad_{\Qsys}$;
\item one trace decision, $\tr\omega=1$ on the scalar
  $\sum_s\tr\omega^{(s)}$, a sum of $|\SecSet|$ block traces accumulated in the
  \exttermref{../foundations/contract.pdf}{term-declared-order}{declared order};
\item in trusted-source mode one marginal decision: form
  $\omega_A=\sum_s\tr_{\Qsys_s}\omega^{(s)}$, which is
  $|\SecSet|$ partial traces onto $A$ and $|\SecSet|-1$ additions of $d_A^2$
  entries, then $\omega_A=\tau_A$ over $d_A^2$ entries;
\item and, per entry $k$, one constraint leg: in trusted-source mode
  $\sum_s\tr_{\Qsys_s}[(\idon{A}\otimes\hat O_k^{(s)})\omega^{(s)}]$, a
  matrix of size $d_A$ assembled from $|\SecSet|$ block terms, compared with
  $\Theta_k$ in the ordering $\bowtie_k$ --- an exact L\"owner comparison,
  i.e.\ the semidefiniteness of one $d_A\times d_A$ difference for an
  inequality, or one exact entrywise equality when $\bowtie_k$ is $=$; in
  untrusted-source mode the scalar
  $\sum_s\tr[(\idon{A}\otimes\hat O_k^{(s)})\omega^{(s)}]$ compared with
  $c_k$ as real numbers, according to the emitted source mode (D(vi)). The
  right-hand side of a trusted-source-mode leg is $c_k\tau_A$ for an entry
  tagged homogeneous and the declared $\Theta_k$ for one tagged general
  (Definition~\ref{def:homog}); the compiler forms $c_k\tau_A$ here in order
  to \emph{test} the candidate, and does not emit it --- the consuming engine
  constructs its own from the same declared $c_k$ and $\tau_A$, so the two
  form the same matrix.
\end{itemize}

\emph{Step 4 (one decision for an equality).} The equality decision is
equivalent to the two inequality conditions. For Hermitian operands $X,Y$, set
$D=X-Y$. If $D\succeq0$ and $-D\succeq0$, each real eigenvalue of $D$ is both
nonnegative and nonpositive, hence zero; Hermitian diagonalization gives
$D=0$. Conversely $D=0$ implies both inequalities. Thus entrywise $X=Y$
decides their conjunction exactly, in one predicate. The argument applies to
every operator row of the general format below as well.

All of these are finitely many arithmetic operations on the declared data, and
the largest object any of them forms is one $n_s\times n_s$ block or one
$d_A\times d_A$ matrix.

\emph{Step 5 (the general format, (i)).} Write
$\chi_s\coloneqq\idon{\Qsys_s}/d_{\Qsys}$ as above, so
$\sum_s\tr\chi_s=1$. The sector decomposition, now along the joint sectors,
is again the presentation and not a condition: \eqref{eq:witness-general} is
one array per joint sector. \emph{Positivity, per joint block}: in
trusted-source mode $\tau^{(t)}_A\otimes\chi_s$ is Hermitian and positive
semidefinite, each $\tau^{(t)}_A$ being so by D(ii) in its general-format
reading; in untrusted-source mode the block is positive definite outright.
\emph{Trace}: $\sum_{(t,s)}\tr\omega^{\star(t,s)}
=\big(\sum_t\tr\tau^{(t)}_A\big)\big(\sum_s\tr\chi_s\big)=1$ in
trusted-source mode by D(ii), and $\sum_t(d_{A_t}/d_A)\sum_s\tr\chi_s=1$ in
untrusted-source mode by $d_A=\sum_td_{A_t}$ (D(i)). \emph{Marginal}, per
Alice sector and in trusted-source mode:
$\sum_s\tr_{\Qsys_s}\omega^{\star(t,s)}
=\tau^{(t)}_A\sum_s\tr\chi_s=\tau^{(t)}_A$, which is the second line of
\eqref{eq:feas-general}; in untrusted-source mode there is no marginal
condition beyond the trace. \emph{Constraints}: with $\CList=\varnothing$ the
three constraint lines of \eqref{eq:feas-general} are empty conjunctions. So
$\omega^{\star}\in\Feas$, and the two points of Step 2 apply unchanged: every
step reads D(ii) as an identity of the declared numbers, and the entries of
\eqref{eq:witness-general} need not be binary64 because the candidate is
formed and decided in exact rational arithmetic and is never emitted.

\emph{Step 6 (the general format, (ii)).} The equivalence is
Definition~\ref{def:feasible} in its general-format reading, read literally.
Writing $n_{t,s}\coloneqq d_{A_t}d_s$ for the size of the joint block, the
verification of a candidate consists of
\begin{itemize}[leftmargin=1.6em]
\item $|\mathbb{T}|\,|\SecSet|$ positivity decisions, one Hermiticity and one
  exact symmetric elimination per joint block, the largest on
  $\max_{(t,s)}n_{t,s}$;
\item one trace decision on $\sum_{(t,s)}\tr\omega^{(t,s)}$, accumulated in
  the declared order of the joint sectors;
\item in trusted-source mode $|\mathbb{T}|$ marginal decisions, one per Alice
  sector: $|\SecSet|$ partial traces onto $A_t$, $|\SecSet|-1$ additions and
  one comparison, each over $d_{A_t}^2$ entries;
\item and, per entry $k$, its own leg. For an entry read through its
  \termref{term-lifted-operator}{lifted operator}, in trusted-source mode one exact L\"owner
  comparison per Alice sector on a $d_{A_t}\times d_{A_t}$ difference for an
  inequality, or one exact entrywise equality when $\bowtie_k$ is $=$, the
  right-hand side $c_k\tau^{(t)}_A$ for a homogeneous entry and the declared
  $\Theta^{(t)}_k$ for a general one; in untrusted-source mode one comparison
  of real numbers against $c_k$. For an entry declared party-local and scoped
  to $f$, one exact L\"owner comparison per pair $(t,s_{f^c})$ for an
  inequality, or one exact entrywise equality per pair when $\bowtie_k$ is
  $=$, on a difference of size
  $d_{A_t}\,d_{f^c,s_{f^c}}=n_{t,s}/d_{f,s_f}$, assembled from the
  $|\SecSet^{f}|$ joint blocks of that row by partial traces over
  $\Qsys_{f,s_f}$ --- which is the last line of \eqref{eq:feas-general}, and
  is \emph{never larger than one joint block}. The same decision rule applies
  in either source mode; a $1\times1$ row uses one real-number comparison as
  its scalar specialization.
\end{itemize}
\emph{No lifted operator is formed at any of these steps.} The value
$\tr_{\Qsys_s}[(\idon{A_t}\otimes\hat O^{\uparrow(s)}_k)\,\omega^{(t,s)}]$ is
computed by contracting the joint block against the declared block
$\hat O^{(s_{f_k})}_k$ on the tensor factor $\Qsys_{f_k,s_{f_k}}$ of that
block and tracing out the others, which reads one joint block and one block of
the scope and forms no array on $\Qsys$; the identity of that contraction with
the displayed trace is the definition of the lifted operator in (DX6).
\emph{So the largest object any decision of this verification forms is one
$n_{t,s}\times n_{t,s}$ joint block, one $d_{A_t}\times d_{A_t}$ matrix, or one
party-local difference, and each of the last two is at most one joint block};
at $|\mathbb{T}|=1$ and one measured factor the list is the list of Step 3,
item for item, with $n_{t_0,s}=n_s$.
\end{proof}

% claim protocol-compiler/rem:nosearch
\claimsummary{protocol-compiler/rem:nosearch}{Deciding non-emptiness of the feasible set for a general list is a semidefinite feasibility problem, while whoever produced the entries can exhibit a reachable state. A supplied witness turns D(vii) into a check, and refusing when no candidate verifies avoids an infimum over the empty set.}
\claimbookmark{protocol-compiler/rem:nosearch}{Remark}{rem:nosearch}{Feasibility by witness rather than by search}
\begin{remark}[Feasibility by witness rather than by search]\label{rem:nosearch}
Deciding $\Feas\ne\varnothing$ for a general list is a semidefinite
feasibility problem. Whoever produced the entries produced them from a
reachable-state argument and can exhibit a reachable state, and a supplied
witness turns D(vii) into a check. Refusing when no candidate verifies is the
safe direction: a declaration whose feasible set is empty makes the consuming
engine's infimum an infimum over the empty set, which is exactly the pathology
the refusal prevents.
\end{remark}
\claimmeta{newly derived
\metasep accepted}

\section{The emitted validator}\label{group-protocol-compiler/def:validator}
\sectionintro{The validator's eight checks are decided exactly on the candidate declaration. The main result, \hyperlink{claim-protocol-compiler/thm:validator-soundness}{Theorem~\ref*{thm:validator-soundness}}, shows that passing every check establishes all eight target conditions of the declared numbers, that every check except the feasibility check is complete, so that a \termref{term-well-formed-declaration}{well-formed declaration} can be refused only for want of a verified witness, and that every decision is finite, exact and taken block by block. Such a refusal means that no verified feasible state was shown, not that the protocol is defective.}

% claim protocol-compiler/def:validator
\claimsummary{protocol-compiler/def:validator}{Eight exact checks on a candidate declaration: sector maps, the marginal, the operators, the key alphabet, sifting, constraint typing, feasibility and party computability. Feasibility passes only on an empty list, a verified default candidate or a verified supplied witness.}
\claimbookmark{protocol-compiler/def:validator}{Definition}{def:validator}{The emitted validator}
\claimleadin{%
The validator is the list of checks the compiler runs on a candidate
declaration before it emits one. A candidate failing a check is refused with a
reason naming the check (Requirement~\ref{req:validator-refusals}); what
passing all eight establishes is Theorem~\ref{thm:validator-soundness}, and which conditions
are decided on the chain rather than here, and which objects check V1 does
not compare, is Remark~\ref{rem:chain-and-validator}. That check V7 decides
on quantities formed in exact arithmetic is
\xref{foundations/req:exact-island} at work, and that no product block,
\termref{term-lifted-operator}{lifted operator} or denoted family of an \termref{term-ensemble-form}{ensemble form} is a
datum is
Requirement~\ref{req:derived-objects}.
}
\begin{definition}[The emitted validator]\label{def:validator}
Given a candidate declaration \eqref{eq:declaration} with its emitted
\termref{term-index-map}{index map}, the checks are as follows. \emph{Every operator condition
below is decided exactly}, in the sense of Definition~\ref{def:exact}, and
every combinatorial condition exactly as well.
\begin{description}[leftmargin=3.0em,style=nextline]
\item[V1 (finiteness, totality, and the \exttermref{../foundations/contract.pdf}{term-sector-map}{sector map})] every declared alphabet is
  finite and non-empty and every declared map is tabulated on all of its stated
  domain; every declared system carries a finite dimension;
  $d_{\Qsys}=\sum_{s\in\SecSet}d_s$; the restriction of $\idxm$ to $\SecSet$ is
  the total order the sector map fixes; \emph{and the sector maps
  agree}: there is one emitted map $\Smap$, and the map carried by each
  supplied $\Qsys$-operator family is compared against it in \emph{block
  kind}, in the \emph{list of block identifiers}, in their \emph{order} and in
  their \emph{dimensions}. A disagreement on any of the four is refused,
  naming the first position at which the lists differ. The families compared
  are exactly the $\Qsys$-operator families of \eqref{eq:declaration}, which
  is the list D(i) quantifies over and no wider one: the emitted family
  $\{F_v\}$, which the last link of $\Chrec$ delivers as its \exttermref{../reductions/contract.pdf}{term-target-povm}{target POVM}
  against the map the chain's producer wrote it against, and each constraint
  entry of $\CList$, declared in $\Spec_{\mathrm{con}}$ or transported by a
  link of either segment. Those two slots are the whole of
  \eqref{eq:declaration} that carries a $\Qsys$-operator
  (Definition~\ref{def:declaration}), so the list is exhaustive by
  construction rather than by enumeration; the objects outside it are compared
  where they belong. The entries of
  $\Spec_{\mathrm{con}}$ are compared as well as the transported ones:
  Definition~\ref{def:compile} fills $\CList$ from both sources,
  Definition~\ref{def:spec}\,(e) has each declared $\hat O_k$ carry its own
  sector map, and an entry written against a different map is the error this
  check exists for, whichever source it arrived from. The comparison is
  combinatorial --- it compares labels and integers, never matrix entries ---
  and it is the whole of the structural verification the emitted declaration
  receives: block-diagonality is carried by the type
  (Theorem~\ref{thm:measurements-and-statistics}\,(ii)), so the one structural error that can still
  arise is two parties disagreeing about what the sectors are, and V1 is the
  check for it.
\item[V2 (\termref{term-marginal-slot}{marginal slot})] the slot holds $\nomarg$, or a matrix $\tau_A$ on $A$
  --- one block, $A$ carrying no sector decomposition --- satisfying
  $\tau_A=\tau_A^{\dagger}$, $\tau_A\succeq0$ and $\tr\tau_A=1$, each decided
  exactly, and with the trace accumulated over the row index verified to
  \exttermref{../foundations/contract.pdf}{term-rounds-nowhere}{round nowhere}; a failure of the last is refused.
\item[V3 (operators, block by block)]
  $M_x=M_x^{\dagger}$ and $M_x\succeq0$ for each $x\in\Xset$, and
  $\sum_xM_x=\idon{A}$ with the sum, in the emitted index order, verified to
  round nowhere, a failure of which is refused;
  $F_v^{(s)}=F_v^{(s)\dagger}$ and $F_v^{(s)}\succeq0$ for each
  $(v,s)\in\Vset\times\SecSet$, and, for each $s$,
  $\sum_vF_v^{(s)}=\idon{\Qsys_s}$ with that sum verified to round nowhere, a
  failure being refused with the sector and the outcome index at which the
  first inexact addition occurred; and
  $\hat O_k^{(s)}=\hat O_k^{(s)\dagger}$ for each $(k,s)$. Every one of these
  is decided exactly. There is no off-block condition: an emitted
  $\Qsys$-operator is a \exttermref{../foundations/contract.pdf}{term-block-list}{block list}, so there is no $\Pi_sF_v\Pi_{s'}$ or
  $\Pi_s\hat O_k\Pi_{s'}$ to evaluate and nothing that could be non-zero.
\item[V4 (key alphabet)] every tabulated value of $r$ lies in
  $\Zval\sqcup\{\perp\}$.
\item[V5 (sifting consistency)] for all $(x,v)$ and $(x',v')$ with
  $a(x,v)=a(x',v')$: $r(x,v)=\perp\iff r(x',v')=\perp$.
\item[V6 (constraint typing, including the tag)] when the marginal
  slot is filled: $\bowtie_k\in\{\preceq,=,\succeq\}$, and the entry's
  right-hand side is of one of the two admitted types, which is what fixes its
  tag (Definition~\ref{def:homog}): \termref{term-homogeneous}{homogeneous}, whose datum is of
  real scalar type, or \termref{term-general}{general}, whose datum is a matrix on $A$ with
  $\Theta_k=\Theta_k^{\dagger}$ decided exactly. When the slot holds
  $\nomarg$: $\bowtie_k\in\{\le,=,\ge\}$ and the right-hand side is of real
  scalar type. A real scalar right-hand side is accepted in
  \exttermref{../transmitter/contract.pdf}{term-trusted-source-mode}{trusted-source mode} at every $d_A$, and is the licensed form there
  (Definition~\ref{def:homog}). The relation symbol, the type of the
  right-hand side and the tag that type fixes are compared exactly; an entry
  whose right-hand-side slot holds neither admitted object, or a complex-typed
  right-hand side in \exttermref{../transmitter/contract.pdf}{term-untrusted-source-mode}{untrusted-source mode}, is refused by type and not
  measured. The check also reads an object \eqref{eq:declaration} does not
  carry, the tag field of an \emph{arriving} entry, in $\Spec_{\mathrm{con}}$
  or transported by a link: where an arriving entry carries an explicit tag
  field, V6 compares that field with the type of the datum it names, and a
  disagreement is refused. An entry of \eqref{eq:declaration} in
  either of the two forms above is the triple
  $(\hat O_k,\bowtie_k,\mathrm{rhs}_k)$ and holds no tag field independent of
  the type of its datum, the emitted tag being read off that type
  (Definition~\ref{def:homog}); so on a candidate declaration this clause has
  no subject, refuses nothing, and is not among the conjuncts D(vi) is paired
  with (Theorem~\ref{thm:validator-soundness}, proof). An entry carrying no tag field is not
  refused: its tag is the type of its datum and is read off it, which is what
  lets an entry declared as
  a real constant, carrying no tag field, be read here.
\item[V7 (feasibility; one-sided, three branches)] V7 passes in exactly three
  cases:
  \begin{enumerate}[label=(\alph*),leftmargin=2.2em]
  \item $\CList=\varnothing$; or
  \item $\CList\ne\varnothing$ and the default candidate
    \eqref{eq:witness-empty}, which the compiler forms and does not emit,
    passes the per-block verification of Theorem~\ref{thm:feasibility}\,(ii); or
  \item $\CList\ne\varnothing$, the default candidate fails, and a witness
    $\omega^{\star}$ is supplied \emph{as a block list against $\Smap$} --- the
    map it carries compared with $\Smap$ in block kind, in identifiers, in
    their order and in their dimensions, exactly as V1 compares the map of a
    declared $\Qsys$-operator family --- and passes that same verification.
    This is the one place a witness's sector map is compared, and it is
    here rather than at V1 because $\omega^{\star}$ is a block list on
    $A\Qsys$, is no slot of \eqref{eq:declaration}, and is quantified over by
    no target condition (Definition~\ref{def:target}).
  \end{enumerate}
  Otherwise V7 fails and the candidate is refused, the refusal naming, where a
  witness was supplied and rejected on its map, the first position at which
  the two lists differ. The verification in (b) and (c) is the $|\SecSet|$
  block positivity decisions --- Hermiticity and semidefiniteness per block
  --- the trace, the marginal equality in trusted-source mode, and one leg per
  constraint entry, all decided exactly, on quantities formed in exact
  arithmetic (Theorem~\ref{thm:feasibility}\,(ii)). It is taken on a
  candidate satisfying D(i)--D(iii) and D(vi), which V1--V3 and V6 decide; a
  candidate failing one of them is refused there. V7 is the one check of the
  list that
  is \emph{not} equivalent to its target condition: it implies D(vii) and is
  implied by nothing weaker than a candidate that verifies
  (Theorem~\ref{thm:validator-soundness}\,(iii)).
\item[V8 (party computability)] $a_A$ and $a_B$ are tabulated on
  $\Oset{\mathsf A}$, $\Oset{\mathsf B}$; and for all $(x,v)$, $(x',v')$ with
  $\rd{\hold}(x,v)=\rd{\hold}(x',v')$ and $a(x,v)=a(x',v')$:
  $r(x,v)=r(x',v')$.
\end{description}

\emph{The general format.} On a candidate declaration in the general format
of Definition~\ref{def:declaration}\,(DX1)--(DX6) the eight checks read as
follows, every operator condition again being decided exactly. In the
\termref{term-basic-format}{basic format} --- one Alice sector, one \termref{term-measured-factor}{measured factor},
\termref{term-povm-form}{POVM form}, two messages neither reading the other and no party-local
entry --- each reading is the check of the
same name above, conjunct for conjunct, so on such a candidate the two
readings admit and refuse alike.

\emph{V1 in the general format.} Finiteness and totality as above, every
table of the transcript included: each $a_m$ is tabulated on all of
$\Oset{\pi_m}\times\mathcal{C}_1\times\dots\times\mathcal{C}_{m-1}$, which is
also where the signature clause of D(viii) lives: a table whose stated
domain names a later message's alphabet is not of that type and is refused
here. \emph{The sector map is one factorised map}: at most one factor tagged
$A$ and, where there is one, first in the list; at least one factor tagged
$\Qsys$; \emph{one block kind across the factors}, a disagreement being
refused, naming the first factor at which the kind differs;
$d_{\Qsys_f}=\sum_{s_f}d_{f,s_f}$ for every factor and
$d_{\Qsys}=\prod_fd_{\Qsys_f}$; and, where an $A$ factor is declared,
$d_A=\sum_td_{A_t}$. The restriction of $\idxm$ to each factor's \exttermref{../reductions/contract.pdf}{term-sector-set}{sector set}
is the total order that factor's own map fixes, and its restriction to
$\mathbb{T}$ is the order the $A$ factor's map fixes; the product orders on
$\SecSet$ and $\Vset$ and the order on the \termref{term-joint-sector}{joint sectors} are derived
by (DX1) and (DX3) and are not compared, there being no second datum to
compare them with. \emph{And the maps agree, factor by factor}: the map
carried by each declared family $\{F^{f}_{v_f}\}_{v_f}$ is compared against
the map of factor $f$, and the map carried by each entry's operator
$\hat O_k$ against the map of that entry's scope, in the same four
components, with the same refusal naming the first differing position.
\emph{Alice's \exttermref{../foundations/contract.pdf}{term-declared-data}{declared data} are compared too, against the map of the
factor tagged $A$}: the map carried by $\tau_A$, the map carried by each
operator of Alice's measurement slot --- $M_x$ in POVM form, $N_x$ in
\termref{term-ensemble-form}{ensemble form} --- and the map carried by each $\Theta_k$ of an entry declared
general are each compared against that factor's map in the same four
components (block kind, block identifiers, their order, dimensions), with the
same refusal naming the first differing position. That comparison is what
makes condition~(T) of Definition~\ref{def:exact} decided for Alice's data,
which its general-format paragraph calls non-empty for them; it is decided
here and nowhere else, and without it a candidate whose Alice block list
carried a different block kind, a different identifier list or a different
order would pass every check and violate D(i). Where no factor is tagged $A$
there is one Alice block and no map to compare, as in V1 above. \emph{The list
of compared objects is again exhaustive by construction}: the objects D(i)
quantifies over are the per-factor families, the entries' operators and
Alice's declared data. The product family is no datum, and the
\termref{term-lifted-operator}{lifted operator} of a scoped entry is no datum; the descriptors that the
last links of a chain produce, decided per factor at
Definition~\ref{def:chaincond}\,(C1), and a supplied witness, compared at
branch~(c) of V7, stay outside the list; and the physical families stay outside
it too, one per factor.

\emph{V2 in the general format.} The slot holds $\nomarg$, or a block list
$(\tau^{(t)}_A)_t$ against the $A$ factor's map with
$\tau^{(t)}_A=\tau^{(t)\dagger}_A$ and $\tau^{(t)}_A\succeq0$ for every $t$
and $\sum_t\tr\tau^{(t)}_A=1$, each trace accumulated over the row index and
the block traces then added in the \exttermref{../foundations/contract.pdf}{term-declared-order}{declared order} of $\mathbb{T}$, with
\emph{every} one of those additions verified to round nowhere; a failure of
the last is refused. \emph{Where Alice's measurement slot is in ensemble form,
additionally $\tau^{(t)}_A\ne0$ for every $t$}, a comparison of declared
entries with the datum $0$, a failure being refused naming the sector. At
$|\mathbb{T}|=1$ this is V2 above, the one block being the whole of $A$.

\emph{V3 in the general format.} \emph{Per measured factor}: for every
$f\in\mathbb{F}$, every $v_f\in\Vset^{f}$ and every $s_f\in\SecSet^{f}$,
$F^{f,(s_f)}_{v_f}$ Hermitian and positive semidefinite, and for each $s_f$
the sum $\sum_{v_f}F^{f,(s_f)}_{v_f}=\idon{\Qsys_{f,s_f}}$ with its
declared-order sum verified to round nowhere, a failure being refused with
the \emph{factor}, the sector and the outcome index at which the first
inexact addition occurred. \emph{On Alice's data, block by block along
$\mathbb{T}$}: in \termref{term-povm-form}{POVM form} $M^{(t)}_x$ Hermitian and positive
semidefinite for every $(x,t)$ and $\sum_xM^{(t)}_x=\idon{A_t}$ for every
$t$, the sum in the emitted index order verified to round nowhere, a failure
being refused with the Alice sector; in \termref{term-ensemble-form}{ensemble form} $N^{(t)}_x$
Hermitian and positive semidefinite for every $(x,t)$ and, \emph{in place of
the completeness identity}, the identity $\sum_xN^{(t)}_x=\tau^{(t)}_A$ of
\eqref{eq:ensembleform} decided entry by entry against the declared block of
$\tau_A$, a failure being refused with the sector and the first differing
entry, with that declared-order sum verified to round nowhere, a failure
being refused with the sector and the index of the first inexact addition.
\emph{An ensemble-form slot beside a marginal slot holding $\nomarg$ is
refused}; this is D(iii)'s general-format clause, and it is a decision on
which of two slots is filled, not on a number. And $\hat O_k^{(s_{f_k})}$ is
Hermitian for every entry and every block of its scope. \emph{Nothing is
decided on a product block \eqref{eq:productfamily}, on a lifted operator or
on the denoted family \eqref{eq:ensemblepovm}}: none of the three is a datum,
and that the first is a POVM on
$\Qsys$ and the third a POVM on $A$ are theorems
(Definition~\ref{def:declaration}\,(DX3), Theorem~\ref{thm:sifting-and-key}\,(v)).

\emph{V4, V5, V7 in the general format.} V4 and V5 keep their wording, with
$a$ the function \eqref{eq:transcript} derives from the transcript's tables;
V5's pairs are compared on the values of that function and on nothing about
how it was built. V7 keeps its three branches and its refusal, and what its
branches (b) and (c) verify is the per-\emph{joint}-sector verification of
Theorem~\ref{thm:feasibility}\,(ii) in the general format, on a candidate
satisfying D(i)--D(iii) and D(vi), in their general-format reading, which
V1--V3 and V6 decide (a candidate failing one of them is refused there): the
default candidate is
\eqref{eq:witness-empty} read per joint sector, and a supplied
$\omega^{\star}$ is required to be a block list over the joint sectors, its
map compared in the same four components at branch~(c) and there alone.

\emph{V6 in the general format.} When the marginal slot is filled:
$\bowtie_k\in\{\preceq,=,\succeq\}$, and the entry's right-hand side is of
one of the \emph{three} admitted forms, which with the arriving tag field is
what fixes its tag (Definition~\ref{def:homog}): homogeneous with a real
scalar, general with a block list $(\Theta^{(t)}_k)_t$ along $\mathbb{T}$
decided Hermitian block by block, or \termref{term-party-local}{party-local} with a real scalar. When the
slot holds $\nomarg$: the right-hand side is of real scalar type, either
untagged or tagged party-local. \emph{The relation symbol is decided by form
and not by mode alone}: for an \emph{untagged} entry
$\bowtie_k\in\{\preceq,=,\succeq\}$ when the slot is filled and
$\bowtie_k\in\{\le,=,\ge\}$ when it holds $\nomarg$, exactly as in V6 above;
for an entry declared party-local, $\bowtie_k\in\{\preceq,=,\succeq\}$
\emph{in either case}, its reading being an operator relation
(Definition~\ref{def:homog}, the relation-symbol clause of the \exttermref{../reductions/contract.pdf}{term-party-local-form}{party-local
form}; Definition~\ref{def:target}, D(vi) in the general format), with the
scalar symbols admitted as well where every row of that entry is $1\times1$
--- every Alice block, and every sector block of every measured factor other
than the entry's scope, of dimension one, $d_A=1$ at one measured factor
included --- the three operator symbols being the three scalar symbols there.
This is a decision on a declared symbol beside a declared tag and measures
nothing, and at an entry not declared party-local it is the clause of V6
above word for word, so nothing that clause admits or refuses moves.
\emph{Every entry's scope names a measured factor of the emitted sector map},
a comparison of a label against the factor list, a failure being refused
naming the entry; this is the general-format half of D(vi) and is decided and
never measured, as the relation symbol is. \emph{A party-local entry is the
one that arrives carrying that tag value}, and it carries a tag field the
type of its datum does not determine, that type not separating it from a
homogeneous one. On a candidate carrying one, the arriving-tag clause of V6
above therefore has a subject, and it is this comparison: the datum is of the
type \emph{that} tag admits, a real scalar, and the scope names a measured
factor; a party-local tag beside a matrix right-hand side is refused. The
sentence of V6 above that an entry holds no tag field independent of the type
of its datum is true of the two untagged forms, homogeneous and general, and
of them alone. \emph{Nothing is inferred}: an entry arriving without the
party-local tag is never read party-locally. So the arriving-tag clause
refuses no party-local entry beside a real scalar, which a literal reading of it
would otherwise do, and wherever no entry is tagged the reading is that of V6
above (Definition~\ref{def:homog}).

\emph{V8 in the general format.} The announcement clause reads the
transcript: every $a_m$ is tabulated on all of its stated domain, which is
V1's clause restated where D(viii) pairs with it; the key clause is unchanged,
with $a$ the derived function.
\end{definition}
\claimmeta{\metalabel{Uses}\mref{Def.}{def:declaration}, \mref{Def.}{def:exact}, \mref{Def.}{def:target}, \mref{Def.}{def:homog}, \mref{Def.}{def:spec}, \mref{Def.}{def:compile}, \mref{Thm.}{thm:feasibility}
\metasep \metalabel{Used by}\mref{Thm.}{thm:typing}, \mref{Thm.}{thm:measurements-and-statistics}, \mref{Rem.}{rem:condfamily}, \mref{Thm.}{thm:validator-soundness}, \mref{Def.}{def:chaincond}, \mref{Rem.}{rem:chain-and-validator}, \mref{Thm.}{thm:chain-consumption}, \mref{Req.}{req:tag-reading}, \mref{Req.}{req:no-window}, \mref{Req.}{req:validator-refusals}, \mref{Req.}{req:feasibility-witness}
\metasep newly derived
\metasep implemented (effective: accepted)
\metasep \metalabel{Code}\coderef{https://github.com/fraunhofersg/certiqs-key/blob/v0.1.0/src/certiqs_key/protocol_compiler/validation.py}{certiqs_key.protocol_compiler.validation:validate}
\metasep \metalabel{Oracles}\coderef{https://github.com/fraunhofersg/certiqs-key/blob/v0.1.0/tests/oracles/registry_protocol_compiler.py}{registry_protocol_compiler.py}}

% claim protocol-compiler/thm:validator-soundness
\claimsummary{protocol-compiler/thm:validator-soundness}{Passing the validator's eight checks establishes every target condition of a well-formed declaration, and every check but the feasibility check is complete, so a well-formed declaration is refused only when no default or supplied candidate is verified feasible. Every decision is finite and exact.}
\claimbookmark{protocol-compiler/thm:validator-soundness}{Theorem}{thm:validator-soundness}{Soundness, completeness and decidability of the validator}
\claimleadin{%
The consuming engine takes D(i)--D(viii) as hypotheses without re-deriving
them, so what a pass of the validator buys has to be exactly those
conditions. The theorem pairs each check of Definition~\ref{def:validator}
with its target condition of Definition~\ref{def:target}, names the one check
that is not equivalent to its condition, and counts what the checks compute.
Remark~\ref{rem:nosearch} records why an empty \termref{term-feasible-set}{feasible set} is the dangerous
error and a spurious refusal is not. For the clause of (iv) that the
validator assembles no array larger than one matrix of size $d_A$, $d_s$ or
$d_Ad_s$ for a single sector $s$, see also
Requirement~\ref{req:operator-form}, by which the engine holds every operator
on $\Qsys$ as its block list.
}
\begin{mainclaim}
\begin{theorem}[Soundness, completeness and decidability of the validator]
\label{thm:validator-soundness}
For a candidate declaration with finite tabulated data:
\begin{enumerate}[label=(\roman*),leftmargin=2.2em]
\item \emph{(soundness, on all eight conditions)} if V1--V8 all
pass then D(i)--D(viii) all hold, of the declared numbers;
\item \emph{(completeness, on seven of the eight)} if D(i)--D(vi)
and D(viii) hold then V1--V6 and V8 all pass;
\item \emph{(V7 is one-sided, and the gap is exactly its third branch)} the
converse of (i) fails, and it fails only at V7: D(vii) does \emph{not} imply
V7. For a witness, take \exttermref{../transmitter/contract.pdf}{term-untrusted-source-mode}{untrusted-source mode} with $\Qsys$ a qubit
carrying the single sector $\SecSet=\{s_0\}$, every other slot well-formed,
and the one-entry list
$\CList=\{(\dyad{0}_{\Qsys},\ \ge,\ 1)\}$, whose entry reads
$\tr[(\idon{A}\otimes\dyad{0})\,\omega]\ge1$. Then $\hat O_1$ is a
one-block list, Hermitian, the typing is that of D(vi), and
$\omega=\eta\otimes\dyad{0}$ is feasible for any state $\eta\in\Seq{A}$ --- it
is Hermitian and positive semidefinite outright, has unit trace, and its
constraint value is exactly $1$ --- so $\Feas\ne\varnothing$ and D(i)--D(viii)
all hold. But $\CList\ne\varnothing$, and the default candidate
\eqref{eq:witness-empty} that branch (b) of V7 tries first gives
$\tr[(\idon{A}\otimes\dyad{0})(\idon{A}/d_A\otimes\idon{\Qsys}/2)]=\tfrac12$,
which fails $\ge1$; so with no witness supplied branch (c) is unavailable, V7
fails and the \termref{term-well-formed-declaration}{well-formed declaration} is refused. Exactly the same
happens for every declaration with $\CList\ne\varnothing$ presented without a
witness whose \termref{term-feasible-set}{feasible set} is non-empty but does not contain
\eqref{eq:witness-empty}. The refusal errs on the safe side --- a refused
declaration is never consumed --- and the incompleteness is deliberate: deciding D(vii) in general is a
semidefinite feasibility problem this module does not solve;
\item \emph{(decidability, and what is computed)} each of V1--V6 and V8 is
decided by finite enumeration over the declared alphabets together with
finitely many arithmetic operations on the declared matrices; V7 is decided in
the same sense on the candidate of its branch (b) or (c), and is otherwise a
refusal. The arithmetic part consists of \emph{exact decisions}: an
identity, a semidefiniteness, or the exactness of an addition. The following
counts give the complete applicable predicate inventory for well-typed data,
not the number an early-terminated execution necessarily performs.
Write $I_{\mathrm{tr}}=1$ in trusted-source mode and $I_{\mathrm{tr}}=0$ when
the marginal slot holds $\nomarg$. The complete arithmetic inventory of V2,
V3 and V6 is
\begin{equation}\label{eq:predcount}
  \underbrace{3I_{\mathrm{tr}}}_{\text{V2}}
  \;+\;\underbrace{2|\Xset|+1}_{\text{V3 on }A}
  \;+\;\underbrace{|\SecSet|\,\big(2|\Vset|+1\big)}_{\text{V3 on }\Qsys}
  \;+\;\underbrace{m\,|\SecSet|}_{\text{V3 on }\CList}
  \;+\;\underbrace{m_{\text{gen}}}_{\text{V6}}
\end{equation}
predicate decisions, where $m_{\text{gen}}\le m$ counts the entries declared
general --- the only ones carrying a matrix whose Hermiticity is
decided; V6's remaining content is $m$ comparisons of a relation symbol, a tag
and a type, which are decided and not computed. \emph{Beside these stand the
summation decisions of condition~(S)}: one on $\tr\tau_A$ in trusted
mode and none on the absent marginal in untrusted mode; one on
$\sum_xM_x$; and one per sector on $\sum_vF_v^{(s)}$, so
$|\SecSet|+1+I_{\mathrm{tr}}$ of them,
each a pass over the additions of one declared-order sum. And V7, on the
candidate it verifies, takes $2|\SecSet|+1$ decisions --- Hermiticity and
semidefiniteness per block, and the trace --- one further in trusted mode for
the marginal equality, and one decision per constraint entry: semidefiniteness
for an inequality, exact entrywise equality for the relation $=$
(Theorem~\ref{thm:feasibility}\,(ii)); in untrusted mode a constraint leg is one
comparison of real numbers instead. Every one of them is taken on a matrix of
size $d_A$, $d_s$ or $d_Ad_s$ for a single $s$; \emph{the validator
assembles no array larger than one such matrix}, which where $|\SecSet|=1$ is
$d_{\Qsys}$ or $d_Ad_{\Qsys}$ square, the one block being the whole space
there (Theorem~\ref{thm:band-a}\,(i)).
In trusted mode the V2 term is $3$ and not $4$: D(ii)'s named Hermiticity
condition and the Hermiticity half of its positivity condition are the same
decision, taken once. In untrusted mode V2 decides only the marginal slot's
$\nomarg$ branch, with no marginal arithmetic or marginal summation predicate.
\end{enumerate}

\emph{The general format.} All four parts hold for a candidate declaration in
the general format of Definition~\ref{def:declaration}\,(DX1)--(DX6), with the
checks read as in the general-format paragraph of
Definition~\ref{def:validator} and the conditions as in that of
Definition~\ref{def:target}: (i) soundness on all eight and (ii) completeness
on D(i)--D(vi) and D(viii); (iii) with the \emph{same} witness, which is a
general-format declaration at $|\mathbb{T}|=1$, one \termref{term-measured-factor}{measured factor} and
no party-local entry; and (iv) with the predicate decisions of V2, V3 and V6
numbering
\begin{equation}\label{eq:predcount-general}
  \underbrace{I_{\mathrm{tr}}\big(2|\mathbb{T}|+1
    +|\mathbb{T}|\,\mathbb 1[\text{ensemble form}]\big)}_{\text{V2}}
  \;+\;\underbrace{|\mathbb{T}|\,\big(2|\Xset|+1\big)}_{\text{V3 on }A}
  \;+\;\underbrace{\sum_{f\in\mathbb{F}}|\SecSet^{f}|
    \big(2|\Vset^{f}|+1\big)}_{\text{V3 per measured factor}}
  \;+\;\underbrace{\sum_{k=1}^{m}|\SecSet^{f_k}|}_{\text{V3 on }\CList}
  \;+\;\underbrace{m_{\text{gen}}\,|\mathbb{T}|}_{\text{V6}} ,
\end{equation}
beside them the summation decisions of condition~(S) --- one on
$\tr\tau_A$ in trusted mode only, one per Alice sector on Alice's family,
and one per factor and sector on that factor's family, so
$|\mathbb{T}|+\sum_f|\SecSet^{f}|+I_{\mathrm{tr}}$ of them --- and V7, on the
candidate it verifies, $2|\mathbb{T}|\,|\SecSet|+1$ decisions, $|\mathbb{T}|$
further in trusted mode for the marginal equalities, and per entry $k$ either
$|\mathbb{T}|$ operator comparisons, one comparison of real numbers in
untrusted mode, or, for a party-local entry,
$|\mathbb{T}|\,|\SecSet|/|\SecSet^{f_k}|$ operator comparisons.
Each operator comparison is one semidefiniteness decision for an inequality
or one exact entrywise equality decision for $=$
(Theorem~\ref{thm:feasibility}\,(ii), general format). \emph{Every one of them is
taken on a matrix of size $d_{A_t}$, $d_s$ or $d_{A_t}d_s$ for a single joint
sector, or of size $d_{A_t}d_s/d_{f,s_f}$ for a single party-local row;
the validator assembles no array larger than one joint block.} In the
\termref{term-basic-format}{basic format} --- $|\mathbb{T}|=1$, one measured factor, POVM form, two
messages neither reading the other and no party-local entry ---
\eqref{eq:predcount-general} is
\eqref{eq:predcount} term for term, and every other count of this paragraph is
the corresponding count of (iv).
\end{theorem}
\end{mainclaim}
\claimmeta{\metalabel{Uses}\mref{Def.}{def:declaration}, \mref{Def.}{def:exact}, \mref{Def.}{def:target}, \mref{Def.}{def:feasible}, \mref{Def.}{def:homog}, \mref{Thm.}{thm:measurements-and-statistics}, \mref{Thm.}{thm:feasibility}, \mref{Def.}{def:validator}
\metasep \metalabel{Used by}\mref{Rem.}{rem:condfamily}, \mref{Rem.}{rem:completeness}, \mref{Thm.}{thm:chain-consumption}, \mref{Def.}{def:eligible}, \mref{Thm.}{thm:gates}, \mref{Req.}{req:validation-transcript}, \mref{Req.}{req:witness-refusal-meaning}
\metasep newly derived
\metasep implemented (effective: accepted)
\metasep \metalabel{Code}\coderef{https://github.com/fraunhofersg/certiqs-key/blob/v0.1.0/src/certiqs_key/protocol_compiler/validation.py}{certiqs_key.protocol_compiler.validation:decidability_counts}
\metasep \metalabel{Oracles}\coderef{https://github.com/fraunhofersg/certiqs-key/blob/v0.1.0/tests/oracles/registry_protocol_compiler.py}{registry_protocol_compiler.py}}

\begin{proof}
\emph{Idea.} Parts (i) and (ii) are proved together, by pairing each check
with its target condition. Three kinds of conjunct arise: combinatorial ones,
which both sides state as the same comparison of labels, integers or
tabulated values; numerical ones, which both sides state as the same exact
predicate on the same declared numbers; and structural ones, which the block
type makes true and neither side states. Two pairs, V5/D(v) and V8/D(viii),
are combinatorial but need an argument. V7 implies D(vii), which completes
(i), and (iii) is the counterexample stated with it. Part (iv) counts the
decisions. The general format takes the same steps, the point being that on
each side the quantifiers still run over one list.

\emph{Step 1 (combinatorial conjuncts: the restatement is literal).} V1's
finiteness, totality, dimension-sum, index-order and sector-map clauses; V4's
key-alphabet clause; V6's relation-symbol, tag and type clauses. Each compares
labels, integers or tabulated values, so V1/D(i), V4/D(iv) and the
combinatorial half of V6/D(vi) are equivalences in both directions. For V1's
sector-map clause we \emph{name} both lists rather than identify them: D(i)
asks that there be one map and that every declared $\Qsys$-operator be a
block list against \emph{it}, and the declared $\Qsys$-operators of
\eqref{eq:declaration} are exactly the family $\{F_v\}_{v\in\Vset}$ together
with the $m$ entries $\hat O_1,\dots,\hat O_m$ of $\CList$
(Definition~\ref{def:declaration}); V1 compares the four components of the
map carried by \emph{each of those same objects} against the emitted one, and
compares the map of nothing else (Definition~\ref{def:validator}). The two
quantifiers therefore run over one list, and two finite lists of labels and
integers agree or they do not. \emph{Neither the descriptor that the chain's
last link produces nor a supplied feasibility witness is on either side}: the
first is a
component of the chain record, decided at
Definition~\ref{def:chaincond}\,(C1), and the second is a block list on
$A\Qsys$ tested at branch~(c) of V7 --- so a candidate meeting D(i)--D(vi)
and D(viii) cannot be refused at V1 on account of either, which is what (ii)
asserts and what a wider V1 would have made false. \emph{V1's index-order
clause is the last clause of D(i)}, word for word: both ask that the
restriction of $\idxm$ to $\SecSet$ be the total order $\Smap$ fixes, and
that comparison is again of two finite lists of labels. Both directions
therefore hold of it as well, which is what makes (ii) true of a
\emph{candidate}: a candidate meeting D(i)--D(vi) and D(viii) meets that
clause too, so V1 cannot refuse it there.

\emph{V6's combinatorial clauses pair with D(vi) on the same $m$ objects.}
D(vi) asks of each entry of $\CList$ that its relation symbol lie in the set
its emitted source mode admits and that its right-hand-side slot hold one of
the two admitted data, whose type \emph{is} the tag
(Definition~\ref{def:homog}); V6 decides those two things, of those same $m$
entries, keyed on whether the marginal slot is filled --- which is the
emitted source mode (Definition~\ref{def:declaration}). V6's remaining
clause, that an explicit tag field on an \emph{arriving} entry agree with the
type of the datum it names, is not a further conjunct on
\eqref{eq:declaration}: no entry of a declaration carries such a field, so on
a candidate that clause has no subject and cannot refuse one
(Definition~\ref{def:validator}, V6). The pairing is therefore between equal
lists here as well.

\emph{Step 2 (numerical conjuncts: the two sides are the same exact
predicate).} V2's three, V3's five and V6's Hermiticity of a $\Theta_k$
declared general are, by construction, the same conditions on the same
declared numbers as the corresponding conjuncts of D(ii), D(iii) and D(vi):
both Definition~\ref{def:target} and Definition~\ref{def:validator} state
them through Definition~\ref{def:exact}, exactly, with no window on either
side. So V2/D(ii), the five-predicate half of V3/D(iii) and the Hermiticity
half of V6/D(vi) are equivalences in both directions, a literal restatement.
Condition~(S) pairs the same way: it is a clause of D(ii) and of D(iii)
(Definition~\ref{def:target}) and a clause of V2 and V3, stated once in
Definition~\ref{def:exact} and read by both.

\emph{Step 3 (structural conjuncts: there is nothing left to pair).} D(iii)'s
last two conjuncts, $\Pi_sF_v\Pi_{s'}=0$ and $\Pi_s\hat O_k\Pi_{s'}=0$, hold
by the emitted type (Theorem~\ref{thm:measurements-and-statistics}\,(ii)) and have no counterpart in
V3. They are therefore vacuous in the soundness direction --- V3 need not
establish what cannot fail --- and vacuous in the completeness direction ---
a declaration satisfying them is not thereby subjected to any check. The
pairing is not broken by their absence; it is one conjunct shorter on each
side. One direction of (ii) does fail for a reason that is not numerical, and
it is a \emph{typing} failure: an operator presented outside the block type,
as a full $d_{\Qsys}\times d_{\Qsys}$ array, is refused at V1 whatever its
off-block entries are, exactly zero included, because it is not a declaration
in the sense of Definition~\ref{def:declaration}. That is why (ii) is stated
for declarations rather than for operators.

\emph{Step 4 (V5 $\Leftrightarrow$ D(v) and V8 $\Leftrightarrow$ D(viii)).}
These two pairs are combinatorial but need an argument.

V5 $\Leftrightarrow$ D(v). ($\Leftarrow$) If $\varsigma$ exists as in
D(v) and $a(x,v)=a(x',v')=\kappa$, then $r(x,v)=\perp$ and $r(x',v')=\perp$ are
both equivalent to $\varsigma(\kappa)=\mathrm{fail}$, so V5 holds.
($\Rightarrow$) Assume V5 and define
$\varsigma(\kappa)\coloneqq\mathrm{fail}$ if there exists
$(x,v)\in a^{-1}(\kappa)$ with $r(x,v)=\perp$, and
$\varsigma(\kappa)\coloneqq\mathrm{pass}$ otherwise (in particular on
$\kappa\notin\operatorname{ran}a$, where the choice is immaterial because the
condition of D(v) quantifies over pairs $(x,v)$). This $\varsigma$ is total.
For $(x,v)$ with $\kappa=a(x,v)$: if $r(x,v)=\perp$ then
$\varsigma(\kappa)=\mathrm{fail}$ by construction; if
$\varsigma(\kappa)=\mathrm{fail}$ then some $(x',v')\in a^{-1}(\kappa)$ has
$r(x',v')=\perp$, and V5 gives $r(x,v)=\perp$. Hence D(v).

V8 $\Leftrightarrow$ D(viii). ($\Leftarrow$) If \eqref{eq:partycomp}
holds with a total $r_{\hold}$, then equal holder readouts and equal
announcements give equal key values, which is V8; the announcement clause of
D(viii) is the first clause of V8. ($\Rightarrow$) Assume V8 and define
$r_{\hold}(o,\kappa)\coloneqq r(x,v)$ for any $(x,v)$ with
$\rd{\hold}(x,v)=o$ and $a(x,v)=\kappa$ if such a pair exists --- well defined
precisely by V8's implication --- and $r_{\hold}(o,\kappa)\coloneqq\perp$ on
pairs $(o,\kappa)$ that are not realized. This $r_{\hold}$ is total on
$\Oset{\hold}\times\Kann$ and satisfies \eqref{eq:partycomp} for every $(x,v)$,
since every such pair realizes $(\rd{\hold}(x,v),a(x,v))$. Hence D(viii). The
value chosen on unrealized pairs is immaterial: \eqref{eq:partycomp} quantifies
over $(x,v)$ only.

Steps 1 to 4 account for (ii) in full, and for (i) except at V7.

\emph{Step 5 (V7 $\Rightarrow$ D(vii), which completes (i); and (iii)).}
Assume V1--V8 all pass. By the equivalences of Steps 1 to 3, D(i)--D(iii)
hold. If V7 passed by branch (a), $\CList=\varnothing$ and
Theorem~\ref{thm:feasibility}\,(i) applies --- its hypotheses are exactly
D(i)--D(iii) --- and exhibits $\omega^{\star}\in\Feas$. If V7 passed by
branch (b) or (c), the candidate it verified satisfies every condition of
Definition~\ref{def:feasible}, decided exactly (Theorem~\ref{thm:feasibility}\,(ii),
whose further hypothesis D(vi) holds because V6 passed, by Steps 1 and 2),
so that candidate lies in $\Feas$. Either way D(vii) holds. Statement (iii) is
proved by the counterexample given with it, which is checked there against
Definition~\ref{def:target} condition by condition.

\emph{Step 6 (decidability, (iv)).} V1 enumerates the declared tables and
compares finitely many declared integers and labels, the sector-map
comparison included: it is $|\SecSet|$ identifier comparisons, $|\SecSet|$
dimension comparisons, one block-kind comparison and one length comparison,
per supplied family, plus the comparison of $\idxm$'s sector order with
$\Smap$'s. V2, V3 and V6 take the decisions \eqref{eq:predcount} counts, on
matrices of the declared finite dimensions: in trusted mode $3$ for V2 ---
$\tau_A=\tau_A^{\dagger}$, $\tau_A\succeq0$, $\tr\tau_A=1$, the first serving
both the Hermiticity condition of D(ii) and the Hermiticity half of its
positivity condition; in untrusted mode none, because V2 takes its
$\nomarg$ branch; $2|\Xset|+1$ for V3's Alice half, one Hermiticity and one
semidefiniteness per $x$ and one completeness identity on $A$;
$|\SecSet|(2|\Vset|+1)$ for its Bob half, the same per $(v,s)$ and one
completeness identity per $s$; $m\,|\SecSet|$ Hermiticity decisions for its
constraint half; and $m_{\text{gen}}$ for V6. Each is finite: an identity is
a comparison of entries, and a semidefiniteness on a Hermitian matrix of size
$n$ is decided by exact symmetric elimination with the zero-pivot row test,
$O(n^3)$ exact rational operations. The summation decisions of (S) are one
pass over the additions of each applicable declared-order sum; the marginal
trace sum exists only in trusted mode. Early termination can omit later
decisions within a check, which is why this complete inventory need not
equal the length of a failed execution's transcript. V4, V5 and V8 enumerate
$|\Xset||\Vset|$ pairs, respectively $(|\Xset||\Vset|)^2$ ordered pairs, and
compare tabulated values. V7 is Theorem~\ref{thm:feasibility}\,(ii), which is likewise
per block, and \emph{on both of its branches the decision is taken in exact
arithmetic}: on branch (b) the candidate is formed by the compiler and decided
in exact rational arithmetic, and on branch (c) the supplied witness's
membership is decided exactly on the dyadic numbers it declares --- not
because those numbers are exactly representable, which by
Definition~\ref{def:exact} warrants nothing about the additions and products
a decision forms, but because every decision this module takes on a computed
quantity is taken exactly (\xref{foundations/req:exact-island}).

\emph{The general format.} The three kinds of conjunct are the same three, and
\emph{no fourth kind arises}; what has to be checked is that on each side the
quantifiers still run over one list. The steps are those above.

\emph{Step 1 in the general format (combinatorial conjuncts).} V1's
finiteness, totality, factor-list, dimension, index-order and map-comparison
clauses; V4's; V6's relation-symbol, tag, type and \emph{scope} clauses. Each
compares labels, integers or tabulated values, so each is an equivalence in
both directions once the two quantifiers are seen to run over one list, which
needs saying in four places. \emph{First, the factor list.} D(i) asks that
the sector map be one factorised map --- at most one factor tagged $A$ and
then first, at least one tagged $\Qsys$, one block kind, the dimension
identities --- and V1 decides those same clauses on the same list
(Definition~\ref{def:validator}, the general format), so they pair literally.
\emph{Second, the objects whose maps are compared.} In the general format the
declared $\Qsys$-operators are the per-factor families
$\{F^{f}_{v_f}\}_{v_f}$, one per $f\in\mathbb{F}$, together with the $m$
operators $\hat O_1,\dots,\hat O_m$, each against the map of its scope
(Definition~\ref{def:declaration}\,(DX1), (DX3), (DX6)); \emph{and Alice's
declared data are on the list beside them} --- $\tau_A$, every operator of her
measurement slot, and every $\Theta_k$ of an entry declared general, each
against the map of the factor tagged $A$, which is where
Definition~\ref{def:exact} states condition~(T) non-empty for them. D(i) asks
that there be one factorised map and that \emph{each} of those objects be a
block list against the map of its own factor; V1 compares the four
components of the map carried by each of \emph{those same} objects and of
nothing else. \emph{So the Alice-side conjunct is on both sides and pairs like
the others}, which is what makes (i) sound and (ii) complete for it: a
candidate whose Alice block list carries a different block kind, a different
identifier list or a different order is refused at V1, and a candidate
meeting D(i) is not refused there, both sides comparing two finite lists of
labels and integers. Where no factor is tagged $A$, Alice's data are single
arrays carrying no map and both sides are empty of them, as in Step 1 above.
\emph{A product block \eqref{eq:productfamily} and the lifted operator of a
scoped entry are on neither side}: neither is a \exttermref{../reductions/contract.pdf}{term-declared-datum}{declared datum}, so D(i) does
not quantify over either and V1 compares neither --- which is what keeps (ii)
true, a wider V1 refusing candidates D(i)--D(vi) and D(viii) admit. The
descriptors that the chain's last links produce, one per factor, and a
supplied witness stay outside both sides for the reasons of Step 1, now read
per factor and over the joint sectors. \emph{Third, the index order.} D(i) asks
that the restriction of $\idxm$ to each factor's sectors be the order that
factor's map fixes, and records that the product order is \emph{derived}; V1
decides exactly that and compares nothing about the derived orders, there
being no second datum to compare them against. Both directions hold, so a
candidate meeting D(i)--D(vi) and D(viii) cannot be refused there.
\emph{Fourth, the transcript.} D(viii) asks that the announcement slot be a
\termref{term-transcript-protocol}{transcript protocol} with every table $a_m$ total on
$\Oset{\pi_m}\times\mathcal{C}_1\times\dots\times\mathcal{C}_{m-1}$ --- so that
a message reads its sender's readout and the earlier messages \emph{by its
signature} --- and V1's totality clause together with V8's announcement clause
decide that same thing of those same $L$ tables. A table whose stated domain
names a later message's alphabet is not of the declared type and is refused at
V1; so on a candidate the signature clause is not a numerical conjunct and
needs no pairing beyond this one.

\emph{And V6 pairs with D(vi) on the same $m$ entries, now with three forms and
a scope.} D(vi) asks of each entry that its relation symbol lie in the set
\emph{its form and its emitted mode} admit --- the mode's set for an untagged
entry, $\{\preceq,=,\succeq\}$ in either mode for an entry declared
party-local, with the scalar symbols also admitted where every row of that
entry is $1\times1$ --- that its right-hand-side slot hold one of the three
admitted data --- whose type, together with the arriving tag field for the
party-local form, \emph{is} the tag (Definition~\ref{def:homog}) --- and that
its scope name a measured factor of the map; V6 decides those three things of
those same $m$ entries, keyed on whether the marginal slot is filled and on
the declared tag, and its symbol clause is D(vi)'s symbol clause set for set,
so that pair is an equivalence in both directions as before. The clause of V6
about an arriving tag field is again not a further conjunct on
\eqref{eq:declaration} for a homogeneous or general entry, whose tag is the
type; \emph{for a party-local entry it is the tag, because the type of the
datum does not separate that form from the homogeneous one}
(Definition~\ref{def:homog}), and D(vi) in its general reading asks for that
declared form, so here too the two sides ask the same thing of the same
objects.

\emph{Step 2 in the general format (numerical conjuncts).} V2's, V3's and
V6's Hermiticity of a $\Theta^{(t)}_k$ declared general are, by construction,
the same exact predicates on the same declared numbers as the corresponding
conjuncts of D(ii), D(iii) and D(vi), both sides stating them through
Definition~\ref{def:exact} in its general-format reading, exactly and with no
window: per measured factor and per factor block on the measured families,
and per Alice sector on Alice's data. Condition~(S) pairs the same way, per
factor and per Alice sector. \emph{Two pairs are new and both are
equivalences.} The ensemble form: V3's identity $\sum_xN^{(t)}_x=\tau^{(t)}_A$
with its summation condition, and V2's $\tau^{(t)}_A\ne0$, are exactly the
conditions Definition~\ref{def:exact} states for that form and D(ii), D(iii)
ask; and V3's refusal of an ensemble-form slot beside a marginal slot holding
$\nomarg$ pairs with D(iii)'s clause that ensemble form is refused in
untrusted-source mode --- a decision on \emph{which of two slots is filled},
hence combinatorial, hence an equivalence in both directions.

\emph{Step 3 in the general format (structural conjuncts: again nothing is
left to pair, and the list of them is longer).} D(iii)'s typing conjuncts
hold by the emitted type along $\mathbb{T}$ as well as along $\SecSet$ and
have no counterpart in V3; and \emph{three further statements that a naive
reading would put on the data are theorems and are conjuncts of neither
side}: that the product family $\{F_v\}_{v\in\Vset}$ is a POVM on $\Qsys$,
block diagonal for the product map (Definition~\ref{def:declaration}\,(DX3));
that the family \eqref{eq:ensemblepovm} an ensemble form denotes is a POVM on
$A$ --- proved at Theorem~\ref{thm:sifting-and-key}\,(v), \emph{named here as a pointer
and used by no step below}, what this step needs being only that neither
D(iii) nor V3 states it; and that a lifted operator is Hermitian and block
diagonal for the product map, which follows from the Hermiticity of
$\hat O^{(s_{f_k})}_k$ that V3 does decide. So each is vacuous in the
soundness direction --- V3 need not establish what follows --- and vacuous in
the completeness direction, a declaration satisfying it being subjected to no
check for it. The typing failure that makes (ii) a statement about
declarations rather than operators is unchanged, and has two more instances: a
family presented as one array over the product alphabet, and an Alice operator
presented as one array on $A$ where an $A$ factor is declared, are refused at
V1 whatever their entries are, because neither is a declaration in the sense
of Definition~\ref{def:declaration}.

\emph{Steps 4 and 5 in the general format (V5, V8 and V7; and (iii)).} The
proofs of V5 $\Leftrightarrow$ D(v) and V8 $\Leftrightarrow$ D(viii) in
Step 4 use of $a$ only that it is a total map into $\Kann$, and construct
$\varsigma$ and $r_{\hold}$ on that alphabet; both hold verbatim with
$\Kann=\prod_m\mathcal{C}_m$ and $a$ the function \eqref{eq:transcript}
derives, \emph{whose totality is derived here and not imported}:
$\kappa_1=a_1\circ\rd{\pi_1}$ is total because $a_1$ and $\rd{\pi_1}$ are, and
if $\kappa_1,\dots,\kappa_{m-1}$ are total on $\Xset\times\Vset$ then so is
$\kappa_m$, $a_m$ being tabulated on all of
$\Oset{\pi_m}\times\mathcal{C}_1\times\dots\times\mathcal{C}_{m-1}$ (V1) and
$\rd{\pi_m}$ total; induction on $m$ gives $a$ total into $\Kann$, which is
all the two proofs consume of it. For V7 $\Rightarrow$ D(vii): if V7 passed by
branch~(a) then $\CList=\varnothing$ and Theorem~\ref{thm:feasibility}\,(i) in its
general-format reading exhibits \eqref{eq:witness-general} in $\Feas$, its
hypotheses being exactly D(i)--D(iii); if by branch~(b) or (c), the candidate
it verified satisfies every line of \eqref{eq:feas-general}, decided exactly
(Theorem~\ref{thm:feasibility}\,(ii), general format, D(vi) holding because
V6 passed), so it lies in $\Feas$. Either
way D(vii) holds, and (i) is complete. \emph{And (iii) holds in the general
format with the witness already given}: that declaration has one Alice sector,
one measured factor, POVM form, two messages and no party-local entry, at which
values D(i)--D(viii) in their general reading are the eight conditions read of
it in (iii) (Definition~\ref{def:target}, the general format) and V7's three
branches are the three branches read of it there, so it is a general-format
declaration satisfying D(i)--D(viii) that V7 refuses. The incompleteness is
therefore at V7 and, by the equivalences just listed, only there.

\emph{Step 6 in the general format ((iv)).} V1 again enumerates declared
tables and compares declared integers and labels: per factor,
$|\SecSet^{f}|$ identifier and $|\SecSet^{f}|$ dimension comparisons, one
block-kind and one length comparison per supplied family against that
factor's map --- \emph{Alice's $\tau_A$, her measurement operators and each
$\Theta_k$ of an entry declared general against the $A$ factor's map among
them}, the same four comparisons each --- plus the comparisons of $\idxm$'s
components with the factors' orders and with $\mathbb{T}$'s. V2, V3 and V6
take the decisions \eqref{eq:predcount-general} counts: $2$ per Alice sector
for V2 --- Hermiticity, which serves both D(ii)'s named condition and the
Hermiticity half of its positivity condition, and semidefiniteness --- and one
trace decision, with one non-zero comparison per sector in ensemble form, all
of these V2 decisions and the marginal trace-summation decision being present
only in trusted mode; when the slot holds $\nomarg$, V2 returns without taking
any of them. V3 has $2|\Xset|+1$ per Alice sector for its Alice half, the
$+1$ being the completeness identity in POVM form and the marginal identity
in ensemble form; $|\SecSet^{f}|(2|\Vset^{f}|+1)$ per measured factor;
$|\SecSet^{f_k}|$ Hermiticity decisions per entry; and
$m_{\text{gen}}|\mathbb{T}|$ for V6. Each is finite and each is taken on a
matrix of one of the four sizes named in the statement, a semidefiniteness on
size $n$ costing $O(n^3)$ exact rational operations as before. V4, V5 and V8
enumerate $|\Xset||\Vset|$ pairs, respectively $(|\Xset||\Vset|)^2$ ordered
pairs, of \emph{tabulated} values, and $a$'s table is finite because each
$a_m$'s is. V7 is Theorem~\ref{thm:feasibility}\,(ii) in its general-format reading,
which is per joint block and in exact arithmetic on both of its branches, for
the reasons given there.
\end{proof}

% claim protocol-compiler/rem:completeness
\claimsummary{protocol-compiler/rem:completeness}{Soundness keeps an ill-formed declaration from being emitted; completeness, which holds for every target condition but feasibility, makes a refusal name a real defect. A refusal at V7 means only that no verified feasible state was shown.}
\claimbookmark{protocol-compiler/rem:completeness}{Remark}{rem:completeness}{The two directions of the validator}
\begin{remark}[The two directions of the validator]
\label{rem:completeness}
The two directions do different jobs and it matters which one is available
where. \emph{Soundness} --- Theorem~\ref{thm:validator-soundness}\,(i), and it holds for all
eight conditions --- is the safety direction: without it the compiler could
emit a declaration that passes every check and is nonetheless ill-formed, which
is the failure this module exists to prevent, since a downstream engine
consumes D(i)--D(viii) as hypotheses without re-deriving them. \emph{Completeness}
is what makes a \emph{refusal} informative: where it holds, a well-formed
declaration is never rejected, so a refusal names a real defect and not a
limitation of the checker. It holds for D(i)--D(vi) and D(viii)
(Theorem~\ref{thm:validator-soundness}\,(ii)); both directions of each of those pairs are
proved, not just the one that keeps bad declarations out. Both directions
are read on the declared numbers: a pass buys D(i)--D(viii) of the declared
numbers, which is what the consuming engine requires of them, and a
declaration well-formed of those numbers passes every check but V7.

Completeness does \emph{not} hold for D(vii), and that is a limitation of
this checker (Theorem~\ref{thm:validator-soundness}\,(iii)). \emph{Its scope is the third
branch of V7 and no wider}: a declaration with an empty constraint list
passes on branch (a), and one whose feasible set contains the default
candidate \eqref{eq:witness-empty} passes on branch (b), with nothing supplied
in either case. What is refused is a well-formed declaration with a non-empty
constraint list whose feasible set does not contain the default candidate and
for which no witness was supplied. Three consequences follow. A refusal at V7
means ``I was not shown a feasible state, and the one I tried is not one'',
never ``there is none'' (Requirement~\ref{req:witness-refusal-meaning}).
Whoever presents a declaration refused for that reason has a remedy inside
this interface --- supply $\omega^{\star}$
(Requirement~\ref{req:feasibility-witness}) --- and no other check has that
character. And the asymmetry is the deliberate one: the direction that is
missing costs a rejected good declaration, while the direction that is present
is the one whose absence would cost a consumed bad one.
\end{remark}
\claimmeta{\metalabel{Uses}\mref{Thm.}{thm:feasibility}, \mref{Thm.}{thm:validator-soundness}
\metasep \metalabel{Used by}\mref{Req.}{req:witness-refusal-meaning}
\metasep newly derived
\metasep accepted}

\section{Consuming the reduction chain}\label{group-protocol-compiler/def:chaincond}
\sectionintro{The chain arrives as a record of two segments, source links and receiver links, and the compiler decides what the record makes decidable: descriptor agreement, the outcome correspondence, the exact admissibility of the target measurement, and the pullback and source identities where channels are supplied as data. A non-empty source segment adds registry checks and two admission rules. \hyperlink{claim-protocol-compiler/thm:chain-consumption}{Theorem~\ref*{thm:chain-consumption}} states what a verified identity buys and names what is never discharged here: that each link is a sound reduction and that any restricted attack set is licensed. Where the reductions module discharges them is recorded, and no claim here relies on that discharge (\hyperlink{claim-protocol-compiler/rem:obligations-discharged}{Remark~\ref*{rem:obligations-discharged}}).}
\sectionimports{\extref{../transmitter/contract.pdf}{transmitter/def:source-segment}}

% claim protocol-compiler/def:chaincond
\claimsummary{protocol-compiler/def:chaincond}{A presented chain must match descriptors, map physical outcomes onto emitted ones, deliver an exactly admissible target measurement and, where channels are data, satisfy the pullback and source identities exactly. A non-empty source segment adds registry checks and two admission rules.}
\claimbookmark{protocol-compiler/def:chaincond}{Definition}{def:chaincond}{Interface conditions on a chain}
\claimleadin{%
The chain is the compiler's third input. A \emph{squashing} link replaces
Bob's measurement and therefore acts on the receiver space; a
\exttermref{../transmitter/contract.pdf}{term-source-link}{source link} replaces Alice's signal ensemble and therefore acts on the
\exttermref{../transmitter/contract.pdf}{term-transmitted-system}{transmitted system} $\Asig$, not on the receiver space at all. Every
link carries two descriptors, of what it starts from and of what it
produces: for a source link these are the transmitter's
\exttermref{../transmitter/contract.pdf}{term-source-descriptor}{source descriptor} and \exttermref{../transmitter/contract.pdf}{term-target-descriptor}{target descriptor}; for a receiver link
they are receiver-side, and the last receiver link produces the emitted
system with its \exttermref{../foundations/contract.pdf}{term-sector-map}{sector map}. One list of links cannot carry both kinds,
and a single descriptor rule cannot type both, so the record is a pair of
segments with a condition set each, and the composition-order rule is a
property of the pair rather than a tag. The warrant for that order is the
published one, \cite[Rem.~9.2]{tupkary2026decoy}, and on the module side
\xref{reductions/rem:ordering}, which states why a source link after a
receiver link is not free; this module types the record and states neither.

That the identity decided at (C4) and (S3) is the identity written is
\xref{foundations/req:exact-island} at work, and a producer repairs its own
emission under \xref{foundations/req:who-snaps}. (C1)--(C5), (S1)--(S4),
the licensed-only rule and the source-mode rule are decided where the chain
is consumed, before a declaration exists, so a failure produces a refusal and
no declaration (Requirement~\ref{req:chain-refusals}), and the \termref{term-gate-outcome-table}{gate outcome
table} of Definition~\ref{def:gate-outcome-table} needs no row for them. The grounds of the
licensed-only rule and of the source-mode rule are
Remark~\ref{rem:source-segment-rules}.
}
\begin{definition}[Interface conditions on a chain]\label{def:chaincond}
A \exttermref{../reductions/contract.pdf}{term-reduction-chain}{reduction chain} $\pcChain=(\Chsrc,\Chrec)$ presented to the compiler, with
$\Chsrc=(\Xi_1,\dots,\Xi_{L_{\mathrm s}})$ and
$\Chrec=(\Lambda_1,\dots,\Lambda_L)$, satisfies the
\termdef{term-interface-conditions-on-a-chain}{interface conditions on a chain}
when the following hold.

\emph{\exttermref{../reductions/contract.pdf}{term-receiver-segment}{Receiver segment}.}
\begin{enumerate}[label=(C\arabic*),leftmargin=3.0em]
\item \emph{(descriptor match)} the target descriptor of $\Lambda_i$ equals the
source descriptor of $\Lambda_{i+1}$ for $1\le i<L$; the source descriptor of
$\Lambda_1$ is the physical receiver space of $\Spec_{\mathrm{rec}}$; and the
target descriptor of $\Lambda_L$ is $(\Qsys,\Smap)$ --- that is, the last
link's target descriptor \emph{is} the sector map. That equality is
decided by the same four-component comparison V1 makes on the map of a
declared $\Qsys$-operator family (Definition~\ref{def:validator});
\item \emph{(outcome correspondence)} the declared outcome correspondence is a
tabulated total surjection $\gamma:\Vph\to\Vset$ from the \emph{physical}
alphabet onto the \emph{emitted} one --- the coarse-graining direction, of
which a bijection is the special case;
\item \emph{(\exttermref{../reductions/contract.pdf}{term-target-povm}{target POVM})} $\{F_v\}_{v\in\Vset}$ is a family of
\exttermref{../foundations/contract.pdf}{term-block-list}{block lists} against $\Smap$ --- \emph{typing}, held by the
emitted form and not checked --- and for each $s$ the block family
$\{F_v^{(s)}\}_v$ is a POVM on $\Qsys_s$, decided \emph{exactly} on the
emitted blocks themselves: $F_v^{(s)}$ Hermitian and positive semidefinite
for each $v$, $\sum_vF_v^{(s)}=\idon{\Qsys_s}$, and the declared-order sum
forming it rounding nowhere (Definition~\ref{def:exact}). This condition is
check V3 restricted to $\{F_v\}$ and states the same list of conjuncts
(Theorem~\ref{thm:chain-consumption}\,(ii));
\item \emph{(pullback, where checkable)} if every link of $\Chrec$ supplies its
channel as data then, with
$\Lambda\coloneqq\Lambda_L\circ\dots\circ\Lambda_1$,
\begin{equation}\label{eq:pullback}
  \Lambda^{\dagger}(F_v)=\sum_{w\in\gamma^{-1}(v)}\Fph_w
  \qquad\text{for every }v\in\Vset ,
\end{equation}
an \emph{exact} matrix identity, decided entry by entry on the emitted matrices
and on the supplied channel data, with the fibre sum accumulated in the
\exttermref{../foundations/contract.pdf}{term-declared-order}{declared order} of $\Vph$. Where some link supplies no channel as data,
only a theorem that constructs one, the identity is not decidable from the
emitted data and is carried instead as an obligation of the chain's
\exttermref{../foundations/contract.pdf}{term-producer}{producer}, Assumption~\ref{ass:link};
\item \emph{(slots)} each link carries its constraint entries and its licence
slot, both possibly empty, and its status tag.
\end{enumerate}

\emph{\exttermref{../transmitter/contract.pdf}{term-source-segment}{Source segment}.} The source segment satisfies the
\exttermref{../transmitter/contract.pdf}{term-interface-conditions}{interface conditions} (S1)--(S4) of
\xref{transmitter/def:source-segment} --- descriptor match with the
\exttermref{../transmitter/contract.pdf}{term-virtual-transmitted-system}{virtual transmitted system} $\Avirt$ as the source descriptor of
$\Xi_1$, on which the certificate of \xref{transmitter/def:srcert} is
written, and every descriptor Alice-side; index preservation; the source
identity $\Xi(\rvirt_x)=\rho_x$ for every $x$, that definition's equation
(S3), where every link supplies its channel as data; and the slots ---
\emph{decided by this module on the presented data}, exactly as that
definition states them. When $L_{\mathrm s}=0$ the segment is empty, $\Xi$ is
the identity, $\Avirt=\Asig$ and $\rvirt_x=\rho_x$, so (S3) holds by
inspection.

\emph{Between the segments.}
\begin{enumerate}[label=(X\arabic*),leftmargin=3.0em]
\item \emph{(order)} the chain is composed source segment first and receiver
segment second. This is not a tag and is not checked: it is the shape of the
record. A source-map link has no slot in $\Chrec$ to be written into, and a
link presented in the wrong segment fails that segment's descriptor rule ---
(C1) for a source map, since $\Asig$ is not the physical receiver space, and
(S1) for a squashing map, since $\Qsys$ is not an Alice-side descriptor.
\end{enumerate}

\emph{At which reading (C4) and (S3) are decided.} Both are \emph{matrix
equalities on \exttermref{../foundations/contract.pdf}{term-declared-data}{declared data}} --- a composite channel supplied as
binary64 data, applied to binary64 POVM elements on one side and to binary64
signal states on the other --- and both are decided \emph{exactly}, entry by
entry, on the numbers the compiler holds. \emph{The composite, the adjoint
and the fibre sum are formed in exact arithmetic as well}, so the identity
decided is the identity written; that
the operands are binary64 is not what makes the decision exact, a product of
two \exttermref{../foundations/contract.pdf}{term-binary64-numbers}{binary64 numbers} need not be one. Exactness here is an obligation of
the chain's producer (Assumption~\ref{ass:link}): a producer whose composite
lands off the identity by a rounding has not delivered the reduction the
interface asks for, and repairs its own emission; no window is opened here. Everything else
in this definition --- descriptors, alphabets, fibres, slots, segments --- is
combinatorial and is read exactly, as D(vi)'s relation symbols are.

\emph{The general format: one receiver segment per \termref{term-measured-factor}{measured factor}.}
For a chain presented as Definition~\ref{def:spec}\,(c$'$), with
$\Chrec=(\pcChain^{f}_{\mathrm{rec}})_{f\in\mathbb{F}}$, the conditions
(C1)--(C5) are decided \emph{of each segment separately, and of nothing
across two segments}. (C1): within segment $f$ the target descriptor of a link
is the source descriptor of the next, the source descriptor of
$\Lambda^{f}_1$ is factor $f$'s physical receiver system, and the target
descriptor of $\Lambda^{f}_{L_f}$ is
$(\Qsys_f,\SecSet^{f},(d_{f,s_f})_{s_f})$ --- \emph{that factor's own sector
map, and not the product map}, which is derived by
Definition~\ref{def:declaration}\,(DX1) and is the target descriptor of
nothing; the same four-component comparison decides it. (C2): one tabulated total surjection
$\gamma_f$ from factor $f$'s physical alphabet onto the alphabet factor $f$'s
last link delivers. (C3): the \exttermref{../reductions/contract.pdf}{term-target-family}{target family} $\{F^{f}_{v_f}\}$ is a family of
block lists against factor $f$'s map and, for each $s_f$, a POVM on
$\Qsys_{f,s_f}$ decided exactly on that factor's blocks --- which is check V3
restricted to that factor, conjunct for conjunct, in its general-format
reading. (C4): where every link of segment $f$ supplies its channel as data,
with $\Lambda_f\coloneqq\Lambda^{f}_{L_f}\circ\dots\circ\Lambda^{f}_1$,
\begin{equation}\label{eq:pullback-factor}
  \Lambda_f^{\dagger}\big(F^{f}_{v_f}\big)
  =\sum_{w_f\in\gamma_f^{-1}(v_f)}F^{\mathrm{ph},f}_{w_f}
  \qquad\text{for every }v_f ,
\end{equation}
an exact matrix identity on that factor's matrices, the fibre sum accumulated
in the declared order of that factor's physical alphabet. \emph{The product
composite $\bigotimes_f\Lambda_f$ and the product identity on $\{F_v\}$ are
formed and decided by no one}: the identity
$\big(\bigotimes_f\Lambda_f\big)^{\dagger}(F_v)
=\bigotimes_f\Lambda_f^{\dagger}(F^{f}_{v_f})$ is a consequence of
\eqref{eq:pullback-factor} and of the product rule
\eqref{eq:productfamily} --- it is \xref{reductions/thm:emitted-product}\,(ii),
named here and decided here by nothing --- and (C4) is the per-factor
identity and nothing wider. (C5) is per link, as above. Where some link of
segment $f$ supplies no channel as data, only a theorem that constructs one,
\eqref{eq:pullback-factor}
for \emph{that} segment is not decidable from the emitted data and is carried
as Assumption~\ref{ass:link}, exactly as above and per factor.

\emph{A non-empty source segment, and the conditions it brings.} \emph{A
presented chain may carry one.} The rule by which a non-empty source segment
composes with a receiver segment is a theorem of the reductions module
(\xref{reductions/lem:precomposition}, \xref{reductions/thm:two-sided}, with
\xref{reductions/def:chain}), in
\exttermref{../transmitter/contract.pdf}{term-trusted-source-mode}{trusted-source mode}; and the module that admits source-link kinds may
emit an artefact carrying a link of an admitted kind. What follows is what
this module decides of such a segment, on the presented data and exactly;
whether a presented chain carries one is data of the record and is asserted
here of no artefact.
\begin{itemize}[leftmargin=2.6em]
\item \emph{(S1)--(S4) as \xref{transmitter/def:source-segment}
states them}, which is what the source-segment paragraph above already says;
and within (S3), \emph{the split by kind of witness}: for a link supplying its
channel as data the per-link identity is decided here, entry by entry, with
the composite formed in exact arithmetic; for a link supplying a
\exttermref{../transmitter/contract.pdf}{term-theorem-constructor}{theorem constructor} it is \emph{not decided here at all} but is the conclusion of the
claim that link's \exttermref{../transmitter/contract.pdf}{term-obligation-record}{obligation record} names, and the composite identity
is the composition of the per-link ones. So what this module decides of (S3)
is the data links' share of it, and the constructor links' share travels as
the obligation Assumption~\ref{ass:link} names ---
Theorem~\ref{thm:chain-consumption}\,(iv), (v) --- and is discharged nowhere in this
contract.
\item \emph{(A1)--(A3) of that definition, which are decided on data and are
not part of (S1)--(S4)}: that the triple a link's obligation record carries
--- its kind, the claim id it names, and the assumption short ids it names ---
appears \emph{entry by entry} in the \exttermref{../transmitter/contract.pdf}{term-soundness-registry}{soundness registry} the artefact
carrying the segment emits, both components of that registry entry's
assumption list compared, so that a record naming a claim's closure and
omitting the \exttermref{../transmitter/contract.pdf}{term-source-model}{source model}'s attestation fails as one naming neither
does; that every assumption of that entry is in the artefact's assumption
ledger; and that the segment is presented inside an artefact of the module
that admits its links. These are comparisons of finitely many declared
strings and of a finite list against a finite list, and \emph{they are decided
here because they are decidable on the presented record}, which is why they
are named beside (S1)--(S4) rather than left to Assumption~\ref{ass:link}.
\emph{What they do not decide} is that the named claim's conclusion applies to
these descriptors --- family membership is the emitting module's assertion and
no consumer re-proves a theorem --- and a link supplying its channel as data
carries no obligation record and is decided by (S3) instead.
\item \emph{the licensed-only rule.} Beside a non-empty source segment
\emph{every} entry of the emitted $\CList$ is required to carry a real scalar
as its right-hand side --- homogeneous, party-local, or an untrusted-mode
scalar (Definition~\ref{def:homog}).
\item \emph{the declared source mode.} \emph{A declaration carrying a
non-empty source segment is in trusted-source mode}: the rule fails
where the \termref{term-marginal-slot}{marginal slot} of \eqref{eq:declaration} holds $\nomarg$
(Definition~\ref{def:declaration}), the \exttermref{../transmitter/contract.pdf}{term-untrusted-source-mode}{untrusted-source mode}, while
$L_{\mathrm s}\ge1$.
\end{itemize}
\emph{(X1) in the general format.} The order rule is unchanged and is still
the type of the record: a source-map link has no slot in any receiver segment
to be written into, and a link presented in the wrong list fails that list's
descriptor rule --- (C1) of the factor whose segment it was written into, or
(S1). In the reductions module the composite of a chain is a pair, and the
reason a source link may not follow a receiver link is a theorem of that
module (\xref{reductions/lem:precomposition},
\xref{reductions/thm:two-sided}) beside the published warrant
\cite[Rem.~9.2]{tupkary2026decoy}; this module types the record, decides the
descriptors, and states neither.
\end{definition}
\claimmeta{\metalabel{Uses}\mref{Def.}{def:declaration}, \mref{Def.}{def:exact}, \mref{Def.}{def:validator}, \extref{../transmitter/contract.pdf}{transmitter/def:source-segment}
\metasep \metalabel{Used by}\mref{Thm.}{thm:choice-encoding}, \mref{Rem.}{rem:relabel}, \mref{Rem.}{rem:alphabet}, \mref{Rem.}{rem:chaincond-reductions}, \mref{Rem.}{rem:source-segment-rules}, \mref{Rem.}{rem:chain-and-validator}, \mref{Thm.}{thm:chain-consumption}, \mref{Rem.}{rem:coarse}, \mref{Thm.}{thm:band-a}, \mref{Ass.}{ass:link}, \mref{Req.}{req:operator-form}, \mref{Req.}{req:chain-refusals}, \mref{Req.}{req:band-a-identity-link}
\metasep newly derived
\metasep implemented (effective: accepted)
\metasep \metalabel{Code}\coderef{https://github.com/fraunhofersg/certiqs-key/blob/v0.1.0/src/certiqs_key/protocol_compiler/chain.py}{certiqs_key.protocol_compiler.chain:check_chain}
\metasep \metalabel{Oracles}\coderef{https://github.com/fraunhofersg/certiqs-key/blob/v0.1.0/tests/oracles/registry_protocol_compiler.py}{registry_protocol_compiler.py}}

% claim protocol-compiler/rem:chaincond-reductions
\claimsummary{protocol-compiler/rem:chaincond-reductions}{The compiler's (C1)\textendash{}(C3) are conditions on the record's fields and differ from the reductions module's composability, tuple chaining and stage order, whose numbers they share. That module proves the chain composes; this one decides on the record what can be decided there and names the rest.}
\claimbookmark{protocol-compiler/rem:chaincond-reductions}{Remark}{rem:chaincond-reductions}{Chain conditions here and in the reductions module}
\begin{remark}[Chain conditions here and in the reductions module]\label{rem:chaincond-reductions}
The two modules that meet at the chain record describe it from two sides: the
reductions module \emph{defines} the chain and \emph{proves} that it composes
(\xref{reductions/def:chain}, \xref{reductions/def:link},
\xref{reductions/thm:chain}); this module \emph{decides} what it can decide
on the presented record and \emph{names} what it cannot. Nothing is asserted
twice: the numbered conditions of Definition~\ref{def:chaincond} are this
module's, are stated on the record's fields alone, and are not the
composability (C1), tuple chaining (C2) and stage order (C3) of that module's
definition, whose numbering they happen to share --- ``(C1) of
Definition~\ref{def:chaincond}'' and ``reductions' composability (C1)'' are
different conditions on the same object.
\end{remark}
\claimmeta{\metalabel{Uses}\mref{Def.}{def:chaincond}
\metasep newly derived
\metasep accepted}

% claim protocol-compiler/rem:source-segment-rules
\claimsummary{protocol-compiler/rem:source-segment-rules}{The licensed-only rule and the source-mode rule are the two halves of the reductions module's admission rule for a non-empty source segment. Both refuse where the chain is consumed; only the first refuses a case an artefact of this repository can present.}
\claimbookmark{protocol-compiler/rem:source-segment-rules}{Remark}{rem:source-segment-rules}{Grounds of the source-segment rules}
\claimleadin{%
Definition~\ref{def:chaincond} states two rules for a non-empty source
segment, the licensed-only rule and the rule on the declared source mode, and
Requirement~\ref{req:chain-refusals} names their refusals. This remark gives
their grounds and compares what each of them refuses.
}
\begin{remark}[Grounds of the source-segment rules]\label{rem:source-segment-rules}
\emph{Where the two rules come from.} The licensed-only rule is the consuming
side of the admission rule of \xref{reductions/def:chain}, whose reason is
\xref{reductions/lem:precomposition} and its counterexample
(\xref{reductions/rem:ordering}); this module states neither and decides the
type of the entry. The rule on the declared source mode is the other half of
that admission rule, which admits a non-empty source segment beside a
constraint list only where every entry is licensed \emph{and only in
\exttermref{../transmitter/contract.pdf}{term-trusted-source-mode}{trusted-source mode}}: the entry half is the licensed-only rule and the
mode half is this one, and both are decided in
Definition~\ref{def:chaincond} on the presented data, the entry half by
reading the declared tag of each entry and the mode half by reading one
declared slot.

\emph{The licensed-only rule refuses a case the \termref{term-basic-format}{basic format} admits.}
A non-empty \exttermref{../transmitter/contract.pdf}{term-source-segment}{source segment} is declarable there --- the source-segment
paragraph of Definition~\ref{def:chaincond} admits $L_{\mathrm s}\ge1$ ---
and, \emph{were this rule absent}, an entry declared general beside one
would, where no other condition or gate refused it, reach the gates and be
emitted by row 5 of the gate outcome table of
Definition~\ref{def:gate-outcome-table} as a
declaration that is not \termref{term-key-length-eligible}{key-length-eligible}
(Requirement~\ref{req:gate-outcomes}). Under the rule the case is refused where the
chain is consumed, before any declaration exists, so row 5 is never reached
for it. The two outcomes are not the same, and the second is strictly the more
conservative. \emph{The case is reachable}, and the rule is then simply a
refusal, which is the direction a consumer may always take.

\emph{Why the source-mode rule refuses rather than weakening the outcome}:
nothing licenses a source map in \exttermref{../transmitter/contract.pdf}{term-untrusted-source-mode}{untrusted-source mode}. The statements
that make a non-empty segment sound --- \xref{reductions/thm:two-sided} and
\xref{reductions/lem:source-transfer} --- are stated in trusted-source mode
alone, because a source map replaces \emph{Alice's preparation} and only a
trusted specification has one: an untrusted specification carries no
certificate and no marginal, its Alice-side system is an output of the attack,
and A-SRC is asserted in trusted-source mode alone
(\xref{transmitter/ass:src}). So a declaration of that shape would carry a
replaced ensemble with no statement transferring anything to it, which is the
one thing a consumer may not emit. \emph{No artefact of this repository can
present the pairing, and the rule is stated so that it is decided rather than
left to provenance}: a non-empty segment is accepted only inside an artefact
of the module that admits its links, by (A3) of
\xref{transmitter/def:source-segment} as Definition~\ref{def:chaincond}
decides it, and the one submodule that admits a kind
declares its data in trusted-source mode
(\xref{transmitter/decoy/cor:emitted-decoy}\,(i)). Unlike the licensed-only
rule, which refuses a case the basic format declares and which, were the
licensed-only rule absent and nothing else refused it, row 5 would emit as not
key-length-eligible, the source-mode rule
refuses nothing that an artefact of this repository can present; what it
removes is the possibility that a later kind emitted in another mode would
pass every condition stated in Definition~\ref{def:chaincond}.
\end{remark}
\claimmeta{\metalabel{Uses}\mref{Def.}{def:chaincond}
\metasep newly derived
\metasep accepted}

% claim protocol-compiler/rem:chain-and-validator
\claimsummary{protocol-compiler/rem:chain-and-validator}{The last clause of (C1), (C4) and the source identity are decided on the chain alone; V3 re-runs (C3) on the emitted family, weighted where a weighting follows the chain. The chain's last target descriptor, a supplied witness and the physical family stay outside V1's list, keeping V1 complete.}
\claimbookmark{protocol-compiler/rem:chain-and-validator}{Remark}{rem:chain-and-validator}{Chain conditions and validator checks}
\claimleadin{%
Some conditions on what the compiler emits are decided where the chain is
consumed and some by the validator. This remark says which, and which objects
each of them compares.
}
\begin{remark}[Chain conditions and validator checks]\label{rem:chain-and-validator}
\emph{What is decided on the chain.} The last clause of
Definition~\ref{def:chaincond}\,(C1), that the target descriptor of
$\Lambda_L$ is $(\Qsys,\Smap)$, is decided at (C1) and nowhere else: it is a
condition on the chain record and not on a slot of \eqref{eq:declaration}, so
no target condition quantifies over it and check V1 does not read it
(Definition~\ref{def:validator}). Nothing is lost by keeping it there alone:
the emitted $\Smap$ \emph{is} that descriptor by construction
(Definition~\ref{def:compile}), and what V1 decides on the emitted data is
that the map carried by this same link's target POVM $\{F_v\}$ is $\Smap$.
(C4) and (S3) are conditions on the chain and not checks of the validator
either: they are decided where the chain is consumed, and they are \emph{not}
among the decisions counted by \eqref{eq:predcount} of
Theorem~\ref{thm:validator-soundness}, a count of V2, V3 and V6 alone; V7's decisions are
counted in the prose that follows that equation and are not terms of it. What
V3 re-runs on the emitted data is (C3), or, where the Bob-side weighting of
Definition~\ref{def:weighted} follows the chain, (C3)'s predicates on the
weighted family the chain's family yields
(Theorem~\ref{thm:chain-consumption}\,(ii)); (C1)'s last clause is re-run by no
check, V1 reading the maps of the declared $\Qsys$-operator families and no
others; and nothing re-runs (C4) or (S3).

\emph{Which objects V1 does not compare.} At a non-trivial chain two families
answer to the receiver's measurement, the emitted family $\{F_v\}$ and the
physical family $\{\Fph_w\}$, and V1 compares the first only. \emph{Two
objects are outside V1's list, and each is compared where it belongs},
because carrying either in V1 would make V1 refuse a candidate that
D(i)--D(viii) admit, and so falsify Theorem~\ref{thm:validator-soundness}\,(ii). \emph{The
chain's last target descriptor} is required to be $(\Qsys,\Smap)$ by (C1),
which decides that equality by the same four-component comparison and is
where a chain is refused for it, as above; what V1 decides instead, that the
map carried by that link's target POVM $\{F_v\}$ is $\Smap$, is a condition
on a declared $\Qsys$-operator and is D(i)'s. \emph{A supplied feasibility
witness} is not a $\Qsys$-operator at all: $\omega^{\star}$ is a
\exttermref{../foundations/contract.pdf}{term-block-list}{block list} on $A\Qsys$, an input consumed and never emitted
(Theorem~\ref{thm:feasibility}\,(ii)), and branch~(c) of V7 already requires it to be
a block list against $\Smap$, which is the one place its map is compared.
\emph{The physical family $\{\Fph_w\}$ is not in the list either}: at a
non-trivial receiver segment it lives on the physical receiver space, which
is not $\Qsys$, so it is not a $\Qsys$-operator family and has no map to
compare against $\Smap$; its descriptor is matched instead by (C1), against
the source descriptor of $\Lambda_1$. The two coincide exactly when the
receiver segment does not change the space --- in \termref{term-band-a}{Band A}, where
$\Lambda_1=\idmap$ and $F_v=\Fph_v$ --- and there the comparison is the same
one either way. Including the physical family at a non-trivial chain would
not be a useful refusal: it would refuse every such chain, and V1 is claimed
complete for D(i), which asks only that every declared
\emph{$\Qsys$-operator} be a block list against $\Smap$.
\end{remark}
\claimmeta{\metalabel{Uses}\mref{Def.}{def:compile}, \mref{Thm.}{thm:feasibility}, \mref{Def.}{def:validator}, \mref{Def.}{def:chaincond}
\metasep newly derived
\metasep accepted}

% claim protocol-compiler/thm:chain-consumption
\claimsummary{protocol-compiler/thm:chain-consumption}{Chain conditions are decided by finite computation; a verified pullback identity reproduces the coarse-grained physical statistics and a verified source identity maps the virtual source onto the specified one. Link soundness and attack-set licences stay obligations this theorem does not discharge.}
\claimbookmark{protocol-compiler/thm:chain-consumption}{Theorem}{thm:chain-consumption}{Decided identities and named obligations of a consumed chain}
\claimleadin{%
Consuming a chain buys only what its record makes decidable. The theorem
states what each decided identity implies and lists what stays an obligation
of the chain's producer; where those obligations are discharged for the chains
the reductions module emits is Remark~\ref{rem:obligations-discharged}.
The attack-set licence that (v) names is Assumption~\ref{ass:atk}. It occurs
in no hypothesis of this theorem: it is spent at the gate that tests for its
record, Theorem~\ref{thm:gates}\,(iii), and is named here and not declared.

How the compiler applies the Bob-side weighting of (ii) after the chain is
Remark~\ref{rem:relabel}. Remark~\ref{rem:coarse} reads the clause of (iii)
on a coarse-graining in the opposite direction as the reason (C2) runs from
the physical alphabet onto the emitted one. That (A1)--(A3) do not decide whether the claim an obligation
record names applies to the descriptors of the link carrying that record,
as (v$'$) says, is
\xref{transmitter/rem:segment-scope}.
}
\begin{theorem}[Decided identities and named obligations of a consumed chain]
\label{thm:chain-consumption}
Let a chain be presented to the compiler, with the
\termref{term-interface-conditions-on-a-chain}{interface conditions} of
Definition~\ref{def:chaincond}. Then:
\begin{enumerate}[label=(\roman*),leftmargin=2.2em]
\item (C1)--(C5) and (S1)--(S4) are decided by finite enumeration and finitely
many arithmetic operations on the emitted data, and a chain failing any of
them is refused. (X1) is not decided,
because it is not a condition on data: it is the type of the record;
\item (C3) is exactly the part of D(iii) concerning $\{F_v\}$, in both halves
--- typing against $\Smap$, and the per-block predicates and summation
condition of Definition~\ref{def:exact} --- hence a chain passing the interface
conditions yields an emitted $\{F_v\}$ for which Theorem~\ref{thm:measurements-and-statistics}
applies. Where the Bob-side weighting of Definition~\ref{def:weighted} is
applied after the chain (Definition~\ref{def:compile}), (C3) is decided on the family the chain delivered,
the emitted family is its weighting $N_{v,f}=P(f\,|\,t(v))F_v$, check V3
decides the emitted family itself (Theorem~\ref{thm:measurements-and-statistics}\,(i)), and
Lemma~\ref{lem:weighted}\,(i), (ii) is what says the weighting cannot break
what (C3) established; the identity of conjunct lists asserted here is between
(C3) and D(iii) on the chain's family;
\item \emph{(what a verified (C4) buys)} when the composite receiver channel
$\Lambda$ is given as data, (C4) is \emph{equivalent} to the statement that
measuring the emitted family $\{F_v\}$ after $\Lambda$ reproduces the
\emph{$\gamma$-coarse-grained} physical measurement statistics:
\begin{equation}\label{eq:pullbackstats}
  \tr[F_v\,\Lambda(\rho)]=\sum_{w\in\gamma^{-1}(v)}\tr[\Fph_w\,\rho]
\end{equation}
for every $v\in\Vset$ and every operator $\rho$ on the physical receiver space.
Consistency with completeness holds \emph{under one hypothesis this contract
names and does not verify}: if the supplied channel data are trace-preserving,
so that $\Lambda^{\dagger}$ is unital, then summing \eqref{eq:pullback} over
$v$ returns completeness of the physical family and nothing more, whereas the
same identity read with a coarse-graining in the opposite direction --- a
total surjection $\gamma':\Vset\to\Vph$, with
$\Lambda^{\dagger}(F_v)=\Fph_{\gamma'(v)}$ --- does not in general, forcing
every physical element covered more than once by $\gamma'$ to vanish. That
trace preservation is
Assumption~\ref{ass:link}, not a check: nothing here evaluates it;
\item \emph{(what a verified (S3) buys)} when the composite source channel
$\Xi$ is given as data, (S3) of
\xref{transmitter/def:source-segment} --- the source identity,
that definition's equation, decided by this module on the presented data
(Definition~\ref{def:chaincond}) --- says that the virtual source, pushed
through $\Xi$, is the specified source: $\Xi(\rvirt_x)=\rho_x$ for every $x$,
hence
\begin{equation}\label{eq:srcresid}
  (\idmap_X\otimes\Xi)\Big(\sum_{x\in\Xset}p_x\dyad{x}_X\otimes\rvirt_x\Big)
  =\sum_{x\in\Xset}p_x\dyad{x}_X\otimes\rho_x ,
\end{equation}
and, averaging, $\sum_xp_x\Xi(\rvirt_x)=\sum_xp_x\rho_x$: the
prepare-and-measure states of the emitted virtual description map onto the
specified ones, index by index, with the weights unchanged by (S2). The
compiler decides the source identity (S3) and nothing beyond it; and
\item the following are \emph{named and never discharged here}: that each link
of either segment is a sound reduction --- that its channel exists as claimed
for a link supplying only a theorem constructor, that a link supplying
its channel as data supplies a completely positive and trace-preserving map,
and that every reachable state after the reduction satisfies every transported
constraint entry --- and that any restricted \exttermref{../foundations/contract.pdf}{term-attack-set}{attack set} claimed by a
link is licensed. These are obligations of the chain's producer, the
soundness of each link being Assumption~\ref{ass:link}. No claim of this contract discharges them for a
presented chain in general --- for the \termref{term-band-a}{Band A} identity link, which
Definition~\ref{def:banda}\,(b) fixes with its channel supplied as data, its
outcome correspondence and its target POVM,
Theorem~\ref{thm:band-a}\,(v) discharges every clause of
Assumption~\ref{ass:link} by exhibition --- and
no proof below invokes them.
\end{enumerate}
\emph{The general format.} For a chain presented as
Definition~\ref{def:spec}\,(c$'$), with the interface conditions read as in the
general-format paragraph of Definition~\ref{def:chaincond}, (i)--(v) hold with
the following readings, of which the five above are the case of one
\termref{term-measured-factor}{measured factor} and an empty \exttermref{../transmitter/contract.pdf}{term-source-segment}{source segment}.
\begin{enumerate}[label=(\roman*$'$),leftmargin=2.6em]
\item (C1)--(C5) are decided per measured factor and (S1)--(S4) as before, all
by finite enumeration and finitely many arithmetic operations, and a chain
failing any of them is refused.
\emph{And (A1)--(A3), the licensed-only rule and the declared-source-mode
rule are decided as well}: the first three by comparing finitely many declared
strings and one finite list against another, the fourth by reading $m$
declared tags, and the fifth by reading one declared slot. (X1) is again not
decided, being the type of the record.
\item (C3) is exactly the part of D(iii) concerning the declared measured
families, \emph{factor by factor}: for each $f$, (C3) read for $f$ and
D(iii)'s clauses on $\{F^{f}_{v_f}\}$ are the same list of conjuncts, so a
chain passing the interface conditions yields per-factor families for which
Theorem~\ref{thm:measurements-and-statistics}\,(i$'$) applies. \emph{That the product family is a POVM on
$\Qsys$ is on neither side of that pairing}, being a consequence of the
factors' conditions and no condition on data
(Definition~\ref{def:declaration}\,(DX3), Theorem~\ref{thm:validator-soundness}, general
format). The Bob-side weighting sentence applies per factor, to that factor's
delivered family.
\item \emph{(what a verified (C4) buys, per factor and jointly)} when every
link of segment $f$ supplies its channel as data,
\eqref{eq:pullback-factor} is equivalent to
$\tr[F^{f}_{v_f}\Lambda_f(\rho_f)]
=\sum_{w_f\in\gamma_f^{-1}(v_f)}\tr[F^{\mathrm{ph},f}_{w_f}\rho_f]$ for every
$v_f$ and every operator $\rho_f$ on factor $f$'s physical receiver space.
\emph{And the joint statement follows for every operator, product or not}: if
every segment supplies its channels as data then, with
$\Lambda\coloneqq\bigotimes_f\Lambda_f$, $\gamma\coloneqq\prod_f\gamma_f$ and
$F^{\mathrm{ph}}_w\coloneqq\bigotimes_fF^{\mathrm{ph},f}_{w_f}$,
\begin{equation}\label{eq:pullback-product}
  \Lambda^{\dagger}(F_v)=\sum_{w\in\gamma^{-1}(v)}F^{\mathrm{ph}}_w ,
  \qquad\text{hence}\qquad
  \tr[F_v\,\Lambda(\rho)]=\sum_{w\in\gamma^{-1}(v)}\tr[F^{\mathrm{ph}}_w\,\rho]
\end{equation}
for every $v\in\Vset$ and every operator $\rho$ on the product physical space.
\emph{Neither side of \eqref{eq:pullback-product} is formed or decided by
anything}: it is a consequence of the per-factor identities, and the objects in
it --- the product channel, the product physical family, the product target
family --- are data of no one. The consistency reading holds under the same
hypothesis read per factor: if every factor's supplied channel data are trace
preserving, so that each $\Lambda_f^{\dagger}$ is unital, then so is
$\Lambda^{\dagger}$ and summing \eqref{eq:pullback-product} over $v$ returns
completeness of the product physical family and nothing more. That trace
preservation is Assumption~\ref{ass:link}, per factor, and is evaluated
nowhere.
\item (S3) is as in (iv), with the split by kind of witness: what a
\emph{decided} (S3) buys is \eqref{eq:srcresid} and its averaged form, for the
data links' share of the composite; for a link supplying a theorem constructor
this module decides nothing and what its share buys is what the claim its
obligation record names concludes, which this contract neither states nor
asserts.
\item the list of obligations named and never discharged here gains three
items, each per the general format: that each link of \emph{each} receiver
segment is a sound reduction, factor by factor; that a non-empty source
segment composes with the receiver segments as a reduction, which is
\xref{reductions/thm:two-sided}, used by no step here; and that the claim an
obligation record names applies to that link's descriptors, which (A1)--(A3)
do \emph{not} decide. These are again the obligations of (v),
Assumption~\ref{ass:link} and the attack-set licence, and the emitting
module's own assertion; no claim of this contract discharges them, save as (v) says of
the \termref{term-band-a}{Band A} identity link, and no proof below invokes them.
\end{enumerate}
\end{theorem}
\claimmeta{\metalabel{Uses}\mref{Def.}{def:declaration}, \mref{Def.}{def:exact}, \mref{Def.}{def:target}, \mref{Def.}{def:spec}, \mref{Def.}{def:compile}, \mref{Thm.}{thm:measurements-and-statistics}, \mref{Def.}{def:validator}, \mref{Thm.}{thm:validator-soundness}, \mref{Def.}{def:chaincond}, \extref{../transmitter/contract.pdf}{transmitter/def:source-segment}
\metasep \metalabel{Assumes}\aref{A-LINK}{ass:link}
\metasep \metalabel{Used by}\mref{Rem.}{rem:coarse}, \mref{Req.}{req:chain-refusals}
\metasep newly derived
\metasep accepted
\metasep \metalabel{Code}\coderef{https://github.com/fraunhofersg/certiqs-key/blob/v0.1.0/src/certiqs_key/protocol_compiler/chain.py}{certiqs_key.protocol_compiler.chain:check_chain}
\metasep \metalabel{Oracles}\coderef{https://github.com/fraunhofersg/certiqs-key/blob/v0.1.0/tests/oracles/registry_protocol_compiler.py}{registry_protocol_compiler.py}}

\begin{proof}
\emph{Idea.} Items (i) and (i$'$) count what each condition decides. Items
(ii) and (ii$'$) compare two lists of conjuncts. Items (iii) and (iii$'$) are
the Hilbert--Schmidt duality between an adjoint identity and a statistics
identity, both exact, with trace preservation of the supplied channels as the
one hypothesis of the consistency reading. Items (iv) and (iv$'$) are the
linearity of a channel, and (v) and (v$'$) are statements about what this
contract asserts.

\emph{Step 1 ((i), decidability).} (C1) and (S1) compare finitely many
declared descriptors --- for (C1)'s last clause, the four components of a
sector map, the same combinatorial comparison V1 makes on the map of a
declared $\Qsys$-operator family and taken here on the descriptor, at (C1);
(C2) enumerates a finite table and checks totality and surjectivity by
counting preimages; (C3) is check V3 of Definition~\ref{def:validator}
restricted to $\{F_v\}$, namely $|\SecSet|(2|\Vset|+1)$ exact decisions on
blocks of size $d_s$ together with $|\SecSet|$ summation decisions, finite by
Theorem~\ref{thm:validator-soundness}\,(iv); (C4) decides $|\Vset|$ matrix identities on
matrices of the declared dimensions --- the adjoint of a channel given as data
(by Kraus operators or as a matrix acting on operators) is computed by finitely
many arithmetic operations, and each right-hand side of \eqref{eq:pullback} is
a finite sum over the tabulated fibre $\gamma^{-1}(v)$; (S2) compares two finite
index sets and $|\Xset|$ declared weights; (S3) decides $|\Xset|$ matrix
identities on matrices of the declared dimensions; (C5) and (S4) check the
presence of fields. (X1) is a statement about which of the two lists a link was
written into, which the record answers by construction.

\emph{Step 2 ((ii), the conjunct lists).} (ii) is the statement of (C3)
against D(iii), conjunct by conjunct. D(iii) asks two things of the emitted
$\{F_v\}$: that it be a family of block lists against $\Smap$, and that each
block family $\{F_v^{(s)}\}_v$ satisfy the predicates and the summation
condition of Definition~\ref{def:exact}. (C3) states the first as the type of
the object the last link delivers and the second as those same conditions, on
the emitted blocks. The two lists of conjuncts are identical, so (C3) is
D(iii) restricted to $\{F_v\}$ and Theorem~\ref{thm:measurements-and-statistics} applies to what the
chain delivered. Nothing in (ii) is inferred from the physical family, so a
chain that merely \emph{claims} a POVM property cannot pass one off as
checked: the predicates are decided on the matrices the compiler emits. The
block-diagonality half, by contrast, is a property the type carries and no
check confirms. An entry that the chain's producer has \emph{mis-assigned} to
a block is a well-typed object no predicate can catch. V1 refuses a
mis-assignment that also disagrees about the sectors; one that does not is a
defect of the chain's producer under Assumption~\ref{ass:link}, named there
and not discharged here.

\emph{Step 3 ((iii), what a verified (C4) buys).} The \exttermref{../reductions/contract.pdf}{term-measurement-channel}{measurement channel} of
$\{F_v\}$ sends $\sigma$ to $\sum_v\tr[F_v\sigma]\dyad{v}$; the
$\gamma$-coarse-graining of the physical measurement channel sends $\rho$ to
$\sum_v\big(\sum_{w\in\gamma^{-1}(v)}\tr[\Fph_w\rho]\big)\dyad{v}$, the fibres
$\gamma^{-1}(v)$ partitioning $\Vph$ because $\gamma$ is a total surjection
(C2). Two maps into the same classical register agree on all inputs iff their
diagonal entries agree, i.e.\ iff \eqref{eq:pullbackstats} holds; and by the
definition of the adjoint with respect to the Hilbert--Schmidt inner product,
$\tr[F_v\Lambda(\rho)]=\tr[\Lambda^{\dagger}(F_v)\rho]$, so
\eqref{eq:pullbackstats} holds for all $\rho$ iff \eqref{eq:pullback} holds for
all $v$ --- the second direction because an operator whose trace against every
$\rho$ vanishes is the zero operator, the trace pairing being non-degenerate.
Both directions are available here because both sides are exact identities.
For the consistency claim, summing \eqref{eq:pullback} over $v$ gives
$\Lambda^{\dagger}\big(\sum_vF_v\big)=\sum_{w\in\Vph}\Fph_w$; with
$\sum_vF_v=\idon{\Qsys}$ from (C3) and $\Lambda^{\dagger}$ unital --- which is
where trace preservation of the supplied data is used, and it is
Assumption~\ref{ass:link} and not a check --- the left-hand side is
$\idon{}$, so what (C4) returns is $\sum_w\Fph_w=\idon{}$: completeness of the
physical family, a condition that family already satisfies, so (C4) adds no
restriction beyond the two completeness conditions.
The opposite direction, under the same hypothesis, is the identity
$\Lambda^{\dagger}(F_v)=\Fph_{\gamma'(v)}$ for every $v\in\Vset$, with
$\gamma':\Vset\to\Vph$ a total surjection. Summing it over $v$, the left-hand
side is again $\Lambda^{\dagger}\big(\sum_vF_v\big)=\idon{}$, by linearity of
$\Lambda^{\dagger}$, (C3) and unitality. On the right, the fibres
$\gamma'^{-1}(w)$ partition $\Vset$ because $\gamma'$ is total, so grouping
the terms by the value of $\gamma'$ gives
$\sum_{w\in\Vph}|\gamma'^{-1}(w)|\,\Fph_w=\idon{}$. The physical family is a
POVM (Definition~\ref{def:spec}\,(b)): $\sum_w\Fph_w=\idon{}$ and every
$\Fph_w\succeq0$. Subtracting,
$\sum_{w\in\Vph}\big(|\gamma'^{-1}(w)|-1\big)\Fph_w=0$, and every coefficient
is a non-negative integer because $\gamma'$ is surjective, so every term is
positive semidefinite. A vanishing sum of positive semidefinite operators has
every term zero: its trace is a sum of non-negative traces, so each term has
trace zero, and a positive semidefinite operator of trace zero is zero. Hence
$\Fph_w=0$ for every $w$ with $|\gamma'^{-1}(w)|\ge2$. That is a restriction
the physical family need not satisfy, so this reading does not in general
return completeness and nothing more. It does exactly when every physical
element covered more than once already vanishes, in particular when $\gamma'$
is a bijection: then every non-zero term of
$\sum_w|\gamma'^{-1}(w)|\,\Fph_w$ has coefficient one, and the summed
identity is $\sum_w\Fph_w=\idon{}$.
In the bijective case $\gamma$ is a bijection,
\eqref{eq:pullback} reads $\Lambda^{\dagger}(F_v)=\Fph_{\gamma^{-1}(v)}$, and
that is the published form of the condition,
\cite[Lemma~9.3, Eq.~(146)]{tupkary2026decoy}, read at the level of operators
--- there the two families are indexed by one alphabet.

\emph{Step 4 ((iv), what a verified (S3) buys).} By linearity,
$\idmap_X\otimes\Xi$ sends $\sum_xp_x\dyad{x}_X\otimes\rvirt_x$ to
$\sum_xp_x\dyad{x}_X\otimes\Xi(\rvirt_x)$; substituting
$\Xi(\rvirt_x)=\rho_x$ from the source identity (S3) gives
\eqref{eq:srcresid}, and tracing out $X$ gives its averaged form. The step is
linear in each $\rvirt_x$ and uses no property of the declared weights. That
the converse implication --- security of the virtual protocol transfers to the
real one --- requires more than the source identity is exactly the content of
Assumption~\ref{ass:link}: the published template is
\cite[Lemma~9.1, Eqs.~(138)--(139)]{tupkary2026decoy}, whose hypothesis is the
existence of such a $\Xi$ and whose conclusion this module does not assert.

\emph{Step 5 ((v), what is asserted).} (v) is a statement about what this
contract asserts. Its content is that Assumption~\ref{ass:link} and the
attack-set licence are named as obligations of the chain's producer, and
that no proof above invokes them: no step invokes Assumption~\ref{ass:link}
as an established obligation, its trace-preserving clause entering (iii) only
as the stated hypothesis of the consistency reading; no step invokes the
attack-set licence; and no step uses a theorem of the reductions
module. The discharge (v) names is Theorem~\ref{thm:band-a}'s, for the
Band-A identity link alone, and no step above rests on it. The
trace-preserving clause of Assumption~\ref{ass:link} is the hypothesis under
which (iii)'s consistency reading holds, so this theorem's statement is
conditional on it and declares it. The order rule (X1) is not an exception --- it is a
\emph{typing} property of the emitted record, and its warrant is the published
one: \cite[Rem.~9.2]{tupkary2026decoy} resolves the tension between source
maps, which are formulated for unrestricted attack sets, and squashing maps,
which restrict them, by applying the source-map lemma first, so that
restrictions arise only afterwards. The compiler enforces the order; it does
not claim that either lemma applies to a given device.

\emph{The general format.} The steps are those above, read per factor.

\emph{Step 1 in the general format ((i$'$)).} Each clause of (i) is the
clause of (i) applied to one segment, the data being per factor and finite;
the refusal naming the factor is the record's own field. For (A1): the
registry an artefact emits is a finite list of triples of declared strings,
the obligation record a triple of declared strings, and the comparison is
entry by entry, including the two components of the assumption list; the
ledger membership test is a finite list against a finite list. For (A3): the
provenance test reads which module emitted the artefact the segment is
presented in, a declared field. The licensed-only rule reads $m$ tags,
terminating and exact. The declared-source-mode rule reads one declared slot,
whether the \termref{term-marginal-slot}{marginal slot} holds $\nomarg$, beside whether the source
segment is empty, which the record answers by its shape. None of the five is
an operator predicate and none is decided at a window.

\emph{Step 2 in the general format ((ii$'$)).} For each $f$, D(iii) asks of
$\{F^{f}_{v_f}\}$ that it be a family of block lists against factor $f$'s map
and that each block family satisfy the predicates and the summation condition
of Definition~\ref{def:exact}, and (C3) read for $f$ states the type and those
same conditions on that factor's emitted blocks; the two lists of conjuncts
are identical, factor by factor. What Step 2 says of the physical family and
of the block-diagonality half applies per factor unchanged, the
mis-assignment exposure now being one per factor and V1's per-factor map
comparison what catches an assignment that also disagrees about that factor's
sectors.

\emph{Step 3 in the general format ((iii$'$)).} The per-factor equivalence is
(iii)'s argument with $\Qsys_f$, $\Lambda_f$, $\gamma_f$ and factor $f$'s two
families in place of $\Qsys$, $\Lambda$, $\gamma$ and the two families: the
fibres $\gamma_f^{-1}(v_f)$ partition factor $f$'s physical alphabet because
$\gamma_f$ is a total surjection (C2), the adjoint identity is the
Hilbert--Schmidt pairing, and both directions are available because both
sides are exact identities. For \eqref{eq:pullback-product}: the adjoint of a
tensor product of maps is the tensor product of their adjoints, so
$\Lambda^{\dagger}(F_v)=\bigotimes_f\Lambda_f^{\dagger}(F^{f}_{v_f})$ by
\eqref{eq:productfamily}, this is
$\bigotimes_f\sum_{w_f\in\gamma_f^{-1}(v_f)}F^{\mathrm{ph},f}_{w_f}$ by
\eqref{eq:pullback-factor}, and expanding the product of sums by
multilinearity gives
$\sum_{w\in\gamma^{-1}(v)}\bigotimes_fF^{\mathrm{ph},f}_{w_f}$, the index set
being exactly $\gamma^{-1}(v)=\prod_f\gamma_f^{-1}(v_f)$. The statistics form
follows for \emph{every} operator $\rho$ on the product physical space, product
or not, because $\tr[F_v\Lambda(\rho)]=\tr[\Lambda^{\dagger}(F_v)\rho]$ is
linear in $\rho$ and \eqref{eq:pullback-product} is an identity of operators.
For the consistency reading, a tensor product of unital maps is unital, so
summing over $v$ and using $\sum_{v_f}F^{f}_{v_f}=\idon{\Qsys_f}$ per factor
(C3) gives $\sum_{w}F^{\mathrm{ph}}_w=\idon{}$, which the product physical
family already satisfies; nothing is added and nothing is evaluated.

\emph{Steps 4 and 5 in the general format ((iv$'$) and (v$'$)).} The first is
(iv)'s computation, which is linear in each $\rvirt_x$ and reads only the
identity, together with the split by kind of witness made in (S3) of
\xref{transmitter/def:source-segment}, restated here and not re-derived. The second is
again a statement about what this contract asserts: no step above invokes
Assumption~\ref{ass:link} as an established obligation, per factor or
otherwise, or the attack-set licence, and none uses a theorem of the
reductions module, and the three items named are named and not used --- in particular no
step above composes a source segment with a receiver segment, which is why the
composition rule appears in the list rather than in a hypothesis. What is
declared here is unchanged: the trace-preserving clause of
Assumption~\ref{ass:link} occurs in (iii$'$)'s consistency reading as in
(iii)'s, and the attack-set licence is named here and spent at the gate
(Theorem~\ref{thm:gates}\,(iii)).
\end{proof}

% claim protocol-compiler/rem:coarse
\claimsummary{protocol-compiler/rem:coarse}{With a trace-preserving chain, a correspondence from the emitted alphabet onto the physical one forces every physical element covered more than once to vanish, so it admits no coarse-graining. The correspondence therefore runs from the physical alphabet onto the emitted one.}
\claimbookmark{protocol-compiler/rem:coarse}{Remark}{rem:coarse}{Direction of the outcome correspondence}
\begin{remark}[Direction of the outcome correspondence]
\label{rem:coarse}
The direction in (C2) is forced, and the argument that forces it is a
\emph{typing} argument: (C2) fixes the direction of a tabulated map, it is
combinatorial, and it is read exactly.

Suppose instead a total surjection
$\gamma':\Vset\to\Vph$ from the emitted alphabet onto the physical one, with the
pullback read as $\Lambda^{\dagger}(F_v)=\Fph_{\gamma'(v)}$. With
trace-preserving channel data, Theorem~\ref{thm:chain-consumption}\,(iii)
says that this reading forces $\Fph_w=0$ for every outcome $w$ covered more
than once; the computation is Step~3 of its proof. The only non-injective $\gamma'$ that
survives is one whose extra outcomes are null, i.e.\ no coarse-graining at all.
Coarse-graining is a real operation --- a squashed family is typically
\emph{coarser} than the physical one --- so the correspondence must run from
the finer alphabet to the coarser, which is (C2), and the identity it feeds
must sum the fibre, which is \eqref{eq:pullback}. Emitting the pair
$(\gamma,\eqref{eq:pullback})$ rather than a bare correspondence of indices is
what keeps the statistics identity \eqref{eq:pullbackstats} true outcome by
outcome.
\end{remark}
\claimmeta{\metalabel{Uses}\mref{Def.}{def:chaincond}, \mref{Thm.}{thm:chain-consumption}
\metasep newly derived
\metasep accepted}

% claim protocol-compiler/rem:obligations-discharged
\claimsummary{protocol-compiler/rem:obligations-discharged}{For a chain the reductions module emits, its licence theorem discharges \hyperlink{claim-protocol-compiler/ass:atk}{A-ATK} for the entries it certified and, at status exact, its chain, transfer and emission theorems discharge \hyperlink{claim-protocol-compiler/ass:link}{A-LINK}, also on a physical side presented by restriction or weighted after the link. Only \hyperlink{claim-protocol-compiler/req:link-row-provenance}{Requirement~\ref*{req:link-row-provenance}} uses them.}
\claimbookmark{protocol-compiler/rem:obligations-discharged}{Remark}{rem:obligations-discharged}{Discharge of the chain obligations by the reductions module}
\claimleadin{%
Theorem~\ref{thm:chain-consumption}\,(v) names the soundness of each link and the
attack-set licence as obligations of the chain's producer and discharges
neither. This remark records where they are discharged for the chains the
reductions module emits, also where a factor's family is weighted after its
link, and for the Band-A identity link.
}
\begin{remark}[Discharge of the chain obligations by the reductions module]\label{rem:obligations-discharged}
\emph{Where the two obligations are discharged.} For a chain emitted by the
reductions module with status \emph{exact}, Assumption~\ref{ass:link} is
discharged by \xref{reductions/thm:chain}\,(i)--(iii),
\xref{reductions/thm:transfer} and \xref{reductions/thm:emit}\,(iii), its
link-soundness clause being \xref{reductions/thm:chain}\,(i)--(iii) and
\xref{reductions/thm:emit}\,(iii). For a chain emitted by the reductions
module, the attack-set licence of Assumption~\ref{ass:atk} is discharged by
\xref{reductions/thm:licence}, the restricted \exttermref{../foundations/contract.pdf}{term-attack-set}{attack set} being that of
\xref{reductions/def:attset}, for the entries the chain transports, each
carrying the spectral pull-back certificate that module discharged for it
(\xref{reductions/cor:entries}); an entry of $\Spec_{\mathrm{con}}$ carrying
no such certificate is covered by no theorem of that module, and for a list
containing one the licence stays the obligation Assumption~\ref{ass:atk}
names. For the identity link of \termref{term-band-a}{Band A}, supplied as data, with the
identity outcome correspondence and the declared family as its target POVM
(Definition~\ref{def:banda}\,(b)), Assumption~\ref{ass:link} is discharged
by exhibition (Theorem~\ref{thm:band-a}\,(v)).

\emph{Several measured factors.} Where each \termref{term-measured-factor}{measured factor}'s receiver
segment is the chain of that factor's artefact, emitted by the reductions
module with status \emph{exact}, and each such segment is that factor's
instance link alone, Assumption~\ref{ass:link} is discharged factor by factor
as above, and jointly: the product of the factors' channels squashes the
product of their \exttermref{../reductions/contract.pdf}{term-declared-model}{declared models} to the product target exactly
(\xref{reductions/thm:emitted-product}\,(ii)), and the physical protocol
measures that product by A-MEAS (\xref{reductions/ass:meas})\,(iv), so
\xref{reductions/thm:transfer} applies to the joint protocol, and every entry
whose scope is a factor is sound for the joint reachable states
(\xref{reductions/thm:emitted-product}\,(iii), with
\xref{reductions/thm:licence}\,(v)). Where each factor's physical side is
presented by restriction, as in the next paragraph, the product of the
presented links is the presented link of the product, its inclusion being
$\bigotimes_f\iota_f$ and its channel $\Lambda\circ\bigotimes_f\mathcal{I}_f$,
since $\bigotimes_f(\Lambda_f\circ\mathcal{I}_f)
=\big(\bigotimes_f\Lambda_f\big)\circ\bigotimes_f\mathcal{I}_f$; the next
paragraph applies to it with $\iota=\bigotimes_f\iota_f$, $\iota_f$ being the
identity on a factor that presents its declared model itself. Where a further
link follows a factor's instance link, that theorem states nothing about the
product, and the joint transfer stays the obligation
Assumption~\ref{ass:link} names. The engine omits the row at several
measured factors only where every factor's segment is discharged in this way
(Requirement~\ref{req:link-row-provenance}).

\emph{A physical side presented by restriction.} Where a receiver segment is
its instance link alone, the link these theorems discharge has as its
\exttermref{../reductions/contract.pdf}{term-source-family}{source family} the declared model $\Mdec$ on the physical system
$\Bsys$, with channel $\Lambda$ and $\Lambda^{\dagger}(F_v)=\Mdec_v$
(\xref{reductions/thm:emit}\,(iii)). The conclusion about the device is that
link's: security transfers to the protocol on $\Bsys$ for every input there.
Let a specification present in its place a subspace $V$ of $\Hs{\Bsys}$, with
isometric inclusion $\iota$, as the physical receiver space, the compressions
$\Fph_v=\iota^{\dagger}\Mdec_v\iota$ as the physical receiver family, and the
identity as the outcome correspondence. The discharge carries over to the link
so presented. Its channel is $\Lambda\circ\mathcal{I}$ with
$\mathcal{I}(X)=\iota X\iota^{\dagger}$, a composition of channels since
$\iota^{\dagger}\iota=\idon{V}$; its adjoint is
$Y\mapsto\iota^{\dagger}\Lambda^{\dagger}(Y)\iota$, so
$(\Lambda\circ\mathcal{I})^{\dagger}(F_v)=\Fph_v$, which is
\eqref{eq:pullback} for it, and every carried operator's pull-back is the
compression of its pull-back under $\Lambda$. The presented link is therefore
again a \exttermref{../reductions/contract.pdf}{term-squashing-relation}{squashing relation} (\xref{reductions/def:squashing}), to which
\xref{reductions/thm:transfer} and \xref{reductions/thm:chain}\,(iii) apply as
to any; and an entry certificate $\Lambda^{\dagger}(\hat O_k)\bowtie_k
c_k\,\idon{\Bsys}$ compresses to $(\Lambda\circ\mathcal{I})^{\dagger}(\hat
O_k)\bowtie_kc_k\,\idon{V}$, compression by $\iota$ being positive and
unital. The presented link inherits its conclusions from the link on
$\Bsys$ and adds none: read alone it would restrict the inputs to $V$, which
the link on $\Bsys$ does not. That the presented family is this compression
is a fact of the instance and not of this paragraph. For the flat instance
with $V=\Hs{\Chan_{\leqN}}$ it holds because its declared model is block
diagonal at the cutoff with the \exttermref{../reductions/contract.pdf}{term-declared-datum}{declared datum} below it, and the target's
\exttermref{../reductions/contract.pdf}{term-data-sector}{data sectors} carry that datum
(\xref{reductions/wpfss/prop:flat-instance}, (I1)--(I2) and its declared
model), so the block of $F_v$ on the data sectors is
$\iota^{\dagger}\Mdec_v\iota$. No other instance has a claim stating that its
docked block is such a compression, so on the routes of every other instance
the row stays (Requirement~\ref{req:link-row-provenance}).

\emph{A factor weighted after its link.} Where a setting-dependent weighting
$P(f\,|\,t)$ is applied after a factor's segment
(Definition~\ref{def:weighted}), the chain record is the unweighted one
(Remark~\ref{rem:relabel}), and the discharge above carries over to the
weighted protocol. With $\Pi\big((v',f)\,|\,v\big)\coloneqq
\delta_{vv'}P(f\,|\,t(v'))$, a stochastic matrix from $\Vset$ to
$\Vset\times\mathcal{F}$, \xref{reductions/lem:linear}\,(i) turns the
squashing relation of the source family and $\{F_v\}$ into the squashing
relation, under the same channel, of $\{P(f\,|\,t(v))\,\Mdec_v\}$,
respectively of the identically weighted presented family, and the emitted
$\{P(f\,|\,t(v))\,F_v\}$, so
\xref{reductions/thm:transfer} applies to the weighted protocol as to the
unweighted one, the transported entries being untouched by a weighting. On
the identity link of \termref{term-band-a}{Band A} the same holds at the identity channel:
the delivered family is the declared one entry for entry
(Definition~\ref{def:banda}\,(b)), so the emitted weighted family is the
identically weighted declared family, the same matrices, and
Theorem~\ref{thm:band-a}\,(v) covers the weighted protocol.

\emph{The producer's statement of the same model.} The reductions module
states the same \exttermref{../foundations/contract.pdf}{term-round-model}{round model}, typed by source mode, and the same
restricted set as the round model of \xref{meat-core/def:roundmodel}, in
\xref{reductions/def:roundmodel} and \xref{reductions/def:attset}: the latter
reproduces the attack sets of \xref{meat-core/def:roundmodel} in both source
modes, and the former types the physical \exttermref{../foundations/contract.pdf}{term-attack-sequence}{attack sequence}, its \exttermref{../reductions/contract.pdf}{term-squashed-sequence}{squashed
sequence} being the one on $\Qsys$.

\emph{Used by a requirement only.} These theorems are used by no theorem,
definition or remark of this contract, so they are in the closure of none of
its results; Requirement~\ref{req:link-row-provenance} uses the ones on which
the omission of the row of Assumption~\ref{ass:link} rests, so that what the
omission rests on is computed. What a run inherits from the reductions module
reaches its assumption manifest as the ledger rows of each reductions artefact
it consumed, and not through this module's graph; the manifest's claim
closure is that of the root claims the run's key-rate result names, read from
the generated claim-closure table.
\end{remark}
\claimmeta{\metalabel{Used by}\mref{Req.}{req:link-row-provenance}
\metasep newly derived
\metasep accepted}

\section{The Band-A route}\label{group-protocol-compiler/def:banda}
\sectionintro{On the direct route a finite measurement is declared on a single sector and the chain is one identity link, given as data, that transports no constraint and delivers the declared measurement unchanged, so every chain condition but the exact admissibility of the declared receiver family holds by inspection and the link is sound by exhibition (\hyperlink{claim-protocol-compiler/thm:band-a}{Theorem~\ref*{thm:band-a}}). That admissibility and the other measurement predicates are still decided exactly, and the residual physical content of the receiver is the ideal-receiver assumption. A specification whose declared data already contradict that assumption is refused (\hyperlink{claim-protocol-compiler/req:band-a-ideal-receiver}{Requirement~\ref*{req:band-a-ideal-receiver}}), and so is one whose link is not decided on the record to be that identity link (\hyperlink{claim-protocol-compiler/req:band-a-identity-link}{Requirement~\ref*{req:band-a-identity-link}}).}
\sectionimports{\extref{../transmitter/contract.pdf}{transmitter/def:source-segment}}

% claim protocol-compiler/def:banda
\claimsummary{protocol-compiler/def:banda}{A Band-A specification declares a finite measurement directly on one sector, with an empty source segment and one identity receiver link, given as data, that carries no constraints and delivers the declared measurement. It is the direct route, with no optical characterisation or photon-number step.}
\claimbookmark{protocol-compiler/def:banda}{Definition}{def:banda}{Band-A specification}
\claimleadin{%
The one receiver link of a Band-A chain is the identity link of
\xref{reductions/ex:identity}; the definition fixes that link in its own
words. A link record carries its channel, its outcome correspondence and its
\exttermref{../reductions/contract.pdf}{term-target-povm}{target POVM} as three separate data, so all three are
fixed here: that the channel is the identity says nothing of the measurement
the link delivers.
}
\begin{definition}[Band-A specification]\label{def:banda}
A specification is
\termdef{term-band-a}{Band A}
when (a) its receiver part declares a POVM $\{\Fph_w\}_{w\in\Vph}$ directly on a
finite-dimensional system, with no characterization step and no photon-number
structure, and the emitted \exttermref{../reductions/contract.pdf}{term-sector-set}{sector set} is the singleton $\SecSet=\{s_0\}$ with
$\Pi_{s_0}=\idon{\Qsys}$; and (b) its \exttermref{../reductions/contract.pdf}{term-reduction-chain}{reduction chain} has an \emph{empty
\exttermref{../transmitter/contract.pdf}{term-source-segment}{source segment}}, $L_{\mathrm s}=0$ --- the source segment of
\xref{transmitter/def:source-segment} at length zero, required here of a
Band-A chain and of no other --- and a \exttermref{../reductions/contract.pdf}{term-receiver-segment}{receiver segment} of one link, $L=1$,
$\Lambda_1=\idmap_{\Qsys}$ supplied as data, with status \emph{exact},
empty constraint slot and empty licence slot, whose outcome correspondence is
the identity map $\gamma=\idmap$ on $\Vset=\Vph$ and whose
target POVM is the declared family itself:
$F_v=\Fph_v$, entry for entry, for every $v\in\Vset$.
\end{definition}
\claimmeta{\metalabel{Uses}\extref{../transmitter/contract.pdf}{transmitter/def:source-segment}
\metasep \metalabel{Used by}\mref{Thm.}{thm:choice-encoding}, \mref{Thm.}{thm:band-a}, \mref{Req.}{req:band-a-ideal-receiver}, \mref{Req.}{req:band-a-identity-link}
\metasep newly derived
\metasep implemented
\metasep \metalabel{Code}\coderef{https://github.com/fraunhofersg/certiqs-key/blob/v0.1.0/src/certiqs_key/protocol_compiler/declared_premises.py}{certiqs_key.protocol_compiler.declared_premises:is_band_a}
\metasep \metalabel{Oracles}\coderef{https://github.com/fraunhofersg/certiqs-key/blob/v0.1.0/tests/oracles/registry_instances_bb84.py}{registry_instances_bb84.py}, \coderef{https://github.com/fraunhofersg/certiqs-key/blob/v0.1.0/tests/oracles/registry_protocol_compiler.py}{registry_protocol_compiler.py}}

% claim protocol-compiler/thm:band-a
\claimsummary{protocol-compiler/thm:band-a}{On the direct one-sector route every chain condition but the receiver family's exact admissibility holds by inspection, and the identity link, which delivers the declared measurement, is sound by exhibition. The receiver's residual physical content is the ideal-receiver assumption.}
\claimbookmark{protocol-compiler/thm:band-a}{Theorem}{thm:band-a}{The Band-A route}
\claimleadin{%
On the direct route no reduction runs: the chain is one identity link, given
as data, which transports no constraint entry and delivers the declared
measurement unchanged. The theorem records what that makes trivial and what
it does not. The one link of the receiver segment is the identity link of
\xref{reductions/ex:identity}. That the link delivers the declared
measurement on the identity outcome correspondence is part of the hypothesis,
Definition~\ref{def:banda}\,(b), and not a consequence of its channel being
the identity: a link record carries its target POVM and its outcome
correspondence as data of their own, and the engine decides both on the
record (Requirement~\ref{req:band-a-identity-link}). The attack-set licence
(Assumption~\ref{ass:atk}) is asserted only for a non-empty constraint list;
on this route that list is the specification's own (iii), so where the
specification declares no entry the licence is not invoked and nothing proved
here rests on it.
}
\begin{theorem}[The Band-A route]\label{thm:band-a}
For a Band-A specification:
\begin{enumerate}[label=(\roman*),leftmargin=2.2em]
\item the block structure of a Band-A declaration is the one-block structure:
$\Smap$ is the singleton map $\big((s_0,d_{s_0})\big)$, $d_{\Qsys}=d_{s_0}$,
and every $\Qsys$-operator is a one-element \exttermref{../foundations/contract.pdf}{term-block-list}{block list}, so that its
single block \emph{is} the whole operator. Two consequences follow, and they
differ. The typing half of D(iii) is not vacuous but trivial to satisfy: an
operator is a one-element list, and V1's sector-map comparison is one
identifier and one dimension. The five predicates of D(iii) are \emph{not}
vacuous at all --- with $|\SecSet|=1$ the per-block predicates are the
full-space ones, decided once each on $d_{\Qsys}\times d_{\Qsys}$ matrices, and
a Band-A declaration is decided exactly like any other, condition~(S)
included;
\item the \termref{term-interface-conditions-on-a-chain}{interface conditions} of
Definition~\ref{def:chaincond} hold by inspection, save the predicates and
summation condition of (C3), which are decided: descriptors match because
source and target of $\idmap$ coincide; the outcome correspondence is the
identity bijection $\gamma=\idmap$ on $\Vset=\Vph$, as
Definition~\ref{def:banda}\,(b) requires; (C3)'s typing half holds
because the declared receiver family is already a one-element block list
against the singleton map, and its predicates and summation condition are the
declared receiver POVM's own, so (C3) holds exactly when that family is
\exttermref{../foundations/contract.pdf}{term-exactly-admissible}{exactly admissible}, which is decided on it and is not given by
inspection; the pullback identity (C4) reads
$\idmap^{\dagger}(F_v)=F_v=\Fph_v$ and holds \emph{identically} --- entry for
entry, since $\idmap^{\dagger}=\idmap$, every fibre $\gamma^{-1}(v)$ is a
singleton, and the delivered $F_v$ \emph{is} the declared $\Fph_v$, the same
matrix, by Definition~\ref{def:banda}\,(b); the source conditions (S1)--(S4) of
\xref{transmitter/def:source-segment} are vacuous, quantifying over an empty
list of links --- the empty \exttermref{../transmitter/contract.pdf}{term-source-segment}{source segment} being that definition's at
$L_{\mathrm s}=0$, which Definition~\ref{def:banda}\,(b) requires of a Band-A
chain and which is therefore a hypothesis of this theorem --- and the
\exttermref{../transmitter/contract.pdf}{term-virtual-ensemble}{virtual ensemble} of (S3) is the specified one, so the source identity
reads $\rho_x=\rho_x$; the ordering (X1) holds because the source segment is
empty; and the slots of (C5) are present and empty. The one link of the
receiver segment is the identity link.
\emph{So both conditions that a general chain has to decide on computed data
are met here by exhibition}, and what a verified pullback identity and a
verified source identity buy on a general chain holds on this path without
appeal to any obligation of the chain's \exttermref{../foundations/contract.pdf}{term-producer}{producer}: measuring the
family the link delivers after $\idmap$ reproduces the physical statistics,
$\tr[F_v\,\idmap(\rho)]=\tr[\Fph_v\,\rho]$ for every $v$ and every operator
$\rho$, and $\idmap$ is exactly trace-preserving \emph{by exhibition}, the
supplied channel data being exactly that map rather than binary64 data
trusted to be trace-preserving, so its adjoint is unital and summing the
pullback identity over $v$ returns $\sum_v\Fph_v=\idon{}$, which on this path
is the same statement as (C3)'s own completeness, $\{\Fph_v\}$ being
$\{F_v\}$: the identity holds as a triviality and carries no content this
path did not already have;
\item the link transports no constraint entry, so $\CList$ is the
specification's own list $\Spec_{\mathrm{con}}$ and the constraint-soundness
obligation of the link (the transported-constraint clause of
Assumption~\ref{ass:link}) is vacuous; where $\Spec_{\mathrm{con}}$ is empty,
$\CList=\varnothing$ and the two-part gate of Theorem~\ref{thm:gates} is not
engaged;
\item consequently no emitted slot, and no claim of this contract used in the
Band-A path, depends on an artefact of the receiver module or of the
reductions module. The residual physical content of the receiver is
Assumption~\ref{ass:ideal}, and it is an idealization of the receiver, not a
reduction of it; an entry the specification declares itself engages the
gate, and with it the attack-set licence of Assumption~\ref{ass:atk}, as on
every route; and
\item \emph{(the identity link is sound)} the Band-A chain satisfies every
condition of Assumption~\ref{ass:link}: (i)(a) the channel of its one link,
$\idmap_{\Qsys}$, exists, the link supplying it as data; (i)(b) that link
satisfies the pullback identity \eqref{eq:pullback}; (i)(c) the protocol
using the measurement the link delivers is the protocol using the declared
receiver family $\{\Fph_v\}$, the delivered $\{F_v\}$ being that family, so
security of the one is security of the other, and where a setting-dependent
weighting $P(f\,|\,t)$ is applied after the link
(Definition~\ref{def:weighted}), the emitted family
$\{P(f\,|\,t(v))\,F_v\}$ is the identically weighted declared family
$\{P(f\,|\,t(v))\,\Fph_v\}$, so the same holds of the weighted protocol; (ii) its source segment is empty; (iii) its link transports no
constraint entry; and (iv) $\idmap_{\Qsys}$ is completely positive and
trace-preserving. No clause of Assumption~\ref{ass:link} is therefore a
premise on this path; that the declared family is the measurement the
receiver performs is Assumption~\ref{ass:ideal}, and no condition on the link.
\end{enumerate}
\end{theorem}
\claimmeta{\metalabel{Uses}\mref{Def.}{def:exact}, \mref{Def.}{def:compile}, \mref{Def.}{def:weighted}, \mref{Def.}{def:chaincond}, \mref{Def.}{def:banda}, \extref{../transmitter/contract.pdf}{transmitter/def:source-segment}
\metasep \metalabel{Assumes}\aref{A-IDEAL}{ass:ideal}
\metasep \metalabel{Used by}\mref{Req.}{req:band-a-ideal-receiver}, \mref{Req.}{req:band-a-identity-link}
\metasep newly derived
\metasep accepted
\metasep \metalabel{Oracles}\coderef{https://github.com/fraunhofersg/certiqs-key/blob/v0.1.0/tests/oracles/registry_instances_bb84.py}{registry_instances_bb84.py}, \coderef{https://github.com/fraunhofersg/certiqs-key/blob/v0.1.0/tests/oracles/registry_protocol_compiler.py}{registry_protocol_compiler.py}}

\begin{proof}
(i) $\sum_sd_s=d_{s_0}$ is the one-term sum, so $d_{\Qsys}=d_{s_0}$. The
typing half: a block list against a singleton map is a one-element list, and
$\bigoplus_{s\in\{s_0\}}F^{(s_0)}=F^{(s_0)}$, so the operator and its block
coincide; V1's comparison is then a comparison of two one-element lists. The
predicates: at $|\SecSet|=1$ the family $\{F_v^{(s_0)}\}_v$ is $\{F_v\}_v$ and
$\Qsys_{s_0}=\Qsys$, so each per-block condition of
Definition~\ref{def:exact} --- Hermiticity, positivity, completeness, and the
summation condition (S) on the completeness sum --- is literally the full-space
one, and none of them is empty.

(ii) Each condition is evaluated at $L_{\mathrm s}=0$ and $L=1$,
$\Lambda_1=\idmap$: (C1) both descriptors of $\idmap_{\Qsys}$, at its source
and at its target, are $(\Qsys,\{s_0\},(d_{s_0}))$ and the physical receiver
space is $\Qsys$ by Definition~\ref{def:banda}\,(a), and the two descriptors
carry the same singleton \exttermref{../foundations/contract.pdf}{term-sector-map}{sector map}, so (C1) holds --- and the emitted
$\{F_v\}$ carries that same singleton map, so V1's comparison passes as well;
(C2) the outcome correspondence is the identity map on $\Vset=\Vph$
(Definition~\ref{def:banda}\,(b)), a tabulated bijection, in particular a
total surjection with singleton fibres; (C3): its typing half is (i), the
family being a one-element list against the singleton map, and by (i) its
blocks are the declared receiver family itself, the same matrices, so its
predicates and summation condition are that family's exact admissibility,
which is decided exactly on those matrices and which nothing in this theorem
supplies; (C4) for every operator $Y$ and every
$\rho$, $\tr[Y\idmap(\rho)]=\tr[Y\rho]$, so $\idmap^{\dagger}=\idmap$ and
\eqref{eq:pullback}, whose fibre $\gamma^{-1}(v)$ is $\{v\}$ for
$\gamma=\idmap$, reads $F_v=\Fph_v$; this holds entry for entry because
Definition~\ref{def:compile} fills $\{F_v\}$ from the last link's
\exttermref{../reductions/contract.pdf}{term-target-povm}{target POVM}, and that target POVM is the declared family
itself, $F_v=\Fph_v$ for every $v$, by Definition~\ref{def:banda}\,(b), so
the delivered matrix \emph{is} the declared one --- a clause of the
hypothesis, and not a consequence of the link's channel being $\idmap$, which
fixes nothing of the target POVM; (C5) is Definition~\ref{def:banda}\,(b),
which gives the one link its status tag, \emph{exact}, and its constraint and
licence slots, both empty; (S1)--(S4) quantify over the links of
$\Chsrc$ and that list is empty, and the source identity (S3) holds with $\Xi$
the identity and $\rvirt_x=\rho_x$, as \xref{transmitter/def:source-segment}
says of the empty segment; (X1) is the shape of the record, with nothing in
the first component.

(iii) $\CList$ is the concatenation of the specification's entry list
$\Spec_{\mathrm{con}}$ with the entries transported by the links
(Definition~\ref{def:compile}); the source segment has no links and the
single receiver link has an empty constraint slot
(Definition~\ref{def:banda}\,(b)), so the links transport no entry and
$\CList$ is $\Spec_{\mathrm{con}}$'s list. The transported-constraint clause
of Assumption~\ref{ass:link} quantifies over the entries a link transports
and is vacuous. Where $\Spec_{\mathrm{con}}$ is empty, $m=0$ and the gate's
premise "$m\ge1$" fails, so Theorem~\ref{thm:gates}\,(ii) is not engaged, and
no step of this proof invokes that theorem's statement.

(iv) By (ii) and (iii) every emitted slot is filled from
$\Spec_{\mathrm{src}}$, $\Spec_{\mathrm{rec}}$, $\Spec_{\mathrm{pub}}$ and
$\Spec_{\mathrm{con}}$, and the only chain-derived slots ($\Qsys$, $\SecSet$,
$(d_s)$, $\{F_v\}$, $\CList$) take the values of the directly declared
receiver data and of the specification's own entry list. The
channel-existence clause of Assumption~\ref{ass:link} --- that the link's
channel exists as claimed --- is discharged by exhibition, since $\idmap$ is
a channel; what remains of the receiver is Assumption~\ref{ass:ideal}. An
entry of $\Spec_{\mathrm{con}}$ is declared by the specification, as on any
route, and the gate it engages is Theorem~\ref{thm:gates}'s, whose licence is
Assumption~\ref{ass:atk}; nothing in this proof rests on it.

(v) Each clause is evaluated at $L_{\mathrm s}=0$ and $L=1$,
$\Lambda_1=\idmap_{\Qsys}$ (Definition~\ref{def:banda}\,(b)). (i)(a) and
(iv): $\idmap_{\Qsys}$ is completely positive and trace-preserving, the
identity on $\Lin{\Qsys}$ tensored with any identity being an identity; it is
a channel that exists, and the link supplies it as data, exactly that map
(Definition~\ref{def:banda}\,(b)). (i)(b) is the pullback identity of (ii),
$\idmap^{\dagger}(F_v)=F_v=\Fph_v$ entry for entry. (i)(c): by (ii), that
is by the target-POVM clause of Definition~\ref{def:banda}\,(b), the
delivered $\{F_v\}$ is the declared $\{\Fph_v\}$, the same matrices, and
nothing else of the protocol depends on the link, so the protocol using the
delivered measurement and the protocol using the declared one are one
protocol, and a security statement about it implies itself. Where a
setting-dependent weighting is applied after the link, the emitted family is
$\{P(f\,|\,t(v))\,F_v\}_{(v,f)\in\Vset_P}$ (Definition~\ref{def:compile},
Definition~\ref{def:weighted}), and with $F_v=\Fph_v$ entry for entry it is
the declared family weighted by the same $P$ through the same setting map,
the same matrices again; so the weighted protocol using the emitted
measurement and the weighted protocol using the declared one are again one
protocol. Clause (ii) of
the assumption quantifies over the links of the source segment, and there are
none. Clause (iii) quantifies over the constraint entries the link transports,
and its constraint slot is empty (Definition~\ref{def:banda}\,(b)).
\end{proof}

\section{The emission gates}\label{group-protocol-compiler/def:eligible}
\sectionintro{A \termref{term-well-formed-declaration}{well-formed declaration} may back a key-length claim only when every check passes, the transmitter's source-replacement licence record is present and re-decided wherever the \termref{term-marginal-slot}{marginal slot} is filled and Alice's dimension is at least two, and a non-empty constraint list is licensed and carries an attack-set licence. The main result, \hyperlink{claim-protocol-compiler/thm:gates}{Theorem~\ref*{thm:gates}}, shows that every gate is decidable or a presence test and that every gate failure has one fixed outcome, a refusal or an emission that is not key-length-eligible; passing a gate asserts nothing about the apparatus or about attack-set inclusion. Dropping a trusted marginal, and reading homogeneous constraints as scalars, enlarges the \termref{term-feasible-set}{feasible set} and so can only lower a minimised quantity (\hyperlink{claim-protocol-compiler/prop:conservative}{Proposition~\ref*{prop:conservative}}); no reachable row of the \termref{term-gate-outcome-table}{gate outcome table} invokes it.}
\sectionimports{\extref{../transmitter/contract.pdf}{transmitter/def:source-spec}, \extref{../transmitter/contract.pdf}{transmitter/def:srcert}, \extref{../transmitter/contract.pdf}{transmitter/thm:certificate-sufficiency}}

% claim protocol-compiler/def:eligible
\claimsummary{protocol-compiler/def:eligible}{An emitted declaration may back a key-length claim when every check passes, the transmitter's source-replacement licence record is present and re-decided wherever a marginal is declared at dimension two or more, and a non-empty constraint list is licensed and carries an attack-set licence.}
\claimbookmark{protocol-compiler/def:eligible}{Definition}{def:eligible}{Key-length eligibility}
\claimleadin{%
A declaration that passes validation is well-formed; whether it may back a
key-length claim is a further question, decided by three gates. The first is
stated on the checks rather than on the target conditions, because the checks
are what a program evaluates and the two are not interchangeable at V7
(Theorem~\ref{thm:validator-soundness}\,(iii)). The emitted flag records the outcome
(Requirement~\ref{req:eligibility-flag}), the \exttermref{../reductions/contract.pdf}{term-licence-record}{licence record} the second gate
reads is passed on to the consuming engine
(Requirement~\ref{req:source-licence-passed-on}), and why that gate admits an
\exttermref{../transmitter/contract.pdf}{term-ensemble-presentation}{ensemble presentation} but refuses a block-form certificate is
Remark~\ref{rem:gate-by-certificate-form}.

Gate (G2) re-decides the transmitter's conditions because a consumer decides
and never infers (\xref{foundations/req:consumer-decides-exactly}). A
trusted-mode specification at $d_A\ge2$ without the record does not reach
this gate at all, the transmitter refusing it before this module sees it;
that, and that no record is demanded at $d_A=1$, is
\xref{transmitter/req:licence-scope} (the mode itself is defined in
\xref{transmitter/def:source-spec}). What the attack-set licence record
attests for a party-local entry, and why that entry needs it, is
Remark~\ref{rem:party-local-soundness}.
}
\begin{definition}[Key-length eligibility]\label{def:eligible}
An emitted declaration is \termdef{term-key-length-eligible}{key-length-eligible}
when all of the following hold:
\begin{enumerate}[label=(G\arabic*),leftmargin=3.0em]
\item V1--V8 all pass --- whereupon D(i)--D(viii) all hold of the declared
numbers, by the soundness half of Theorem~\ref{thm:validator-soundness}\,(i);
\item \emph{(a consumption gate on the transmitter artefact)} if the
\termref{term-marginal-slot}{marginal slot} is filled and $d_A\ge2$: the transmitter artefact
carries the source-replacement licence record of
\xref{transmitter/def:srcert} --- the certificate
$(\tau_A,\Psi,\{p_x,\rvirt_x\})$ with its attestation (SR3) --- whose
\emph{five decidable conditions} the transmitter decided exactly on the same
data (\xref{transmitter/thm:certificate-sufficiency}\,(i)) and which \emph{this
module re-decides on the emitted data}: $\Psi=\Psi^{\dagger}$; $\Psi^2=\Psi$;
$\tr\Psi=1$, the trace accumulated over the row index and
\exttermref{../foundations/contract.pdf}{term-rounds-nowhere}{rounding nowhere}; (SR1); and (SR2) --- each
\emph{exactly}, as matrix and scalar identities. The transmitter's decision is
not taken on its word. \emph{No condition of (G2) is stated on
a vector.} The certificate's \exttermref{../reductions/contract.pdf}{term-declared-datum}{declared datum} is the operator $\Psi$; the unit
vector $\ket{\psi}$ with $\Psi=\dyad{\psi}$ exists by $\Psi^2=\Psi$ and
$\tr\Psi=1$ and is never itself declared (\xref{transmitter/def:srcert}).
Clause (SR3) carries no number and is not decidable
(Theorem~\ref{thm:gates}\,(i)): (G2) asks for its presence only. At $d_A=1$
no record is demanded, the licence being vacuous there;
\item if $m\ge1$: (a) $\CList$ lies in the licensed class of the emitted source
mode (Definition~\ref{def:homog}) --- \emph{a type check on each entry's
right-hand-side tag, comparing no numbers} --- \emph{and}
(b) the attack-set licence record, the attestation of
Assumption~\ref{ass:atk}, is present.
\end{enumerate}

\emph{(G2) by form of the certificate.} The five conditions above are those
of a certificate declared in \exttermref{../transmitter/contract.pdf}{term-pure-form}{pure form}: a pure-form certificate is
re-decided exactly as stated, and refused when one of the five fails.

A certificate declared in \exttermref{../transmitter/contract.pdf}{term-block-form}{block form} is refused at this gate by its
declared form, whatever its blocks, by row~3 of the gate outcome table of
Definition~\ref{def:gate-outcome-table}. A block-form
certificate (\xref{transmitter/def:srcert}) is a positive semidefinite
unit-trace \exttermref{../foundations/contract.pdf}{term-block-list}{block list} on which that definition imposes no purity, so
one with more than one non-zero block, or whose one non-zero block is not a
rank-one projector, also fails $\Psi^2=\Psi$; but the refusal does not rest on
that, since one whose one non-zero block is a rank-one projector can meet all
five conditions above.

An \exttermref{../transmitter/contract.pdf}{term-ensemble-presentation}{ensemble presentation} declares no $\Psi$ and no $\{M_x\}$ to
re-decide. For it the gate decides, in place of the five conditions, what the
presentation declares: (E1), (E2) and (E3) of \xref{transmitter/def:srcert}
on the declared weighted ensemble, and that the emitted declaration is in
\termref{term-ensemble-form}{ensemble form} with its Alice slots \emph{the} ones the presentation
defines --- the \exttermref{../transmitter/contract.pdf}{term-sector-structure}{sector structure} of the factor tagged $A$ is the
matched structure, and for every $t$ and $x$ the block $N^{(t)}_x$ is entry
for entry the transpose of the declared block $\sigma^{\mathrm v,(t)}_x$ and
$\tau^{(t)}_A$ the transpose of their declared-order sum. \emph{A presentation
accompanying a POVM-form declaration is refused}, nothing tying the declared
$\{M_x\}$ to it, which is the transmitter's own rule for a POVM-form consumer.
In every form, (G2) asks for the presence of (SR3).

\emph{(G3)(a) with the \exttermref{../reductions/contract.pdf}{term-party-local-form}{party-local form}.} The licensed class of
Definition~\ref{def:homog} contains, in either source mode, the entries
declared \termref{term-party-local}{party-local}; the decision remains a type check on each
entry's tag and compares no numbers. A list carrying such an entry is
key-length-eligible only with the attack-set licence record of (b), as every
non-empty list is, and what that record attests for a tagged entry is a
certificate of the module that produces the chain, on the scoped factor's own
channel (Definition~\ref{def:homog}).
\end{definition}
\claimmeta{\metalabel{Uses}\mref{Def.}{def:declaration}, \mref{Def.}{def:exact}, \mref{Def.}{def:homog}, \mref{Thm.}{thm:validator-soundness}, \extref{../transmitter/contract.pdf}{transmitter/def:srcert}, \extref{../transmitter/contract.pdf}{transmitter/thm:certificate-sufficiency}
\metasep \metalabel{Used by}\mref{Rem.}{rem:gate-by-certificate-form}, \mref{Def.}{def:gate-outcome-table}, \mref{Thm.}{thm:gates}, \mref{Req.}{req:source-licence-passed-on}, \mref{Req.}{req:eligibility-flag}, \mref{Req.}{req:gate-outcomes}, \mref{Req.}{req:attestation-presence}, \mref{Req.}{req:vacuous-gate-report}
\metasep newly derived
\metasep implemented (effective: accepted)
\metasep \metalabel{Code}\coderef{https://github.com/fraunhofersg/certiqs-key/blob/v0.1.0/src/certiqs_key/protocol_compiler/gates.py}{certiqs_key.protocol_compiler.gates:GateId}
\metasep \metalabel{Oracles}\coderef{https://github.com/fraunhofersg/certiqs-key/blob/v0.1.0/tests/oracles/registry_protocol_compiler.py}{registry_protocol_compiler.py}, \coderef{https://github.com/fraunhofersg/certiqs-key/blob/v0.1.0/tests/oracles/registry_transmitter.py}{registry_transmitter.py}}

% claim protocol-compiler/rem:gate-by-certificate-form
\claimsummary{protocol-compiler/rem:gate-by-certificate-form}{Gate (G2) refuses every block-form certificate by its form, pure or not, and admits an ensemble presentation, whose blocks are rank one by theorem; the refusal rests on scope, not on that difference. At several Alice sectors a presentation is not globally pure, and the gate does not demand it.}
\claimbookmark{protocol-compiler/rem:gate-by-certificate-form}{Remark}{rem:gate-by-certificate-form}{The source gate on each certificate form}
\claimleadin{%
Gate (G2) of Definition~\ref{def:eligible} re-decides a pure-form
certificate, refuses a block-form one and decides an ensemble presentation by
what it declares. What purity adds to a certificate, and why the certificate
has a block form and an ensemble presentation at all, is
\xref{transmitter/rem:certificate-forms}; this remark says what separates the
forms at the gate.
}
\begin{remark}[The source gate on each certificate form]\label{rem:gate-by-certificate-form}
\emph{The block form.} Gate (G2) decides $\Psi^2=\Psi$ of a pure-form
certificate and refuses a \exttermref{../transmitter/contract.pdf}{term-block-form}{block-form} certificate by its
declared form, whatever its blocks (Definition~\ref{def:eligible}). Most
block-form certificates fail $\Psi^2=\Psi$ as well --- every one with two or
more non-zero blocks, and every one whose one non-zero block is not a
rank-one projector --- but one whose one non-zero block is a rank-one
projector is pure and is refused all the same. So (G2) is not the consumer that
\xref{transmitter/rem:certificate-forms} describes, which demands purity of
every certificate and refuses a block form for want of it: the ground of the
refusal is the scope rule stated at the end of this remark.

\emph{What separates a general block form from an ensemble presentation} is
not the word ``purity'', which is waived on both sides once Alice carries
several sectors. A general block form declares blocks that are positive
semidefinite and \emph{not known to be rank one}, so the ``only if'' half of
\xref{transmitter/thm:certificate-sufficiency}\,(ii) --- proved for a block
form exactly where every non-zero block is a positive multiple of a rank-one
projector --- is lost with it, and a source/POVM mismatch would then be
excluded only by the POVM condition decided on its own. A presentation's
blocks are rank one \emph{by theorem}, so that half is not lost.

\emph{An ensemble presentation at one Alice sector.} For an \exttermref{../transmitter/contract.pdf}{term-ensemble-presentation}{ensemble
presentation} there is no declared $\Psi$ and no declared $\{M_x\}$ to
re-decide: purity of the blocks, (SR1) and (SR2) are \emph{theorems} of the
presentation (\xref{transmitter/lem:ensemble-presentation}\,(i), (ii)), that
lemma's (ii) giving every non-zero block as a positive multiple of a rank-one
projector. At one Alice sector no purity requirement is waived: there the
$\Psi$ the presentation defines is that one block, of unit trace, hence a
rank-one projector, so the five conditions of (G2) hold of it as theorems
rather than as decisions, and the gate demands of such a presentation
everything it demands of a pure-form certificate.

\emph{At two or more Alice sectors.} There the defined operator is
$\Psi=\bigoplus_t\Psi_t$ with $\tr\Psi_t=g_t$ and $\sum_tg_t=1$, so
$\Psi^2=\bigoplus_tg_t\Psi_t\ne\Psi$ as soon as two blocks are non-zero, by
the same arithmetic by which a block form with two non-zero blocks fails
$\Psi^2=\Psi$. Global purity is therefore not
demanded of a presentation with several Alice sectors. What the gate rests on
there instead is two statements of the transmitter's, both available for a
presentation, of which a general block form has only the second.
\begin{enumerate}[label=(\alph*),leftmargin=2.2em]
\item Every block is rank one by theorem
  (\xref{transmitter/lem:ensemble-presentation}\,(ii)), which is exactly the
  hypothesis under which \xref{transmitter/thm:certificate-sufficiency}\,(ii)
  holds in \emph{both} directions.
\item \xref{transmitter/thm:certificate-sufficiency}\,(iii), the sufficiency
  statement at the \exttermref{../foundations/contract.pdf}{term-module-boundary}{module boundary}, assumes (SR1)--(SR3) and that
  module's sequential assumption and \emph{reads purity nowhere}.
\end{enumerate}
The asymmetry in (a) is not relied on. The ``only if'' direction is what a
consumer that does \emph{not} decide the POVM condition separately would need,
and this module decides it at check V3 in either form ---
$\sum_xM^{(t)}_x=\idon{A_t}$ in \termref{term-povm-form}{POVM form}, the ensemble identity of
\eqref{eq:ensembleform} in ensemble form (D(iii)) --- which is stronger than
the compressed condition that converse delivers; and for a presentation there
is in any case no mismatch to detect, $\{M_x\}$ being \emph{defined} by the
presentation rather than declared and gate (G2) pinning the declaration's
Alice slots to the transposed declared blocks entry for entry, so that
direction holds vacuously. Clause (a) records what the two forms differ in;
it does not carry the refusal of a general block form, which rests on scope:
admitting one would widen (G2) in the unsafe direction, to certificates that
can accompany a declaration in the \termref{term-basic-format}{basic format}.

\emph{Where a presentation can occur.} A presentation is admitted only beside
an ensemble-form declaration, and no declaration in the basic format is in
ensemble form.

\emph{No form's (SR1)--(SR2) is assumed.} In none of its forms are the
conditions (SR1)--(SR2) of a \exttermref{../transmitter/contract.pdf}{term-source-replacement-certificate}{source-replacement certificate} an
assumption of this module: what gate (G2) decides of each form is
Definition~\ref{def:eligible}'s and Theorem~\ref{thm:gates}'s, and
(SR1)--(SR2) are named as a hypothesis at the one place a derivation uses
them, \xref{transmitter/thm:certificate-sufficiency}.
\end{remark}
\claimmeta{\metalabel{Uses}\mref{Def.}{def:eligible}
\metasep newly derived
\metasep accepted}

% claim protocol-compiler/def:gate-outcome-table
\claimsummary{protocol-compiler/def:gate-outcome-table}{The gate outcome table gives each gate failure a refusal naming what failed or an emission that is not key-length-eligible; its two certificate-absent rows give the compiler no outcome. Where a declaration meets a refusal row and an emission row, the refusal wins.}
\claimbookmark{protocol-compiler/def:gate-outcome-table}{Definition}{def:gate-outcome-table}{The gate outcome table}
\claimleadin{%
Rows 2 and 4 describe a case that does not reach the compiler: no
\termref{term-protocol-specification}{protocol specification} meets their premise (Theorem~\ref{thm:gates}\,(ii)), and the
transmitter refuses such a request before any artefact exists
(Requirement~\ref{req:gate-outcomes}). Their cell says so; the table gives
the compiler no outcome there.
}
\begin{definition}[The gate outcome table]\label{def:gate-outcome-table}
The
\termdef{term-gate-outcome-table}{gate outcome table}
is the following table of the failures of the gates of
Definition~\ref{def:eligible} and their outcomes; where the failures of one declaration meet the
premises of a refusal row and of an emission row both, the outcome is the
refusal.

\medskip
\noindent\emph{Gate failure to emitted outcome.} No row substitutes a default
value, drops an entry, changes the source mode or synthesizes an attestation.
\phantomsection\label{tab:gateoutcome}

\noindent{\small
\begin{tabular}{@{}>{\raggedright\arraybackslash}p{0.58\textwidth}
                  >{\raggedright\arraybackslash}p{0.36\textwidth}@{}}
\hline
\emph{What failed} & \emph{Outcome} \\
\hline
(G1): any of V1--V8 & refusal, naming the check \\
(G2): trusted mode at $d_A\ge2$, certificate \emph{absent}, and
$m=0$ &
none (no outcome of the compiler) \\
(G2): trusted mode \emph{at any $d_A$}, certificate \emph{present}
but one
of the five conditions of Definition~\ref{def:eligible}\,(G2) fails ---
$\Psi$ Hermitian, $\Psi^2=\Psi$, $\tr\Psi=1$ rounding nowhere, (SR1), (SR2) &
refusal, naming the condition \\
(G2): trusted mode at $d_A\ge2$, certificate absent, and $m\ge1$ &
none (no outcome of the compiler) \\
(G3)(a): $\CList$ not in the licensed class --- some entry declared
\termref{term-general}{general} in \exttermref{../transmitter/contract.pdf}{term-trusted-source-mode}{trusted-source mode} &
emission, not key-length-eligible \\
(G3)(b): the attack-set licence record absent &
emission, not key-length-eligible \\
\hline
\end{tabular}}

\noindent The cell of rows 2 and 4 assigns the compiler no outcome.
\end{definition}
\claimmeta{\metalabel{Uses}\mref{Def.}{def:eligible}
\metasep \metalabel{Used by}\mref{Thm.}{thm:gates}, \mref{Req.}{req:gate-outcomes}
\metasep newly derived
\metasep implemented (effective: accepted)
\metasep \metalabel{Code}\coderef{https://github.com/fraunhofersg/certiqs-key/blob/v0.1.0/src/certiqs_key/protocol_compiler/gates.py}{certiqs_key.protocol_compiler.gates:resolve_gates}}

% claim protocol-compiler/thm:gates
\claimsummary{protocol-compiler/thm:gates}{Each gate is decidable or a presence test, and each failure has one fixed outcome, a refusal or an emission flagged as backing no key length, never a default, dropped entry, changed source mode or licence the compiler writes. Passing asserts nothing about the apparatus or attack-set inclusion.}
\claimbookmark{protocol-compiler/thm:gates}{Theorem}{thm:gates}{The emission gates and their outcomes}
\claimleadin{%
One principle fixes every outcome. A check that fails, so that no declaration
is shown to be well-formed, or a certificate that does not verify is a
refusal, and no declaration is emitted; a \termref{term-well-formed-declaration}{well-formed declaration} on
verified data whose key-length claim is unsupported is emitted, not
key-length-eligible, and may be explored numerically. The refusal strings and
the flag that carry these outcomes are Requirement~\ref{req:gate-outcomes}'s.
}
\begin{mainclaim}
\begin{theorem}[The emission gates and their outcomes]\label{thm:gates}
\begin{enumerate}[label=(\roman*),leftmargin=2.2em]
\item \emph{(decidability)} (G1) is decidable by Theorem~\ref{thm:validator-soundness}; (G2)'s
conditions --- the five of Definition~\ref{def:eligible}\,(G2), namely $\Psi$
Hermitian, $\Psi^2=\Psi$, $\tr\Psi=1$ with the trace rounding nowhere, (SR1)
and (SR2) --- are finitely many exact identities on the emitted data
(\xref{transmitter/thm:certificate-sufficiency}\,(i)); (G3)(a) is decided by
\emph{reading the type of the datum each entry's right-hand-side slot holds,
which is its tag}, and by nothing else --- it compares no matrix against any
other and constructs no right-hand side; (G3)(b) is a presence test. (G2)'s
clause (SR3) and (G3)(b)'s record are attestations and are \emph{not}
decidable; they are tested for presence only;
\item \emph{(outcome of a failure)} If any gate fails, the compiler emits no
key-length-eligible declaration, and which of the \emph{two} permitted
outcomes it produces --- refusal, or emission of a declaration that is not
key-length-eligible --- is fixed by the \termref{term-gate-outcome-table}{gate outcome table}
(Definition~\ref{def:gate-outcome-table}) and is not an
implementer's choice; where the failures of one declaration meet the premises
of a refusal row and of an emission row both, the outcome is the refusal. No
protocol specification meets the premise of rows 2 and 4, the two rows of the
gate outcome table that describe a certificate-absent case: by
Definition~\ref{def:spec}\,(a) a transmitter artefact in trusted-source mode at
$d_A\ge2$ carries a source-replacement certificate. The compiler
never drops a constraint entry silently, never
substitutes a default, never changes the source mode it was handed, and never
synthesizes a licence.
\item \emph{(no assertion)} Passing (G3)(b) records that an attestation was
supplied. It asserts nothing about attack-set inclusion, and no claim of this
contract does. Likewise passing (G2) records that the transmitter's licence
record is present and that its five decided conditions hold of the emitted
data: it asserts nothing about the apparatus, and what it buys for the run is
bought under Assumption~\ref{ass:rhom} --- the certificate is decided once, for
the one declaration that describes every round, which is where this module
spends round homogeneity (\xref{transmitter/fact:sr} drops the per-round
superscript of the source state under exactly that assumption of the
consumer's).
\end{enumerate}

\emph{The general format.} For a declaration in the general format of
Definition~\ref{def:declaration}\,(DX1)--(DX6) the three parts hold with the
following readings, of which the three above are the case of one Alice sector,
one \termref{term-measured-factor}{measured factor}, \termref{term-povm-form}{POVM form}, a pure-form certificate and no
party-local entry.
\begin{enumerate}[label=(\roman*$'$),leftmargin=2.6em]
\item \emph{(decidability)} (G1) is decidable by Theorem~\ref{thm:validator-soundness} in its
general-format reading. (G2) keeps its five conditions and is decidable in
\emph{one further declared form} (Definition~\ref{def:eligible}): the five
conditions of a pure-form certificate, as above, a certificate declared
in \exttermref{../transmitter/contract.pdf}{term-block-form}{block form} being refused at this gate; and, for an \exttermref{../transmitter/contract.pdf}{term-ensemble-presentation}{ensemble
presentation}, (E1), (E2), (E3) on the declared weighted ensemble together with
the entry-for-entry identity of the declaration's Alice slots with the
transposes of the declared blocks and of their declared-order sum --- finitely
many exact identities, comparisons of labels and integers, and finitely many
summation decisions, the number fixed by $|\mathbb{T}|$, $|\Xset|$ and the
declared dimensions. \emph{Purity is decided by nothing there}: a presentation
declares no $\Psi$, and purity of its blocks, (SR1) and (SR2) are theorems of
the presentation (\xref{transmitter/lem:ensemble-presentation}\,(i), (ii)), the
second giving every non-zero block as a positive multiple of a rank-one
projector. \emph{At one Alice sector it is not thereby waived}, the defined
$\Psi$ being that one block, of unit trace, hence a rank-one projector of which
the five conditions hold. \emph{At two or more it is waived}, and the gate
rests there on two statements --- rank-one blocks by theorem, which is the
hypothesis under which \xref{transmitter/thm:certificate-sufficiency}\,(ii) holds in both directions,
and that theorem's (iii), which reads purity nowhere --- the first of which a
general block form does not supply. (G3)(a) is decided by reading the type of the
datum each entry's right-hand-side slot holds and, for the third form, the tag
field it arrives with --- three forms being in the licensed class, and still
\emph{nothing computed and nothing compared numerically}. (G3)(b) is a presence
test, and (SR3) is tested for presence in every form.
\item \emph{(no row is added)} The outcome of every failure of the general
format is fixed by a \emph{row the gate outcome table already has}. A
general-format check
of Definition~\ref{def:validator} that fails is (G1) and is row 1, the refusal
naming the check. A certificate that is \emph{present} and fails a condition
of (G2) is row 3: a certificate that does not verify is a defect in the
supplied data, not a weaker claim. \emph{Row 3's premise is read of the
declared form}, which reads the cell and adds no row: for a pure-form
certificate the five conditions the cell enumerates are the conditions, and
the cell is literal; for a certificate declared in block form the condition
it fails is the form itself, (G2) admitting no block-form certificate
(Definition~\ref{def:eligible}), and the cell's premise --- a certificate
present and failing a condition of (G2) --- covers it, whether or not it also
fails $\Psi^2=\Psi$; for an ensemble presentation the conditions are (E1) to (E3) and the
identity of the declaration's Alice slots with the transposed declared blocks,
and the cell's premise --- a certificate present and failing a condition of
(G2) --- covers them, the enumeration in it being the pure form's list and not
a restriction of the premise. \emph{An ensemble presentation accompanying a
POVM-form declaration is a failure of that kind and is refused} (row 3),
because nothing ties the declared
$\{M_x\}$ to the presentation (Definition~\ref{def:eligible}). A list
carrying an entry declared general in trusted-source mode is row 5 and is
emitted, not key-length-eligible; \emph{a list carrying a \termref{term-party-local}{party-local}
entry is not that row}, that form being licensed in either mode.
\item \emph{(no assertion)} Item (iii) holds as stated, and passing (G3)(b)
records, rather than asserts, one further item: for a list carrying a
party-local entry, what the attack-set licence record attests is a certificate
of the module that produces the chain, on the scoped factor's own channel
(Definition~\ref{def:homog}), and
\emph{no function of the emitted matrices decides it}, exactly as for the
untagged forms. Passing (G2) on an ensemble presentation likewise records that
the transmitter's record is present and that the conditions of \emph{that
form} hold of the emitted data, and asserts nothing about the apparatus; what
it buys for the run is bought under Assumption~\ref{ass:rhom} as in (iii), the
certificate being decided once for the one declaration that describes every
round.
\end{enumerate}
\end{theorem}
\end{mainclaim}
\claimmeta{\metalabel{Uses}\mref{Def.}{def:declaration}, \mref{Def.}{def:exact}, \mref{Def.}{def:homog}, \mref{Def.}{def:spec}, \mref{Def.}{def:compile}, \mref{Thm.}{thm:validator-soundness}, \mref{Def.}{def:eligible}, \mref{Def.}{def:gate-outcome-table}, \extref{../transmitter/contract.pdf}{transmitter/def:source-spec}, \extref{../transmitter/contract.pdf}{transmitter/thm:certificate-sufficiency}
\metasep \metalabel{Assumes}\aref{A-ATK}{ass:atk}, \aref{A-RHOM}{ass:rhom}
\metasep \metalabel{Used by}\mref{Rem.}{rem:notfallback}, \mref{Req.}{req:gate-outcomes}, \mref{Req.}{req:attestation-presence}
\metasep newly derived
\metasep implemented (effective: accepted)
\metasep \metalabel{Code}\coderef{https://github.com/fraunhofersg/certiqs-key/blob/v0.1.0/src/certiqs_key/protocol_compiler/gates.py}{certiqs_key.protocol_compiler.gates:resolve_gates}
\metasep \metalabel{Oracles}\coderef{https://github.com/fraunhofersg/certiqs-key/blob/v0.1.0/tests/oracles/registry_protocol_compiler.py}{registry_protocol_compiler.py}, \coderef{https://github.com/fraunhofersg/certiqs-key/blob/v0.1.0/tests/oracles/registry_transmitter.py}{registry_transmitter.py}}

\begin{proof}
\emph{Idea.} Each gate condition is either a finite list of exact identities,
a type read off each constraint entry, or a presence test; the outcome of a
failure is the one the gate outcome table of
Definition~\ref{def:gate-outcome-table} fixes, no row of it changes the
source mode or drops an entry, and by the definition of a protocol
specification no specification meets its certificate-absent rows; and
the general format adds decisions of the same kinds and failures that land on
existing rows.

\emph{Step 1 (decidability, (i)).} Each clause is as stated:
Theorem~\ref{thm:validator-soundness}\,(iv) for (G1);
\xref{transmitter/thm:certificate-sufficiency}\,(i) for the (G2) identities.
For (G3)(a): by Definition~\ref{def:homog} the licensed class is defined by
which of the two forms each entry \emph{declares}, and the tag is a field of
the datum (Definition~\ref{def:declaration}), so deciding the class is reading
$m$ tags and checking that the datum each names has the type it names --- $m$
comparisons of a label and a type, terminating and exact. \emph{Nothing is
computed and nothing is compared numerically.} Presence tests are decidable
trivially. That (SR3) and the licence record are attestations is
\xref{transmitter/def:srcert}\,(SR3) and Assumption~\ref{ass:atk}: neither is
a property of the emitted matrices, so no function of them decides it.

\emph{Step 2 (the marginal slot and the certificate-absent rows, for (ii)).}
$\tau_A$ reaches the emitted
data only through the transmitter artefact's marginal slot ---
Definition~\ref{def:compile} passes that slot through, and no input of this
module carries a marginal or a mode separately (Definition~\ref{def:spec}\,(a):
the mode is read off the artefact's marginal slot and is never a second
datum). The premise of
rows 2 and 4 is an artefact in trusted mode at $d_A\ge2$ whose certificate is
absent, and by Definition~\ref{def:spec}\,(a) a transmitter artefact in
trusted-source mode at $d_A\ge2$ additionally carries a source-replacement
certificate; so no protocol specification meets that premise. That the engine
refuses such a request is Requirement~\ref{req:gate-outcomes}'s, and no step
here uses it. At
$d_A=1$, where (G2) does not demand the certificate at all, the slot holds
$\tau_A=[1]$, the value the transmitter filled it with
(Definition~\ref{def:spec}\,(a)), and this
module passes the slot through --- that is not a default substituted for a
missing datum here, which item (ii) forbids, because this module fills
nothing: $\Seq{A}=\{[1]\}$ is a singleton, so D(ii) admits that value and no
other, and at $d_A=1$ the two source modes impose the same condition
(\xref{transmitter/def:source-spec}), so the filled slot excludes nothing.

\emph{Step 3 ((ii) and (iii)).} (ii) and (iii) are rules of the
\termref{term-compilation-map}{compilation map}, stated here because they are the behaviour the gate exists to produce.
(ii)'s precedence clause is part of that rule and is what makes the table fix
\emph{one} outcome for each declaration: one declaration can meet the premise
of a refusal row and of an emission row at once --- a certificate failing a
condition of (G2), row~3, beside a list carrying an entry declared general,
row~5 --- and the clause takes the refusal, so no declaration whose supplied
data failed a check or a certificate is emitted.
(ii) is consistent: the one non-refusal outcome is exactly the one for which
this contract proves the result sound --- an ineligible declaration is not
consumed by a key-length statement at all --- and no row changes the source
mode, re-reads a constraint entry or drops one. No row invokes
Proposition~\ref{prop:conservative}: the rows on which it could bear, the
certificate-absent ones, are met by no protocol specification (Step~2).
(iii)'s second sentence is
Assumption~\ref{ass:rhom} read where it is spent: (G2) decides the certificate
once and \xref{transmitter/thm:certificate-sufficiency}\,(iii) is a statement
about one round's source state, so the flag says of the run only what one
declaration can say of every round.

\emph{Step 4 (the general format, (i$'$)).} Theorem~\ref{thm:validator-soundness}\,(iv) in its
general-format reading gives (G1). For (G2): each condition
Definition~\ref{def:eligible} lists for a form is either an identity or an
order relation between declared numbers, decided as
Definition~\ref{def:exact} decides such a predicate, or a comparison of labels,
integers and structures, or one summation decision --- and each list is finite,
its length fixed by $|\mathbb{T}|$, $|\Xset|$ and the declared dimensions. That
purity is decided by nothing \emph{for a presentation} is read off
\xref{transmitter/def:srcert}: a presentation declares no $\Psi$, and purity,
(SR1) and (SR2) are conclusions of the transmitter's lemma, so no function of
the emitted data is asked for them. Three separate statements describe what
the gate then demands. \emph{First}, the five conditions are decided of every
certificate declared in pure form, and at one Alice sector --- the only case a
declaration in the \termref{term-basic-format}{basic format} can present --- they \emph{hold} of the
$\Psi$ a presentation defines, by
\xref{transmitter/lem:ensemble-presentation}\,(ii) with unit \exttermref{../reductions/contract.pdf}{term-total-weight}{total weight}.
\emph{Second}, a presentation is admitted only beside an
\emph{ensemble-form} declaration, which no basic-format declaration is; and a
certificate declared in \emph{block form} is refused here
(Definition~\ref{def:eligible}). \emph{Third}, at $|\mathbb{T}|\ge2$ the
defined $\Psi=\bigoplus_t\Psi_t$ has $\Psi^2=\bigoplus_tg_t\Psi_t\ne\Psi$ once
two blocks are non-zero, so global purity is not demanded of a presentation
there, and what the gate rests on in its place is what
Remark~\ref{rem:gate-by-certificate-form} names --- rank-one blocks by theorem,
giving both directions of \xref{transmitter/thm:certificate-sufficiency}\,(ii),
and that theorem's (iii), which reads purity nowhere. A general block form
lacks the first and has the second, that theorem being stated in every form;
the refusal of a block form rests on scope and not on that difference, and
the distinction is not the word ``purity''
(Remark~\ref{rem:gate-by-certificate-form}). For (G3)(a): by
Definition~\ref{def:homog} the licensed class is defined by which of the three
forms an entry \emph{declares}, the first two by the type of the datum and the
third by the tag field beside a real scalar, so deciding the class is reading
$m$ tags and checking that the datum each names has the type it names --- $m$
comparisons of a label and a type, terminating and exact, and again nothing
computed: no form of the general format introduces a numerical comparison.

\emph{Step 5 (the general format, (ii$'$)).} (ii$'$) is a statement about
where each failure lands, and it is proved by exhibiting the row: a failing
check is (G1) by the definition of that gate, and (G1)'s row is a refusal
naming the check, so every failing check of the general format lands on that
row; a certificate present and failing a condition of (G2) is the premise of
row 3, its outcome a refusal naming the condition. \emph{A block-form
certificate meets that premise by its form}: the record is present, and the
condition of (G2) it fails is that (G2) admits no certificate declared in
block form (Definition~\ref{def:eligible}), which holds of every block-form
certificate whatever its blocks --- including one whose one non-zero block is
a rank-one projector and which therefore meets $\Psi^2=\Psi$ --- so the
outcome is a refusal naming the form. \emph{An ensemble presentation
beside a POVM-form declaration meets the premise too}: the record is present,
and the condition of its form that it fails is the one requiring the
declaration's Alice slots to be \emph{the} slots the presentation defines,
which a POVM-form declaration has not got --- so the outcome is a refusal, and
the refusal names which condition failed. No premise of any other row is met
by a general-format failure, and each of the six rows keeps its outcome; in
particular no general-format failure yields a key-length-eligible
declaration, and none yields an emission on a row whose outcome is a refusal.

\emph{Step 6 (the general format, (iii$'$)).} The added item is the
\termref{term-soundness-condition-of-a-party-local-entry}{soundness condition of a party-local entry} (Definition~\ref{def:homog};
that the party-local form is sound only under it is
Remark~\ref{rem:party-local-soundness}), read where the record is tested:
every attack reaches the scoped factor through a channel on that factor alone
carrying the spectral certificate, which is a certificate of the module that
produces the chain (Definition~\ref{def:eligible}), and the gate tests for
the presence of the record that attests it. \emph{So admitting the third form
at (G3)(a) moves no decision from an attestation to this module and none from
this module to an attestation}: it adds one licensed tag and one clause to
what the same record attests. The rest of (iii) holds as stated, including
where Assumption~\ref{ass:rhom} is spent.
\end{proof}

% claim protocol-compiler/rem:degenerate
\claimsummary{protocol-compiler/rem:degenerate}{In untrusted-source mode gate (G3)(a) holds vacuously, so the constraint gate rests on the attack-set licence alone. At dimension one (G3)(a) is still a type check, so a general-tagged 1\ensuremath{\times}1 entry is unlicensed, although it could always have been declared homogeneous.}
\claimbookmark{protocol-compiler/rem:degenerate}{Remark}{rem:degenerate}{The licensed-class gate in untrusted-source mode and at d\_A=1}
\claimleadin{%
Gate (G3) of Definition~\ref{def:eligible} reads as two checks, a type check
on the list and a presence test for the licence. This remark says what the
type check amounts to in \exttermref{../transmitter/contract.pdf}{term-untrusted-source-mode}{untrusted-source mode} and at $d_A=1$; what the gate
report records of a vacuous pass is Requirement~\ref{req:vacuous-gate-report}.
}
\begin{remark}[The licensed-class gate in untrusted-source mode and at $d_A=1$]\label{rem:degenerate}
Gate (G3)(a) is a genuine restriction in \exttermref{../transmitter/contract.pdf}{term-trusted-source-mode}{trusted-source mode}: an entry
declared \termref{term-general}{general} fails it, and such an entry is exactly what an
upstream module emits when it has an operator right-hand side rather than a
scalar one. It is \emph{vacuously true} in \exttermref{../transmitter/contract.pdf}{term-untrusted-source-mode}{untrusted-source mode} at any
$d_A$, where every right-hand side is already a scalar by D(vi), so the gate
degenerates there from a restriction into a type check that a
\termref{term-well-formed-declaration}{well-formed declaration} passes by construction, and rests entirely on
(G3)(b), which no program can verify.

\emph{In trusted-source mode at $d_A=1$ the gate is not vacuous.} Every
$1\times1$ Hermitian $\Theta_k$ is $\Theta_k\cdot[1]$, but under the declared
tag an entry tagged general declares a matrix and fails a type check at every
dimension: (G3)(a) is a type check at every $d_A$, so such an entry is
unlicensed (Definition~\ref{def:homog}), and the consuming engine reads its
gate the same way (\xref{meat-core/rem:cgate}). A declaration carrying one is
emitted, not key-length-eligible. Its producer could have declared the entry
homogeneous with the same condition, since at $d_A=1$ the constructed
$c_k\tau_A=c_k[1]$ is the scalar itself, so there the gate rejects nothing a
producer cannot trivially avoid.
\end{remark}
\claimmeta{\metalabel{Used by}\mref{Req.}{req:vacuous-gate-report}
\metasep newly derived
\metasep accepted}

% claim protocol-compiler/prop:conservative
\claimsummary{protocol-compiler/prop:conservative}{Dropping a trusted marginal and reading homogeneous constraints as scalars enlarges the feasible set, for every admissible marginal, so it can only lower a minimised objective. It compares the two source modes and is no rule of emission.}
\claimbookmark{protocol-compiler/prop:conservative}{Proposition}{prop:conservative}{Dropping a trusted marginal enlarges the feasible set}
\claimleadin{%
A trusted marginal is one more condition on the states the adversary may
hold. Dropping it, and reading each homogeneous entry as its scalar, can only
admit more states; the proposition states this as an inclusion of
\termref{term-feasible-set}{feasible sets}, and Remark~\ref{rem:notfallback} says what it licenses and what it does
not.
}
\begin{proposition}[Dropping a trusted marginal enlarges the feasible set]\label{prop:conservative}
Let a specification request trusted-source mode with marginal $\tau_A$
satisfying D(ii), and a licensed list
$\{(\hat O_k,\bowtie_k,c_k)\}_{k=1}^{m}$, $m\ge0$, satisfying D(iii) --- every
$\hat O_k$ a \exttermref{../foundations/contract.pdf}{term-block-list}{block list} against $\Smap$ and Hermitian block by block,
which the proof spends through Lemma~\ref{lem:hermitian-constraint-values} ---
with every entry declared \termref{term-homogeneous}{homogeneous}, so that the trusted-mode
right-hand side it denotes is the constructed $c_k\tau_A$
(Definition~\ref{def:homog}). Let $\Feas^{\mathrm{tr}}$ be the
\termref{term-feasible-set}{feasible set} of the trusted-mode declaration and $\Feas^{\mathrm{un}}$
that of the untrusted-mode declaration obtained by replacing the marginal
slot with $\nomarg$ and reading each entry's right-hand side as the scalar
$c_k$, all other slots unchanged. Then
\begin{equation}\label{eq:inclusion}
  \Feas^{\mathrm{tr}}\subseteq\Feas^{\mathrm{un}} ,
\end{equation}
and consequently, for \emph{any} functional $J$ on $\Lin{A\Qsys}$,
$\inf_{\Feas^{\mathrm{un}}}J\le\inf_{\Feas^{\mathrm{tr}}}J$. The inclusion is
\emph{uniform in $\tau_A$}: it holds for every $\tau_A$ satisfying D(ii), and
the untrusted set does not depend on which one the request carried.
\end{proposition}
\claimmeta{\metalabel{Uses}\mref{Def.}{def:exact}, \mref{Def.}{def:feasible}, \mref{Lem.}{lem:hermitian-constraint-values}, \mref{Def.}{def:homog}
\metasep \metalabel{Used by}\mref{Rem.}{rem:notfallback}
\metasep newly derived
\metasep accepted}

\begin{proof}
Let $\omega\in\Feas^{\mathrm{tr}}$. The block presentation and $\tr\omega=1$
are shared by the two definitions, as are the per-block positivity conditions;
the untrusted marginal condition is empty, and a set defined by a conjunction
is contained in the set defined by any sub-conjunction, which settles $m=0$
outright and for every $\tau_A$.

For $m\ge1$, fix $k$ and consider first
$\bowtie_k$ equal to $\succeq$, so that
$\tr_{\Qsys}[(\idon{A}\otimes\hat O_k)\omega]\succeq c_k\tau_A$. Both sides are
Hermitian --- the left by Lemma~\ref{lem:hermitian-constraint-values}, using
$\hat O_k$ Hermitian from D(iii) and $\omega$ Hermitian from \eqref{eq:feas},
and the right because $c_k$ is real and $\tau_A$ is Hermitian by D(ii) --- so
the L\"owner comparison is well posed and the traces below are real. Taking
traces of both sides and using that $Y\succeq Y'$ implies $\tr Y\ge\tr Y'$
(the trace of a positive semidefinite operator is non-negative, applied to
$Y-Y'$),
\begin{equation*}
  \tr\!\big[(\idon{A}\otimes\hat O_k)\,\omega\big]
  =\tr\Big(\tr_{\Qsys}\big[(\idon{A}\otimes\hat O_k)\,\omega\big]\Big)
  \;\ge\;\tr(c_k\tau_A)=c_k\,\tr\tau_A=c_k ,
\end{equation*}
using $\tr\tau_A=1$ from D(ii) --- \emph{an identity of the declared numbers,
which is what Definition~\ref{def:exact} delivers and what condition~(S) on the
trace is there to make true} --- and the fact that a partial trace followed by
a trace is the full trace. This is the untrusted entry with $\ge$. The case
$\bowtie_k$ equal to $\preceq$ is identical with the inequality reversed, and
the case $=$ follows by applying both. Hence $\omega\in\Feas^{\mathrm{un}}$,
which is \eqref{eq:inclusion}. The infimum
statement is immediate: an infimum over a superset is not larger than an
infimum over a subset. \emph{Uniformity in $\tau_A$}: the only property of
$\tau_A$ used is $\tr\tau_A=1$, which D(ii) gives for every admissible
marginal, and no $\tau_A$ appears in $\Feas^{\mathrm{un}}$ at all.

\emph{At $d_A=1$ the inclusion is an equality.} There $\tau_A=[1]$ is the only
marginal and $\omega_A=\tau_A$ is the condition $\tr\omega=1$, which
$\Feas^{\mathrm{un}}$ already imposes; and every licensed entry's constructed
right-hand side $c_k[1]$ is the scalar $c_k$ itself, so the two constraint
conditions coincide entry by entry. Hence the two feasible sets are equal and
the two emissions are the same declaration (\xref{transmitter/rem:da1} states
the device-side half of this).
\end{proof}

% claim protocol-compiler/rem:notfallback
\claimsummary{protocol-compiler/rem:notfallback}{\hyperlink{claim-protocol-compiler/prop:conservative}{Proposition~\ref*{prop:conservative}} licenses dropping a trusted marginal, which can only lower a minimised quantity, and never adding or weakening one; no reachable gate row uses it. At d\_A \ensuremath{\ge} 2 the maximally mixed marginal is strictly stronger than no marginal, and at d\_A = 1 it is [1], which excludes nothing.}
\claimbookmark{protocol-compiler/rem:notfallback}{Remark}{rem:notfallback}{Dropping a marginal, never replacing it}
\begin{remark}[Dropping a marginal, never replacing it]\label{rem:notfallback}
\emph{The direction of the licence.} Proposition~\ref{prop:conservative}
licenses \emph{dropping} a marginal, never adding or weakening one. Emitting
the untrusted-mode declaration when the source-replacement attestation is
absent is therefore a change in the safe direction for any downstream quantity
obtained by minimizing over the \termref{term-feasible-set}{feasible set}: dropping the marginal
\emph{lowers} any minimized quantity, so it can only cost rate; it can never
turn an unsupported declaration into a supported-looking one. A consumer that
drops a declared marginal --- the transmitter's \exttermref{../transmitter/contract.pdf}{term-untrusted-source-mode}{untrusted-source mode},
which is always available --- is doing exactly what the proposition settles,
in the direction it says.

\emph{Where the compiler uses it.} The compiler exercises that licence on no
row of the gate outcome table (Definition~\ref{def:gate-outcome-table}) that
fixes the outcomes of Theorem~\ref{thm:gates}\,(ii): the
certificate-absent rows are unreachable through the transmitter artefact, no
protocol specification meeting their premise (Theorem~\ref{thm:gates}\,(ii)),
the transmitter refusing a trusted-mode request at $d_A\ge2$ without a
certificate before this module sees it (Requirement~\ref{req:gate-outcomes}),
and the proposition stands as a statement about
the feasible sets, not as a rule of emission. On the certificate-absent row
with $m\ge1$ the inclusion is \emph{not} what is missing, since it holds at
every $m$: what is missing is a referent for the attestation, the module that
produces the chain having certified its entries against $c_k\tau_A$, and with
no declared $\tau_A$ that attestation (Assumption~\ref{ass:link}, named and
not used here) has none.

\emph{No neutral marginal.} In particular, at $d_A\ge2$,
$\tau_A=\idon{A}/d_A$ is not a neutral way to say ``nothing is known about
the source'': it fixes the marginal to maximally mixed and excludes every
$\omega$ whose $A$-marginal differs from it at all, which is strictly
stronger than $\nomarg$, which imposes no marginal condition. The compiler
never substitutes it, filling no slot (Requirement~\ref{req:pass-through}),
and at $d_A\ge2$ neither does the transmitter (\xref{transmitter/ass:src},
\xref{transmitter/rem:source-mode}). At $d_A=1$ the maximally mixed
marginal is $[1]$, the only state on $A$, and it excludes no normalized
$\omega$, so it imposes the same condition as $\nomarg$
(\xref{transmitter/rem:da1}); it is the one value the transmitter fills into
an empty marginal slot (\xref{transmitter/req:marginal-fill}), and the
compiler passes it through unchanged.
\end{remark}
\claimmeta{\metalabel{Uses}\mref{Thm.}{thm:gates}, \mref{Prop.}{prop:conservative}
\metasep newly derived
\metasep accepted}

\section{Assumptions}\label{group-protocol-compiler/ass:spec}
\sectionintro{The declared data, the declared measurements included, describe the apparatus: the physical receiver measurements, in every round and with no memory of the others, and, where no source-replacement certificate is demanded, Alice's are the ones performed. The key rule yields a key value on every input; where no certificate is demanded Alice measures as declared in every round; one declaration describes every round; public choices are private coins of declared bias, and the Band-A receiver is an ideal qubit measurement. The soundness of each reduction link and the attack-set licence are obligations of the module that presents the chain. The compiler decides every condition of a \termref{term-well-formed-specification}{well-formed specification}, and its presence checks establish none of these assumptions.}
\sectionimports{\extref{../foundations/contract.pdf}{foundations/def:round-homogeneity}, \extref{../transmitter/contract.pdf}{transmitter/def:source-segment}}

% claim protocol-compiler/ass:spec
\claimsummary{protocol-compiler/ass:spec}{Declared dimensions, alphabets, tables and measurements describe the apparatus: the receiver's in every round whatever the others did and, where no source certificate is demanded, Alice's of two or more letters; every input yields a key. The compiler checks only the last, in part.}
\claimbookmark{protocol-compiler/ass:spec}{Assumption}{ass:spec}{Specification faithful to the apparatus}
\claimleadin{%
The compiler decides, on the declared numbers, every condition of a
\termref{term-well-formed-specification}{well-formed specification}. What no function of the
\exttermref{../foundations/contract.pdf}{term-declared-data}{declared data} decides is whether those data are the apparatus: a
mis-declared dimension, or a measurement that is complete as written but is
not the one performed, passes every check. Where a
\exttermref{../transmitter/contract.pdf}{term-source-replacement-certificate}{source-replacement certificate} is demanded, in
\exttermref{../transmitter/contract.pdf}{term-trusted-source-mode}{trusted-source mode} at $d_A\ge2$, Alice's measurement is tied to the
source through that certificate, and clause~(i) says nothing of it. On the
receiver side clause~(i) includes that the measurement keeps no memory from
round to round, on every route, a directly declared family and a chain from
elsewhere included: no \exttermref{../reductions/contract.pdf}{term-declared-datum}{declared datum} describes more than one round. The
key rule's codomain is decided only in part: check V5 refuses a key rule that
withholds a key on part of a passing class, and not one that withholds it on
the whole of such a class.
}
\SetShortId{A-SPEC}
\begin{assumption}[Specification faithful to the apparatus]\label{ass:spec}
For a \termref{term-protocol-specification}{protocol specification} $\Spec$ (Definition~\ref{def:spec}):
(i) the declared dimensions, alphabets and tables of $\Spec$ describe the
apparatus, and so do its declared measurements, namely: the physical
receiver family $\{\Fph_w\}_{w\in\Vph}$ (Definition~\ref{def:spec}\,(b); in
the general format, each \termref{term-measured-factor}{measured factor}'s family of (b$'$)) is the
measurement performed, in every round, on that round's physical receiver
system alone, whatever happens in the other rounds, their outcomes and public
messages included; and
where no source-replacement certificate is demanded of $\Spec$ (in
\exttermref{../transmitter/contract.pdf}{term-untrusted-source-mode}{untrusted-source mode}, or at $d_A=1$), $\{M_x\}$
(Definition~\ref{def:spec}\,(a)) is the measurement Alice performs on $A$;
(ii) the key rule $r^{0}_{\hold}$ is valued in $\Zval$.
\end{assumption}
\claimverification{The engine does not verify this assumption; a run that rests on it records it as unverified.}
\claimreading{Your declared numbers describe your apparatus: each system has its declared dimension, every alphabet and table (announcements, sifting, the raw key bit) is what your devices and software do, and the key rule yields a key value, never "no key", on every input. Each measuring station, of either party, performs exactly its declared measurement in every round, on that round's input alone; afterpulsing or a dead time violates this. A declared receiver instrument's measurement is the model the engine emits from it; the emission step is recorded, not bounded. Where no source-replacement certificate is demanded (an untrusted source, or a one-dimensional A), Alice's declared measurement of her letter is the one she performs; \hyperlink{claim-protocol-compiler/ass:alice-meas}{A-ALICE-MEAS} adds every round, alone. The engine checks that the declaration is well formed, not that it matches your hardware.}
\claimmeta{\metalabel{Uses}\mref{Def.}{def:spec}
\metasep \metalabel{Category}protocol structure
\metasep \metalabel{Affects}secrecy
\metasep \metalabel{Assumed by}\mref{Thm.}{thm:typing}, \mref{Thm.}{thm:measurements-and-statistics}, \mref{Thm.}{thm:sifting-and-key}, \extref{../meat-core/contract.pdf}{meat-core/thm:keylength}
\metasep \metalabel{Machine key}\iface{specification_describes_apparatus}
\metasep newly derived
\metasep accepted}

% claim protocol-compiler/ass:alice-meas
\claimsummary{protocol-compiler/ass:alice-meas}{Where no source-replacement certificate is demanded, in untrusted-source mode or at d\_A = 1, and at least two letters are declared, Alice's letter in every round is the outcome of her declared measurement on that round's system A alone, whatever happens in the other rounds.}
\claimbookmark{protocol-compiler/ass:alice-meas}{Assumption}{ass:alice-meas}{Alice's declared measurement where no source certificate is demanded}
\claimleadin{%
The condition narrows Assumption~\ref{ass:spec} to Alice's measurement on
$A$ where no \exttermref{../transmitter/contract.pdf}{term-source-replacement-certificate}{source-replacement certificate} is demanded: clause~(i)
there makes $\{M_x\}$ the measurement Alice performs on $A$, and this
condition makes it so in every round separately. It constrains the
measurement and not the systems measured: in \exttermref{../transmitter/contract.pdf}{term-untrusted-source-mode}{untrusted-source mode} the
systems $A$ of different rounds may be correlated, the adversary supplying
them.
}
\SetShortId{A-ALICE-MEAS}
\begin{assumption}[Alice's declared measurement where no source certificate is demanded]\label{ass:alice-meas}
(i) Where no source-replacement certificate is demanded of the
\termref{term-protocol-specification}{protocol specification} (in untrusted-source mode, or at $d_A=1$), the
letter $x\in\Xset$ Alice obtains in every round is the outcome of the
declared measurement $\{M_x\}$ (Definition~\ref{def:spec}\,(a)) applied to
that round's system $A$ alone, the measurement she performs on $A$ being
$\{M_x\}$ whatever happens in the other rounds, their outcomes and public
messages included.
\end{assumption}
\claimverification{The engine does not verify this assumption; a run that rests on it records it as unverified.}
\claimreading{Applies where no source-replacement certificate is demanded and at least two letters are declared. Where the system A is one-dimensional, the declared measurement is only the probabilities of Alice's letter, which is then her local randomness; with a trusted source the engine refuses letter probabilities that differ from the source's and a sent state that depends on the letter. With an untrusted source and A of dimension two or more, Alice measures at her station the system A she receives. In every round she obtains her letter by the same declared measurement on that round's system alone, whatever happens in other rounds, their outcomes and the public messages included. Systems from an untrusted source, being the adversary's, may be correlated across rounds. A device of Alice's declared as a measured factor is a receiver measurement of the specification, not covered here.}
\claimmeta{\metalabel{Uses}\mref{Def.}{def:spec}
\metasep \metalabel{Category}protocol structure
\metasep \metalabel{Affects}secrecy
\metasep \metalabel{Refines}\aref{A-SPEC}{ass:spec}
\metasep \metalabel{Assumed by}\extref{../meat-core/contract.pdf}{meat-core/lem:honest-attacks}, \extref{../meat-core/contract.pdf}{meat-core/thm:keylength}
\metasep \metalabel{Machine key}\iface{alice_measurement_as_declared}
\metasep newly derived
\metasep accepted}

% claim protocol-compiler/ass:rhom
\claimsummary{protocol-compiler/ass:rhom}{The same specification, source mode and operators describe every round: the sameness clause of round homogeneity, not the absence of memory across rounds. One compiled declaration and one certificate decision then serve the whole run.}
\claimbookmark{protocol-compiler/ass:rhom}{Assumption}{ass:rhom}{Round homogeneity of the declaration}
\claimleadin{%
Rounds are not averaged and per-round variation is not modelled: one
declaration stands for every round. Only sameness is asserted; that the
devices have no memory across rounds, the second clause of
\exttermref{../foundations/contract.pdf}{term-round-homogeneous}{round homogeneity}, is not part of this assumption.
}
\SetShortId{A-RHOM}
\begin{assumption}[Round homogeneity of the declaration]\label{ass:rhom}
(i) Every round of the protocol is described by the same specification, hence
by the same declaration, the same source mode and the same operators: the
family of the specification, the declaration, the source mode and the
operators satisfies clause~(i), sameness, of
\exttermref{../foundations/contract.pdf}{term-round-homogeneous}{round homogeneity}
(\xref{foundations/def:round-homogeneity}).
\end{assumption}
\claimverification{The engine does not verify this assumption; a run that rests on it records it as unverified.}
\claimreading{One protocol description covers the whole run. Every round uses the same measurements, the same setting choices and their probabilities, and the same classical processing, and, where the source is trusted, the same source, so one declaration and one certificate decision apply to all n rounds. A run in which you change settings, or your devices drift so that later rounds need different operators, violates this. Only parts of a run fixed before it starts, each processed as its own execution and composed under \hyperref{../meat-core/contract.pdf}{}{claim-meat-core/ass:separate-executions}{A-SEPARATE}, which also asks that no device carry a state from one part into the next, can be analysed separately; a split chosen after seeing the data is not covered. That no device remembers earlier rounds is not part of this premise: for the measuring stations it is \hyperlink{claim-protocol-compiler/ass:spec}{A-SPEC}, for Alice's letter \hyperlink{claim-protocol-compiler/ass:alice-meas}{A-ALICE-MEAS}, and where the source is trusted the source premise.}
\claimmeta{\metalabel{Uses}\extref{../foundations/contract.pdf}{foundations/def:round-homogeneity}
\metasep \metalabel{Category}protocol structure
\metasep \metalabel{Affects}secrecy
\metasep \metalabel{Assumed by}\mref{Thm.}{thm:gates}, \extref{../meat-core/contract.pdf}{meat-core/thm:accumulated-entropy}
\metasep \metalabel{Machine key}\iface{round_homogeneity_declaration}
\metasep newly derived
\metasep accepted}

% claim protocol-compiler/ass:coin
\claimsummary{protocol-compiler/ass:coin}{Every public choice is drawn with its declared bias; those of fixed bias are jointly independent of each other, the quantum systems and the adversary, and one weighted by a station's setting depends on the round through that setting only. It lets a weight inside the measurement describe the round.}
\claimbookmark{protocol-compiler/ass:coin}{Assumption}{ass:coin}{Public choices as private coins of declared bias}
\claimleadin{%
A public choice of known bias enters the declaration as a weight inside the
measurement operators, and that weight describes the round only if the choice
is drawn as declared. For a choice $f$ weighted through the private setting
of the station that measures a system, the conditional independence given $t$
in clause~(ii) is what makes the only dependence of $f$ on the round its
dependence through $t$: a value whose weight varied with the device outcome
would be an inference from that outcome and not a choice, and such a
weighting is refused by type.
}
\SetShortId{A-COIN}
\begin{assumption}[Public choices as private coins of declared bias]\label{ass:coin}
(i) The declared public choices $t\in\Tpub$ are drawn from private random
sources with their declared weights $q_t$, mutually independent, whichever
party draws each, and, taken together, independent of the quantum systems and
of the adversary;
(ii) a public choice $f$ declared with a setting-dependent weighting
$P(f\,|\,t)$ --- the \termref{term-setting-dependent-form}{setting-dependent form} of
Definition~\ref{def:weighted}, which may be declared at any
\termref{term-measured-factor}{measured factor} --- is drawn in each round by a party whose
\termref{term-readout-map}{readout} records the \termref{term-setting-component}{setting component} $t$ of the
outcome, from a private random source with the declared conditional
distribution $P(\,\cdot\,|\,t)$ given that round's setting $t$; given $t$,
$f$ is independent of the quantum systems, of the letter and every outcome of
the round (the device outcome $y$ included), of every other declared public
choice and of the adversary, taken together;
(iii) every such choice is announced, or withheld, exactly as the public
structure declares.
\end{assumption}
\claimverification{The engine does not verify this assumption; a run that rests on it records it as unverified.}
\claimreading{Every public choice of your protocol (for example an actively chosen basis, or whether a round is a test round; a basis set passively by a beam splitter is part of the measurement, not a choice), whichever party draws it and at whichever of its stations, comes from a local random number generator with exactly the declared probabilities, and each is announced or kept secret exactly as declared. The choices with fixed probabilities are independent of each other and, all together, of the quantum signals and of anything an attacker does. A choice whose probabilities depend on a station's private setting is drawn by a party whose recorded outcome includes that setting, and may depend on that setting only, never otherwise on Alice's letter, another choice or a detection result.}
\claimmeta{\metalabel{Uses}\mref{Def.}{def:declaration}, \mref{Def.}{def:weighted}
\metasep \metalabel{Category}local randomness
\metasep \metalabel{Affects}secrecy
\metasep \metalabel{Assumed by}\mref{Thm.}{thm:choice-encoding}, \mref{Rem.}{rem:exactness}, \mref{Rem.}{rem:relabel}
\metasep \metalabel{Machine key}\iface{public_choice_weights}
\metasep newly derived
\metasep accepted}

% claim protocol-compiler/rem:coin-scope
\claimsummary{protocol-compiler/rem:coin-scope}{A state-dependent marginal of a setting-conditioned public choice is not a violation of \hyperlink{claim-protocol-compiler/ass:coin}{A-COIN}: it is the output-state model of the reductions module. A violation is repaired by no encoding, only by a different protocol description.}
\claimbookmark{protocol-compiler/rem:coin-scope}{Remark}{rem:coin-scope}{Scope of the private-coin assumption}
\begin{remark}[Scope of the private-coin assumption]\label{rem:coin-scope}
No encoding repairs a violation of Assumption~\ref{ass:coin}; the remedy is a
different protocol description. Under a violation the weighted POVM
misdescribes the round: for the \termref{term-setting-independent-form}{setting-independent form}
\eqref{eq:choicemass} no longer gives the physical choice statistics, and for
the setting-dependent form the relabelled family of Remark~\ref{rem:relabel}
no longer describes how $f$ was drawn; a choice correlated with the adversary,
or with the device outcome beyond its setting, is exactly the attack the
encoding assumes away. The state-dependent \emph{marginal} of $f$
that the setting-dependent form of Definition~\ref{def:weighted} produces is
not a violation: it is the output-state model of
\xref{reductions/ass:tsm}, and
Lemma~\ref{lem:weighted}\,(iii) is asserted for the setting-independent form
only (Theorem~\ref{thm:choice-encoding}).
\end{remark}
\claimmeta{\metalabel{Uses}\mref{Def.}{def:weighted}
\metasep newly derived
\metasep accepted}

% claim protocol-compiler/ass:link
\claimsummary{protocol-compiler/ass:link}{Every link of the reduction chain is sound: its channel exists and is a channel, the reduced protocol's security implies the physical one's, and every transported constraint holds on reachable states. A decided pullback or source identity is a necessary condition only, never this obligation.}
\claimbookmark{protocol-compiler/ass:link}{Assumption}{ass:link}{Soundness of every reduction link}
\claimleadin{%
The compiler decides the pullback identity and the source identity (S3)
exactly where the data make them decidable; those are necessary conditions,
and not this obligation. The obligation belongs to whoever presents the
chain, a channel supplied as data included. On the Band-A identity link,
given as data with the identity outcome correspondence and the directly
declared family as its target, every clause holds by exhibition, and on a single instance
link the reductions module emits with status exact its own theorems prove
every clause; there nothing is left to attest. The published templates are
\cite[Lemma~9.3, Eqs.~(146)--(147)]{tupkary2026decoy} for the squashing case
and \cite[Lemma~9.1, Eqs.~(138)--(139)]{tupkary2026decoy} for the source
case, composed in the order of \cite[Rem.~9.2]{tupkary2026decoy}.
}
\SetShortId{A-LINK}
\begin{assumption}[Soundness of every reduction link]\label{ass:link}
Each link of either segment of the \exttermref{../reductions/contract.pdf}{term-reduction-chain}{reduction chain} is sound:
(i) for a \emph{receiver} link, (a) where the link supplies no channel as
data, only a theorem that constructs one, its channel exists as claimed;
(b) it satisfies the pullback identity \eqref{eq:pullback}; and (c) security
of the protocol using the emitted measurement implies security of the
protocol using the physical one;
(ii) for a \emph{source} link, (a) the channel of the source identity (S3) of
\xref{transmitter/def:source-segment} exists as claimed; and (b) security of
the protocol preparing the \exttermref{../transmitter/contract.pdf}{term-virtual-ensemble}{virtual ensemble} implies security of the
protocol preparing the specified one;
(iii) for a link of either segment, every reachable state after the reduction
satisfies every constraint entry the link transports;
(iv) a link of either segment that supplies its channel \emph{as data}
supplies a completely positive and trace-preserving map.
\end{assumption}
\claimverification{The engine does not verify this assumption; a run that rests on it records it as unverified.}
\claimreading{Every reduction step between your physical devices and the finite model the engine solves is sound, and so is their combination. Each step's channel exists, a channel supplied as numbers is a genuine quantum channel, and a receiver step maps the model's measurement back to the physical one. Security of the reduced protocol implies security of the physical one, and every constraint a step passes on holds for every state an attacker can reach. The run lists this premise wherever it cannot point to a proof: it is not listed on the idealised Band-A route, whose one step is checked to be the identity, nor where every measuring device's single reduction step was built in this run by the engine's reductions module with status exact, with no source step and outcomes kept as they are, where that module's theorems prove it. A chain or channel from elsewhere must come with the proof.}
\claimmeta{\metalabel{Uses}\mref{Def.}{def:chaincond}, \extref{../transmitter/contract.pdf}{transmitter/def:source-segment}
\metasep \metalabel{Category}supplied artefact
\metasep \metalabel{Affects}secrecy
\metasep \metalabel{Assumed by}\mref{Thm.}{thm:choice-encoding}, \mref{Thm.}{thm:chain-consumption}, \extref{../meat-core/contract.pdf}{meat-core/thm:keylength}
\metasep \metalabel{Machine key}\iface{reduction_chain_sound}
\metasep newly derived
\metasep accepted}

% claim protocol-compiler/ass:atk
\claimsummary{protocol-compiler/ass:atk}{For a licensed non-empty constraint list that backs a key-length claim, its producer's reduction represents every physical attack with finite registers and every signal at its first channel by an admissible attack sequence. It lets constraints restrict the attacks considered.}
\claimbookmark{protocol-compiler/ass:atk}{Assumption}{ass:atk}{Attack-set licence}
\claimleadin{%
The licence is the attack-set licence of the module that produces the chain,
and it is that module's obligation. The published template for such a
certificate is \cite[Lemma~9.3, Eqs.~(146)--(147)]{tupkary2026decoy}
(squashing maps: the squashed POVM pulls back to the physical one and every
composed physical attack lands in the restricted set), applied in the order of
\cite[Rem.~9.2]{tupkary2026decoy} (source maps before squashing
restrictions). No part of the condition is specific to \exttermref{../transmitter/contract.pdf}{term-trusted-source-mode}{trusted-source mode}:
the licensed class and the admissibility format follow the declaration's
source mode.
}
\SetShortId{A-ATK}
\begin{assumption}[Attack-set licence]\label{ass:atk}
For a non-empty constraint list $\CList$ that is
\termref{term-licensed-for-the-declared-source-mode}{licensed for the declared source mode} (Definition~\ref{def:homog})
and backs a key-length claim:
(i) the reduction chain that produced the declaration certifies the
\emph{licensed reduction}: security of the physical protocol reduces to
security of the consuming engine's generated \exttermref{../foundations/contract.pdf}{term-round-model}{round model}
(\xref{meat-core/def:roundmodel}) with \exttermref{../foundations/contract.pdf}{term-attack-sequence}{attack
sequences} restricted to the $\CList$-admissible ones, that is, every attack
available against the physical protocol whose registers are finite-dimensional
and whose first channel receives every \exttermref{../transmitter/contract.pdf}{term-transmitted-system}{transmitted system} is
represented, after the reduction, by a sequence whose per-round channel $A_j$ lies, for every round $j$, in the
restricted set $\AttSet_{\CList}(E_{j-1},\Qsys E'_j)$ of that definition, with
no $A$ output, in trusted-source mode, and in its untrusted-source-mode
counterpart, with $A$ supplied by the adversary, in untrusted-source mode.
\end{assumption}
\claimverification{The engine does not verify this assumption; a run that rests on it records it as unverified.}
\claimreading{Applies only where constraints restrict the attacks the analysis considers, each in the form your declared source mode admits, which the engine checks. Whoever produced the reduction chain has proved that every attack on your physical system becomes, after the reduction, an attack that respects every listed constraint in every round. The proof is asked of attacks with a finite-dimensional memory, of any size, that receive every signal at the start; the engine extends the security statement to every other attack by a limit (\xref{meat-core/cor:composed-security}). For constraints the reductions module produces and certifies, that proof is its licence theorem. A constraint you add yourself, or one supplied without such a certificate, needs that proof from you, or must be removed; the analysis without constraints stays valid.}
\claimmeta{\metalabel{Uses}\mref{Def.}{def:homog}
\metasep \metalabel{Category}supplied artefact
\metasep \metalabel{Affects}secrecy
\metasep \metalabel{Assumed by}\mref{Thm.}{thm:gates}, \extref{../meat-core/contract.pdf}{meat-core/cor:constrained-key-length}, \extref{../meat-core/contract.pdf}{meat-core/ex:fss}, \extref{../meat-core/contract.pdf}{meat-core/lem:range}, \extref{../meat-core/contract.pdf}{meat-core/rem:cgate}, \extref{../meat-core/contract.pdf}{meat-core/rem:no-constraints}, \extref{../meat-core/contract.pdf}{meat-core/rem:not-licensed}, \extref{../meat-core/contract.pdf}{meat-core/thm:constrained-accumulation}, \extref{../meat-core/contract.pdf}{meat-core/thm:keylength}
\metasep \metalabel{Machine key}\iface{attack_set_license}
\metasep newly derived
\metasep accepted}

% claim protocol-compiler/ass:ideal
\claimsummary{protocol-compiler/ass:ideal}{On the Band-A route the declared receiver is an exactly projective qubit measurement with unit efficiency and no loss, dark counts, multi-photon events or misalignment, on a transmitted system that is exactly a qubit. The route is an explicit idealization, not a squashing argument.}
\claimbookmark{protocol-compiler/ass:ideal}{Assumption}{ass:ideal}{Ideal qubit receiver on the Band-A route}
\claimleadin{%
Used to reproduce known comparison values, a Band-A declaration is sound,
and its idealization is visible in this assumption. It is an idealization,
not a reduction: no squashing argument is made and none is needed, because
nothing infinite-dimensional is being reduced.
}
\SetShortId{A-IDEAL}
\begin{assumption}[Ideal qubit receiver on the Band-A route]\label{ass:ideal}
In a \termref{term-band-a}{Band A} specification:
(i) the declared receiver is exactly an ideal projective qubit measurement;
(ii) it has unit detection efficiency, no dark counts, no loss, no
multi-photon events and no misalignment;
(iii) the \exttermref{../transmitter/contract.pdf}{term-transmitted-system}{transmitted system} is exactly a qubit.
\end{assumption}
\claimverification{The engine does not verify this assumption; a run that rests on it records it as unverified.}
\claimreading{Only for an idealised Band-A analysis: your receiver is a perfect qubit measurement, with unit detection efficiency, no loss, no dark counts, no multi-photon events and no misalignment, and what reaches it in each round is exactly one qubit. The engine refuses a Band-A declaration whose own numbers contradict this: more than one measuring station, a measured system (or, with a trusted source, a transmitted system) of dimension other than two, or a measurement not made of projectors (per basis, where a basis choice is declared). No real optical receiver meets this; use the route to reproduce textbook comparison values. Declare a real device through the receiver module and a squashing reduction.}
\claimmeta{\metalabel{Category}receiver
\metasep \metalabel{Affects}secrecy
\metasep \metalabel{Assumed by}\mref{Thm.}{thm:band-a}
\metasep \metalabel{Machine key}\iface{ideal_qubit_receiver}
\metasep newly derived
\metasep accepted}

\clearpage
\part{Engine requirements}
\sectionintro{What an implementation of this module must do, record or refuse; each requirement names the result that motivates it.}

\section{Requirements: the emitted declaration}\label{group-protocol-compiler/req:declaration-keys}
\sectionintro{The engine emits the declaration under its machine keys, with the \termref{term-index-map}{index map} beside it, and holds every operator on the measured system as its block list. It passes the declared data it reads through unchanged, forms none of the derived objects of the general format as data, and decides every predicate per \termref{term-measured-factor}{measured factor} and per Alice sector on declared blocks. It reads each constraint tag off its datum or off the field it arrives with, infers none, and decides nothing at a window. With the declaration it passes on the transmitter's source-replacement licence record where one is demanded, and records the choice weighting beside the chain record. At a one-dimensional A in trusted-source mode it refuses a specification whose source emission depends on the letter or whose letter weights differ from Alice's measurement.}
\sectionimports{\extref{../transmitter/contract.pdf}{transmitter/def:source-segment}, \extref{../transmitter/contract.pdf}{transmitter/def:srcert}}

% claim protocol-compiler/req:declaration-keys
\claimsummary{protocol-compiler/req:declaration-keys}{The sector map is emitted as three keys, its identifiers, their dimensions and its block kind; the index map goes beside the declaration and each constraint entry's tag with the entry. The sifting predicate and the key-live set are derived from the specification and are not emitted.}
\claimbookmark{protocol-compiler/req:declaration-keys}{Requirement}{req:declaration-keys}{Machine keys of the emitted declaration}
\claimleadin{%
Definition~\ref{def:declaration} makes the \exttermref{../foundations/contract.pdf}{term-sector-map}{sector map} one object and
derives the sifting predicate and the key-live set from the specification;
this requirement says which of these the emitted declaration carries.
}
\begin{requirement}[Machine keys of the emitted declaration]\label{req:declaration-keys}
The engine emits the \exttermref{../foundations/contract.pdf}{term-sector-map}{sector map} $\Smap$ of a
\termref{term-protocol-declaration}{protocol declaration} as three machine keys,
\iface{sector_index_set} for its ordered block identifiers,
\iface{sector_dims} for their dimensions and \iface{sector_block_kind} for
its block kind, so that a consumer comparing sector maps can compare all
three and not two of them. It emits the label-to-index map $\idxm$ of
Definition~\ref{def:compile} alongside the declaration, under
\iface{label_index_map}, and with each constraint entry the tag of its
right-hand side, under \iface{constraint_rhs_tag}, so that the consuming
engine's type check has a datum (Definition~\ref{def:homog}). The sifting
predicate $\varsigma$ and the set $\Chkey$ of \termref{term-key-live}{key-live} values are not
emitted.
\end{requirement}
\claimmeta{\metalabel{Motivated by}\mref{Def.}{def:declaration}, \mref{Def.}{def:homog}, \mref{Def.}{def:compile}
\metasep accepted
\metasep \metalabel{Code}\coderef{https://github.com/fraunhofersg/certiqs-key/blob/v0.1.0/src/certiqs_key/protocol_compiler/declaration.py}{certiqs_key.protocol_compiler.declaration:CandidateDeclaration}}

% claim protocol-compiler/req:operator-form
\claimsummary{protocol-compiler/req:operator-form}{Every operator on the measured system is held and emitted as its block list. Where there are several sectors its direct sum is formed only as the argument of chain condition (C4), decided there and discarded.}
\claimbookmark{protocol-compiler/req:operator-form}{Requirement}{req:operator-form}{Operators on the measured system held as block lists}
\claimleadin{%
By Definition~\ref{def:declaration} an operator on the \exttermref{../reductions/contract.pdf}{term-measured-system}{measured system} is its
\exttermref{../foundations/contract.pdf}{term-block-list}{block list}, so sector block-diagonality holds by type and the declaration
holds no array larger than one sector block.
}
\begin{requirement}[Operators on the measured system held as block lists]\label{req:operator-form}
The engine holds and emits every operator on $\Qsys$ as its
\exttermref{../foundations/contract.pdf}{term-block-list}{block list} against the sector map $\Smap$
(Definition~\ref{def:declaration}). Where $|\SecSet|>1$ it does not
materialize the direct sum $\bigoplus_sF^{(s)}$ of such a list, with one
exception: condition (C4) of Definition~\ref{def:chaincond} forms it as the
argument of that one decision, decides it and discards it. The convention
\iface{operator_form} states this rule and its scope.
\end{requirement}
\claimmeta{\metalabel{Motivated by}\mref{Def.}{def:declaration}, \mref{Def.}{def:chaincond}
\metasep accepted}

% claim protocol-compiler/req:pass-through
\claimsummary{protocol-compiler/req:pass-through}{The engine passes the declared data it reads into a declaration through unchanged: the marginal slot, Alice's measurement slot in its declared form, the transcript, and each entry's scope and right-hand side. It fills no slot, constructs no right-hand side and converts no form.}
\claimbookmark{protocol-compiler/req:pass-through}{Requirement}{req:pass-through}{Declared data passed through unchanged}
\claimleadin{%
The \termref{term-compilation-map}{compilation map} of Definition~\ref{def:compile} reads
\exttermref{../foundations/contract.pdf}{term-declared-data}{declared data} from the transmitter artefact, the chain and the
specification, and constructs only the raw-key map, the \termref{term-index-map}{index map} and
the weighted families. Data passed on unchanged keep their \exttermref{../foundations/contract.pdf}{term-declarer}{declarer}
(\xref{foundations/def:declared-data}\,(iii)), and a consumer repairs nothing
it was handed (\xref{foundations/req:consumer-decides-exactly}).
}
\begin{requirement}[Declared data passed through unchanged]\label{req:pass-through}
The engine passes through unchanged the declared data that the compilation map
reads into a declaration: the transmitter artefact's \termref{term-marginal-slot}{marginal slot}, the
value $[1]$ the transmitter filled it with at $d_A=1$ included; Alice's
measurement slot, in the form the artefact declares, which it changes only by
weighting it in that form (Definition~\ref{def:weighted}) where a choice
family is presented; the \termref{term-transcript-protocol}{transcript}; and each
constraint entry's right-hand-side datum and, where the entry carries one, its
\termref{term-scope}{scope}. It fills no slot and demands no certificate of its own.
It constructs no right-hand side of an emitted constraint entry: it forms no
$c_k\tau_A$ for one and reads no $c_k$ off any matrix. It converts neither
form of Alice's measurement slot into the other: it forms neither
\eqref{eq:ensemblepovm} from an ensemble slot nor
$\tau_A^{1/2}M_x\tau_A^{1/2}$ from a POVM slot.
\end{requirement}
\claimmeta{\metalabel{Motivated by}\mref{Def.}{def:declaration}, \mref{Def.}{def:compile}, \mref{Def.}{def:weighted}, \extref{../foundations/contract.pdf}{foundations/def:declared-data}, \extref{../foundations/contract.pdf}{foundations/req:consumer-decides-exactly}
\metasep accepted
\metasep \metalabel{Code}\coderef{https://github.com/fraunhofersg/certiqs-key/blob/v0.1.0/src/certiqs_key/protocol_compiler/compile.py}{certiqs_key.protocol_compiler.compile:compile_specification}}

% claim protocol-compiler/req:derived-objects
\claimsummary{protocol-compiler/req:derived-objects}{The product family, the lifted operators and the family an ensemble form denotes are rules of the general format and never data: nothing forms, rounds or decides them. Every exact predicate is decided per measured factor and per Alice sector on declared blocks.}
\claimbookmark{protocol-compiler/req:derived-objects}{Requirement}{req:derived-objects}{Per-factor decisions on the declared blocks}
\claimleadin{%
In the general format of Definition~\ref{def:declaration} the product
family, the \termref{term-lifted-operator}{lifted operators} and the POVM an
\termref{term-ensemble-form}{ensemble form} denotes are defined by rules from declared data; what is
declared, and decided, is the data the rules read.
}
\begin{requirement}[Per-factor decisions on the declared blocks]\label{req:derived-objects}
For a declaration in the general format of
Definition~\ref{def:declaration}: no entry of a product family
\eqref{eq:productfamily} is formed as a datum, rounded or decided on, here or
downstream; the \termref{term-lifted-operator}{lifted operator} of a scoped entry is never formed as a
datum; and the family \eqref{eq:ensemblepovm} that an \termref{term-ensemble-form}{ensemble form}
denotes is formed by no one and decided by nothing. Every exact predicate is
decided per \termref{term-measured-factor}{measured factor}, on that factor's blocks and in that
factor's \exttermref{../foundations/contract.pdf}{term-declared-order}{declared order}, and on Alice's data per Alice sector
(Definition~\ref{def:exact}). The engine refuses by type a presentation of a
\termref{term-choice-weighted-family}{choice-weighted family} indexed by a product alphabet
(Definition~\ref{def:weighted}).
\end{requirement}
\claimmeta{\metalabel{Motivated by}\mref{Def.}{def:declaration}, \mref{Def.}{def:exact}, \mref{Def.}{def:weighted}
\metasep accepted}

% claim protocol-compiler/req:tag-reading
\claimsummary{protocol-compiler/req:tag-reading}{The tag of an untagged entry is read off the type of its datum, and an arriving tag field is checked against that type; the party-local tag is read off the field it arrives with. No tag is inferred and none is taken on the upstream module's word.}
\claimbookmark{protocol-compiler/req:tag-reading}{Requirement}{req:tag-reading}{Reading the tag of a constraint entry}
\claimleadin{%
For the two untagged forms of Definition~\ref{def:homog} the tag is the type
of the right-hand side; the \exttermref{../reductions/contract.pdf}{term-party-local-form}{party-local form} is the one form the type does
not separate from another.
}
\begin{requirement}[Reading the tag of a constraint entry]\label{req:tag-reading}
The engine reads the tag of an entry of one of the two untagged forms of
Definition~\ref{def:homog} off the type of its right-hand-side datum. Where
such an entry arrives carrying a tag field, check V6 of
Definition~\ref{def:validator} compares that field with the type of the
datum it names and refuses a disagreement; the tag is never taken on the
upstream module's word against the object the entry holds. An entry is read
as \termref{term-party-local}{party-local} exactly when it arrives carrying that tag; check V6
then requires its right-hand-side slot to hold a real scalar and its scope to
name a measured factor of the emitted sector map, and refuses otherwise. No
tag is inferred: an entry arriving without the party-local tag is never read
party-locally, whatever module emitted it.
\end{requirement}
\claimmeta{\metalabel{Motivated by}\mref{Def.}{def:homog}, \mref{Def.}{def:validator}
\metasep implemented (effective: accepted)
\metasep \metalabel{Code}\coderef{https://github.com/fraunhofersg/certiqs-key/blob/v0.1.0/src/certiqs_key/protocol_compiler/validation.py}{certiqs_key.protocol_compiler.validation:check_v6}}

% claim protocol-compiler/req:no-window
\claimsummary{protocol-compiler/req:no-window}{The engine declares no numerical tolerance, so no check of the validator reads a window and nothing in this contract is decided at one. Licensing and the admission predicates are decided exactly.}
\claimbookmark{protocol-compiler/req:no-window}{Requirement}{req:no-window}{Exact decisions without a tolerance}
\claimleadin{%
A consumer decides every admission predicate on declared data exactly and
reads no tolerance (\xref{foundations/req:consumer-decides-exactly},
\xref{foundations/req:tolerance-scope}).
}
\begin{requirement}[Exact decisions without a tolerance]\label{req:no-window}
The engine declares no numerical tolerance, and there is no window in this
contract at which anything could be decided: no check of the validator
(Definition~\ref{def:validator}) reads a window or is decided at one, and
licensing (Definition~\ref{def:homog}) and the predicates of
Definition~\ref{def:exact} are decided exactly.
\end{requirement}
\claimmeta{\metalabel{Motivated by}\mref{Def.}{def:exact}, \mref{Def.}{def:homog}, \mref{Def.}{def:validator}, \extref{../foundations/contract.pdf}{foundations/req:consumer-decides-exactly}, \extref{../foundations/contract.pdf}{foundations/req:tolerance-scope}
\metasep implemented (effective: accepted)
\metasep \metalabel{Code}\coderef{https://github.com/fraunhofersg/certiqs-key/blob/v0.1.0/src/certiqs_key/protocol_compiler/validation.py}{certiqs_key.protocol_compiler.validation:validate}}

% claim protocol-compiler/req:source-licence-passed-on
\claimsummary{protocol-compiler/req:source-licence-passed-on}{A declaration carries the transmitter's source-replacement licence record, re-decided by gate (G2), to the consuming engine exactly when it is emitted in trusted-source mode with Alice's dimension at least two. Its certificate and attestation are the ones supplied.}
\claimbookmark{protocol-compiler/req:source-licence-passed-on}{Requirement}{req:source-licence-passed-on}{The source-replacement licence passed on with the declaration}
\claimleadin{%
The record gate (G2) consumes is what licenses a declared marginal at
$d_A\ge2$, and the consuming engine receives it with the declaration rather
than from the transmitter.
}
\begin{requirement}[The source-replacement licence passed on with the declaration]\label{req:source-licence-passed-on}
The engine passes the transmitter's source-replacement \exttermref{../reductions/contract.pdf}{term-licence-record}{licence record}
\iface{source_replacement_license} (\xref{transmitter/def:srcert}) on to the
consuming engine with a declaration exactly when it emits that declaration in
\exttermref{../transmitter/contract.pdf}{term-trusted-source-mode}{trusted-source mode} at $d_A\ge2$, after gate (G2) of
Definition~\ref{def:eligible} has re-decided its conditions on the emitted
data; the certificate and the attestation (SR3) it carries are the ones
supplied, and the engine asserts no attestation of its own
(\xref{transmitter/req:attestation-record}). At $d_A=1$ the record is neither
demanded nor emitted (\xref{transmitter/req:licence-scope}).
\end{requirement}
\claimmeta{\metalabel{Motivated by}\mref{Def.}{def:eligible}, \extref{../transmitter/contract.pdf}{transmitter/def:srcert}, \extref{../transmitter/contract.pdf}{transmitter/req:attestation-record}, \extref{../transmitter/contract.pdf}{transmitter/req:licence-scope}
\metasep accepted
\metasep \metalabel{Code}\coderef{https://github.com/fraunhofersg/certiqs-key/blob/v0.1.0/src/certiqs_key/protocol_compiler/compile.py}{certiqs_key.protocol_compiler.compile:compile_specification}}

% claim protocol-compiler/req:da1-source-identity
\claimsummary{protocol-compiler/req:da1-source-identity}{In trusted-source mode at d\_A = 1 each weighted member of the source's ensemble must equal the letter's measurement weight times the summed ensemble, decided exactly, or the specification is refused. Alice's letter weights then equal the source's, and no emitted state depends on the letter.}
\claimbookmark{protocol-compiler/req:da1-source-identity}{Requirement}{req:da1-source-identity}{Letter-independent emission at a one-dimensional trusted A}
\claimleadin{%
At $d_A=1$ no \exttermref{../transmitter/contract.pdf}{term-source-replacement-certificate}{source-replacement certificate} is demanded
(\xref{transmitter/rem:da1}), and the consuming \exttermref{../foundations/contract.pdf}{term-round-model}{round model} takes Alice's
system to be trivial (\xref{meat-core/prop:trivial-alice-system},
\xref{meat-core/def:roundmodel}): her letter is the outcome of $\{[m_x]\}$ on
a one-dimensional register, independent of every system the adversary holds.
In \exttermref{../transmitter/contract.pdf}{term-trusted-source-mode}{trusted-source mode} the source premise
(\xref{transmitter/ass:ens}, or the premise a \exttermref{../transmitter/contract.pdf}{term-source-model}{source model} declares in
its place) has the source emit $\rho_x$ given the letter, and those systems
reach the adversary. The two descriptions agree only where the emission does
not depend on the letter and the letter weights agree, which is condition
(SR2) of \xref{transmitter/def:srcert} read at $d_A=1$: with $A$ trivial,
$\sigma_x=p_x\rvirt_x$ and $\Psi=\sum_y\sigma_y$, the condition
$\tr_A[(M_x\otimes\idon{\Avirt})\Psi]=\sigma_x$ reads $m_x\Psi=\sigma_x$.
Taking the trace gives $m_x=p_x$, the identity on which
Assumption~\ref{ass:alice-meas} and the source premise agree on the letter
distribution; on a letter with $p_x>0$ it gives $\rvirt_x=\sum_yp_y\rvirt_y$,
one common state. Both sides are declared numbers, so the engine decides the
identity. In \exttermref{../transmitter/contract.pdf}{term-untrusted-source-mode}{untrusted-source mode} no source premise is emitted and
there is nothing to compare.
}
\begin{requirement}[Letter-independent emission at a one-dimensional trusted A]\label{req:da1-source-identity}
For a specification in \exttermref{../transmitter/contract.pdf}{term-trusted-source-mode}{trusted-source mode} at $d_A=1$, the engine
decides, for every $x\in\Xset$ in the \exttermref{../foundations/contract.pdf}{term-declared-order}{declared order} and exactly on the
declared numbers, the identity
\begin{equation*}
  \sigma_x \;=\; m_x\sum_{y\in\Xset}\sigma_y ,
\end{equation*}
where $m_x$ is the one entry of Alice's measurement element $M_x$, read as the
declaration carries it (Definition~\ref{def:compile}: for a
\termref{term-choice-weighted-family}{choice-weighted family} the weighted element the engine forms; in the
\termref{term-ensemble-form}{ensemble form}, $N_x$, which equals $M_x$ at $d_A=1$ by
\eqref{eq:ensemblepovm}), and $\sigma_x=p_x\rvirt_x$ is the weighted member
for $x$ of the \exttermref{../transmitter/contract.pdf}{term-virtual-ensemble}{virtual ensemble} the transmitter artefact emits, which
at an empty \exttermref{../transmitter/contract.pdf}{term-source-segment}{source segment} is the specified one
(\xref{transmitter/def:source-segment}); where the artefact emits the
ensemble by its \exttermref{../transmitter/contract.pdf}{term-weighted-members}{weighted members}, $\sigma_x$ is the declared member. Products
and sums are formed in exact arithmetic, entry by entry. A specification for
which the identity fails at some letter is refused, before a declaration is
emitted, with \iface{da1_source_identity_failed}, naming the first such letter
in the declared order and the first entry at which the two sides differ.
\end{requirement}
\claimreading{With a trusted source and a one-dimensional retained system A, Alice's letter is local randomness that the analysis treats as independent of everything the source sends out. The engine therefore requires two things of the declaration and refuses it otherwise: Alice's letter probabilities equal the source's, number for number, and every letter of positive probability is sent as the same state. A source whose emitted state depends on the letter is declared with a retained system A of dimension two or more and a source-replacement certificate.}
\claimmeta{\metalabel{Motivated by}\mref{Def.}{def:declaration}, \mref{Def.}{def:spec}, \mref{Def.}{def:compile}, \mref{Ass.}{ass:alice-meas}, \extref{../meat-core/contract.pdf}{meat-core/def:roundmodel}, \extref{../meat-core/contract.pdf}{meat-core/prop:trivial-alice-system}, \extref{../transmitter/contract.pdf}{transmitter/ass:ens}, \extref{../transmitter/contract.pdf}{transmitter/def:source-segment}, \extref{../transmitter/contract.pdf}{transmitter/def:srcert}, \extref{../transmitter/contract.pdf}{transmitter/rem:da1}
\metasep accepted
\metasep \metalabel{Code}\coderef{https://github.com/fraunhofersg/certiqs-key/blob/v0.1.0/src/certiqs_key/protocol_compiler/declared_premises.py}{certiqs_key.protocol_compiler.declared_premises:decide_da1_source_identity}}

% claim protocol-compiler/req:choice-weighting-record
\claimsummary{protocol-compiler/req:choice-weighting-record}{Both forms of the choice-weighted family are formed, and emitted wherever they are the emitted family. The setting-dependent weighting adds no link and no register to the chain record, which stays unchanged, and is recorded beside it, one entry per weighted factor.}
\claimbookmark{protocol-compiler/req:choice-weighting-record}{Requirement}{req:choice-weighting-record}{Emission of the weighted families and the weighting record}
\claimleadin{%
A relabelling composes no channel and so has no place in a chain
(Remark~\ref{rem:relabel}); what the compiler applied is recorded beside the
chain instead.
}
\begin{requirement}[Emission of the weighted families and the weighting record]\label{req:choice-weighting-record}
The engine emits both forms of the \termref{term-choice-weighted-family}{choice-weighted family} of
Definition~\ref{def:weighted}, each where it is the emitted family: the
\termref{term-setting-independent-form}{setting-independent form} \eqref{eq:weighted} on either party's side,
and the \termref{term-setting-dependent-form}{setting-dependent form} \eqref{eq:weighted-t} at a measured
system, of whichever party, after that system's \exttermref{../reductions/contract.pdf}{term-receiver-segment}{receiver segment}. A
setting-independent presentation of a measuring party's physical receiver
family is the emitted family only where that party's receiver segment
delivers it as its \exttermref{../reductions/contract.pdf}{term-target-povm}{target POVM}; in front of any other
segment it is handed to the chain as the physical receiver family of
Definition~\ref{def:spec}\,(b), and the emitted family is the chain's. For the
setting-dependent relabelling it emits no link and no register
for the choice value, and the chain record it consumes and re-emits is
unchanged by the weighting; the weightings are recorded beside the chain in
the emitted artefact, in \iface{choice_weighting_record}, one entry per
\termref{term-measured-factor}{measured factor} at which a weighting is applied, each entry naming
that factor; like the setting-independent choice-weighted
family, they are this module's own construction.
\end{requirement}
\claimmeta{\metalabel{Motivated by}\mref{Def.}{def:spec}, \mref{Def.}{def:weighted}, \mref{Rem.}{rem:relabel}
\metasep accepted}

\section{Requirements: choice weighting}\label{group-protocol-compiler/req:conditional-family-decision}
\sectionintro{The engine decides each declared conditional family, per choice value, before it forms any product, and rejects a weight that is not a finite binary64 number as malformed input; decimal or rational user weights may be prepared as binary64 before the module boundary. It refuses a setting-dependent weighting read through anything but the station's setting, decides the weight list and the stochastic matrix exactly, and emits a weighted family only after deciding every product representable and the emitted family exactly admissible. It declares no lattice exponent, snaps nothing, and never emits a fixed fraction outside the measurement operators.}

% claim protocol-compiler/req:conditional-family-decision
\claimsummary{protocol-compiler/req:conditional-family-decision}{Each conditional family of a choice-weighted presentation is decided exactly admissible, per choice value, before any product is formed or any chain condition decided. A failure is refused with its own reason, naming the first failure.}
\claimbookmark{protocol-compiler/req:conditional-family-decision}{Requirement}{req:conditional-family-decision}{Decision on each declared conditional family}
\claimleadin{%
A \termref{term-choice-weighted-family}{choice-weighted family} can satisfy every condition of exact
admissibility while none of the conditional families it was built from is a
measurement (Example~\ref{ex:condfamily}), so each conditional family is
decided on its own, on the declared input (Remark~\ref{rem:condfamily}).
}
\begin{requirement}[Decision on each declared conditional family]\label{req:conditional-family-decision}
Wherever a choice-weighted presentation is read, on either party's side, the
engine decides, of each choice value $t\in\Tpub$ separately and on the
declared numbers, the condition Definition~\ref{def:weighted} places on the
conditional family $\{N_{u|t}\}_u$ --- in \termref{term-povm-form}{POVM form}, its exact
admissibility in the reading of that definition; in \termref{term-ensemble-form}{ensemble form}, the
condition that definition states in its place: (H) and (P) per block, the
ensemble identity in place of (C), and (T) --- before it forms any product and
before any condition of Definition~\ref{def:chaincond} is decided
on the chain: wherever the presented physical family is the weighted one,
that decision and the representability decision on the products it permits
precede the consumption of the chain, so that
no condition on the chain is decided on a product that was not decided
first. A failure is refused with
\iface{conditional_family_not_admissible}, naming the declaring party, and
the choice value, the predicate, the block and the index of the \emph{first}
failure --- for (S), the index of the first inexact addition; in ensemble
form the predicate named is the ensemble identity where that is the one that
failed. The decision stops at the first failure and does not enumerate the
rest.
\end{requirement}
\claimmeta{\metalabel{Motivated by}\mref{Def.}{def:compile}, \mref{Def.}{def:weighted}, \mref{Rem.}{rem:condfamily}, \mref{Ex.}{ex:condfamily}, \mref{Def.}{def:chaincond}
\metasep implemented (effective: accepted)
\metasep \metalabel{Code}\coderef{https://github.com/fraunhofersg/certiqs-key/blob/v0.1.0/src/certiqs_key/protocol_compiler/weights.py}{certiqs_key.protocol_compiler.weights:decide_conditional_families}}

% claim protocol-compiler/req:binary64-weights
\claimsummary{protocol-compiler/req:binary64-weights}{A declared weight in either form that is not exactly a finite binary64 number is rejected as malformed input before it enters any sum or product. This is the declared input type, not a failed predicate, and it has no refusal string of its own.}
\claimbookmark{protocol-compiler/req:binary64-weights}{Requirement}{req:binary64-weights}{Declared weights as binary64 input}
\claimleadin{%
The conditions of Definition~\ref{def:weighted} are decided on declared
numbers, so the declared weights have to be numbers of the type they are
decided on.
}
\begin{requirement}[Declared weights as binary64 input]\label{req:binary64-weights}
The weights of either form of Definition~\ref{def:weighted} --- each $q_t$ of
the \termref{term-setting-independent-form}{setting-independent form} and each entry of $P$ of the
\termref{term-setting-dependent-form}{setting-dependent form} --- are \exttermref{../foundations/contract.pdf}{term-binary64-numbers}{binary64 numbers}. The engine
rejects as malformed input a declared weight, in either form, that is not
exactly representable as a finite binary64 number, before that weight is used
in an accumulation or a product. This is the declared input type: the
rejection is not a failed summation predicate, and it adds no row of the gate
report and no refusal string.
\end{requirement}
\claimmeta{\metalabel{Motivated by}\mref{Def.}{def:weighted}
\metasep implemented (effective: accepted)
\metasep \metalabel{Code}\coderef{https://github.com/fraunhofersg/certiqs-key/blob/v0.1.0/src/certiqs_key/protocol_compiler/weights.py}{certiqs_key.protocol_compiler.weights:weighted_family}}

% claim protocol-compiler/req:weight-preparation
\claimsummary{protocol-compiler/req:weight-preparation}{Decimal and rational user weights may be converted to binary64 by the declaring entry point before they reach this module, with the moves recorded and normalisation and bound directions preserved. The preparation asserts no bound on the resulting change of key rate.}
\claimbookmark{protocol-compiler/req:weight-preparation}{Requirement}{req:weight-preparation}{Preparation of user weights before the module boundary}
\claimleadin{%
This module reads only binary64 weights
(Requirement~\ref{req:binary64-weights}); ordinary user inputs reach that
type through a preparation that belongs to the declarer
(\xref{foundations/req:who-snaps}).
}
\begin{requirement}[Preparation of user weights before the module boundary]\label{req:weight-preparation}
Ordinary decimal and rational user inputs for the weights of
Definition~\ref{def:weighted} may be converted and prepared by the declaring
entry point before they cross the \exttermref{../foundations/contract.pdf}{term-module-boundary}{module boundary} into this module. An
entry point that does so records the supplied and the prepared values and
their displacement,
preserves normalisation and every required bound direction, and leaves this
module's exact tests unchanged. The preparation asserts no bound on a
resulting change of the key rate.
\end{requirement}
\claimmeta{\metalabel{Motivated by}\mref{Def.}{def:weighted}, \mref{Req.}{req:binary64-weights}, \extref{../foundations/contract.pdf}{foundations/req:who-snaps}
\metasep implemented (effective: accepted)
\metasep \metalabel{Code}\coderef{https://github.com/fraunhofersg/certiqs-key/blob/v0.1.0/src/certiqs_key/runs/inputs.py}{certiqs_key.runs.inputs:prepare_probabilities}}

% claim protocol-compiler/req:setting-only-weighting
\claimsummary{protocol-compiler/req:setting-only-weighting}{A setting-dependent weighting presented through anything but the station's setting, such as the outcome alphabet or the device outcome, is refused by type. A value whose weight varied with the device outcome would be inferred from it and would not be a choice.}
\claimbookmark{protocol-compiler/req:setting-only-weighting}{Requirement}{req:setting-only-weighting}{Setting-dependent weighting read through the setting only}
\claimleadin{%
A value $f$ whose weight varied with the device outcome would be inferred from
that outcome and would not be a choice (Assumption~\ref{ass:coin},
Remark~\ref{rem:relabel}); so the weighting reaches an outcome only through
its \termref{term-setting-component}{setting component}, as Theorem~\ref{thm:choice-encoding}\,(viii) records.
}
\begin{requirement}[Setting-dependent weighting read through the setting only]\label{req:setting-only-weighting}
The engine refuses by type, with \iface{choice_weight_depends_on_outcome}, a
weighting of the \termref{term-setting-dependent-form}{setting-dependent form} of
Definition~\ref{def:weighted} that is presented indexed by the \exttermref{../reductions/contract.pdf}{term-outcome-alphabet}{outcome
alphabet} $\Vset$, by the device outcome, or by any component of the outcome
other than its setting; in the general format, one presented indexed by a
\termref{term-measured-factor}{measured factor}'s whole outcome, by another factor's outcome, or by any
component other than that factor's setting.
\end{requirement}
\claimmeta{\metalabel{Motivated by}\mref{Def.}{def:weighted}, \mref{Thm.}{thm:choice-encoding}, \mref{Rem.}{rem:relabel}, \mref{Ass.}{ass:coin}
\metasep implemented (effective: accepted)
\metasep \metalabel{Code}\coderef{https://github.com/fraunhofersg/certiqs-key/blob/v0.1.0/src/certiqs_key/protocol_compiler/weights.py}{certiqs_key.protocol_compiler.weights:refuse_outcome_dependent_weighting}}

% claim protocol-compiler/req:weight-list
\claimsummary{protocol-compiler/req:weight-list}{The declared weights of the setting-independent form are decided exactly: each positive by its sign, and their sum equal to one in the declared order without rounding. A weight that is not positive is refused by type, and a list whose sum is not one or rounds is refused with one named reason.}
\claimbookmark{protocol-compiler/req:weight-list}{Requirement}{req:weight-list}{Decision on the declared weight list}
\claimleadin{%
The weights of the \termref{term-setting-independent-form}{setting-independent form} are declared scalars, and
Definition~\ref{def:weighted} reads each of its conditions on them exactly.
}
\begin{requirement}[Decision on the declared weight list]\label{req:weight-list}
The engine decides the conditions of Definition~\ref{def:weighted} on the
declared weights $(q_t)_{t\in\Tpub}$ of a choice-weighted presentation: each
$q_t>0$ by its sign, and $\sum_tq_t=1$ exactly, with the sum accumulated in
the \exttermref{../foundations/contract.pdf}{term-declared-order}{declared order} of $\Tpub$ --- the order in which the declared
public-choice record lists it (\xref{foundations/req:declared-order}) --- and
required to \exttermref{../foundations/contract.pdf}{term-rounds-nowhere}{round nowhere}. A declared weight that is
zero or negative is refused by type. A weight list whose exact sum is not $1$,
or whose declared-order accumulation rounds at some addition, is refused with
\iface{choice_weights_not_normalized}, and no declaration is emitted.
\end{requirement}
\claimmeta{\metalabel{Motivated by}\mref{Def.}{def:weighted}, \extref{../foundations/contract.pdf}{foundations/req:declared-order}
\metasep implemented (effective: accepted)
\metasep \metalabel{Code}\coderef{https://github.com/fraunhofersg/certiqs-key/blob/v0.1.0/src/certiqs_key/protocol_compiler/weights.py}{certiqs_key.protocol_compiler.weights:weighted_family}}

% claim protocol-compiler/req:stochastic-weights
\claimsummary{protocol-compiler/req:stochastic-weights}{The declared matrix of a setting-dependent weighting is decided exactly: every entry non-negative by its sign and every row sum equal to one in the declared order without rounding. A negative entry or a row not summing exactly to one is refused with its own reason.}
\claimbookmark{protocol-compiler/req:stochastic-weights}{Requirement}{req:stochastic-weights}{Stochastic matrix of a setting-dependent weighting}
\claimleadin{%
The setting-dependent form is built from a stochastic matrix $P$, and
Theorem~\ref{thm:choice-encoding}\,(viii) states that a $P$ that is not one is
refused; its proof gives the reason.
}
\begin{requirement}[Stochastic matrix of a setting-dependent weighting]\label{req:stochastic-weights}
The engine decides the conditions of Definition~\ref{def:weighted} on the
declared matrix $P$ of the \termref{term-setting-dependent-form}{setting-dependent form}: every entry
$P(f\,|\,t)\ge0$ by its sign, and every row sum $\sum_fP(f\,|\,t)=1$ exactly,
accumulated in the \exttermref{../foundations/contract.pdf}{term-declared-order}{declared order} of the \termref{term-choice-alphabet}{choice alphabet}
(\xref{foundations/req:declared-order}) and required to
\exttermref{../foundations/contract.pdf}{term-rounds-nowhere}{round nowhere}; the \termref{term-support-of-the-weighting}{support of the weighting} is
decided by the sign of each entry. A negative entry, or a row whose
declared-order sum is not exactly $1$, is refused with
\iface{choice_weight_not_stochastic}.
\end{requirement}
\claimmeta{\metalabel{Motivated by}\mref{Def.}{def:weighted}, \mref{Thm.}{thm:choice-encoding}, \extref{../foundations/contract.pdf}{foundations/req:declared-order}
\metasep implemented (effective: accepted)
\metasep \metalabel{Code}\coderef{https://github.com/fraunhofersg/certiqs-key/blob/v0.1.0/src/certiqs_key/protocol_compiler/weights.py}{certiqs_key.protocol_compiler.weights:setting_weighted_family}}

% claim protocol-compiler/req:weighted-family-emission
\claimsummary{protocol-compiler/req:weighted-family-emission}{Every product a weighting forms is decided binary64, emitted or handed on, and the emitted family exactly admissible, before anything is emitted. A failure refuses and emits nothing; an exact product that is binary64 but off the instance's lattice is emitted as it is.}
\claimbookmark{protocol-compiler/req:weighted-family-emission}{Requirement}{req:weighted-family-emission}{Emission decisions on the weighted products}
\claimleadin{%
A product of two binary64 numbers need not be one, and a family whose entries
are all binary64 may still have a declared-order sum that rounds; both are
decided before a weighted family leaves the compiler
(\xref{foundations/req:exact-emission}).
}
\begin{requirement}[Emission decisions on the weighted products]\label{req:weighted-family-emission}
The engine forms the products $q_tN_{u|t}$ and $P(f\,|\,t(v))F_v$ of
Definition~\ref{def:weighted} exactly, on the declared dyadic numbers, and
takes the two emission decisions of that definition before emitting anything.
The first, that every entry of every product it forms is a binary64 number,
is taken on every product the map forms, whether the product is emitted or
handed on: the intermediate families of an iterated presentation are
products, and a presentation on Bob's side before the chain produces the
physical family that (C2) and (C4) of Definition~\ref{def:chaincond} read and
that, in \termref{term-band-a}{Band A}, (C3) decides, so no condition of the chain is decided on a
product whose entries were not decided first. The second is that the emitted
family satisfies Definition~\ref{def:exact}. If either decision fails, the
engine refuses --- with \iface{emitted_value_not_representable} in the first
case and with the predicate that failed in the second --- and emits no
declaration. \iface{emitted_value_not_representable} covers exactly an exact
product of declared numbers that is not a binary64 number: an entry that is
binary64 but lies off whatever lattice the \exttermref{../foundations/contract.pdf}{term-instance}{instance} declared its own
data on is emitted, and the five conditions the consumer states are decided
on it exactly and on it alone.
\end{requirement}
\claimmeta{\metalabel{Motivated by}\mref{Def.}{def:exact}, \mref{Def.}{def:weighted}, \mref{Thm.}{thm:choice-encoding}, \mref{Def.}{def:chaincond}, \extref{../foundations/contract.pdf}{foundations/req:exact-emission}
\metasep implemented (effective: accepted)
\metasep \metalabel{Code}\coderef{https://github.com/fraunhofersg/certiqs-key/blob/v0.1.0/src/certiqs_key/protocol_compiler/weights.py}{certiqs_key.protocol_compiler.weights:decide_emission}}

% claim protocol-compiler/req:no-lattice-exponent
\claimsummary{protocol-compiler/req:no-lattice-exponent}{The compiler declares no lattice exponent for what it emits, so it never rounds or snaps its emission and never declares a nearby family: a failed emission condition refuses. It owes its consumer the five exact conditions, not a lattice.}
\claimbookmark{protocol-compiler/req:no-lattice-exponent}{Requirement}{req:no-lattice-exponent}{No lattice exponent and no snapping of the emission}
\claimleadin{%
\xref{foundations/req:exact-emission} binds every producer, and the
choice-weighted families are this module's products. Five of its items bind
every producer, and for these families they are the five conditions of
Definition~\ref{def:exact} --- the \exttermref{../foundations/contract.pdf}{term-sector-map}{sector map}, exact Hermiticity, exact
positive semidefiniteness, exact completeness per block, and condition~(S);
a sixth, (D-LAT) of \xref{foundations/req:declared-lattice}, binds only a
producer that declares a \exttermref{../foundations/contract.pdf}{term-lattice-exponent}{lattice exponent} $k\le53$. That requirement
makes a producer that cannot establish the conditions that bind it refuse,
naming the predicate; only a producer that declares $k$ may instead declare a
nearby family established under all six, carrying three disclosures, and
\xref{foundations/req:who-snaps} lets only such a producer snap its own
emission. (D-LAT) is a \emph{sufficient} condition for the five to be
establishable and is not among the conditions the consuming engine states,
so this module states no (D-LAT) obligation of its own.
}
\begin{requirement}[No lattice exponent and no snapping of the emission]\label{req:no-lattice-exponent}
The engine declares no lattice exponent for the families it emits. It
establishes the five conditions of Definition~\ref{def:exact} of those
families directly, on the numbers it is about to emit, and refuses where one
fails (Requirement~\ref{req:weighted-family-emission}); it neither rounds nor
snaps its own emission, and it never declares a nearby family
(\iface{no_snapping}). A specification whose weighting does not land exactly
on the lattice is a specification to be re-declared at the \exttermref{../foundations/contract.pdf}{term-instance}{instance}.
What the engine owes its consumer is the five conditions, and it decides them
rather than inheriting them from a lattice.
\end{requirement}
\claimmeta{\metalabel{Motivated by}\mref{Def.}{def:exact}, \mref{Def.}{def:weighted}, \mref{Req.}{req:weighted-family-emission}, \extref{../foundations/contract.pdf}{foundations/req:declared-lattice}, \extref{../foundations/contract.pdf}{foundations/req:exact-emission}, \extref{../foundations/contract.pdf}{foundations/req:who-snaps}
\metasep implemented (effective: accepted)
\metasep \metalabel{Code}\coderef{https://github.com/fraunhofersg/certiqs-key/blob/v0.1.0/src/certiqs_key/protocol_compiler/weights.py}{certiqs_key.protocol_compiler.weights:decide_emission}}

% claim protocol-compiler/req:no-fixed-fraction
\claimsummary{protocol-compiler/req:no-fixed-fraction}{The emitted declaration carries the output-state test-statistics model and never a fixed fraction. A specification requesting a state-independent fraction applied outside the operators is refused.}
\claimbookmark{protocol-compiler/req:no-fixed-fraction}{Requirement}{req:no-fixed-fraction}{No fixed fraction outside the measurement operators}
\claimleadin{%
A state-independent fraction applied outside the operators is unsound, since
the terms that carry the weight depend on the state
(Remark~\ref{rem:exactness}), and by
Theorem~\ref{thm:choice-encoding}\,(ii) the weight inside the operators
reproduces the declared choice statistics exactly wherever the weighted
family is the emitted one, so nothing is lost by refusing the fraction; in
front of a \exttermref{../reductions/contract.pdf}{term-receiver-segment}{receiver segment} that squashes, the weight sits inside the
physical family handed to the chain, and no fraction is emitted either. This
is the rule the encoding exists to enforce.
}
\begin{requirement}[No fixed fraction outside the measurement operators]\label{req:no-fixed-fraction}
The emitted declaration carries the output-state test-statistics model
(\xref{reductions/ass:tsm}) and never the fixed-fraction model: no path of the
engine produces an artefact carrying a fixed-fraction model, and no weight
$q_t$ is emitted outside a POVM element. A specification requesting the
fixed-fraction encoding of Theorem~\ref{thm:choice-encoding}\,(iv) --- the conditional
families together with a separate state-independent fraction $q_t$ to be
applied outside the operators --- is refused with
\iface{fixed_fraction_requested}.
\end{requirement}
\claimmeta{\metalabel{Motivated by}\mref{Thm.}{thm:choice-encoding}, \mref{Rem.}{rem:exactness}, \extref{../reductions/contract.pdf}{reductions/ass:tsm}
\metasep implemented (effective: accepted)
\metasep \metalabel{Code}\coderef{https://github.com/fraunhofersg/certiqs-key/blob/v0.1.0/src/certiqs_key/protocol_compiler/weights.py}{certiqs_key.protocol_compiler.weights:refuse_fixed_fraction}}

\section{Requirements: validation, feasibility and the chain}\label{group-protocol-compiler/req:validator-refusals}
\sectionintro{The engine refuses a failing check with a reason naming it, repairs nothing, and reports in a check's transcript only the decisions it took. It tests the default feasibility candidate before a supplied witness, emits neither, and never reports a refusal for want of a witness as a defect of the protocol. It decides the presented physical receiver family exactly admissible on every route, and refuses a chain failing one of its conditions where the chain is consumed, before a declaration exists. It takes the Band-A route only for one \termref{term-measured-factor}{measured factor} whose receiver system, and transmitted system where the source is trusted, is two-dimensional and whose declared family is projective, and only where that factor's one link supplies the identity channel as data, keeps the outcomes as they are and delivers the declared family. It emits both segments of the chain record and both slots of every link, empty or not, and omits the link-soundness row from its ledger only on \termref{term-band-a}{Band A} and on a single exact instance link the reductions module emitted in the same run, a provenance it never reads from the specification. Where the chain carries a reference companion, it compiles a companion declaration beside the declaration by the same map, decides none of its family's predicates and emits it evaluation-only, with no run of its own.}
\sectionimports{\extref{../foundations/contract.pdf}{foundations/def:factor-enclosure}, \extref{../reductions/contract.pdf}{reductions/cor:entries}, \extref{../reductions/contract.pdf}{reductions/cor:reference-companion}, \extref{../reductions/contract.pdf}{reductions/lem:linear}, \extref{../reductions/contract.pdf}{reductions/lem:row-relaxation}, \extref{../reductions/contract.pdf}{reductions/thm:chain}, \extref{../reductions/contract.pdf}{reductions/thm:emit}, \extref{../reductions/contract.pdf}{reductions/thm:emitted-product}, \extref{../reductions/contract.pdf}{reductions/thm:licence}, \extref{../reductions/contract.pdf}{reductions/thm:transfer}}

% claim protocol-compiler/req:validator-refusals
\claimsummary{protocol-compiler/req:validator-refusals}{A candidate declaration failing a check of the validator is refused, and no declaration is emitted, with a reason naming the check and what that check names. Nothing is repaired and no table is completed with a default value.}
\claimbookmark{protocol-compiler/req:validator-refusals}{Requirement}{req:validator-refusals}{Refusals of the validator's checks}
\claimleadin{%
A pass of every check of Definition~\ref{def:validator} is what establishes
the target conditions (Theorem~\ref{thm:validator-soundness}\,(i)), and the compiler, as a
\exttermref{../foundations/contract.pdf}{term-consumer}{consumer} of what it is handed, decides and never repairs
(\xref{foundations/req:consumer-decides-exactly}).
}
\begin{requirement}[Refusals of the validator's checks]\label{req:validator-refusals}
The engine runs checks V1 to V8 of Definition~\ref{def:validator} on every
candidate declaration and refuses a candidate that fails one, emitting no
declaration, with a reason naming the failing check. It repairs nothing and completes no declared table with a default
value. Where a check has a reason of its own, the refusal carries that reason
and what the check names with it:
\begin{enumerate}[label=(\roman*),leftmargin=2.2em]
\item V1: \iface{sector_map_mismatch}, naming the first position at which the
  compared lists differ, and in the general format, for the block kind across
  the factors, the first factor at which the kind differs;
\item V2: \iface{alice_marginal_trace_rounds} where an addition of the
  marginal's trace rounds; in the general format
  \iface{alice_sector_marginal_zero}, naming the sector, where a block of the
  marginal is zero beside a measurement slot in \termref{term-ensemble-form}{ensemble form};
\item V3: \iface{alice_povm_summation_rounds}, in the general format carrying
  the Alice sector; \iface{bob_povm_summation_rounds}, carrying the sector and
  the outcome index at which the first inexact addition occurred, and in the
  general format the factor as well; in the general format
  \iface{alice_ensemble_marginal_mismatch}, carrying the sector and the first
  differing entry, \iface{alice_ensemble_sum_rounds}, carrying the sector and
  the index of the first inexact addition, and
  \iface{ensemble_form_in_untrusted_mode} for a slot in ensemble form beside a
  \termref{term-marginal-slot}{marginal slot} holding $\nomarg$;
\item V6: in the general format \iface{constraint_scope_not_a_factor},
  naming the entry;
\item V7: \iface{no_feasibility_witness}, naming, where a supplied witness was
  rejected on its map, the first position at which the two lists differ.
\end{enumerate}
\end{requirement}
\claimmeta{\metalabel{Motivated by}\mref{Def.}{def:validator}, \mref{Thm.}{thm:validator-soundness}, \extref{../foundations/contract.pdf}{foundations/req:consumer-decides-exactly}
\metasep accepted
\metasep \metalabel{Code}\coderef{https://github.com/fraunhofersg/certiqs-key/blob/v0.1.0/src/certiqs_key/protocol_compiler/validation.py}{certiqs_key.protocol_compiler.validation:validate}}

% claim protocol-compiler/req:validation-transcript
\claimsummary{protocol-compiler/req:validation-transcript}{A check may stop at a failed prerequisite or predicate. Its transcript then reports only the decisions it took, in their execution order, and no outcome is inferred for a predicate it did not execute.}
\claimbookmark{protocol-compiler/req:validation-transcript}{Requirement}{req:validation-transcript}{Transcript of a check that stops early}
\claimleadin{%
The counts of Theorem~\ref{thm:validator-soundness}\,(iv) are the complete applicable
inventory of the decisions the checks take on well-typed data, not the number
an execution that stops early performs.
}
\begin{requirement}[Transcript of a check that stops early]\label{req:validation-transcript}
A check of Definition~\ref{def:validator} may stop after a failed
prerequisite or predicate. Its transcript reports only the decisions the check
actually took, in their execution order, and the engine infers no outcome for
a predicate it did not execute.
\end{requirement}
\claimmeta{\metalabel{Motivated by}\mref{Def.}{def:validator}, \mref{Thm.}{thm:validator-soundness}
\metasep implemented (effective: accepted)
\metasep \metalabel{Code}\coderef{https://github.com/fraunhofersg/certiqs-key/blob/v0.1.0/src/certiqs_key/protocol_compiler/validation.py}{certiqs_key.protocol_compiler.validation:CheckOutcome}}

% claim protocol-compiler/req:feasibility-witness
\claimsummary{protocol-compiler/req:feasibility-witness}{At a non-empty constraint list the engine tests the default candidate first and, if it fails, a supplied witness, emits neither, and refuses exactly when the default candidate fails and no supplied witness verifies.}
\claimbookmark{protocol-compiler/req:feasibility-witness}{Requirement}{req:feasibility-witness}{Default candidate before a supplied witness}
\claimleadin{%
Theorem~\ref{thm:feasibility}\,(ii) reduces D(vii) at a non-empty constraint list to
exact decisions on at most two candidates, and check V7 of
Definition~\ref{def:validator} is where they are taken.
}
\begin{requirement}[Default candidate before a supplied witness]\label{req:feasibility-witness}
Where the constraint list is non-empty, on a candidate declaration
satisfying D(i)--D(iii) and D(vi), which checks V1--V3 and V6 of
Definition~\ref{def:validator} decide (a candidate failing one of them is
refused there), the engine decides the
conditions of
\eqref{eq:feas}, in the general format those of \eqref{eq:feas-general},
first on the default candidate \eqref{eq:witness-empty}, in the general format
\eqref{eq:witness-general}, and, if it fails, on a
supplied witness $\omega^{\star}$, read from the input
\iface{feasibility_witness}. It forms the default candidate only to test it,
and it emits neither candidate. It refuses with \iface{no_feasibility_witness}
exactly when the default candidate fails and either no witness is supplied or
the supplied one fails.
\end{requirement}
\claimmeta{\metalabel{Motivated by}\mref{Def.}{def:feasible}, \mref{Thm.}{thm:feasibility}, \mref{Def.}{def:validator}
\metasep implemented (effective: accepted)
\metasep \metalabel{Code}\coderef{https://github.com/fraunhofersg/certiqs-key/blob/v0.1.0/src/certiqs_key/protocol_compiler/feasibility.py}{certiqs_key.protocol_compiler.feasibility:check_feasibility}}

% claim protocol-compiler/req:witness-refusal-meaning
\claimsummary{protocol-compiler/req:witness-refusal-meaning}{A refusal at the feasibility check is never reported to a user as a defect of the protocol: it means only that no feasible state was shown, and a well-formed declaration can be refused there.}
\claimbookmark{protocol-compiler/req:witness-refusal-meaning}{Requirement}{req:witness-refusal-meaning}{Reporting a refusal for want of a feasibility witness}
\claimleadin{%
Check V7 is the one check that can refuse a \termref{term-well-formed-declaration}{well-formed declaration}
(Theorem~\ref{thm:validator-soundness}\,(iii)): its refusal means that no feasible state was
shown, not that none exists (Remark~\ref{rem:completeness}).
}
\begin{requirement}[Reporting a refusal for want of a feasibility witness]\label{req:witness-refusal-meaning}
The engine never reports a refusal with \iface{no_feasibility_witness} to a
user as a defect in the protocol.
\end{requirement}
\claimmeta{\metalabel{Motivated by}\mref{Thm.}{thm:validator-soundness}, \mref{Rem.}{rem:completeness}
\metasep accepted}

% claim protocol-compiler/req:physical-family-decision
\claimsummary{protocol-compiler/req:physical-family-decision}{The physical receiver family as presented is decided exactly admissible on every route, as one block on its physical receiver space, before any condition on the chain is decided. A failure is refused with its own reason, naming the first failure.}
\claimbookmark{protocol-compiler/req:physical-family-decision}{Requirement}{req:physical-family-decision}{Decision on the presented physical receiver family}
\claimleadin{%
Clause~(iv) of Definition~\ref{def:well-formed-spec} reads the physical
receiver family as presented like every other declared family, so nothing
of its exact admissibility is left to whoever presents the chain. On the
Band-A route the same family is the one the link delivers, emitted as it is
unless a setting-dependent choice weighting is declared, and check V3 decides it, or its
weighting, again (Theorem~\ref{thm:band-a}\,(ii)).
}
\begin{requirement}[Decision on the presented physical receiver family]\label{req:physical-family-decision}
On every route the engine decides, on the declared numbers, the condition
clause~(iv) of Definition~\ref{def:well-formed-spec} places on the physical
receiver family as presented, in the general format on each
\termref{term-measured-factor}{measured factor}'s: presented directly, (S), (H), (P) and (C) of
Definition~\ref{def:exact} on its physical receiver space as one block, in
the \exttermref{../foundations/contract.pdf}{term-declared-order}{declared order} of its alphabet; presented as a
\termref{term-choice-weighted-family}{choice-weighted family}, condition~(S) on its own declared-order sum,
after the decision of Requirement~\ref{req:conditional-family-decision} on
its conditional families and the representability decision on the products
formed from them. It decides it before any condition of
Definition~\ref{def:chaincond} is decided on the chain, and refuses a failure
with \iface{physical_family_not_admissible}, naming the measured factor, the
predicate and the index of the first failure (for (S), the index of the
first inexact addition). The decision stops at the first failure.
\end{requirement}
\claimmeta{\metalabel{Motivated by}\mref{Def.}{def:exact}, \mref{Def.}{def:well-formed-spec}, \mref{Def.}{def:weighted}, \mref{Def.}{def:chaincond}, \mref{Thm.}{thm:band-a}, \mref{Req.}{req:conditional-family-decision}
\metasep accepted
\metasep \metalabel{Code}\coderef{https://github.com/fraunhofersg/certiqs-key/blob/v0.1.0/src/certiqs_key/protocol_compiler/compile.py}{certiqs_key.protocol_compiler.compile:decide_physical_family}}

% claim protocol-compiler/req:chain-refusals
\claimsummary{protocol-compiler/req:chain-refusals}{A chain failing a condition on it is refused where it is consumed, before a declaration exists, with a reason naming the condition; the licensed-only and source-mode rules have reasons of their own. No refusal emits a declaration.}
\claimbookmark{protocol-compiler/req:chain-refusals}{Requirement}{req:chain-refusals}{Refusals of the interface conditions on a chain}
\claimleadin{%
Theorem~\ref{thm:chain-consumption}\,(i) and (i$'$) make (C1)--(C5),
(S1)--(S4), (A1)--(A3), the licensed-only rule and the declared-source-mode
rule decidable by finite computation on the presented record. The grounds of
the two source-segment rules are Remark~\ref{rem:source-segment-rules}.
}
\begin{requirement}[Refusals of the interface conditions on a chain]\label{req:chain-refusals}
The engine decides the \termref{term-interface-conditions-on-a-chain}{interface conditions on a chain} of
Definition~\ref{def:chaincond} where the chain is consumed, before a
declaration exists, and refuses a chain that fails one; no refusal emits a
declaration.
\begin{enumerate}[label=(\roman*),leftmargin=2.2em]
\item A chain failing one of (C1)--(C5) or (S1)--(S4) is refused with
  \iface{chain_interface_condition_failed}, naming the condition and, for
  (C1)--(C5) in the general format, the factor, and, for (C4) and (S3), the outcome or
  signal index at which the first entry disagreed. A non-empty
  \exttermref{../transmitter/contract.pdf}{term-source-segment}{source segment} failing one of (A1)--(A3) is refused.
\item Beside a non-empty source segment, an entry declared general is refused
  with \iface{general_entry_beside_source_segment}, naming the entry.
\item A declaration carrying a non-empty source segment in
  \exttermref{../transmitter/contract.pdf}{term-untrusted-source-mode}{untrusted-source mode} is refused with
  \iface{source_segment_in_untrusted_mode}.
\end{enumerate}
\end{requirement}
\claimmeta{\metalabel{Motivated by}\mref{Def.}{def:chaincond}, \mref{Rem.}{rem:source-segment-rules}, \mref{Thm.}{thm:chain-consumption}
\metasep accepted
\metasep \metalabel{Code}\coderef{https://github.com/fraunhofersg/certiqs-key/blob/v0.1.0/src/certiqs_key/protocol_compiler/chain.py}{certiqs_key.protocol_compiler.chain:check_chain}}

% claim protocol-compiler/req:band-a-ideal-receiver
\claimsummary{protocol-compiler/req:band-a-ideal-receiver}{A specification on the Band-A route is refused unless it has one measured factor, a two-dimensional receiver system and transmitted system, and a projective declared receiver measurement. The route then never rests on an ideal-receiver premise its own declared data contradict.}
\claimbookmark{protocol-compiler/req:band-a-ideal-receiver}{Requirement}{req:band-a-ideal-receiver}{Band A only within the ideal qubit receiver}
\claimleadin{%
On the Band-A route the residual physical content is
Assumption~\ref{ass:ideal}, an ideal projective measurement on a qubit
(Theorem~\ref{thm:band-a}\,(iv)), while
Theorem~\ref{thm:band-a} itself holds for any declared dimension and any
\exttermref{../foundations/contract.pdf}{term-exactly-admissible}{exactly admissible} family. Part of that assumption is a property of the
\exttermref{../foundations/contract.pdf}{term-declared-data}{declared data}: the number of measured factors, the dimensions and the
projectivity of the declared family. The engine decides that part, so that the
route is never taken for a declaration that contradicts the premise it rests
on. Exact admissibility of the same family is decided first, on every route
(Requirement~\ref{req:physical-family-decision}). For a family presented as
a \termref{term-choice-weighted-family}{choice-weighted family}, Assumption~\ref{ass:ideal}\,(i) is read of
each conditional family: a public choice of known bias followed by an ideal
projective measurement. The weighted family itself is not projective when a
weight is below one, and the same measurement presented directly, as a
passive basis choice, is refused on purpose: a passive choice is made by an
optical element, which the idealisation does not describe.
}
\begin{requirement}[Band A only within the ideal qubit receiver]\label{req:band-a-ideal-receiver}
The engine recognises the \termref{term-band-a}{Band A} route on the record, reading the
shape that condition (b) of Definition~\ref{def:banda} gives the chain and
the sector clause of its (a), for a specification whose \exttermref{../transmitter/contract.pdf}{term-source-segment}{source segment} is
empty, whose emitted \exttermref{../reductions/contract.pdf}{term-sector-set}{sector set} (in the general format, the joint one) is a singleton, and
in which, for some \termref{term-measured-factor}{measured factor}, the \exttermref{../reductions/contract.pdf}{term-receiver-segment}{receiver segment} is one
link with status exact, equal source and
\exttermref{../transmitter/contract.pdf}{term-target-descriptor}{target descriptors}, a channel that, where the link
supplies one as data, is exactly the identity map on its system, an empty
constraint slot and an empty licence slot; it decides the link condition on
the chain record of every measured factor. The clauses of (b) this
recognition leaves open --- that the link supplies its channel as data, its
outcome correspondence and its \exttermref{../reductions/contract.pdf}{term-target-povm}{target POVM} --- are
decided by Requirement~\ref{req:band-a-identity-link}. It refuses such a
specification, before a declaration is emitted, with
\iface{band_a_outside_ideal_receiver}, naming the first clause below that
fails and, for (iii), the choice value where there is one and the outcome
index of the first element that is not a projector, unless
\begin{enumerate}[label=(\roman*),leftmargin=2.2em]
\item the specification declares exactly one measured factor;
\item that factor's physical receiver system has dimension two and, in
  \exttermref{../transmitter/contract.pdf}{term-trusted-source-mode}{trusted-source mode}, the \exttermref{../transmitter/contract.pdf}{term-transmitted-system}{transmitted system} of the
  transmitter artefact (Definition~\ref{def:spec}\,(a)) has dimension two; and
\item that factor's physical receiver family as presented is projective,
  decided exactly on the declared numbers: presented directly, every element
  satisfies $\Fph_w\Fph_w=\Fph_w$; presented as a
  \termref{term-choice-weighted-family}{choice-weighted family}, every element of every declared conditional
  family satisfies $N_{u|t}N_{u|t}=N_{u|t}$
  (Definition~\ref{def:weighted}).
\end{enumerate}
\end{requirement}
\claimreading{The idealised Band-A route is taken only for a receiver that can be the ideal qubit measurement its premise \hyperlink{claim-protocol-compiler/ass:ideal}{A-IDEAL} describes: one measuring station, a two-dimensional measured system and, with a trusted source, a two-dimensional transmitted system, and a declared measurement made of projectors (per basis, where a basis choice is declared). Anything else is refused rather than analysed under a premise its own numbers already contradict; such a receiver is declared through the receiver module and a squashing reduction.}
\claimmeta{\metalabel{Motivated by}\mref{Def.}{def:spec}, \mref{Def.}{def:weighted}, \mref{Def.}{def:banda}, \mref{Thm.}{thm:band-a}, \mref{Ass.}{ass:ideal}, \mref{Req.}{req:physical-family-decision}, \mref{Req.}{req:band-a-identity-link}
\metasep \metalabel{Used by}\mref{Req.}{req:band-a-identity-link}
\metasep accepted
\metasep \metalabel{Code}\coderef{https://github.com/fraunhofersg/certiqs-key/blob/v0.1.0/src/certiqs_key/protocol_compiler/declared_premises.py}{certiqs_key.protocol_compiler.declared_premises:decide_band_a_ideal_receiver}}

% claim protocol-compiler/req:band-a-identity-link
\claimsummary{protocol-compiler/req:band-a-identity-link}{On the Band-A route the one receiver link must supply the identity channel as data and keep the outcomes as they are, and the pullback identity then decides that the measurement the link delivers is the declared one. A specification failing any of this is refused before a declaration exists.}
\claimbookmark{protocol-compiler/req:band-a-identity-link}{Requirement}{req:band-a-identity-link}{Band A only on an identity link that delivers the declared family}
\claimleadin{%
Theorem~\ref{thm:band-a}\,(ii) and~(v) read the family the link delivers as the
declared one and the outcome correspondence as the identity, and
Requirement~\ref{req:link-row-provenance}\,(i) omits the link-soundness row of
Assumption~\ref{ass:link} on that strength. A link record carries its
channel, its outcome correspondence and its target POVM as separate data, and
where the link only names its channel the pullback identity (C4) is not
decidable (Definition~\ref{def:chaincond}); on such a record nothing would
tie the measurement the key length is computed for to the one declared. So
the route is taken only where all three are decided on the record.
}
\begin{requirement}[Band A only on an identity link that delivers the declared family]\label{req:band-a-identity-link}
For a specification that Requirement~\ref{req:band-a-ideal-receiver}
recognises on the \termref{term-band-a}{Band A} route and does not refuse, the engine
decides on the chain record of its one \termref{term-measured-factor}{measured factor}, before a
declaration is emitted, that
\begin{enumerate}[label=(\roman*),leftmargin=2.2em]
\item the one link of that factor's \exttermref{../reductions/contract.pdf}{term-receiver-segment}{receiver segment} supplies its
  channel as data, which that recognition has then decided to be exactly the
  identity map on its system; and
\item the outcome correspondence of that segment is the identity: the
  physical receiver alphabet $\Vph$ and the alphabet $\Vset$ of the link's
  \exttermref{../reductions/contract.pdf}{term-target-povm}{target POVM} are one set, and $\gamma(w)=w$ for every
  $w\in\Vph$;
\end{enumerate}
and refuses a specification failing either, before a declaration is emitted,
with \iface{band_a_link_not_identity}, naming the first clause that fails. On
a specification passing both, condition (C4) of
Definition~\ref{def:chaincond} is decided, as on every chain whose links
supply their channel as data, and reads $F_v=\Fph_v$ for every $v\in\Vset$,
entry for entry, $\{\Fph_w\}$ being the physical receiver family as presented
and, for a \termref{term-choice-weighted-family}{choice-weighted family}, the family of products
Definition~\ref{def:weighted} forms; a failure of it is refused as every
failure of (C4) is (Requirement~\ref{req:chain-refusals}). A declaration
emitted from a specification so recognised therefore satisfies
Definition~\ref{def:banda}\,(b) on its record --- its one link the identity,
supplied as data, its outcome correspondence the identity and its target
POVM the declared family --- and the rules that read the route, the ledger's
among them (Requirement~\ref{req:link-row-provenance}), read it of such a
declaration alone.
\end{requirement}
\claimreading{The idealised Band-A route is taken only when its one reduction step is the identity map given as numbers, keeps your outcomes as they are, and hands on exactly the measurement you declared. The engine checks all three on the record and refuses anything else, so on this route the key length is always computed for the measurement your declaration names, with any declared public choice applied to it, and the premise that the reduction step is sound is proved rather than left to you.}
\claimmeta{\metalabel{Motivated by}\mref{Def.}{def:weighted}, \mref{Def.}{def:chaincond}, \mref{Def.}{def:banda}, \mref{Thm.}{thm:band-a}, \mref{Ass.}{ass:link}, \mref{Req.}{req:chain-refusals}, \mref{Req.}{req:band-a-ideal-receiver}, \mref{Req.}{req:link-row-provenance}
\metasep implemented (effective: accepted)
\metasep \metalabel{Code}\coderef{https://github.com/fraunhofersg/certiqs-key/blob/v0.1.0/src/certiqs_key/protocol_compiler/declared_premises.py}{certiqs_key.protocol_compiler.declared_premises:decide_band_a_identity_link}}

% claim protocol-compiler/req:chain-record-slots
\claimsummary{protocol-compiler/req:chain-record-slots}{The emitted chain record carries the source segment, the receiver segment and, on every link, the constraint slot and the licence slot, whether or not they are empty, Band A included.}
\claimbookmark{protocol-compiler/req:chain-record-slots}{Requirement}{req:chain-record-slots}{Emission of both chain segments and every link slot}
\claimleadin{%
The order rule (X1) of Definition~\ref{def:chaincond} is carried by the shape
of the record, two segments, and the constraint entries each link transports
are read off its slots (C5); the reason a \exttermref{../transmitter/contract.pdf}{term-source-link}{source link} may not follow a
receiver link is \xref{reductions/rem:ordering}.
}
\begin{requirement}[Emission of both chain segments and every link slot]
\label{req:chain-record-slots}
The engine emits the chain record (\iface{reduction_chain_record}) with both
of its segments, the \exttermref{../transmitter/contract.pdf}{term-source-segment}{source segment} and the \exttermref{../reductions/contract.pdf}{term-receiver-segment}{receiver segment}, and on
every link with both its constraint slot and its licence slot, each whether
or not it is empty. This includes \termref{term-band-a}{Band A}, where the source segment is
empty and the receiver segment is one identity link with both slots empty.
\end{requirement}
\claimmeta{\metalabel{Motivated by}\mref{Def.}{def:chaincond}, \extref{../reductions/contract.pdf}{reductions/rem:ordering}
\metasep implemented
\metasep \metalabel{Code}\coderef{https://github.com/fraunhofersg/certiqs-key/blob/v0.1.0/src/certiqs_key/protocol_compiler/compile.py}{certiqs_key.protocol_compiler.compile:chain_record}}

% claim protocol-compiler/req:link-row-provenance
\claimsummary{protocol-compiler/req:link-row-provenance}{The ledger omits \hyperlink{claim-protocol-compiler/ass:link}{A-LINK} only on Band A and where every measured factor's receiver segment is the single exact instance link of a reductions artefact of the same run, with the identity outcome correspondence. That provenance is the run's, never read from the specification.}
\claimbookmark{protocol-compiler/req:link-row-provenance}{Requirement}{req:link-row-provenance}{Omission of the link-soundness row only on a proved chain}
\claimleadin{%
Theorem~\ref{thm:chain-consumption}\,(v) names the soundness of each link as
an obligation of the chain's producer, and Remark~\ref{rem:obligations-discharged}
records the chains for which the producer's own theorems discharge it. A
specification is a record (Definition~\ref{def:spec}\,(c)): its link types
and status tags are data that a specification written by hand can carry
equally, so whether a chain is one the reductions module emitted cannot be
read off it. Condition~(ii) below is the shape that the theorems this
requirement uses discharge. An instance link followed by a further link
squashes another model (a link of \xref{reductions/domination/def:domination-link}),
and on such a chain the row stays. At several
\termref{term-measured-factor}{measured factors} whose segments are each that
factor's instance link alone, the joint transfer is
\xref{reductions/thm:emitted-product}\,(ii) and~(iii), the apparatus
measuring the product of the factors' families by A-MEAS
(\xref{reductions/ass:meas})\,(iv). A setting-dependent weighting applied
after a factor's segment leaves the chain record as it is
(Remark~\ref{rem:relabel}), so such a factor meets condition~(ii) as an
unweighted one does, and the weighted relation is
\xref{reductions/lem:linear}\,(i). At Band A the discharge is
Theorem~\ref{thm:band-a}\,(v), whose hypothesis, Definition~\ref{def:banda}\,(b),
is decided on the record (Requirement~\ref{req:band-a-identity-link}).
}
\begin{requirement}[Omission of the link-soundness row only on a proved chain]\label{req:link-row-provenance}
The engine omits the row of Assumption~\ref{ass:link} from the assumption
ledger of a compilation only
(i) on the \termref{term-band-a}{Band A} route, for a declaration whose record satisfies
Definition~\ref{def:banda}\,(b) as Requirements~\ref{req:band-a-ideal-receiver}
and~\ref{req:band-a-identity-link} decide it --- the one link the identity,
supplied as data, its outcome correspondence the identity and its target
POVM, by the decided (C4), the presented physical family --- where
Theorem~\ref{thm:band-a}\,(v) proves every clause of it; or
(ii) where the \exttermref{../transmitter/contract.pdf}{term-source-segment}{source segment} is empty and, for every
\termref{term-measured-factor}{measured factor} of the specification, all of the following hold: its
\exttermref{../reductions/contract.pdf}{term-receiver-segment}{receiver segment} consists of one link, the instance link of an
artefact that the reductions module emitted in the same run, with status
\emph{exact} and no channel supplied as data; the outcome correspondence of
the segment is the identity on that factor's target alphabet; and that
factor's physical receiver family as presented is exactly the compression
$\{\iota^{\dagger}\Mdec_v\iota\}_v$ of that artefact's \exttermref{../reductions/contract.pdf}{term-declared-model}{declared model}
onto the presented physical receiver space, $\iota$ its isometric inclusion
into the physical system, or the declared model itself. At several measured
factors the omission rests in addition on A-MEAS
(\xref{reductions/ass:meas})\,(iv), that the round's physical measurement is
the product of the factors' measurements. On either route, at a factor
whose delivered family is weighted after its segment, in the
\termref{term-setting-dependent-form}{setting-dependent form} of Definition~\ref{def:weighted}, the omission
rests in addition on \xref{reductions/lem:linear}\,(i): the weighting of the
family the segment's link delivers and the identically weighted physical
family stand in the \exttermref{../reductions/contract.pdf}{term-squashing-relation}{squashing relation} under the same channel, so the
transfer applies to the weighted protocol as to the unweighted one; on
route~(i) the channel is the identity and the two weighted families are the
same matrices (Theorem~\ref{thm:band-a}\,(v)).
That each link is that artefact's instance link and that each presented
family is that compression is established by the part of the engine that ran
the reductions module and docked its artefacts, only for an instance for which
Remark~\ref{rem:obligations-discharged} names the claim showing the docked
family to be that compression, and is passed to the compiler
beside the specification (\iface{chain_provenance}); the compiler never reads
it from the specification. The compiler decides, on the presented chain, the
parts of~(ii) that are data: for every measured factor the length of its
segment, the status tag, the channel slot and the outcome correspondence
(Definition~\ref{def:chaincond}, (C2)), and the length of the source segment.
Where condition~(ii) is not passed, or a part of it decided on the chain
fails, the row is emitted.
\end{requirement}
\claimmeta{\metalabel{Motivated by}\mref{Def.}{def:spec}, \mref{Def.}{def:weighted}, \mref{Rem.}{rem:relabel}, \mref{Def.}{def:chaincond}, \mref{Thm.}{thm:chain-consumption}, \mref{Rem.}{rem:obligations-discharged}, \mref{Def.}{def:banda}, \mref{Thm.}{thm:band-a}, \mref{Ass.}{ass:link}, \mref{Req.}{req:band-a-ideal-receiver}, \mref{Req.}{req:band-a-identity-link}, \extref{../reductions/contract.pdf}{reductions/cor:entries}, \extref{../reductions/contract.pdf}{reductions/domination/def:domination-link}, \extref{../reductions/contract.pdf}{reductions/lem:linear}, \extref{../reductions/contract.pdf}{reductions/thm:chain}, \extref{../reductions/contract.pdf}{reductions/thm:emit}, \extref{../reductions/contract.pdf}{reductions/thm:emitted-product}, \extref{../reductions/contract.pdf}{reductions/thm:licence}, \extref{../reductions/contract.pdf}{reductions/thm:transfer}
\metasep \metalabel{Assumes}\extaref{../reductions/contract.pdf}{reductions/ass:meas}{A-MEAS}
\metasep implemented (effective: accepted)
\metasep \metalabel{Code}\coderef{https://github.com/fraunhofersg/certiqs-key/blob/v0.1.0/src/certiqs_key/protocol_compiler/compile.py}{certiqs_key.protocol_compiler.compile:link_row_discharged}}

% claim protocol-compiler/req:reference-companion-compiled
\claimsummary{protocol-compiler/req:reference-companion-compiled}{Where the chains of one or more measured factors carry a reference companion, the compiler forms beside the declaration one companion replacing each, by the same map, weighted exactly with its factor enclosures, decides none of its predicates and emits it evaluation-only, with no run of its own.}
\claimbookmark{protocol-compiler/req:reference-companion-compiled}{Requirement}{req:reference-companion-compiled}{The reference companion compiled beside the declaration}
\claimleadin{%
The companion a reductions emission carries beside its artefact
(\xref{reductions/req:reference-companion}) is the same receiver segment run
on the reference. Compiling it by the same map, from the same specification,
gives a tuple that differs from the declaration in the measured families of
the factors it replaces and the entries of their last links and in nothing
else (\xref{meat-core/prop:reference-companion}). Those families are known
through enclosures and are off the binary64 grid, so the checks this module
decides on emitted \exttermref{../foundations/contract.pdf}{term-binary64-numbers}{binary64 numbers}, the validator's
(Definition~\ref{def:validator}) and the chain conditions'
(Definition~\ref{def:chaincond}), do not apply to them; that each is a POVM
is imported from the reductions corollary, and the consumer that evaluates it
admits it under its own rule (\xref{meat-core/req:companion-admission}).
}
\begin{requirement}[The reference companion compiled beside the declaration]\label{req:reference-companion-compiled}
Where the last link of the receiver segment of one or more
\termref{term-measured-factor}{measured factors} was emitted with a reference
companion (\xref{reductions/req:reference-companion}), the engine emits,
beside the declaration $\Decl$ it compiles, one \emph{compiled reference
companion} $\Decl^{\mathrm{ref}}$, in \iface{reference_companion}:
\begin{enumerate}[label=(\roman*),leftmargin=2.4em]
\item it is the tuple the \termref{term-compilation-map}{compilation map} (Definition~\ref{def:compile})
  gives from the same specification with, for every measured factor whose
  last link was so emitted, that link's \exttermref{../reductions/contract.pdf}{term-target-family}{target family} replaced by its
  companion's and that link's entries replaced by its companion's relaxed
  entries (\xref{reductions/lem:row-relaxation}), those of the entries
  requested of the link included where the companion carries them, and none
  in place of a requested entry the companion does not carry
  (\xref{reductions/req:reference-companion}\,(iii)); the engine forms it by
  the code path that forms $\Decl$, every other slot read from the same
  data, the family and the entries of every other measured factor's last
  link among them;
\item where a \termref{term-measured-factor}{measured factor}'s family is weighted after its chain,
  by that factor's weighting (Definition~\ref{def:weighted}), every product
  $P(f\,|\,t(v))\,F_v$ is
  formed exactly from the companion's enclosure: the centre multiplied by
  the weight and every radius by its modulus, so that the result encloses
  the product of the weight and the companion's element. Where the companion
  carries a \exttermref{../foundations/contract.pdf}{term-factor-enclosure}{factor enclosure} (\xref{foundations/def:factor-enclosure})
  of a block of $F_v$ (\xref{reductions/req:reference-companion}\,(vii)),
  the block of the product carries, beside its \exttermref{../foundations/contract.pdf}{term-entrywise-enclosure}{entrywise enclosure}, the one of
  Lemma~\ref{lem:weighted-factor}: the same box, its centres and radii
  unaltered, and the core multiplied by the weight in exact rational
  arithmetic. Where the family is not weighted after its chain, every factor
  enclosure the companion carries is passed to $\Decl^{\mathrm{ref}}$
  unaltered;
\item the engine decides on $\Decl^{\mathrm{ref}}$ none of the predicates
  it decides on the binary64 numbers of $\Decl$: no Hermiticity, positivity
  or completeness of the family, no completeness sum in a \exttermref{../foundations/contract.pdf}{term-declared-order}{declared order}
  (Definition~\ref{def:exact}), no chain condition and no check of the
  validator; it passes the companion's entries through as
  Requirement~\ref{req:pass-through} passes every entry;
\item it emits no bundle of a run of its own from $\Decl^{\mathrm{ref}}$
  and marks it evaluation-only, a declaration with no run of its own, with
  a reference to the declaration $\Decl$ it is the
  companion of and the companion's record carried unchanged.
\end{enumerate}
Where no measured factor's last link carries a companion, nothing of this
requirement applies, and the run forms its first-order number on $\Decl$
alone, its blocks resting as the records of the artefact say.
\end{requirement}
\claimmeta{\metalabel{Motivated by}\mref{Def.}{def:declaration}, \mref{Def.}{def:exact}, \mref{Def.}{def:compile}, \mref{Def.}{def:weighted}, \mref{Lem.}{lem:weighted}, \mref{Lem.}{lem:weighted-factor}, \mref{Def.}{def:validator}, \mref{Def.}{def:chaincond}, \mref{Req.}{req:pass-through}, \extref{../foundations/contract.pdf}{foundations/def:factor-enclosure}, \extref{../meat-core/contract.pdf}{meat-core/prop:reference-companion}, \extref{../meat-core/contract.pdf}{meat-core/req:companion-admission}, \extref{../reductions/contract.pdf}{reductions/cor:reference-companion}, \extref{../reductions/contract.pdf}{reductions/lem:row-relaxation}, \extref{../reductions/contract.pdf}{reductions/req:reference-companion}
\metasep \metalabel{Used by}\extref{../meat-core/contract.pdf}{meat-core/req:companion-admission}
\metasep implemented (effective: accepted)
\metasep \metalabel{Code}\coderef{https://github.com/fraunhofersg/certiqs-key/blob/v0.1.0/src/certiqs_key/protocol_compiler/companion.py}{certiqs_key.protocol_compiler.companion:compile_reference_companion}}

\section{Requirements: the emission gates}\label{group-protocol-compiler/req:eligibility-flag}
\sectionintro{The engine emits the eligibility flag true only for a key-length-eligible declaration and produces, for each gate failure, the outcome the \termref{term-gate-outcome-table}{gate outcome table} fixes. It tests attestations for presence and never for truth, and records a vacuous pass of the licensed-class gate as vacuous.}
\sectionimports{\extref{../transmitter/contract.pdf}{transmitter/req:licence-scope}}

% claim protocol-compiler/req:eligibility-flag
\claimsummary{protocol-compiler/req:eligibility-flag}{The emitted flag reads true only for a declaration that passes all three gates of \hyperlink{claim-protocol-compiler/def:eligible}{Definition~\ref*{def:eligible}}. A declaration marked false may be explored numerically but backs no key-length claim.}
\claimbookmark{protocol-compiler/req:eligibility-flag}{Requirement}{req:eligibility-flag}{The key-length eligibility flag}
\claimleadin{%
Theorem~\ref{thm:gates} fixes, for each gate failure, whether the outcome is a
refusal or an emission that is not key-length-eligible; the flag an emitted
declaration carries never reads true for one that fails a gate.
}
\begin{requirement}[The key-length eligibility flag]\label{req:eligibility-flag}
The engine emits the flag \iface{key_length_eligible} true for a declaration
only if that declaration is \termref{term-key-length-eligible}{key-length-eligible}
(Definition~\ref{def:eligible}).
\end{requirement}
\claimmeta{\metalabel{Motivated by}\mref{Def.}{def:eligible}, \mref{Thm.}{thm:gates}
\metasep implemented (effective: accepted)
\metasep \metalabel{Code}\coderef{https://github.com/fraunhofersg/certiqs-key/blob/v0.1.0/src/certiqs_key/protocol_compiler/gates.py}{certiqs_key.protocol_compiler.gates:resolve_gates}}

% claim protocol-compiler/req:gate-outcomes
\claimsummary{protocol-compiler/req:gate-outcomes}{A failing check or a failing certificate is refused with the check or the condition named, whatever else fails; an unlicensed list or an absent attack-set licence alone is emitted marked not key-length-eligible. The engine never supplies a missing attestation or proceeds as though it had one.}
\claimbookmark{protocol-compiler/req:gate-outcomes}{Requirement}{req:gate-outcomes}{Outcome of each gate failure}
\claimleadin{%
The \termref{term-gate-outcome-table}{gate outcome table} (Definition~\ref{def:gate-outcome-table}) fixes which
of the two outcomes each failure has, as Theorem~\ref{thm:gates}\,(ii) states;
this requirement names the strings and the flag that carry them. A check that
fails, so that no declaration is shown to be well-formed, or a certificate
that does not verify, is a refusal; a \termref{term-well-formed-declaration}{well-formed declaration} on
verified data whose key-length claim is unsupported is emitted and marked. The
transmitter refuses a request whose certificate fails a condition of its form
(\xref{transmitter/req:failing-certificate}); gate (G2) re-decides the
conditions all the same, a consumer deciding the data it is handed
(\xref{foundations/req:consumer-decides-exactly}).
}
\begin{requirement}[Outcome of each gate failure]\label{req:gate-outcomes}
For each gate failure the engine produces the outcome the gate outcome table
(Definition~\ref{def:gate-outcome-table}) fixes for it, and no other; where
the failures of one declaration meet the premises of a refusal row and of an
emission row both, it produces the refusal:
\begin{enumerate}[label=(\roman*),leftmargin=2.2em]
\item row~1, a check that fails: a refusal naming the check, which in the
  general format covers \iface{ensemble_form_in_untrusted_mode},
  \iface{alice_ensemble_marginal_mismatch}, \iface{alice_ensemble_sum_rounds},
  \iface{alice_sector_marginal_zero} and
  \iface{constraint_scope_not_a_factor};
\item row~3, \exttermref{../transmitter/contract.pdf}{term-trusted-source-mode}{trusted-source mode} at any $d_A$ with a certificate present and
  failing a condition of gate (G2) for its form: a refusal naming the
  condition, an \exttermref{../transmitter/contract.pdf}{term-ensemble-presentation}{ensemble
  presentation} beside a POVM-form declaration being refused with
  \iface{ensemble_presentation_without_ensemble_form};
\item rows~5 and~6, a constraint list not in the licensed class or the
  attack-set \exttermref{../reductions/contract.pdf}{term-licence-record}{licence record} \iface{attack_set_license} absent: an emission
  with \iface{key_length_eligible} false.
\end{enumerate}
\emph{A trusted-mode specification at $d_A\ge2$ whose certificate is absent
does not reach this module.} The mode is read off the transmitter artefact's
\termref{term-marginal-slot}{marginal slot} (Definition~\ref{def:compile}), and the transmitter
refuses a trusted-mode request at $d_A\ge2$ that carries no certificate before
any artefact is emitted, with \iface{trusted_mode_without_certificate}
(\xref{transmitter/req:licence-scope}; the modes are defined in
\xref{transmitter/def:source-spec}). So an artefact in trusted mode at
$d_A\ge2$ carries the record gate (G2) tests for, and the two rows of the
table that describe a certificate-absent case --- rows 2 and 4, at $m=0$ and
at $m\ge1$ --- are \emph{unreachable through the transmitter artefact}. At
$d_A=1$ the gate demands no certificate at all, the slot holds the $[1]$ the
transmitter filled it with (\xref{transmitter/req:marginal-fill}), and one
supplied there and failing is refused by row 3, its scope being every $d_A$.
Rows 2 and 4 are unreachable in the general format for the same reason, which
reads the transmitter's refusal and no property of the general format. The
engine never supplies a missing attestation and never proceeds as though it
had one.
\end{requirement}
\claimmeta{\metalabel{Motivated by}\mref{Def.}{def:compile}, \mref{Def.}{def:eligible}, \mref{Def.}{def:gate-outcome-table}, \mref{Thm.}{thm:gates}, \extref{../foundations/contract.pdf}{foundations/req:consumer-decides-exactly}, \extref{../transmitter/contract.pdf}{transmitter/def:source-spec}, \extref{../transmitter/contract.pdf}{transmitter/req:failing-certificate}, \extref{../transmitter/contract.pdf}{transmitter/req:licence-scope}, \extref{../transmitter/contract.pdf}{transmitter/req:marginal-fill}
\metasep accepted
\metasep \metalabel{Code}\coderef{https://github.com/fraunhofersg/certiqs-key/blob/v0.1.0/src/certiqs_key/protocol_compiler/gates.py}{certiqs_key.protocol_compiler.gates:resolve_gates}}

% claim protocol-compiler/req:attestation-presence
\claimsummary{protocol-compiler/req:attestation-presence}{The retention attestation (SR3) behind gate (G2) and the attack-set licence behind gate (G3)(b) are tested for presence and never for truth. The engine never supplies one and never verifies one.}
\claimbookmark{protocol-compiler/req:attestation-presence}{Requirement}{req:attestation-presence}{Attestations tested for presence only}
\claimleadin{%
Neither attestation is a property of the emitted matrices, so no function of
them decides it (Theorem~\ref{thm:gates}\,(i)); what the engine can do is
record whether it was supplied.
}
\begin{requirement}[Attestations tested for presence only]\label{req:attestation-presence}
The engine tests the attestations of Definition~\ref{def:eligible} for
presence only, never for truth: at gate (G2) the clause (SR3) of the
source-replacement licence record \iface{source_replacement_license}
(\xref{transmitter/def:srcert}), in every form of the certificate, and at gate
(G3)(b) the attack-set licence record \iface{attack_set_license}, the
attestation of Assumption~\ref{ass:atk}. It never supplies an attestation and
never verifies one.
\end{requirement}
\claimmeta{\metalabel{Motivated by}\mref{Def.}{def:eligible}, \mref{Thm.}{thm:gates}, \mref{Ass.}{ass:atk}, \extref{../transmitter/contract.pdf}{transmitter/def:srcert}
\metasep accepted
\metasep \metalabel{Code}\coderef{https://github.com/fraunhofersg/certiqs-key/blob/v0.1.0/src/certiqs_key/protocol_compiler/gates.py}{certiqs_key.protocol_compiler.gates:resolve_gates}}

% claim protocol-compiler/req:vacuous-gate-report
\claimsummary{protocol-compiler/req:vacuous-gate-report}{In untrusted-source mode gate (G3)(a) passes by construction, so the gate report records it as passed vacuously and not as passed. A vacuous pass is never reported as evidence of anything.}
\claimbookmark{protocol-compiler/req:vacuous-gate-report}{Requirement}{req:vacuous-gate-report}{A vacuous pass recorded as vacuous}
\claimleadin{%
In \exttermref{../transmitter/contract.pdf}{term-untrusted-source-mode}{untrusted-source mode} the constraint gate rests on the attack-set
licence alone (Remark~\ref{rem:degenerate}); the gate report must not make
its type-check half look like a second check.
}
\begin{requirement}[A vacuous pass recorded as vacuous]\label{req:vacuous-gate-report}
Where gate (G3)(a) of Definition~\ref{def:eligible} holds vacuously, in
untrusted-source mode at any $d_A$, the engine records it in the gate report
\iface{gate_report} as \emph{passed vacuously} rather than \emph{passed}, and
it never reports the vacuous pass as evidence of anything.
\end{requirement}
\claimmeta{\metalabel{Motivated by}\mref{Def.}{def:eligible}, \mref{Rem.}{rem:degenerate}
\metasep accepted
\metasep \metalabel{Code}\coderef{https://github.com/fraunhofersg/certiqs-key/blob/v0.1.0/src/certiqs_key/protocol_compiler/gates.py}{certiqs_key.protocol_compiler.gates:resolve_gates}}
\clearpage
\bookmarksetupnext{level=part}
\section*{Interface}\addcontentsline{toc}{section}{Interface}
\sectionintro{The declared inputs and outputs of the module and its submodules: each key with its symbol, when it is present and the claims that fix it. The domains, types and refusal rules are in the interface files themselves.}
\subsection*{Inputs: Protocol compiler (protocol-compiler)}
\noindent{\small From \href{interface.yaml}{\nolinkurl{interface.yaml}}.}\par\smallskip
{\footnotesize\begin{xltabular}{\linewidth}{@{}>{\raggedright\arraybackslash}p{4.4cm}>{\raggedright\arraybackslash}X>{\raggedright\arraybackslash}p{2.9cm}@{}}
\toprule
Key and symbol & Presence & Fixed by\\
\midrule
\endhead
\iface{transmitter_artifact}\newline Spec\_src & always & \mref{Def.}{def:spec}, \mref{Def.}{def:compile}, \extref{../transmitter/contract.pdf}{transmitter/rem:da1}, \mref{Def.}{def:eligible}, \mref{Thm.}{thm:gates}, \mref{Req.}{req:gate-outcomes}\\
\addlinespace
\iface{receiver_povm}\newline \{F\textasciicircum{}ph\_w\}\_\{w in V\textasciicircum{}ph\} & always; in Band A it IS the family the link delivers, the link's `target\_povm` being decided equal to it entry for entry (req:band-a-identity-link), and so the emitted family unless a `choice\_weighting` is declared, which then weights it; it is declared directly (thm:band-a) & \mref{Def.}{def:spec}, \mref{Def.}{def:well-formed-spec}, \mref{Req.}{req:physical-family-decision}, \mref{Def.}{def:validator}, \mref{Thm.}{thm:chain-consumption}, \mref{Thm.}{thm:band-a}, \mref{Req.}{req:band-a-identity-link}\\
\addlinespace
\iface{public_choices}\newline (T\_pub, q, \{N\_\{u\textbar{}t\}\}) & optional; one such record per party, either possibly absent & \mref{Def.}{def:weighted}, \mref{Def.}{def:compile}, \mref{Ass.}{ass:coin}, \mref{Def.}{def:well-formed-spec}\\
\addlinespace
\iface{choice_weighting}\newline (F, t: V -\textgreater{} T, P(f \textbar{} t)) & OPTIONAL, and per measured factor: at most one per factor, of either party (def:spec (b), (b')); its choice is drawn by a party whose readout records that factor's setting component (ass:coin (ii)). Absent is the ordinary case: with no such record the emitted family is the chain's \{F\_v\} unchanged (rem:relabel, the paragraph "With no setting-dependent choice declared"). & \mref{Def.}{def:weighted}, \mref{Def.}{def:spec}, \mref{Def.}{def:compile}, \mref{Thm.}{thm:choice-encoding}, \mref{Rem.}{rem:relabel}, \mref{Ass.}{ass:coin}\\
\addlinespace
\iface{reduction_chain}\newline Chain = (Ch\_src, Ch\_rec) & always & \mref{Def.}{def:spec}, \mref{Def.}{def:chaincond}, \extref{../transmitter/contract.pdf}{transmitter/def:source-segment}, \mref{Thm.}{thm:chain-consumption}, \mref{Def.}{def:banda}, \mref{Thm.}{thm:band-a}\\
\addlinespace
\iface{target_sectors}\newline Smap = (S, (d\_s)\_s, block kind) & always & \mref{Def.}{def:declaration}, \mref{Def.}{def:validator}, \mref{Def.}{def:chaincond}, \mref{Def.}{def:target}\\
\addlinespace
\iface{target_povm}\newline \{F\_v\}\_\{v in V\} & always & \mref{Def.}{def:spec}, \mref{Def.}{def:chaincond}, \mref{Def.}{def:compile}, \mref{Def.}{def:exact}, \mref{Def.}{def:validator}, \mref{Thm.}{thm:chain-consumption}, \mref{Def.}{def:banda}, \mref{Req.}{req:band-a-identity-link}\\
\addlinespace
\iface{public_structure}\newline (Party, a\_A, a\_B, sigma, r\textasciicircum{}0\_h, h, phi, Chat) & always & \mref{Def.}{def:spec}, \mref{Rem.}{rem:sidecond}, \mref{Thm.}{thm:sifting-and-key}, \mref{Def.}{def:compile}\\
\addlinespace
\iface{constraint_entries}\newline \{(hat O\_k, bowtie\_k, rhs\_k)\} & always; m \textgreater{}= 0, EMPTY IS THE DEFAULT ROUTE & \mref{Def.}{def:spec}, \mref{Def.}{def:homog}, \mref{Def.}{def:declaration}, \mref{Def.}{def:validator}, \mref{Def.}{def:compile}\\
\addlinespace
\iface{attack_set_license} & optional; REQUIRED by gate (G3)(b) whenever m \textgreater{}= 1 and a key-length-eligible declaration is wanted & \mref{Def.}{def:eligible}, \mref{Thm.}{thm:gates}\\
\addlinespace
\iface{feasibility_witness}\newline omega* & optional; required by branch (c) of check V7 when m \textgreater{}= 1 and the default candidate fails & \mref{Def.}{def:validator}, \mref{Thm.}{thm:feasibility}, \mref{Def.}{def:spec}\\
\addlinespace
\iface{chain_provenance} & optional; passed only by the part of the engine that ran the reductions module and docked its artefacts into the chain, and absent on every other route & \mref{Req.}{req:link-row-provenance}, \mref{Rem.}{rem:obligations-discharged}, \mref{Ass.}{ass:link}\\
\addlinespace
\iface{alice_sector_factor}\newline (T, (d\_A\_t)\_t) with a party label & general format only; absent is \textbar{}T\textbar{} = 1 with A\_t0 = A, which is the basic format & \mref{Def.}{def:declaration}, \mref{Def.}{def:spec}, \mref{Def.}{def:validator}, \mref{Def.}{def:target}\\
\addlinespace
\iface{alice_form}\newline `povm` or `ensemble` & general format only; absent is `povm`, which is the basic format & \mref{Def.}{def:declaration}, \mref{Def.}{def:spec}, \mref{Def.}{def:compile}, \mref{Def.}{def:validator}\\
\addlinespace
\iface{alice_ensemble}\newline \{N\_x\}\_\{x in X\} & general format, ensemble form only; in place of `alice\_povm` & \mref{Def.}{def:declaration}, \mref{Def.}{def:exact}, \mref{Def.}{def:validator}, \mref{Thm.}{thm:sifting-and-key}\\
\addlinespace
\iface{receiver_parts}\newline ((party\_f, physical system\_f, \{F\textasciicircum{}ph,f\_w\}, optional P\_f))\_\{f in F\} & general format only; a one-entry list IS the pair `receiver\_povm` plus `choice\_weighting` of the rows above, which is the basic format & \mref{Def.}{def:spec}, \mref{Def.}{def:compile}, \mref{Def.}{def:weighted}\\
\addlinespace
\iface{reduction_chain.receiver_segments}\newline (R\textasciicircum{}f\_rec)\_\{f in F\} & general format only; one segment IS the `reduction\_chain` row's Ch\_rec, which is the basic format & \mref{Def.}{def:spec}, \mref{Def.}{def:chaincond}, \mref{Thm.}{thm:chain-consumption}\\
\addlinespace
\iface{reduction_chain.source_segment}\newline Ch\_src with its obligation records & always present, and EMPTY for every matrix-declared specification; a NON-EMPTY segment may be presented, and is then decided as `domain` and `what\_this\_module\_decides\_on\_it` say & \mref{Def.}{def:spec}, \mref{Def.}{def:chaincond}, \extref{../transmitter/contract.pdf}{transmitter/def:source-segment}, \mref{Thm.}{thm:chain-consumption}\\
\addlinespace
\iface{public_structure.transcript}\newline ((pi\_m, C\_m, a\_m))\_\{m = 1..L\} & general format only; L = 2 with pi = (A, B) and a\_2 independent of kappa\_1 IS the pair `announce\_alice`, `announce\_bob`, which is the basic format & \mref{Def.}{def:declaration}, \mref{Def.}{def:spec}, \mref{Thm.}{thm:sifting-and-key}, \mref{Def.}{def:target}\\
\addlinespace
\iface{constraint_entries.scope_and_tag}\newline (O\_k, rel\_k, rhs\_k; f\_k) & general format only; one measured factor and no `party\_local` tag is the basic format, where the scope is that factor and the two tags are the two of def:homog & \mref{Def.}{def:declaration}, \mref{Def.}{def:homog}, \mref{Def.}{def:spec}, \mref{Def.}{def:validator}, \mref{Def.}{def:eligible}\\
\addlinespace
\iface{feasibility_witness.joint}\newline omega* & general format only; \textbar{}T\textbar{} = 1 is the `feasibility\_witness` row above & \mref{Def.}{def:validator}, \mref{Thm.}{thm:feasibility}, \mref{Def.}{def:spec}\\
\bottomrule
\end{xltabular}}
\subsection*{Outputs: Protocol compiler (protocol-compiler)}
\noindent{\small From \href{interface.yaml}{\nolinkurl{interface.yaml}}.}\par\smallskip
{\footnotesize\begin{xltabular}{\linewidth}{@{}>{\raggedright\arraybackslash}p{4.4cm}>{\raggedright\arraybackslash}X>{\raggedright\arraybackslash}p{2.9cm}@{}}
\toprule
Key and symbol & Presence & Fixed by\\
\midrule
\endhead
\iface{alice_dim}\newline d\_A & always & \mref{Def.}{def:declaration}, \mref{Def.}{def:spec}\\
\addlinespace
\iface{alice_marginal}\newline tau\_A, or null & always & \mref{Def.}{def:declaration}, \mref{Def.}{def:target}, \mref{Def.}{def:validator}, \mref{Def.}{def:compile}\\
\addlinespace
\iface{source_replacement_license} & emitted exactly when trusted mode is emitted at d\_A \textgreater{}= 2 & \mref{Def.}{def:eligible}, \mref{Thm.}{thm:gates}, \extref{../transmitter/contract.pdf}{transmitter/ass:src}, \extref{../transmitter/contract.pdf}{transmitter/def:srcert}, \mref{Req.}{req:source-licence-passed-on}\\
\addlinespace
\iface{party_partition}\newline Party & always & \mref{Def.}{def:declaration}, \mref{Def.}{def:spec}, \mref{Def.}{def:target}\\
\addlinespace
\iface{sector_index_set}\newline S & always & \mref{Def.}{def:declaration}, \mref{Req.}{req:declaration-keys}, \mref{Def.}{def:target}, \mref{Def.}{def:validator}\\
\addlinespace
\iface{sector_dims}\newline (d\_s)\_s & always & \mref{Def.}{def:declaration}, \mref{Def.}{def:validator}\\
\addlinespace
\iface{sector_block_kind} & always & \mref{Def.}{def:declaration}, \mref{Req.}{req:declaration-keys}\\
\addlinespace
\iface{alice_povm}\newline \{M\_x\}\_\{x in X\} & always & \mref{Def.}{def:declaration}, \mref{Def.}{def:target}, \mref{Def.}{def:validator}, \mref{Def.}{def:exact}\\
\addlinespace
\iface{bob_povm}\newline \{F\_v\}\_\{v in V\}, or \{N\_(v,f)\}\_\{(v,f) in V\_P\} & always & \mref{Def.}{def:declaration}, \mref{Def.}{def:target}, \mref{Def.}{def:exact}, \mref{Def.}{def:validator}, \mref{Def.}{def:weighted}, \mref{Thm.}{thm:choice-encoding}, \mref{Def.}{def:compile}\\
\addlinespace
\iface{reference_companion}\newline D\textasciicircum{}ref & exactly where the last link of the receiver segment of one or more measured factors was emitted with a reference companion (reductions/req:reference-companion); one companion for all of them & \mref{Req.}{req:reference-companion-compiled}, \mref{Def.}{def:compile}, \mref{Def.}{def:weighted}, \mref{Lem.}{lem:weighted}, \mref{Lem.}{lem:weighted-factor}, \extref{../meat-core/contract.pdf}{meat-core/prop:reference-companion}\\
\addlinespace
\iface{announce_alice}\newline a\_A & always & \mref{Def.}{def:declaration}, \mref{Def.}{def:target}\\
\addlinespace
\iface{announce_bob}\newline a\_B & always & \mref{Def.}{def:declaration}, \mref{Def.}{def:target}, \mref{Def.}{def:compile}, \mref{Rem.}{rem:alphabet}\\
\addlinespace
\iface{accept_statistic}\newline phi: Kann -\textgreater{} Chat & always & \mref{Def.}{def:phi}, \mref{Def.}{def:declaration}, \mref{Def.}{def:target}, \mref{Rem.}{rem:phitotal}\\
\addlinespace
\iface{key_map}\newline r: X x V -\textgreater{} Z & always & \mref{Def.}{def:compile}, \mref{Thm.}{thm:sifting-and-key}, \mref{Def.}{def:target}, \mref{Def.}{def:validator}\\
\addlinespace
\iface{key_holder}\newline h & always & \mref{Def.}{def:declaration}, \mref{Def.}{def:target}, \mref{Def.}{def:validator}, \mref{Thm.}{thm:sifting-and-key}\\
\addlinespace
\iface{constraint_list}\newline CList & always; EMPTY IS THE DEFAULT ROUTE & \mref{Def.}{def:declaration}, \mref{Def.}{def:compile}, \mref{Def.}{def:target}, \mref{Def.}{def:homog}, \mref{Def.}{def:validator}\\
\addlinespace
\iface{constraint_rhs_tag} & one per entry of CList & \mref{Def.}{def:homog}, \mref{Def.}{def:declaration}, \mref{Def.}{def:validator}\\
\addlinespace
\iface{attack_set_license} & present only if supplied & \mref{Def.}{def:eligible}, \mref{Thm.}{thm:gates}\\
\addlinespace
\iface{acceptance_alphabet}\newline Chat & always & \mref{Def.}{def:phi}, \mref{Def.}{def:declaration}\\
\addlinespace
\iface{label_index_map}\newline idxm & always & \mref{Def.}{def:compile}, \mref{Def.}{def:target}, \mref{Thm.}{thm:typing}, \mref{Def.}{def:validator}, \mref{Def.}{def:exact}\\
\addlinespace
\iface{reduction_chain_record}\newline Chain & always & \mref{Def.}{def:spec}, \mref{Def.}{def:chaincond}, \mref{Thm.}{thm:chain-consumption}, \mref{Rem.}{rem:relabel}, \mref{Req.}{req:chain-record-slots}\\
\addlinespace
\iface{choice_weighting_record}\newline (F, t, P, the orders) & emitted exactly when a `choice\_weighting` was declared at some measured factor & \mref{Def.}{def:weighted}, \mref{Def.}{def:compile}, \mref{Req.}{req:choice-weighting-record}, \mref{Rem.}{rem:relabel}\\
\addlinespace
\iface{key_length_eligible} & always & \mref{Def.}{def:eligible}, \mref{Thm.}{thm:gates}\\
\addlinespace
\iface{gate_report} & always & \mref{Def.}{def:eligible}, \mref{Thm.}{thm:gates}, \mref{Rem.}{rem:degenerate}, \mref{Def.}{def:exact}, \mref{Def.}{def:validator}, \mref{Def.}{def:chaincond}\\
\addlinespace
\iface{refusal} & present exactly when the compiler refuses & \mref{Thm.}{thm:gates}, \mref{Def.}{def:validator}, \mref{Def.}{def:chaincond}\\
\addlinespace
\iface{assumption_ledger}\newline A-... & always & ---\\
\addlinespace
\iface{assumption_bindings} & always; one entry per ledger row, in its order; empty on a refusal & \mref{Ass.}{ass:spec}, \mref{Ass.}{ass:alice-meas}, \mref{Ass.}{ass:rhom}, \mref{Ass.}{ass:coin}, \mref{Ass.}{ass:link}, \mref{Ass.}{ass:atk}, \mref{Ass.}{ass:ideal}, \extref{../transmitter/contract.pdf}{transmitter/ass:ens}, \extref{../transmitter/contract.pdf}{transmitter/ass:src}, \extref{../transmitter/contract.pdf}{transmitter/decoy/ass:wcp}\\
\addlinespace
\iface{mathematical_spec_version} & always & ---\\
\addlinespace
\iface{sector_factorisation}\newline the ordered factor list of (DX1) & general format only; absent is one measured factor and no A factor & \mref{Def.}{def:declaration}, \mref{Def.}{def:compile}, \mref{Def.}{def:validator}\\
\addlinespace
\iface{alice_form}\newline `povm` or `ensemble` & general format only; absent is `povm` & \mref{Def.}{def:declaration}, \mref{Def.}{def:compile}\\
\addlinespace
\iface{alice_ensemble}\newline \{N\_x\}\_\{x in X\} & general format, ensemble form only; in place of `alice\_povm` & \mref{Def.}{def:declaration}, \mref{Def.}{def:exact}\\
\addlinespace
\iface{alice_marginal.blocks}\newline (tau\_A\textasciicircum{}(t))\_t & general format with an A factor; \textbar{}T\textbar{} = 1 is the `alice\_marginal` row & \mref{Def.}{def:target}, \mref{Def.}{def:validator}\\
\addlinespace
\iface{bob_povm.per_factor}\newline (\{F\textasciicircum{}f\_v\_f\})\_\{f in F\} & general format only; one factor is the `bob\_povm` row & \mref{Def.}{def:declaration}, \mref{Def.}{def:compile}, \mref{Thm.}{thm:chain-consumption}\\
\addlinespace
\iface{transcript}\newline ((pi\_m, C\_m, a\_m))\_\{m = 1..L\} & general format only; L = 2 is the pair `announce\_alice`, `announce\_bob`, which stays the emitted form of that case & \mref{Def.}{def:declaration}, \mref{Def.}{def:compile}, \mref{Thm.}{thm:sifting-and-key}\\
\addlinespace
\iface{constraint_scope}\newline f\_k & general format only; one measured factor is the scope of every entry & \mref{Def.}{def:declaration}, \mref{Def.}{def:spec}, \mref{Def.}{def:compile}\\
\addlinespace
\iface{constraint_rhs_tag.party_local}\newline the third tag value & general format only & \mref{Def.}{def:homog}, \mref{Def.}{def:eligible}, \mref{Thm.}{thm:gates}\\
\bottomrule
\end{xltabular}}
\clearpage
\bookmarksetupnext{level=part}
\section*{Claim index}\addcontentsline{toc}{section}{Claim index}
{\small\begin{xltabular}{\linewidth}{@{}>{\raggedright\arraybackslash}p{2.7cm}>{\raggedright\arraybackslash}Xr@{}}
\toprule
Claim & Title and identifier & Page\\
\midrule
\endhead
\mref{Definition}{def:declaration} & Protocol declaration\newline{\footnotesize\color{metaGrey}\claimid{protocol-compiler/def:declaration}} & \pageref{def:declaration}\\
\addlinespace[2pt]
\mref{Remark}{rem:party-label-restriction} & Party labels in the general format\newline{\footnotesize\color{metaGrey}\claimid{protocol-compiler/rem:party-label-restriction}} & \pageref{rem:party-label-restriction}\\
\addlinespace[2pt]
\mref{Definition}{def:exact} & Exact admissibility of declared operator data\newline{\footnotesize\color{metaGrey}\claimid{protocol-compiler/def:exact}} & \pageref{def:exact}\\
\addlinespace[2pt]
\mref{Remark}{rem:summation-condition} & The summation condition beside completeness\newline{\footnotesize\color{metaGrey}\claimid{protocol-compiler/rem:summation-condition}} & \pageref{rem:summation-condition}\\
\addlinespace[2pt]
\mref{Definition}{def:target} & Target conditions of a well-formed declaration\newline{\footnotesize\color{metaGrey}\claimid{protocol-compiler/def:target}} & \pageref{def:target}\\
\addlinespace[2pt]
\mref{Remark}{rem:engine-admission} & Relation to the consuming engine's admission conditions\newline{\footnotesize\color{metaGrey}\claimid{protocol-compiler/rem:engine-admission}} & \pageref{rem:engine-admission}\\
\addlinespace[2pt]
\mref{Definition}{def:feasible} & Feasible set\newline{\footnotesize\color{metaGrey}\claimid{protocol-compiler/def:feasible}} & \pageref{def:feasible}\\
\addlinespace[2pt]
\mref{Lemma}{lem:hermitian-constraint-values} & Hermiticity of the constraint values\newline{\footnotesize\color{metaGrey}\claimid{protocol-compiler/lem:hermitian-constraint-values}} & \pageref{lem:hermitian-constraint-values}\\
\addlinespace[2pt]
\mref{Remark}{rem:wellposed} & Hermiticity as a declared condition\newline{\footnotesize\color{metaGrey}\claimid{protocol-compiler/rem:wellposed}} & \pageref{rem:wellposed}\\
\addlinespace[2pt]
\mref{Remark}{rem:block-restriction} & Restriction of the feasible set to block-diagonal states\newline{\footnotesize\color{metaGrey}\claimid{protocol-compiler/rem:block-restriction}} & \pageref{rem:block-restriction}\\
\addlinespace[2pt]
\mref{Definition}{def:homog} & Licensed constraint list\newline{\footnotesize\color{metaGrey}\claimid{protocol-compiler/def:homog}} & \pageref{def:homog}\\
\addlinespace[2pt]
\mref{Remark}{rem:party-local-soundness} & Soundness of a party-local entry\newline{\footnotesize\color{metaGrey}\claimid{protocol-compiler/rem:party-local-soundness}} & \pageref{rem:party-local-soundness}\\
\addlinespace[2pt]
\mref{Example}{ex:two-party-shape} & Declaration of a two-party protocol with an untrusted source\newline{\footnotesize\color{metaGrey}\claimid{protocol-compiler/ex:two-party-shape}} & \pageref{ex:two-party-shape}\\
\addlinespace[2pt]
\mref{Definition}{def:spec} & Protocol specification\newline{\footnotesize\color{metaGrey}\claimid{protocol-compiler/def:spec}} & \pageref{def:spec}\\
\addlinespace[2pt]
\mref{Definition}{def:well-formed-spec} & Well-formed specification\newline{\footnotesize\color{metaGrey}\claimid{protocol-compiler/def:well-formed-spec}} & \pageref{def:well-formed-spec}\\
\addlinespace[2pt]
\mref{Definition}{def:compile} & The compilation map\newline{\footnotesize\color{metaGrey}\claimid{protocol-compiler/def:compile}} & \pageref{def:compile}\\
\addlinespace[2pt]
\mref{Theorem}{thm:typing} & Typing of an emitted declaration\newline{\footnotesize\color{metaGrey}\claimid{protocol-compiler/thm:typing}} & \pageref{thm:typing}\\
\addlinespace[2pt]
\mref{Remark}{rem:rbuilt} & The raw-key map built from sifting and the key rule\newline{\footnotesize\color{metaGrey}\claimid{protocol-compiler/rem:rbuilt}} & \pageref{rem:rbuilt}\\
\addlinespace[2pt]
\mref{Definition}{def:phi} & Declared statistics functional\newline{\footnotesize\color{metaGrey}\claimid{protocol-compiler/def:phi}} & \pageref{def:phi}\\
\addlinespace[2pt]
\mref{Theorem}{thm:measurements-and-statistics} & Exact measurement families and the statistics functional\newline{\footnotesize\color{metaGrey}\claimid{protocol-compiler/thm:measurements-and-statistics}} & \pageref{thm:measurements-and-statistics}\\
\addlinespace[2pt]
\mref{Remark}{rem:phitotal} & Totality of the acceptance-statistic map\newline{\footnotesize\color{metaGrey}\claimid{protocol-compiler/rem:phitotal}} & \pageref{rem:phitotal}\\
\addlinespace[2pt]
\mref{Definition}{def:weighted} & Choice-weighted families\newline{\footnotesize\color{metaGrey}\claimid{protocol-compiler/def:weighted}} & \pageref{def:weighted}\\
\addlinespace[2pt]
\mref{Lemma}{lem:weighted} & Properties of the choice-weighted families\newline{\footnotesize\color{metaGrey}\claimid{protocol-compiler/lem:weighted}} & \pageref{lem:weighted}\\
\addlinespace[2pt]
\mref{Lemma}{lem:weighted-factor} & A factor enclosure through a setting-dependent weighting\newline{\footnotesize\color{metaGrey}\claimid{protocol-compiler/lem:weighted-factor}} & \pageref{lem:weighted-factor}\\
\addlinespace[2pt]
\mref{Theorem}{thm:choice-encoding} & Encoding of public choices in the measurement operators\newline{\footnotesize\color{metaGrey}\claimid{protocol-compiler/thm:choice-encoding}} & \pageref{thm:choice-encoding}\\
\addlinespace[2pt]
\mref{Remark}{rem:exactness} & Scope of the exact encoding\newline{\footnotesize\color{metaGrey}\claimid{protocol-compiler/rem:exactness}} & \pageref{rem:exactness}\\
\addlinespace[2pt]
\mref{Remark}{rem:condfamily} & Deciding the conditional families\newline{\footnotesize\color{metaGrey}\claimid{protocol-compiler/rem:condfamily}} & \pageref{rem:condfamily}\\
\addlinespace[2pt]
\mref{Example}{ex:condfamily} & An exact weighted family built from incomplete conditional families\newline{\footnotesize\color{metaGrey}\claimid{protocol-compiler/ex:condfamily}} & \pageref{ex:condfamily}\\
\addlinespace[2pt]
\mref{Remark}{rem:relabel} & The setting-dependent relabelling after the chain\newline{\footnotesize\color{metaGrey}\claimid{protocol-compiler/rem:relabel}} & \pageref{rem:relabel}\\
\addlinespace[2pt]
\mref{Remark}{rem:alphabet} & The outcome alphabet is chosen upstream\newline{\footnotesize\color{metaGrey}\claimid{protocol-compiler/rem:alphabet}} & \pageref{rem:alphabet}\\
\addlinespace[2pt]
\mref{Theorem}{thm:sifting-and-key} & Announcement locality, sifting and party computability\newline{\footnotesize\color{metaGrey}\claimid{protocol-compiler/thm:sifting-and-key}} & \pageref{thm:sifting-and-key}\\
\addlinespace[2pt]
\mref{Remark}{rem:sidecond} & Necessity of the key-rule side condition\newline{\footnotesize\color{metaGrey}\claimid{protocol-compiler/rem:sidecond}} & \pageref{rem:sidecond}\\
\addlinespace[2pt]
\mref{Remark}{rem:vacuous} & Party computability read against the readouts\newline{\footnotesize\color{metaGrey}\claimid{protocol-compiler/rem:vacuous}} & \pageref{rem:vacuous}\\
\addlinespace[2pt]
\mref{Example}{ex:transcript} & A three-message transcript\newline{\footnotesize\color{metaGrey}\claimid{protocol-compiler/ex:transcript}} & \pageref{ex:transcript}\\
\addlinespace[2pt]
\mref{Theorem}{thm:feasibility} & Feasibility by an explicit or a verified candidate\newline{\footnotesize\color{metaGrey}\claimid{protocol-compiler/thm:feasibility}} & \pageref{thm:feasibility}\\
\addlinespace[2pt]
\mref{Remark}{rem:nosearch} & Feasibility by witness rather than by search\newline{\footnotesize\color{metaGrey}\claimid{protocol-compiler/rem:nosearch}} & \pageref{rem:nosearch}\\
\addlinespace[2pt]
\mref{Definition}{def:validator} & The emitted validator\newline{\footnotesize\color{metaGrey}\claimid{protocol-compiler/def:validator}} & \pageref{def:validator}\\
\addlinespace[2pt]
\mref{Theorem}{thm:validator-soundness} & Soundness, completeness and decidability of the validator\newline{\footnotesize\color{metaGrey}\claimid{protocol-compiler/thm:validator-soundness}} & \pageref{thm:validator-soundness}\\
\addlinespace[2pt]
\mref{Remark}{rem:completeness} & The two directions of the validator\newline{\footnotesize\color{metaGrey}\claimid{protocol-compiler/rem:completeness}} & \pageref{rem:completeness}\\
\addlinespace[2pt]
\mref{Definition}{def:chaincond} & Interface conditions on a chain\newline{\footnotesize\color{metaGrey}\claimid{protocol-compiler/def:chaincond}} & \pageref{def:chaincond}\\
\addlinespace[2pt]
\mref{Remark}{rem:chaincond-reductions} & Chain conditions here and in the reductions module\newline{\footnotesize\color{metaGrey}\claimid{protocol-compiler/rem:chaincond-reductions}} & \pageref{rem:chaincond-reductions}\\
\addlinespace[2pt]
\mref{Remark}{rem:source-segment-rules} & Grounds of the source-segment rules\newline{\footnotesize\color{metaGrey}\claimid{protocol-compiler/rem:source-segment-rules}} & \pageref{rem:source-segment-rules}\\
\addlinespace[2pt]
\mref{Remark}{rem:chain-and-validator} & Chain conditions and validator checks\newline{\footnotesize\color{metaGrey}\claimid{protocol-compiler/rem:chain-and-validator}} & \pageref{rem:chain-and-validator}\\
\addlinespace[2pt]
\mref{Theorem}{thm:chain-consumption} & Decided identities and named obligations of a consumed chain\newline{\footnotesize\color{metaGrey}\claimid{protocol-compiler/thm:chain-consumption}} & \pageref{thm:chain-consumption}\\
\addlinespace[2pt]
\mref{Remark}{rem:coarse} & Direction of the outcome correspondence\newline{\footnotesize\color{metaGrey}\claimid{protocol-compiler/rem:coarse}} & \pageref{rem:coarse}\\
\addlinespace[2pt]
\mref{Remark}{rem:obligations-discharged} & Discharge of the chain obligations by the reductions module\newline{\footnotesize\color{metaGrey}\claimid{protocol-compiler/rem:obligations-discharged}} & \pageref{rem:obligations-discharged}\\
\addlinespace[2pt]
\mref{Definition}{def:banda} & Band-A specification\newline{\footnotesize\color{metaGrey}\claimid{protocol-compiler/def:banda}} & \pageref{def:banda}\\
\addlinespace[2pt]
\mref{Theorem}{thm:band-a} & The Band-A route\newline{\footnotesize\color{metaGrey}\claimid{protocol-compiler/thm:band-a}} & \pageref{thm:band-a}\\
\addlinespace[2pt]
\mref{Definition}{def:eligible} & Key-length eligibility\newline{\footnotesize\color{metaGrey}\claimid{protocol-compiler/def:eligible}} & \pageref{def:eligible}\\
\addlinespace[2pt]
\mref{Remark}{rem:gate-by-certificate-form} & The source gate on each certificate form\newline{\footnotesize\color{metaGrey}\claimid{protocol-compiler/rem:gate-by-certificate-form}} & \pageref{rem:gate-by-certificate-form}\\
\addlinespace[2pt]
\mref{Definition}{def:gate-outcome-table} & The gate outcome table\newline{\footnotesize\color{metaGrey}\claimid{protocol-compiler/def:gate-outcome-table}} & \pageref{def:gate-outcome-table}\\
\addlinespace[2pt]
\mref{Theorem}{thm:gates} & The emission gates and their outcomes\newline{\footnotesize\color{metaGrey}\claimid{protocol-compiler/thm:gates}} & \pageref{thm:gates}\\
\addlinespace[2pt]
\mref{Remark}{rem:degenerate} & The licensed-class gate in untrusted-source mode and at $d_A=1$\newline{\footnotesize\color{metaGrey}\claimid{protocol-compiler/rem:degenerate}} & \pageref{rem:degenerate}\\
\addlinespace[2pt]
\mref{Proposition}{prop:conservative} & Dropping a trusted marginal enlarges the feasible set\newline{\footnotesize\color{metaGrey}\claimid{protocol-compiler/prop:conservative}} & \pageref{prop:conservative}\\
\addlinespace[2pt]
\mref{Remark}{rem:notfallback} & Dropping a marginal, never replacing it\newline{\footnotesize\color{metaGrey}\claimid{protocol-compiler/rem:notfallback}} & \pageref{rem:notfallback}\\
\addlinespace[2pt]
\mref{Assumption}{ass:spec} & Specification faithful to the apparatus\newline{\footnotesize\color{metaGrey}\claimid{protocol-compiler/ass:spec}} & \pageref{ass:spec}\\
\addlinespace[2pt]
\mref{Assumption}{ass:alice-meas} & Alice's declared measurement where no source certificate is demanded\newline{\footnotesize\color{metaGrey}\claimid{protocol-compiler/ass:alice-meas}} & \pageref{ass:alice-meas}\\
\addlinespace[2pt]
\mref{Assumption}{ass:rhom} & Round homogeneity of the declaration\newline{\footnotesize\color{metaGrey}\claimid{protocol-compiler/ass:rhom}} & \pageref{ass:rhom}\\
\addlinespace[2pt]
\mref{Assumption}{ass:coin} & Public choices as private coins of declared bias\newline{\footnotesize\color{metaGrey}\claimid{protocol-compiler/ass:coin}} & \pageref{ass:coin}\\
\addlinespace[2pt]
\mref{Remark}{rem:coin-scope} & Scope of the private-coin assumption\newline{\footnotesize\color{metaGrey}\claimid{protocol-compiler/rem:coin-scope}} & \pageref{rem:coin-scope}\\
\addlinespace[2pt]
\mref{Assumption}{ass:link} & Soundness of every reduction link\newline{\footnotesize\color{metaGrey}\claimid{protocol-compiler/ass:link}} & \pageref{ass:link}\\
\addlinespace[2pt]
\mref{Assumption}{ass:atk} & Attack-set licence\newline{\footnotesize\color{metaGrey}\claimid{protocol-compiler/ass:atk}} & \pageref{ass:atk}\\
\addlinespace[2pt]
\mref{Assumption}{ass:ideal} & Ideal qubit receiver on the Band-A route\newline{\footnotesize\color{metaGrey}\claimid{protocol-compiler/ass:ideal}} & \pageref{ass:ideal}\\
\addlinespace[2pt]
\mref{Requirement}{req:declaration-keys} & Machine keys of the emitted declaration\newline{\footnotesize\color{metaGrey}\claimid{protocol-compiler/req:declaration-keys}} & \pageref{req:declaration-keys}\\
\addlinespace[2pt]
\mref{Requirement}{req:operator-form} & Operators on the measured system held as block lists\newline{\footnotesize\color{metaGrey}\claimid{protocol-compiler/req:operator-form}} & \pageref{req:operator-form}\\
\addlinespace[2pt]
\mref{Requirement}{req:pass-through} & Declared data passed through unchanged\newline{\footnotesize\color{metaGrey}\claimid{protocol-compiler/req:pass-through}} & \pageref{req:pass-through}\\
\addlinespace[2pt]
\mref{Requirement}{req:derived-objects} & Per-factor decisions on the declared blocks\newline{\footnotesize\color{metaGrey}\claimid{protocol-compiler/req:derived-objects}} & \pageref{req:derived-objects}\\
\addlinespace[2pt]
\mref{Requirement}{req:tag-reading} & Reading the tag of a constraint entry\newline{\footnotesize\color{metaGrey}\claimid{protocol-compiler/req:tag-reading}} & \pageref{req:tag-reading}\\
\addlinespace[2pt]
\mref{Requirement}{req:no-window} & Exact decisions without a tolerance\newline{\footnotesize\color{metaGrey}\claimid{protocol-compiler/req:no-window}} & \pageref{req:no-window}\\
\addlinespace[2pt]
\mref{Requirement}{req:source-licence-passed-on} & The source-replacement licence passed on with the declaration\newline{\footnotesize\color{metaGrey}\claimid{protocol-compiler/req:source-licence-passed-on}} & \pageref{req:source-licence-passed-on}\\
\addlinespace[2pt]
\mref{Requirement}{req:da1-source-identity} & Letter-independent emission at a one-dimensional trusted A\newline{\footnotesize\color{metaGrey}\claimid{protocol-compiler/req:da1-source-identity}} & \pageref{req:da1-source-identity}\\
\addlinespace[2pt]
\mref{Requirement}{req:choice-weighting-record} & Emission of the weighted families and the weighting record\newline{\footnotesize\color{metaGrey}\claimid{protocol-compiler/req:choice-weighting-record}} & \pageref{req:choice-weighting-record}\\
\addlinespace[2pt]
\mref{Requirement}{req:conditional-family-decision} & Decision on each declared conditional family\newline{\footnotesize\color{metaGrey}\claimid{protocol-compiler/req:conditional-family-decision}} & \pageref{req:conditional-family-decision}\\
\addlinespace[2pt]
\mref{Requirement}{req:binary64-weights} & Declared weights as binary64 input\newline{\footnotesize\color{metaGrey}\claimid{protocol-compiler/req:binary64-weights}} & \pageref{req:binary64-weights}\\
\addlinespace[2pt]
\mref{Requirement}{req:weight-preparation} & Preparation of user weights before the module boundary\newline{\footnotesize\color{metaGrey}\claimid{protocol-compiler/req:weight-preparation}} & \pageref{req:weight-preparation}\\
\addlinespace[2pt]
\mref{Requirement}{req:setting-only-weighting} & Setting-dependent weighting read through the setting only\newline{\footnotesize\color{metaGrey}\claimid{protocol-compiler/req:setting-only-weighting}} & \pageref{req:setting-only-weighting}\\
\addlinespace[2pt]
\mref{Requirement}{req:weight-list} & Decision on the declared weight list\newline{\footnotesize\color{metaGrey}\claimid{protocol-compiler/req:weight-list}} & \pageref{req:weight-list}\\
\addlinespace[2pt]
\mref{Requirement}{req:stochastic-weights} & Stochastic matrix of a setting-dependent weighting\newline{\footnotesize\color{metaGrey}\claimid{protocol-compiler/req:stochastic-weights}} & \pageref{req:stochastic-weights}\\
\addlinespace[2pt]
\mref{Requirement}{req:weighted-family-emission} & Emission decisions on the weighted products\newline{\footnotesize\color{metaGrey}\claimid{protocol-compiler/req:weighted-family-emission}} & \pageref{req:weighted-family-emission}\\
\addlinespace[2pt]
\mref{Requirement}{req:no-lattice-exponent} & No lattice exponent and no snapping of the emission\newline{\footnotesize\color{metaGrey}\claimid{protocol-compiler/req:no-lattice-exponent}} & \pageref{req:no-lattice-exponent}\\
\addlinespace[2pt]
\mref{Requirement}{req:no-fixed-fraction} & No fixed fraction outside the measurement operators\newline{\footnotesize\color{metaGrey}\claimid{protocol-compiler/req:no-fixed-fraction}} & \pageref{req:no-fixed-fraction}\\
\addlinespace[2pt]
\mref{Requirement}{req:validator-refusals} & Refusals of the validator's checks\newline{\footnotesize\color{metaGrey}\claimid{protocol-compiler/req:validator-refusals}} & \pageref{req:validator-refusals}\\
\addlinespace[2pt]
\mref{Requirement}{req:validation-transcript} & Transcript of a check that stops early\newline{\footnotesize\color{metaGrey}\claimid{protocol-compiler/req:validation-transcript}} & \pageref{req:validation-transcript}\\
\addlinespace[2pt]
\mref{Requirement}{req:feasibility-witness} & Default candidate before a supplied witness\newline{\footnotesize\color{metaGrey}\claimid{protocol-compiler/req:feasibility-witness}} & \pageref{req:feasibility-witness}\\
\addlinespace[2pt]
\mref{Requirement}{req:witness-refusal-meaning} & Reporting a refusal for want of a feasibility witness\newline{\footnotesize\color{metaGrey}\claimid{protocol-compiler/req:witness-refusal-meaning}} & \pageref{req:witness-refusal-meaning}\\
\addlinespace[2pt]
\mref{Requirement}{req:physical-family-decision} & Decision on the presented physical receiver family\newline{\footnotesize\color{metaGrey}\claimid{protocol-compiler/req:physical-family-decision}} & \pageref{req:physical-family-decision}\\
\addlinespace[2pt]
\mref{Requirement}{req:chain-refusals} & Refusals of the interface conditions on a chain\newline{\footnotesize\color{metaGrey}\claimid{protocol-compiler/req:chain-refusals}} & \pageref{req:chain-refusals}\\
\addlinespace[2pt]
\mref{Requirement}{req:band-a-ideal-receiver} & Band A only within the ideal qubit receiver\newline{\footnotesize\color{metaGrey}\claimid{protocol-compiler/req:band-a-ideal-receiver}} & \pageref{req:band-a-ideal-receiver}\\
\addlinespace[2pt]
\mref{Requirement}{req:band-a-identity-link} & Band A only on an identity link that delivers the declared family\newline{\footnotesize\color{metaGrey}\claimid{protocol-compiler/req:band-a-identity-link}} & \pageref{req:band-a-identity-link}\\
\addlinespace[2pt]
\mref{Requirement}{req:chain-record-slots} & Emission of both chain segments and every link slot\newline{\footnotesize\color{metaGrey}\claimid{protocol-compiler/req:chain-record-slots}} & \pageref{req:chain-record-slots}\\
\addlinespace[2pt]
\mref{Requirement}{req:link-row-provenance} & Omission of the link-soundness row only on a proved chain\newline{\footnotesize\color{metaGrey}\claimid{protocol-compiler/req:link-row-provenance}} & \pageref{req:link-row-provenance}\\
\addlinespace[2pt]
\mref{Requirement}{req:reference-companion-compiled} & The reference companion compiled beside the declaration\newline{\footnotesize\color{metaGrey}\claimid{protocol-compiler/req:reference-companion-compiled}} & \pageref{req:reference-companion-compiled}\\
\addlinespace[2pt]
\mref{Requirement}{req:eligibility-flag} & The key-length eligibility flag\newline{\footnotesize\color{metaGrey}\claimid{protocol-compiler/req:eligibility-flag}} & \pageref{req:eligibility-flag}\\
\addlinespace[2pt]
\mref{Requirement}{req:gate-outcomes} & Outcome of each gate failure\newline{\footnotesize\color{metaGrey}\claimid{protocol-compiler/req:gate-outcomes}} & \pageref{req:gate-outcomes}\\
\addlinespace[2pt]
\mref{Requirement}{req:attestation-presence} & Attestations tested for presence only\newline{\footnotesize\color{metaGrey}\claimid{protocol-compiler/req:attestation-presence}} & \pageref{req:attestation-presence}\\
\addlinespace[2pt]
\mref{Requirement}{req:vacuous-gate-report} & A vacuous pass recorded as vacuous\newline{\footnotesize\color{metaGrey}\claimid{protocol-compiler/req:vacuous-gate-report}} & \pageref{req:vacuous-gate-report}\\
\bottomrule
\end{xltabular}}

\clearpage
\bookmarksetupnext{level=part}
\printindex

\clearpage\phantomsection\bookmarksetupnext{level=part}%
\addcontentsline{toc}{section}{References}
\bibliographystyle{plain}
\bibliography{../../../refs/references}

\end{document}
